% ============================================================
% rk_error_analysis_deep_dive_20260426_zh.tex
% ------------------------------------------------------------
% PDC 靶点动量重建 —— RK 误差分析详细版
% (Companion deep-dive to rk_reconstruction_status_20260416_zh.tex)
%
% Build: cd docs/reports/reconstruction/momentum_reco/rk
%        latexmk -xelatex rk_error_analysis_deep_dive_20260426_zh.tex
%
% 结构:
%   第一部分: 数学基础   (3 个 \input 片段, ~2000 行)
%   第二部分: 系统误差预算 (3 个 \input 片段, ~1600 行)
%   第三部分: 实证诊断   (1 个 \input 片段, ~1450 行)
%   附录:    glossary + 复现 + 公开问题清单
%
% 创建: 2026-04-26
% 作者: Baiting Tian
% ============================================================

\documentclass[a4paper, 11pt]{article}

\usepackage[margin=2.4cm]{geometry}
\usepackage{ctex}
\usepackage{amsmath, amssymb, amsthm}
\usepackage{booktabs}
\usepackage{multirow}
\usepackage{array}
\usepackage{makecell}
\usepackage{siunitx}
\usepackage{xcolor}
\usepackage{hyperref}
\usepackage{enumitem}
\usepackage{float}
\usepackage{graphicx}
\usepackage{tikz}
\usetikzlibrary{arrows.meta, positioning, calc, decorations.markings}
\usepackage{algorithm}
\usepackage{algpseudocode}
\usepackage{listings}

% [EN] xelatex-friendly listings setup with Chinese-safe basic style.
% [CN] xelatex 下兼容中文与代码并存的 listings 样式。
\lstdefinestyle{deepdive}{
  basicstyle=\footnotesize\ttfamily,
  keywordstyle=\color{blue!60!black},
  stringstyle=\color{red!50!black},
  commentstyle=\color{gray!60!black}\itshape,
  showstringspaces=false,
  numbers=left,
  numberstyle=\tiny\color{gray},
  numbersep=6pt,
  frame=single,
  rulecolor=\color{black!30},
  framesep=4pt,
  tabsize=2,
  breaklines=true,
  breakatwhitespace=true,
  columns=fullflexible,
  keepspaces=true,
  extendedchars=true,
  inputencoding=utf8,
}
\lstset{style=deepdive}

\definecolor{goodgreen}{RGB}{34,139,34}
\definecolor{warnred}{RGB}{200,50,50}

\hypersetup{
  colorlinks=true,
  linkcolor=blue!60!black,
  citecolor=blue!60!black,
  urlcolor=blue!60!black,
  pdftitle={PDC 靶点动量重建 —— RK 误差分析详细版},
  pdfauthor={Baiting Tian},
}

% ===== amsthm 环境声明 (供 Part I 使用) =====
\theoremstyle{plain}
\newtheorem{theorem}{定理}[section]
\newtheorem{proposition}[theorem]{命题}
\newtheorem{lemma}[theorem]{引理}
\newtheorem{corollary}[theorem]{推论}

\theoremstyle{definition}
\newtheorem{definition}[theorem]{定义}

\theoremstyle{remark}
\newtheorem{remark}[theorem]{评注}
\newtheorem{example}[theorem]{例}

% [EN] Replace amsthm 'Proof.' marker with Chinese.
% [CN] 把英文 Proof 标头改成中文证明。
\renewcommand\proofname{证明}

% ===== 跨文档宏 (子片段共用) =====
% [EN] Cross-references to the status report (sibling document).
% [CN] 指向同目录下状态报告的引用宏；子片段以 \StatusReport § X 形式引用。
\providecommand{\StatusReport}{\href{rk_reconstruction_status_20260416_zh.pdf}{状态报告~(2026-04-16)}}

% [EN] Path to deep-dive data root (CSV / figure outputs from new analyses).
% [CN] 详细版数据根目录；后续 \DeepDiveData/path/to/file 形式引用。
% [EN] Underscores escaped to avoid math-mode misparse when used in text.
% [CN] 下划线转义，避免在普通文本中触发数学下标解析。
\providecommand{\DeepDiveData}{../../../../../build/test\_output/ensemble\_n50\_targetframe}

% [EN] Path to status-report's analysis run (kept for legacy figure cross-refs).
% [CN] 状态报告分析 run 路径；与状态报告保持一致以便重用图表引用。
\providecommand{\AnaRoot}{../../../../../results/pdc\_rk\_error\_after\_pdcana\_20260411}

\title{\textbf{PDC 靶点动量重建——RK 误差分析详细版\\
       \large 数学基础 \(\cdot\) 系统误差预算 \(\cdot\) 实证诊断}}
\author{Baiting Tian \\ RIKEN 仁科中心 / DPOL 实验}
\date{2026 年 4 月 26 日}

\begin{document}
\maketitle

% ===== Abstract =====
\begin{abstract}
\noindent
本报告是
\StatusReport\
的\textbf{详细版伴随文档}。状态报告承担"执行摘要"角色，给出 RK 后端从
PDC 命中坐标重建靶点动量的完整管线、六种误差分析方法的统一对照、以及
2026-04-19 坐标系修复后的 744 事件 ensemble 覆盖率结论。本报告并不重述这些
内容，而是沿三条轴线深入：

\textbf{第一部分：数学基础}~——~Cramér-Rao 下界、Fisher 信息矩阵在 SAMURAI
2-PDC 几何下的具体形式（含动量空间 Jacobian $M$ 的解析推导）、Wilks 定理与
Profile Likelihood 的有效边界、Laplace 方法与 Bernstein–von Mises 定理、
Metropolis-Hastings 详细平衡与最优 acceptance 0.234、Geweke $z$ 与 ESS、
Glückstern 解析动量分辨极限、Bethe-Bloch 平均能损与 Landau/Vavilov 涨落
（解释当前 $-10.8\,\mathrm{MeV}/c$ 的 $p_z$ 系统偏差）、Highland 多次散射公式。

\textbf{第二部分：系统误差预算}~——~每条物理来源（PDC 位置分辨、空气/PDC
气体/Sn 靶 多次散射、$dE/dx$、磁场图插值、对位与靶倾角、$(u,v,p)$ 参数化、
LM 收敛盆地）独立成章；每章给出物理机制推导、量级估计、与 744 事件残差对比、
预算行（区分 BIAS 与 WIDTH 列）、修复路径与预测后效果。最后通过\textbf{二次方
和与观测残差并表}调和：当前预算 $p_x$ 预测 $3.22\,\mathrm{MeV}/c$ vs
观测半宽 $5.98\,\mathrm{MeV}/c$（缺口 $1.86\times$）；$p_y$ 预测 $0.83$ vs
观测 $1.40\,\mathrm{MeV}/c$；$p_z$ 预测 $3.18$ vs 观测 $5.57$。诊断指向：约
50\% 的缺口源于 Fisher $\Sigma_\mathrm{meas}$ 缺少 MS 过程噪声（Kalman 风格的
通用修复），约 20\% 来自未调研材料、约 10\% 来自 LM 收敛盆地尾部、约 5\% 来
自参数化非线性。

\textbf{第三部分：实证诊断}~——~Clopper-Pearson vs Wilson vs Wald
覆盖率方法学（含 ensemble 大小 $\to$ CP CI 半宽的扫描表）、pull 分布的定义
与解读（744 事件靶系实测：$p_x$ pull 中位 $-0.026$，$\sigma_z = 4.22$；$p_y$
对应 $-0.066$ 与 $2.48$ —— 与第二部分参数化欠覆盖结论一致）、$\chi^2$ 分布
与失败模式分类（Class A/B/C 阈值表）、5--10 个单事件深度示例（含 $(-100, 0)$
病态点 12\% 失败率剖析）、待执行扫描（B-field $\pm5\%$、$\sigma_\mathrm{wire}$
$0.1\to 2.0\,\mathrm{mm}$、$\alpha_\mathrm{tgt}$、ensemble 规模 10/50/200/1000）
的实施模板。

\textbf{执行摘要}：在 RK 后端继续作为 SAMURAI/DPOL 主重建器的前提下，本报告
给出按修复优先级排序的清单——
(1) Bethe-Bloch $dE/dx$ 注入 RK 步进 \(\Rightarrow\) 消除 $-10.8\,\mathrm{MeV}/c$
$p_z$ 偏差；
(2) MS 过程噪声并入 $\Sigma_\mathrm{meas}$（Kalman 化）\(\Rightarrow\) 关闭 $p_y$
Fisher 欠估、补 50\% 残差缺口；
(3) Cartesian $(p_x,p_y,p_z)$ 参数化 \(\Rightarrow\) 边际 $\sigma_{p_y}^\mathrm{Fisher}$
回升 $0.75 \to 0.85\,\mathrm{MeV}/c$；
(4) 非对称扩展粗格 + 反向初值 \(\Rightarrow\) $(-100,0)$ 失败率 12\% \(\to\) <1\%；
(5) tricubic B-field 插值与全材料调研 \(\Rightarrow\) 末位优化。
\end{abstract}

\tableofcontents
\listoffigures
\listoftables
\listofalgorithms
\newpage

% ============================================================
% 0. 前言、记号约定、阅读路径
% ============================================================
\section*{0. 阅读指南、记号约定、与状态报告的关系}
\addcontentsline{toc}{section}{0. 阅读指南、记号约定、与状态报告的关系}

\paragraph{与状态报告的关系。}
本文档是
\StatusReport
的伴随深入版，后者承担执行摘要、当前性能与覆盖率结论的"骨架"角色，
并在 §坐标系约定、§误差分析、§理论分辨极限、§覆盖率验证 等节给出本文档反复引用的
基础数字（例如 744 事件 $|\vec p|$ 半宽 $5.30\,\mathrm{MeV}/c$、$p_x$
Fisher 68\% 覆盖 $0.80$、Profile $|\vec p|$ 半宽中位 $76.91\,\mathrm{MeV}/c$ 等）。
本文不重述这些数字，而是给出推导、调和、扩展。

\paragraph{记号。}
全文采用以下记号约定（与状态报告一致，新引入处显式声明）：
\begin{itemize}[nosep]
  \item 参数矢量 $\vec\alpha = (\delta x, \delta y, u, v, p)$，
        其中 $u = p_x/p_z$、$v = p_y/p_z$；动量模 $p = |\vec p|$。
        在 \texttt{kFixedTargetPdcOnly} 模式下 $\delta x = \delta y = 0$，
        参数压缩为 $(u, v, p)$。
  \item 物理动量分量记 $\vec p = (p_x, p_y, p_z)$；坐标系以\textbf{靶系
        （beam-as-Z）}为默认；当涉及实验室系（lab）时显式标注 lab。
  \item PDC 测量值 $m_i$；前向模型预言 $f_i(\vec\alpha)$；
        原始残差 $r_i = f_i - m_i$（$i = 1, \dots, 4$ 对应
        PDC1$u$、PDC1$v$、PDC2$u$、PDC2$v$）；白化残差 $\vec w$
        通过 Cholesky $C = LL^\top$ 得到 $\vec w = L^{-1}\vec r$。
  \item Jacobian 定义 $J_{ij} = \partial w_i / \partial \alpha_j$，是
        $4\times 3$（或 \texttt{kThreePointFree} 下 $6\times 5$）矩阵。
  \item 动量空间 Jacobian $M = \partial(p_x,p_y,p_z)/\partial(u,v,p)$，
        其解析形式见第~\ref{sec:partIa}~部分 §3。
  \item 跨文档引用：\StatusReport\ § X 表示状态报告章节 X；
        本文档内部章节交叉引用使用 \texttt{\textbackslash{}ref}。
\end{itemize}

\paragraph{数据来源。}
所有出现的数值或者
(a) 在本文档内由首原理推导，或
(b) 引用状态报告并在引文处给出 \StatusReport\ § X 锚点，或
(c) 标记为 \texttt{\% TODO: needs run R-N} 的待执行实验，
绝不杜撰。CSV 路径汇总于附录 A.2，主要 run 包括：
\StatusReport\ 中的 \texttt{ensemble\_n50\_targetframe} (744 事件，靶系，
2026-04-19 坐标系修复后) 与 \texttt{pdc\_rk\_error\_after\_pdcana\_20260411}
(815 事件，缺 truth)。

\paragraph{阅读路径建议。}
\begin{itemize}[nosep]
  \item 首读：本文 §0 + 摘要 + 第三部分 §1, §10 (Part II reconciliation)
        \(\to\) 快速获得修复优先级清单。
  \item 算法/统计同事：第一部分全篇。
  \item 探测器/实验同事：第二部分全篇 + 第三部分 §6 ($(-100,0)$ 病态)。
  \item 数据分析同事：第三部分全篇 + 附录 A 复现章。
\end{itemize}

\newpage

% ============================================================
% 第一部分: 数学基础
% ============================================================
\part{数学基础}\label{part:math}\label{sec:partIa}

本部分给出 RK 重建管线中用到的全部统计推断机器的完整推导，并把每条结论都
落地到 SAMURAI 2-PDC、4 残差、3 (或 5) 参数的具体几何上。九章覆盖：
似然与极大似然估计 \(\to\) Cramér-Rao 下界与正则条件 \(\to\) 我们问题
专属的 Fisher 信息矩阵形式 \(\to\) Wilks 定理与 Profile Likelihood
\(\to\) Laplace 方法与 Bernstein–von Mises \(\to\) MCMC 详细平衡、最优
acceptance、Geweke 与 ESS \(\to\) Glückstern 闭式动量分辨极限 \(\to\)
Bethe-Bloch 与 Landau/Vavilov （解释当前 $p_z$ 系统偏差）\(\to\) Highland
多次散射公式与精度修正。每章末给出与状态报告中相应结论的对照。

% ============================================================
% File: rk_deep_dive_part1a_math_ch1to3_zh.tex
% Purpose: Part I (数学基础) chapters 1-3 — likelihood foundations,
%          Cramér-Rao 下界与正则条件, Fisher information matrix in
%          the SAMURAI 2-PDC + 3-parameter geometry.
% Parent : rk_error_analysis_deep_dive_20260426_zh.tex (wrapper)
% Sibling: rk_deep_dive_part1b_math_ch4to6_zh.tex (Wilks/Laplace/MCMC)
%          rk_deep_dive_part1c_math_ch7to9_zh.tex (Glückstern/Bethe-Bloch/Highland)
% Date   : 2026-04-26
% Author : TBT
%
% Brief TOC:
%   §1 似然函数与极大似然估计基础
%   §2 Cramér-Rao 下界与正则条件
%   §3 Fisher 信息矩阵在 SAMURAI 2-PDC 几何下的具体形式
%
% Notation conventions (inherited from status report):
%   \vec\alpha = (\delta x, \delta y, u, v, p)         5-DoF kThreePointFree
%   \vec\alpha = (u, v, p)                              3-DoF kFixedTargetPdcOnly
%   u = p_x/p_z, v = p_y/p_z, p = |\vec p|
%   \vec w = whitened residual (4-vector for 2-PDC)
%   \vec r = raw residual; \hat u, \hat v wire directions (Eq. eq:wire_dir)
%   Default frame: target frame (beam-as-Z), all truth vs reco comparisons
%   in target frame.
%
% Wrapper macros consumed: \StatusReport (provided in wrapper). No new macros
% are introduced in this fragment.
% ============================================================

\providecommand{\StatusReport}{文献~[1]}

% ============================================================
\section{似然函数与极大似然估计基础}
\label{sec:partIa:likelihood:basics}
% ============================================================

本章建立 PDC 重建问题的统计学语言。后两章 (CRB、Fisher 矩阵) 与 Part II 的系统误差预算 (\StatusReport § 误差分析) 都以本章定义为出发点。所有讨论默认采用靶系 (beam-as-Z)，参数向量沿用状态报告 \StatusReport § 问题设定与符号约定 中的约定。

\subsection{统计模型：白化残差的高斯似然}
\label{sec:partIa:lik:model}

设单条重建轨迹的测量量为 4 维实向量
\begin{equation}
  \vec m = \bigl(u_1^\mathrm{meas},\; v_1^\mathrm{meas},\; u_2^\mathrm{meas},\; v_2^\mathrm{meas}\bigr)^\top \in \mathbb R^4,
  \label{eq:partIa:lik:meas}
\end{equation}
真实分布建模为
\begin{equation}
  \vec m = \vec f(\vec\alpha_\star) + \boldsymbol\epsilon,
  \qquad \boldsymbol\epsilon \sim \mathcal N(\vec 0,\;\Sigma_\mathrm{meas}),
  \label{eq:partIa:lik:gauss}
\end{equation}
其中前向模型 $\vec f:\mathbb R^{n_p}\to\mathbb R^4$ 由 RK4 轨迹积分加丝坐标投影定义 (\StatusReport § 残差构建)，$\Sigma_\mathrm{meas}$ 是测量协方差。在 PDC 几何下两个丝方向 $\hat u, \hat v$ 不正交 (式 \StatusReport eq:wire\_dir) 且测量误差 $\sigma_\mathrm{wire}=0.5\;\mathrm{mm}$ 在丝坐标系内独立、各向同性：单 PDC 内残差协方差由 $\hat u\cdot\hat v=\sin^2\theta$ 引入耦合，需要 Cholesky 白化 (\StatusReport eq:whitening) 才能得到独立单位方差残差。

定义白化残差
\begin{equation}
  \vec w(\vec\alpha) \equiv L^{-1}\bigl[\vec f(\vec\alpha)-\vec m\bigr]
  \;\sim\; \mathcal N(\vec 0, I_4)\quad\text{at }\vec\alpha=\vec\alpha_\star,
  \label{eq:partIa:lik:whiten}
\end{equation}
其中 $L$ 是 $\Sigma_\mathrm{meas}=LL^\top$ 的 Cholesky 因子。等价地，$\vec w_i(\vec\alpha_\star)$ 是 i.i.d.\ 标准正态。这一性质使得我们可以把 $\chi^2$ 写为白化残差平方和 (式 \StatusReport eq:chi2\_recall)，并在似然层面统一处理两个 PDC 平面。

\subsection{似然函数与对数似然}
\label{sec:partIa:lik:loglik}

由式 \eqref{eq:partIa:lik:whiten} 立即得到单条事件的似然函数 (取 $\Sigma_\mathrm{meas}$ 不依赖参数)：
\begin{equation}
  L(\vec\alpha\mid\vec m) \;=\; \prod_{i=1}^{4}\phi\bigl(w_i(\vec\alpha)\bigr)
  \;=\; (2\pi)^{-2}\,\exp\!\Biggl[-\tfrac12\sum_{i=1}^{4} w_i^2(\vec\alpha)\Biggr],
  \label{eq:partIa:lik:def}
\end{equation}
其中 $\phi(z)=(2\pi)^{-1/2}e^{-z^2/2}$ 是标准正态密度。对数似然为
\begin{equation}
  \ln L(\vec\alpha\mid\vec m)
  \;=\; -\tfrac12\sum_{i=1}^{4} w_i^2(\vec\alpha) \;+\; C
  \;=\; -\tfrac12\,\chi^2(\vec\alpha) + C,
  \label{eq:partIa:lik:loglik}
\end{equation}
$C=-2\ln(2\pi)$ 不依赖 $\vec\alpha$。这就把状态报告的 LM 最小二乘 (\StatusReport § 最小二乘拟合) 与统计学的极大似然估计 (MLE) 严格联系起来：
\begin{equation}
  \hat{\vec\alpha}_\mathrm{MLE} \;=\; \arg\max_{\vec\alpha}\,\ln L(\vec\alpha\mid\vec m)
  \;=\; \arg\min_{\vec\alpha}\,\chi^2(\vec\alpha).
  \label{eq:partIa:lik:mle}
\end{equation}

LM 最优化算法返回的中心值即为白化残差高斯模型下的 MLE。状态报告 § Levenberg-Marquardt 优化 中的迭代更新本质上是 Gauss-Newton 法对 $-\nabla^2\ln L$ 加阻尼正则化 (式 \StatusReport eq:lm\_normal) 的具象化。

\subsection{充分统计量}
\label{sec:partIa:lik:sufficiency}

对于式 \eqref{eq:partIa:lik:gauss} 高斯模型，给定 $\vec\alpha$ 与 $\Sigma_\mathrm{meas}$ 不依赖 $\vec\alpha$ 的假设，由因子分解定理可知白化残差向量 $\vec w(\vec\alpha)$ 本身是 $\vec\alpha$ 的充分统计量：
\begin{proposition}[白化残差充分性]\label{prop:partIa:suff}
  在式 \eqref{eq:partIa:lik:gauss} 模型下，给定 $\Sigma_\mathrm{meas}$ 已知且与 $\vec\alpha$ 无关，统计量 $T(\vec m)=\vec w(\vec\alpha=\vec 0)=L^{-1}\vec m$ 对参数族 $\{p(\vec m\mid\vec\alpha):\vec\alpha\in\mathbb R^{n_p}\}$ 是充分统计量。
\end{proposition}
\begin{proof}
  $L(\vec\alpha\mid\vec m)\propto\exp[-\tfrac12\|L^{-1}\vec m-L^{-1}\vec f(\vec\alpha)\|^2]$，把 $L^{-1}\vec m$ 视为 $T$，似然只通过 $T$ 依赖 $\vec m$。Neyman 因子分解定理给出充分性。
\end{proof}

物理上这意味着：所有的统计推断 (Fisher、Profile、MCMC、Laplace) 都可以仅依赖 $\vec w$ 而不需要原始 $\vec m$。状态报告中 \texttt{PDCRkAnalysisInternal::computeWhitenedResidual} 输出的 4 维向量正是后续所有方法的统一输入接口，与本命题的数学结构对应。

\subsection{正则统计族 (regular family)}
\label{sec:partIa:lik:regular}

CRB 与 Fisher 矩阵的所有定理都假设似然属于"正则族" (regular family)。本节按 Lehmann-Casella 第 6 章的定义列出正则条件，并逐项检验本问题。

\begin{definition}[正则族]\label{def:partIa:regular}
  参数族 $\{p(\vec m\mid\vec\alpha):\vec\alpha\in\Theta\}$ 称为正则族，若：
  \begin{enumerate}[label=(R\arabic*),nosep]
    \item $\Theta$ 是 $\mathbb R^{n_p}$ 的开集 (参数在内点)；
    \item 支撑 $\{\vec m: p(\vec m\mid\vec\alpha)>0\}$ 不依赖 $\vec\alpha$；
    \item 对几乎所有 $\vec m$，$\ln p(\vec m\mid\vec\alpha)$ 关于 $\vec\alpha$ 二次连续可微；
    \item 期望与微分可交换：$\nabla_{\vec\alpha}\int p\,d\vec m=\int\nabla_{\vec\alpha}p\,d\vec m$，二阶同；
    \item Score 函数 $S(\vec\alpha)\equiv\nabla_{\vec\alpha}\ln L$ 满足 $\mathbb E[S]=0$、$\mathrm{Cov}(S)=\mathcal I(\vec\alpha)$ 有限且正定。
  \end{enumerate}
\end{definition}

\paragraph{逐项检验。} 在我们的高斯白化模型下：
\begin{itemize}[nosep]
  \item (R1) 参数空间 $\Theta=\mathbb R^{n_p}$ 显式开集；$u,v$ 取实数，$p>0$ 但 reco 中位数 $\sim 635\;\mathrm{MeV}/c$ 远离 $p=0$ 边界，故内点条件满足。
  \item (R2) 支撑为 $\mathbb R^4$，与 $\vec\alpha$ 无关。
  \item (R3) $\ln L=-\tfrac12\|\vec w(\vec\alpha)\|^2+C$，依赖 $\vec\alpha$ 经 $\vec w=L^{-1}[\vec f(\vec\alpha)-\vec m]$。前向模型 $\vec f$ 由 RK4 积分给出，沿 PDC1/PDC2 投影解析；磁场图通过三线性插值 $C^0$ 连续，但插值导数仅 $C^{-1}$ (分段常值)。这是\textbf{潜在违反点}：在磁场图块界面 $\vec f$ 的二阶导不存在；具体磁场图三线性插值的精度影响见 Part~II §\ref{sec:budget-bfield} 的系统误差预算条目。在 RK 步长尺度上做平滑后此违反可忽略。
  \item (R4) 高斯密度衰减快，期望与微分可交换标准成立。
  \item (R5) Score $S=\partial_{\vec\alpha}\ln L=-J^\top\vec w$，其中 $J=\partial\vec w/\partial\vec\alpha$；显然 $\mathbb E_{\vec\alpha_\star}[\vec w]=\vec 0\Rightarrow\mathbb E[S]=\vec 0$。$\mathcal I=\mathbb E[SS^\top]=\mathbb E[J^\top\vec w\vec w^\top J]=J^\top J$ (利用 $\mathbb E[\vec w\vec w^\top]=I_4$)。要求 $J$ 列满秩；这是\emph{可识别性条件}，下节讨论。
\end{itemize}

\subsection{可识别性 (identifiability)}
\label{sec:partIa:lik:identifiability}

正则族成立的最关键前提是 Fisher 矩阵 $\mathcal I=J^\top J$ 正定，等价于 $J\in\mathbb R^{4\times 3}$ 列满秩。直观地说：参数 $\vec\alpha$ 到测量 $\vec m$ 的映射在最优点附近\emph{局部一一对应}。

\begin{lemma}[局部可识别性]\label{lem:partIa:ident}
  若 $\vec f$ 在 $\vec\alpha_\star$ 处 $C^1$ 且 $J(\vec\alpha_\star)\in\mathbb R^{4\times 3}$ 列满秩，则在 $\vec\alpha_\star$ 的某邻域内 $\vec\alpha\mapsto\vec f(\vec\alpha)$ 单射；从而 MLE 在该邻域内唯一。
\end{lemma}
\begin{proof}
  逆函数定理推论：$J$ 列满秩等价于 $J^\top J$ 正定。$\chi^2(\vec\alpha)=\|\vec w(\vec\alpha)\|^2$ 在 $\vec\alpha_\star$ 处的 Hessian 为 $2J^\top J+O(\vec w)$，正定 $\Rightarrow$ $\chi^2$ 局部严格凸 $\Rightarrow$ $\arg\min$ 唯一。
\end{proof}

对 kFixedTargetPdcOnly 模式，4 个测量约束 3 个自由参数，过约束 1 个自由度 ($N_\mathrm{dof}=1$)。可识别性等价于 $J$ 不存在退化方向。状态报告 \StatusReport § 为什么 $\sigma_{p_y}$ 极小？ 给出的物理图像 ($\hat u,\hat v$ 的 $y$ 分量为 $\pm 1/\sqrt 2$) 是 $J$ 第二列大、第一/三列小的代数表现。当 $u\to 0$、$v\to 0$ 同时 $p\to\infty$ (直线极限) 时第三列 (动量灵敏度) $\to 0$；这是潜在的弱可识别性源，对低弯曲事件需关注 SVD 条件数 (§\ref{sec:partIa:fisher:condnum})。

\subsection{MLE 的渐近性质 (前瞻)}
\label{sec:partIa:lik:asymptotic}

经典定理：在正则族下，当样本量 $N\to\infty$ 时
\begin{equation}
  \sqrt N\,(\hat{\vec\alpha}_\mathrm{MLE}-\vec\alpha_\star)
  \;\xrightarrow{d}\; \mathcal N\bigl(\vec 0,\;\mathcal I_1^{-1}(\vec\alpha_\star)\bigr),
  \label{eq:partIa:lik:asymptotic}
\end{equation}
其中 $\mathcal I_1$ 是单样本 Fisher 信息。在本问题中"样本量"是\emph{每条事件的测量数}，单条事件 $N=4$、$N_\mathrm{dof}=1$ —— 距渐近域 ($N\gg n_p$) 极远。这正是状态报告 § 适用条件与局限 中"Fisher 估计在 $N_\mathrm{dof}$ 很小时可能过于乐观"的根本原因。本部分第 4 章 (Wilks 定理) 将给出 $N$ 较小时的修正路径；本章后两节为该讨论奠定 CRB 与 Fisher 几何的基础。

% ============================================================
\section{Cramér-Rao 下界与正则条件}
\label{sec:partIa:crb}
% ============================================================

\subsection{CRB 的陈述与几何含义}
\label{sec:partIa:crb:statement}

\begin{theorem}[向量参数 Cramér-Rao 下界]\label{thm:partIa:crb}
  设似然 $\{p(\vec m\mid\vec\alpha):\vec\alpha\in\Theta\}$ 满足正则条件 (R1-R5)；$\hat{\vec\alpha}(\vec m)$ 为 $\vec\alpha$ 的无偏估计量 (即 $\mathbb E_{\vec\alpha}[\hat{\vec\alpha}]=\vec\alpha$ 对所有 $\vec\alpha\in\Theta$ 成立)。则
  \begin{equation}
    \mathrm{Cov}_{\vec\alpha}\bigl(\hat{\vec\alpha}\bigr) \;\succeq\; \mathcal I^{-1}(\vec\alpha),
    \label{eq:partIa:crb:bound}
  \end{equation}
  其中 $A\succeq B$ 表示 $A-B$ 半正定。等号成立当且仅当 $\hat{\vec\alpha}$ 是有效估计量。
\end{theorem}

几何含义：在参数空间中，无偏估计量的协方差椭球必须包含 Fisher 椭球 $\{\vec\alpha:\vec\alpha^\top\mathcal I\vec\alpha\le 1\}$。Fisher 椭球是任何无偏估计量\emph{不可逾越}的精度极限。

\subsection{CRB 的证明 (向量版本)}
\label{sec:partIa:crb:proof}

\begin{proof}
  分四步。固定 $\vec\alpha$；记 score 函数 $S(\vec m;\vec\alpha)=\nabla_{\vec\alpha}\ln L$。

  \textbf{第 1 步：score 期望为零。}
  由 (R4) 微分与积分可交换：
  \begin{equation}
    \mathbb E[S]
    \;=\; \int (\nabla_{\vec\alpha}\ln p)\,p\,d\vec m
    \;=\; \nabla_{\vec\alpha}\int p\,d\vec m
    \;=\; \nabla_{\vec\alpha}1 \;=\;\vec 0.
    \label{eq:partIa:crb:scoremean}
  \end{equation}

  \textbf{第 2 步：score 与无偏估计量的协方差等于单位阵。}
  对无偏性条件 $\int\hat{\vec\alpha}(\vec m)\,p(\vec m;\vec\alpha)\,d\vec m=\vec\alpha$ 求 $\nabla_{\vec\alpha}$：
  \begin{equation}
    I_{n_p}
    \;=\; \nabla_{\vec\alpha}\!\!\int\!\hat{\vec\alpha}\,p\,d\vec m
    \;=\; \int\!\hat{\vec\alpha}\,(\nabla_{\vec\alpha}p)\,d\vec m
    \;=\; \int\!\hat{\vec\alpha}\,(\nabla_{\vec\alpha}\ln p)\,p\,d\vec m
    \;=\; \mathbb E[\hat{\vec\alpha}\,S^\top].
    \label{eq:partIa:crb:cross}
  \end{equation}
  结合 $\mathbb E[S]=\vec 0$、$\mathbb E[\hat{\vec\alpha}]=\vec\alpha$，
  \begin{equation}
    \mathrm{Cov}(\hat{\vec\alpha},S) \;=\; \mathbb E[\hat{\vec\alpha}\,S^\top]-\mathbb E[\hat{\vec\alpha}]\mathbb E[S]^\top \;=\;I_{n_p}.
    \label{eq:partIa:crb:cov_eqI}
  \end{equation}

  \textbf{第 3 步：联合协方差矩阵半正定。}
  把 $\hat{\vec\alpha}-\vec\alpha$ 与 $S$ 拼成 $2n_p$ 维向量 $Z=\bigl(\hat{\vec\alpha}-\vec\alpha,\;S\bigr)$。其协方差矩阵
  \begin{equation}
    \Sigma_Z \;=\; \begin{pmatrix}\mathrm{Cov}(\hat{\vec\alpha}) & I_{n_p}\\ I_{n_p} & \mathcal I(\vec\alpha)\end{pmatrix}\;\succeq\;0
  \end{equation}
  半正定 (任何随机向量的协方差矩阵半正定)。利用分块矩阵 Schur 补：$\Sigma_Z\succeq 0$ 且 $\mathcal I\succ 0$ 给出
  \begin{equation}
    \mathrm{Cov}(\hat{\vec\alpha}) - I_{n_p}\,\mathcal I^{-1}\,I_{n_p}^\top \;\succeq\;0,
    \label{eq:partIa:crb:schur}
  \end{equation}
  即 $\mathrm{Cov}(\hat{\vec\alpha})\succeq\mathcal I^{-1}$。

  \textbf{第 4 步：等号条件。}
  Cauchy-Schwarz 等式条件要求 $\hat{\vec\alpha}-\vec\alpha=A\,S$ 几乎处处成立，对线性变换 $A$；代回 $\mathrm{Cov}(\hat{\vec\alpha},S)=I$ 得 $A=\mathcal I^{-1}$，即 $\hat{\vec\alpha}=\vec\alpha+\mathcal I^{-1}S$。这是\textbf{有效估计量}的形式，仅在指数族中成立。
\end{proof}

证明对应向量版的 Cauchy-Schwarz / Schur 补论证。该结构在 §\ref{sec:partIa:fisher} 中将以 Fisher 矩阵 $\mathcal I=J^\top J$ 的具体形式重新出现。

\subsection{正则条件的精细化讨论}
\label{sec:partIa:crb:reg}

第 \ref{sec:partIa:lik:regular} 节的检验已经覆盖大部分情形。本节强调三个对本问题重要的细节。

\paragraph{微分与积分可交换 (R4) 的依据。} Lebesgue 控制收敛定理要求 $|\partial_{\alpha_i}p(\vec m;\vec\alpha)|$ 被某个 $\vec\alpha$-邻域内可积函数控制。对高斯密度
\begin{equation}
  p(\vec m;\vec\alpha) \;=\;(2\pi)^{-2}|\det L|^{-1}\exp\!\bigl[-\tfrac12\|\vec w(\vec\alpha)\|^2\bigr],
\end{equation}
$\partial_{\alpha_i}p=p\cdot[-J_i^\top\vec w]$，被 $|\vec w|\,p$ 控制；后者一阶矩有限。同理二阶微分被 $\|\vec w\|^2 p$ 控制。R4 自动成立。

\paragraph{Score 二阶矩有限 (R5)。} $\mathbb E[\|S\|^2]=\mathrm{tr}(\mathcal I)=\mathrm{tr}(J^\top J)=\|J\|_F^2<\infty$ 当且仅当 $J$ 元素有限。若磁场图存在奇异点 (本问题不存在) 或 $\vec f$ 在某 $\vec\alpha$ 不可微，R5 失效；本问题仅在 RK 步与磁场表块界面交叉时有 $C^1$ 但非 $C^2$ 的弱奇性，对 $\mathcal I$ 数值上影响 $O(10^{-4})$，可忽略。

\paragraph{支撑不依赖参数 (R2)。} 对截断分布 (例如丝信号有上下截止) 此条件可能失效；本问题中我们假设 PDC 命中位置可在整个 $\mathbb R^4$ 上 (实际命中在 PDC 灵敏区内但建模时不引入硬截断)。R2 满足。

\subsection{CRB 在本问题中失效或宽松的四种机制}
\label{sec:partIa:crb:fail}

CRB 是\emph{无偏}估计量的下界。本问题的 RK 重建在多个层面偏离这一前提；逐一陈述。

\paragraph{(F1) 渐近域外。} CRB 的可达性 (即 MLE 达到 CRB) 是一个渐近结果 ($N\to\infty$)。在 $N_\mathrm{dof}=1$ 的极小样本下，MLE 协方差与 $\mathcal I^{-1}$ 之间存在 $O(N^{-1})$ 阶修正项，可能使 Fisher 估计偏紧 (低估方差)。Bartlett 修正 (Part I 第 4 章 Wilks 部分讨论) 是修正此偏差的标准工具。

\paragraph{(F2) 边界参数。} CRB 假设 $\vec\alpha$ 在 $\Theta$ 内点。本问题的 $u,v$ 自由实数；$p>0$ 但 reco 中位数 $\sim 635\;\mathrm{MeV}/c$ 远离边界。\textbf{此机制在本问题不触发}。

\paragraph{(F3) 非高斯似然。} 我们的 $\chi^2$ 二次结构假设白化残差严格高斯。前向模型 $\vec f(\vec\alpha)$ 关于 $\vec\alpha$ 高度非线性 ($u,v,p\to p_x,p_y,p_z\to$ RK 轨迹$\to$ 丝坐标投影)。在最优点附近做泰勒展开
\begin{equation}
  w_i(\vec\alpha)\approx w_i(\hat{\vec\alpha})+J_{ij}\Delta\alpha_j+\tfrac12 H_{ijk}\Delta\alpha_j\Delta\alpha_k+\dots
\end{equation}
线性截断后似然才是严格高斯。二阶项 $H_{ijk}$ 引入 \emph{Edgeworth} 修正，对 $\sigma$ 的影响为 $O(H_{ijk}/J_{ij})$。状态报告 § 为什么 $\sigma_{p_y}$ 极小？ 已经定性指出 $(u,v,p)\to p_y$ 的二阶项是 $p_y$ 方向 Fisher 低估的主因；§\ref{sec:partIa:fisher:momentum} 给出 $M$ 的解析二阶项。

\paragraph{(F4) 偏差。} CRB 不约束有偏估计量。状态报告 \StatusReport § 修复前 lab/靶系未对齐 记录了一个标志性事件：修复前 $p_x$ 上存在系统偏置 $\Delta p_x\approx -33\;\mathrm{MeV}/c$，远超 Fisher $\sigma_{p_x}\sim 7\;\mathrm{MeV}/c$；CRB 在偏置存在时无法保证覆盖率。修复后 $p_x$ 偏置降至 $\sim -0.33\;\mathrm{MeV}/c$ ($-0.05\sigma$ 量级)，覆盖率从 $\sim 0.05$ 恢复到 $\sim 0.80$。Part II 第 7 章 (PDC 对齐) 与 §\ref{sec:partIa:fisher:momentum} 共同处理剩余的偏置源。

\subsection{应用：Fisher 在 744-事件 ensemble 中的近似饱和性}
\label{sec:partIa:crb:saturation}

\begin{proposition}[本问题的 CRB 饱和度]\label{prop:partIa:saturation}
  在 \StatusReport 引用的 2026-04-17 ensemble study (744 个 truth-grid 事件，靶系) 中：
  \begin{enumerate}[nosep]
    \item Fisher 在 $|\vec p|$、$p_x$、$p_z$ 三个分量上 \emph{近似饱和} CRB —— 68\% 名义覆盖率的实测值为 $0.761,\,0.801,\,0.772$，对应 Clopper-Pearson 95\% CI 分别为 $[0.728,0.791]$、$[0.771,0.829]$、$[0.740,0.801]$；名义 0.68 落在 CI 边缘外侧，方向是 \emph{保守} (覆盖率偏高，$\sigma$ 估计偏大约 $\sim 25\%$)。
    \item Fisher 在 $p_y$ 分量上 \emph{显著违反} CRB 的实际可达性 —— 实测 68\% 覆盖率 $0.421$，CP 95\% CI $[0.385, 0.457]$，名义 $\sigma_{p_y}$ 被低估约 $1.9\times$。
  \end{enumerate}
  违反机制为 (F3) 非高斯似然 + (F4) 参数化诱导偏置，\emph{非}正则性失效；具体论证见 §\ref{sec:partIa:fisher:momentum} 关于 $M$ 矩阵二阶项的展开。
\end{proposition}

数值来源：\StatusReport § 误差分析覆盖率小节，由独立的 truth-grid 744-事件 run 给出。本命题的关键陈述是：CRB 的失败发生在\emph{参数化}层面 ($u,v,p\mapsto p_x,p_y,p_z$ 的非线性)，而非\emph{似然}层面 (R1-R5 全部成立)。修复路径是使用\emph{无参数化偏置}的 reparam (例如直接以 $p_x,p_y,p_z$ 拟合或在 $p_y$ 方向加二阶项校正)，由 Part II § 8 详述。

% ============================================================
\section{Fisher 信息矩阵在 SAMURAI 2-PDC 几何下的具体形式}
\label{sec:partIa:fisher}
% ============================================================

本章把抽象的 Fisher 矩阵代数翻译成 SAMURAI 谱仪几何下的可观测量，定量解释状态报告中 $\sigma_{p_y}\ll\sigma_{p_x}$ 的反直觉结果，并给出 $(u,v,p)\to(p_x,p_y,p_z)$ 协方差变换的解析形式。

\subsection{Fisher 矩阵的两种等价定义}
\label{sec:partIa:fisher:def}

\begin{definition}[Fisher 信息矩阵]\label{def:partIa:fisher}
  \begin{align}
    \mathcal I_{ij}(\vec\alpha) &\;=\; \mathrm{Cov}\bigl(S_i,S_j\bigr) \;=\;\mathbb E[S_iS_j],
      \label{eq:partIa:fisher:def_cov}\\
    \mathcal I_{ij}(\vec\alpha) &\;=\; -\mathbb E\!\left[\partial_i\partial_j\ln L\right].
      \label{eq:partIa:fisher:def_hess}
  \end{align}
\end{definition}

两个表达式在正则族下相等。证明利用 $\partial_j\partial_i\ln p=\partial_jp^{-1}\cdot\partial_ip+p^{-1}\partial_j\partial_ip$，第一项给 $-\mathcal I_{ij}$，第二项 (取期望后利用 (R4)) 为零。

对式 \eqref{eq:partIa:lik:loglik} 高斯白化模型：
\begin{equation}
  \ln L = -\tfrac12\vec w^\top\vec w + C
  \;\Rightarrow\;
  S = -J^\top\vec w,
  \quad
  \partial_i\partial_j\ln L = -J_{ki}J_{kj}-w_k\partial_i\partial_jw_k.
\end{equation}
取 $\mathbb E_{\vec\alpha_\star}$，第二项消去 (因 $\mathbb E[\vec w]=\vec 0$)，得
\begin{equation}
  \boxed{\;\mathcal I(\vec\alpha) \;=\; J^\top(\vec\alpha)\,J(\vec\alpha),\quad J\equiv\partial\vec w/\partial\vec\alpha\in\mathbb R^{4\times n_p}.\;}
  \label{eq:partIa:fisher:JTJ}
\end{equation}
此即状态报告 \StatusReport eq:fisher\_cov 的根。下节展开 $J$ 的列结构。

\subsection{kFixedTargetPdcOnly 模式下 $J$ 的符号导出}
\label{sec:partIa:fisher:Jderiv}

参数 $\vec\alpha=(u,v,p)$。前向模型分两步：
\begin{enumerate}[nosep]
  \item \textbf{初值映射}：$(u,v,p)\mapsto\vec p_0=p\,\hat n(u,v)$，其中 $\hat n=(u,v,1)/\sqrt{D}$，$D=1+u^2+v^2$。
  \item \textbf{RK 传播 + PDC 投影}：$\vec p_0\mapsto\vec r^\mathrm{traj}_k(\vec\alpha)$ ($k=1,2$)，再投影到丝方向给出 $f_{2k-1}=\vec r_k\cdot\hat u$，$f_{2k}=\vec r_k\cdot\hat v$。
\end{enumerate}

链式法则：
\begin{equation}
  J_{ij} \;=\; \partial w_i/\partial\alpha_j
  \;=\; (L^{-1})_{ii'}\,\sum_{k}\,(\hat u_i \text{ or }\hat v_i)_l\,
        \frac{\partial r_{k,l}^\mathrm{traj}}{\partial p_{0,m}}\,\frac{\partial p_{0,m}}{\partial\alpha_j}.
  \label{eq:partIa:fisher:chainrule}
\end{equation}
其中 $\partial p_{0,m}/\partial\alpha_j$ 解析；$\partial r_{k,l}^\mathrm{traj}/\partial p_{0,m}$ 由 RK4 沿轨迹的传递矩阵 (Jacobi 方程) 决定，无解析闭形，状态报告通过中心差分数值计算 (\StatusReport eq:jacobian\_explicit)。

\paragraph{初值导数解析形式。}
\begin{equation}
  \frac{\partial\vec p_0}{\partial u}
  \;=\;\frac{p}{\sqrt{D}}\!\left[\hat e_x-\frac{u}{D}(u,v,1)^\top\right],\quad
  \frac{\partial\vec p_0}{\partial v}
  \;=\;\frac{p}{\sqrt{D}}\!\left[\hat e_y-\frac{v}{D}(u,v,1)^\top\right],
  \label{eq:partIa:fisher:p0_deriv_uv}
\end{equation}
\begin{equation}
  \frac{\partial\vec p_0}{\partial p}\;=\;\frac{1}{\sqrt{D}}(u,v,1)^\top \;=\;\hat n.
  \label{eq:partIa:fisher:p0_deriv_p}
\end{equation}

物理解读：$\partial\vec p_0/\partial u$ 主分量在 $\hat e_x$ 方向，强度 $\sim p/\sqrt{D}\approx p$ (小角度极限 $D\to 1$)；$\partial\vec p_0/\partial p$ 沿出射方向 $\hat n$。后者是 \emph{动量灵敏度} 的核心：改变 $|\vec p|$ 不改变出射方向，但改变弯曲半径 $\rho=p/(qB_y)$，从而改变 PDC 处的命中点 —— $J$ 第三列的 \emph{所有} 信号通过磁场弯曲产生。

\paragraph{第 $j$ 列的物理含义 (定性表)。}
\begin{table}[H]
\centering
\caption{$J\in\mathbb R^{4\times 3}$ 列的物理含义。"幅度"指典型 $|J_{ij}|$ 量级，归一化到 $\sigma_\mathrm{wire}=0.5\;\mathrm{mm}$ 单位 (来自状态报告 \StatusReport § Fisher 输出 数值)。}
\label{tab:partIa:fisher:Jcols}
\begin{tabular}{c l l c}
\toprule
列 & 物理含义 & 主导路径 & 典型 $|J_{ij}|$ \\
\midrule
1 ($u$) & $p_x/p_z$ 出射方向 & 直线 + 弯曲耦合 & $\sim 10^4\,\hat u_x$ (mm/单位) \\
2 ($v$) & $p_y/p_z$ 出射方向 & 直线 (磁场不弯) & $\sim 10^4\,\hat u_y$ (mm/单位) \\
3 ($p$) & $|\vec p|$ 弯曲半径 & 仅经磁场 & $\sim O(1)$ (mm/(MeV/$c$)) \\
\bottomrule
\end{tabular}
\end{table}

注意第 2 列幅度比第 1 列大\emph{约 $1/\cos\theta\approx 1.84$ 倍}，因为 $\hat u,\hat v$ 在 $y$ 方向投影 $1/\sqrt 2$ 而 $x$ 方向仅 $\cos\theta/\sqrt 2$ ($\theta=57^\circ$)。这是下节"$\sigma_{p_y}$ 极小"的代数源。

\subsection{为什么 $\sigma_{p_y}$ 极小}
\label{sec:partIa:fisher:sigmapy}

由式 \eqref{eq:partIa:fisher:p0_deriv_uv}-\eqref{eq:partIa:fisher:p0_deriv_p} 与丝方向 (\StatusReport eq:wire\_dir)：
\begin{equation}
  \hat u\cdot\hat e_y = \tfrac{1}{\sqrt 2},\quad \hat v\cdot\hat e_y = -\tfrac{1}{\sqrt 2},
  \quad
  \hat u\cdot\hat e_x=\hat v\cdot\hat e_x=\tfrac{\cos\theta}{\sqrt 2}.
\end{equation}
第二列的丝投影 $|J_{i2}|\propto p\,(\hat e_y\cdot\hat u\text{ 或 }\hat v)/\sqrt D = p/(\sqrt{2D})$；第一列为 $p\cos\theta/(\sqrt{2D})$。比值
\begin{equation}
  \frac{|J_{i2}|}{|J_{i1}|} \;=\;\frac{1}{\cos\theta}\approx 1.836\;\text{ at }\theta=57^\circ.
\end{equation}
Fisher 对角元为列范数平方，故 $\mathcal I_{vv}/\mathcal I_{uu}=1/\cos^2\theta\approx 3.37$。在仅考虑对角的近似下 $\sigma_v^2\propto 1/\mathcal I_{vv}$，结合非线性传播 $\sigma_{p_y}\approx p_z\sigma_v$ (式 \eqref{eq:partIa:fisher:M}) 给出
\begin{equation}
  \frac{\sigma_{p_y}}{\sigma_{p_x}}\;\sim\;\cos\theta\approx 0.545,
\end{equation}
即 $\sigma_{p_y}$ 在 \emph{Fisher 估计内} 比 $\sigma_{p_x}$ 小约 1.8 倍。这与状态报告 \StatusReport § 为什么 $\sigma_{p_y}$ 极小？ 给出的事件级数字 $\sigma_{p_x}\approx 2.26\;\mathrm{MeV}/c$、$\sigma_{p_y}\approx 0.35\;\mathrm{MeV}/c$ 量级一致 (后者更小是因为还耦合了第三列 $p$ 的关联以及 truth-frame 旋转效应)。

\paragraph{警告：Fisher 估计正确，但在 $p_y$ 上 Fisher 与真实方差脱节。} 上述代数说的是\emph{Fisher 自身} 的对角项很小 —— 这是精确事实。但 Proposition \ref{prop:partIa:saturation} 表明实测 $\sigma_{p_y}$ 比 Fisher 预言大 $\sim 1.9\times$。差距 \emph{不在} Fisher 矩阵本身，而在 $(u,v,p)\to(p_x,p_y,p_z)$ 的协方差传播 (式 \eqref{eq:partIa:fisher:Sp})；下节展开。

\subsection{条件数与病态性诊断}
\label{sec:partIa:fisher:condnum}

$\mathcal I=J^\top J$ 的奇异值分解 $\mathcal I=V\Lambda V^\top$ 给出主轴方向 (列向量 $V_{\cdot k}$) 与对应方差 $\Lambda_{kk}^{-1}$。条件数
\begin{equation}
  \kappa(\mathcal I) \;=\;\frac{\sigma_\mathrm{max}(J)^2}{\sigma_\mathrm{min}(J)^2}.
  \label{eq:partIa:fisher:cond}
\end{equation}
状态报告 \StatusReport § 完整算法流程 (Algorithm 1) 把 $\kappa$ 作为 Fisher 计算成功/失败的诊断量。物理意义如下。

\begin{itemize}[nosep]
  \item $\kappa\sim 1$：参数空间各方向被均衡约束。
  \item $\kappa\gg 1$：存在退化方向 $V_{\cdot\mathrm{min}}$ —— 沿此方向移动 $\vec\alpha$，$\chi^2$ 几乎不变。物理上对应"参数组合无法被该几何独立测量"。
  \item $\kappa>\kappa_\mathrm{max}$ (代码默认阈值 $\sim 10^{12}$)：返回失败标志，弃事件。
\end{itemize}

在 SAMURAI 2-PDC 几何下，已知病态方向是 \emph{动量-曲率简并}：当 $u,v$ 极小且粒子能量很高时，弯曲半径远大于 PDC 弧长，$J$ 第三列退化到 $J$ 第一/二列张成的子空间。$\kappa$ 在此区域 $\sim 10^4$ 但仍可逆。低能 ($p\to 0$) 区域则反过来：$J$ 第三列变大，$\kappa$ 改善，但 RK 步长稳定性下降。本问题中 reco $|\vec p|\sim 635\;\mathrm{MeV}/c$ 处 $\kappa$ 实测 (状态报告事件 \#399 同口径) 在 $10^2-10^3$ 之间，远低于警戒阈值。

\paragraph{代码实现注。} 状态报告 § 完整算法流程 中标注 "SVD" 的步骤在实际代码 (\texttt{libs/analysis\_pdc\_reco/src/PDCErrorAnalysis.cc}, 行 $\sim 355$) 调用 \texttt{TMatrixDSymEigen}：对 $\mathcal I$ 做对称特征值分解 $\mathcal I=V\Lambda V^\top$，等价于 $J^\top J$ 的右奇异向量与奇异值平方。条件数定义为最大/最小特征值的比 (因为 $\mathcal I$ 已是对称半正定，特征值即为奇异值的平方)。

\subsection{从 $(u,v,p)$ 到 $(p_x,p_y,p_z)$ 的协方差变换}
\label{sec:partIa:fisher:momentum}

物理观测量是 $(p_x,p_y,p_z)$；Fisher 给出的是 $(u,v,p)$ 协方差 $\Sigma_{\vec\alpha}=\mathcal I^{-1}$。需要 Jacobian
\begin{equation}
  M \;\equiv\;\frac{\partial(p_x,p_y,p_z)}{\partial(u,v,p)},\quad
  \Sigma_{\vec p}\;=\;M\,\Sigma_{\vec\alpha}\,M^\top.
  \label{eq:partIa:fisher:Sp}
\end{equation}

由 $(p_x,p_y,p_z)=p\cdot(u,v,1)/\sqrt{D}$，$D=1+u^2+v^2$，逐项求偏导：
\begin{align}
  \frac{\partial p_x}{\partial u}&= p\!\cdot\!\frac{\sqrt D - u\cdot u/\sqrt D}{D}
                                 = \frac{p}{\sqrt D}\!\left(1-\frac{u^2}{D}\right)
                                 = \frac{p_z}{1}\!\left(1-\frac{u^2}{D}\right),\\[2pt]
  \frac{\partial p_x}{\partial v}&= -p\,\frac{uv}{D^{3/2}} = -\frac{p_z\,uv}{D},\\
  \frac{\partial p_x}{\partial p}&= \frac{u}{\sqrt D}.
\end{align}
对 $p_y,p_z$ 分量重复同样代数，得
\begin{equation}
  \boxed{\;M = \begin{pmatrix}
    p_z\!\left(1-u^2/D\right) & -p_z\,uv/D & u/\sqrt D \\
    -p_z\,uv/D & p_z\!\left(1-v^2/D\right) & v/\sqrt D \\
    -p_z\,u/D & -p_z\,v/D & 1/\sqrt D
  \end{pmatrix}.\;}
  \label{eq:partIa:fisher:M}
\end{equation}

\begin{lemma}[$M$ 的行列式]\label{lem:partIa:fisher:detM}
  $\det M = p_z^2/\sqrt D$.
\end{lemma}
\begin{proof}
  分块或直接展开均可。经济做法：把 $\vec p=p\,\hat n(u,v,1)/\sqrt D$ 看作 $p\to|\vec p|$、$(u,v)\to$ 单位球面上的角坐标的极坐标变换；$|\det\partial(p_x,p_y,p_z)/\partial(p,u,v)|=p^2|\sin\theta|$ 形式经 $u=\tan\theta_x,v=\tan\theta_y$ 重参数化即得 $p_z^2/\sqrt D$。
\end{proof}

\paragraph{小角度极限 $u,v\to 0$。} 此时 $D\to 1$，$M$ 退化为
\begin{equation}
  M_0 \;=\;\begin{pmatrix}p_z & 0 & 0\\ 0 & p_z & 0\\ 0 & 0 & 1\end{pmatrix}\quad(u=v=0).
  \label{eq:partIa:fisher:M_smallangle}
\end{equation}
即 \emph{一阶线性化下}：$\sigma_{p_x}=p_z\sigma_u$、$\sigma_{p_y}=p_z\sigma_v$、$\sigma_{p_z}=\sigma_p$。这是状态报告 \StatusReport § 从参数空间传播到动量空间 中默认的"事件 \#399 比值 1.00"代数。在 $|\vec p|\sim 635\;\mathrm{MeV}/c$、$p_y\sim -55\;\mathrm{MeV}/c$ (即 $v\sim -0.087$) 的真实事件下这一近似的二阶残差是
\begin{equation}
  M_{22}-p_z = -p_z v^2/D \;\approx\; -p_z\cdot v^2 \;\sim\;-4.7\;\mathrm{MeV}/c\cdot\,\text{(per unit }v\text{)},
\end{equation}
量级 $\sim 7.6\times 10^{-3}$ 的相对修正 —— 看似小，但与下节"二阶项的偏置贡献"相加后可解释 $p_y$ 上 1.9× 的覆盖率亏损。

\paragraph{二阶展开与覆盖率亏损。}
若把 $\vec p=g(\vec\alpha)$ 二阶展开
\begin{equation}
  \vec p \;\approx\; g(\hat{\vec\alpha}) + M\,\Delta\vec\alpha + \tfrac12 H\,\Delta\vec\alpha^{\otimes 2}
\end{equation}
代入 $\mathbb E[\Delta\vec\alpha]=\vec 0$、$\mathrm{Cov}(\Delta\vec\alpha)=\Sigma_{\vec\alpha}$，得
\begin{equation}
  \mathbb E[\vec p] \;=\; g(\hat{\vec\alpha}) + \tfrac12\mathrm{tr}\!\bigl(H\,\Sigma_{\vec\alpha}\bigr),
  \quad
  \mathrm{Cov}(\vec p)\;=\;M\Sigma_{\vec\alpha}M^\top + O(\Sigma_{\vec\alpha}^2),
  \label{eq:partIa:fisher:nonlin_bias}
\end{equation}
即 \textbf{非线性传播在均值上引入偏置}（项 $\tfrac12\mathrm{tr}(H\Sigma_{\vec\alpha})$），\textbf{在协方差上 $M\Sigma M^\top$ 形式只是一阶}。CRB 的可达性需要 MLE 无偏；非线性偏置直接破坏了这一点。

具体到 $p_y=p\,v/\sqrt D$，$\partial^2 p_y/\partial v^2=-p\,(2v/D-3v^3/D^2)/\sqrt D\sim-2p_z v/D$；在 $v\sim-0.087$、$\sigma_v\sim 5\times 10^{-4}$ (Fisher 量级) 时，二阶项贡献的方差修正约为
\begin{equation}
  \tfrac12\bigl|\partial^2_v p_y\bigr|\,\sigma_v^2 \;\sim\; 0.087\cdot p_z\cdot\sigma_v^2/D
  \;\sim\;\mathcal O(10^{-2})\;\mathrm{MeV}/c
\end{equation}
对均值；但更关键的是 \emph{二阶 Edgeworth 校正} 给方差的修正比例与 $|H|\sigma_\alpha/|M|$ 同阶，在 $p_y$ 方向上估算为 $\sim 50\%$。这直接解释了 Proposition \ref{prop:partIa:saturation} 中 $\sigma_{p_y}$ 被低估 $1.9\times$ 的根因 —— Part II 第 8 章将给出可验证的修复方案 (reparam 到 $p_x,p_y,p_z$ 直接拟合)。

\subsection{kThreePointFree 模式的扩展}
\label{sec:partIa:fisher:5dof}

主拟合模式为 \texttt{kThreePointFree} ($\vec\alpha=(\delta x,\delta y,u,v,p)$，5 自由参数)，$N_\mathrm{res}=4+2=6$ (4 PDC 残差 + 2 靶点软先验残差)，$N_\mathrm{dof}=1$。

Fisher 矩阵扩展为 $5\times 5$：
\begin{equation}
  \mathcal I_5 \;=\;
  \begin{pmatrix}
    \mathcal I_{xx} & \mathcal I_{xy} & \mathcal I_{xu} & \mathcal I_{xv} & \mathcal I_{xp}\\
    \cdot & \mathcal I_{yy} & \mathcal I_{yu} & \mathcal I_{yv} & \mathcal I_{yp}\\
    \cdot & \cdot & \mathcal I_{uu} & \mathcal I_{uv} & \mathcal I_{up}\\
    \cdot & \cdot & \cdot & \mathcal I_{vv} & \mathcal I_{vp}\\
    \cdot & \cdot & \cdot & \cdot & \mathcal I_{pp}
  \end{pmatrix}.
\end{equation}
靶点软先验 ($\sigma_{xy}^\mathrm{tgt}=5\;\mathrm{mm}$) 给 $\mathcal I_{xx},\mathcal I_{yy}$ 各加上 $1/(\sigma_{xy}^\mathrm{tgt})^2=4\times 10^{-2}\;\mathrm{mm}^{-2}$。$3\times 3$ 子块 $(u,v,p)$ 与 §\ref{sec:partIa:fisher:Jderiv} 推导的形式相同；新增的 $\delta x,\delta y$ 行/列描述靶点位置与方向/动量的耦合。

物理上 $\delta x,\delta y$ 与 $u,v$ 的耦合通过 RK 轨迹积分强相关：在 PDC 处的命中位置 $\Delta r\sim\Delta x_\mathrm{tgt}+L_\mathrm{arm}\Delta u$，$L_\mathrm{arm}\sim 5\;\mathrm m$，因此 $\mathcal I_{xu}\sim L_\mathrm{arm}\,\mathcal I_{uu}$；$5\times 5$ 矩阵存在显著的非对角块。状态报告 \StatusReport § 关于 \texttt{kThreePointFree} 的 worked example 展示了实际 $\Sigma_5$。

\subsection{动量空间协方差矩阵 $M$ (5DoF) 的扩展}
\label{sec:partIa:fisher:M5}

5DoF 模式下 $M\in\mathbb R^{3\times 5}$：动量分量不依赖 $\delta x,\delta y$，故前两列为零：
\begin{equation}
  M_5 \;=\;
  \begin{pmatrix}
    0 & 0 & p_z(1-u^2/D) & -p_z uv/D & u/\sqrt D\\
    0 & 0 & -p_z uv/D & p_z(1-v^2/D) & v/\sqrt D\\
    0 & 0 & -p_z u/D & -p_z v/D & 1/\sqrt D
  \end{pmatrix}.
\end{equation}
$\Sigma_{\vec p}=M_5\Sigma_5 M_5^\top$ 仅取出 $\Sigma_5$ 的 $(u,v,p)$ 子块进行变换 —— 数学上 5DoF 的 $\Sigma_{\vec p}$ 与 3DoF 在 $(u,v,p)$ 子块对齐时\emph{相等}。然而，5DoF 拟合中 $\Sigma_{(u,v,p)}^{(5)}$ 与 3DoF $\Sigma_{(u,v,p)}^{(3)}$ \emph{不相等}：5DoF 需要从 5 维 Fisher 中边际化掉 $\delta x,\delta y$，得到的 $(u,v,p)$ 边际方差比 3DoF 大 (因为 5DoF 多了 $\delta x,\delta y$ 的不确定性，通过非对角项漏到 $u,v,p$)。这是 5DoF 模式下 Fisher $\sigma$ 系统性偏大的来源；状态报告 § 拟合模式 已隐含该效应。

\subsection{条件数与几何参数的关系}
\label{sec:partIa:fisher:cond_geom}

把 Fisher 谱直观化：
\begin{itemize}[nosep]
  \item PDC 间距 $\Delta z_{12}=z_2-z_1\approx 1\;\mathrm m$ —— 决定 $\partial r_2/\partial u$ 与 $\partial r_1/\partial u$ 的差，进而决定 $J$ 的"角度"列 (1, 2)。
  \item PDC-靶距 $L_\mathrm{arm}\sim 5\;\mathrm m$ —— 决定 $\delta x_\mathrm{tgt}$ 到 PDC 的杠杆，进而决定 $\delta x$ 与 $u$ 的耦合强度。
  \item 磁场弯曲角 $\sim 25^\circ$ ($p\sim 635\;\mathrm{MeV}/c$, $B\sim 1.15\;\mathrm T$) —— 决定 $J$ 第三列 (动量) 的幅度。
  \item PDC 倾角 $\theta=57^\circ$ —— 决定 $\hat u,\hat v$ 在 $(x,y,z)$ 的投影分配，进而决定 $J_{i1}$ 与 $J_{i2}$ 的相对大小 (本章 §\ref{sec:partIa:fisher:sigmapy})。
\end{itemize}

设 PDC 位置分辨为 $\sigma_\mathrm{wire}=0.5\;\mathrm{mm}$，量纲分析给出：
\begin{equation}
  \sigma_u^\mathrm{Fisher} \;\sim\;\frac{\sigma_\mathrm{wire}\sqrt 2}{\Delta z_{12}\cos\theta/\sqrt 2}\;\sim\;\frac{0.5\;\mathrm{mm}\cdot 2}{1000\;\mathrm{mm}\cdot 0.385}\approx 2.6\times 10^{-3}.
\end{equation}
代回 $\sigma_{p_x}\approx p_z\sigma_u\sim 1.6\;\mathrm{MeV}/c$ —— 与状态报告事件 \#399 的 $\sigma_{p_x}=2.26\;\mathrm{MeV}/c$ 量级吻合 (差异来自非对角项与靶点软先验)。同样估算 $\sigma_{p_y}$ 给出 $\sim 0.9\;\mathrm{MeV}/c$ —— 与事件级 $0.35\;\mathrm{MeV}/c$ 比有 2-3× 偏差，主要因为多事件平均与单事件值相差 (上述估算未考虑 5DoF $\delta x,\delta y$ 漏入)；定性结论 $\sigma_{p_y}\ll\sigma_{p_x}$ 仍稳健。

闭合形 Fisher 公式 (Glückstern 形式) 由 Part I 第 7 章给出，与本节量纲估算严格一致。

\subsection{代码锚点与算法实现}
\label{sec:partIa:fisher:code}

Fisher 矩阵在仓库中的实现路径：
\begin{itemize}[nosep]
  \item \texttt{libs/analysis\_pdc\_reco/src/PDCErrorAnalysis.cc} —— \texttt{computeFisher} 取 LM 收敛态，调用 \texttt{computeWhitenedResidualJacobian} 得 $J$，构造 $\mathcal I=J^\top J$。注意此处实际使用 \texttt{TMatrixDSymEigen}（行 $\sim 355$）做对称特征值分解，\emph{不是} 显式 SVD —— 但二者对正定 $\mathcal I$ 等价 (特征值即奇异值平方)。
  \item 同文件 \texttt{computeMomentumCovariance} —— 数值差分计算 $M=\partial(p_x,p_y,p_z)/\partial\vec\alpha$ (中心差分，与 LM 内部 Jacobian 共享步长策略 \StatusReport § Levenberg-Marquardt)；最终 $\Sigma_{\vec p}=M\Sigma_{\vec\alpha}M^\top$。
  \item 数据结构 \texttt{IntervalEstimate} (\texttt{libs/analysis\_pdc\_reco/include/PDCErrorAnalysis.hh}) 统一存储 Fisher / Laplace / Profile / MCMC 四种方法的中心值与 68\%/95\% 区间，供下游 \texttt{validation\_summary.csv} 统计。
\end{itemize}

\paragraph{数值健壮性。} 实现中两道闸门：
\begin{enumerate}[nosep]
  \item 条件数检查 $\kappa<\kappa_\mathrm{max}\sim 10^{12}$；
  \item 最小特征值检查 $\lambda_\mathrm{min}>\epsilon\sim 10^{-14}$。
\end{enumerate}
触发任一闸门 $\Rightarrow$ Fisher 标记失败，\texttt{validation\_summary.csv} 中本事件 \texttt{cover68/95} 不计入。在 744-事件 ensemble 中 Fisher 失败率 $<1\%$（细节由 Part III 第 6 章 LM 收敛性研究给出）。

\subsection{小结}
\label{sec:partIa:fisher:summary}

本章建立了 Fisher 矩阵从抽象 $\mathcal I=J^\top J$ 到 SAMURAI 谱仪几何下可观测量的完整桥梁。三个关键事实：
\begin{enumerate}[nosep]
  \item Fisher 矩阵在似然层面 (R1-R5) 完全合法；CRB 等号在 $|\vec p|, p_x, p_z$ 三分量上近似饱和 (覆盖率 0.76-0.80)。
  \item $\sigma_{p_y}$ 极小，根源是丝方向 $\hat u,\hat v$ 在 $y$ 方向投影最大，$J$ 第二列幅度比第一列大约 $1/\cos\theta\approx 1.84$ 倍 —— 这是 \emph{Fisher 估计本身正确} 的事实。
  \item $\sigma_{p_y}$ 实测覆盖率 0.42 (低估 $\sim 1.9\times$) 来自 $(u,v,p)\to(p_x,p_y,p_z)$ 协方差变换中被忽略的二阶项 (式 \eqref{eq:partIa:fisher:nonlin_bias})；修复路径见 Part II § 8。
\end{enumerate}

下一章 (本部分第 4 章 Wilks 定理) 将处理小样本 ($N_\mathrm{dof}=1$) 下 Fisher 失效的另一面 —— 似然函数偏离高斯，需要 Profile Likelihood 或 MCMC 才能得到正确分位数。

% TODO: figure - Fisher J 三列在 SAMURAI xz 平面的几何示意 (status report figures/geometry/samurai_geometry_xz.png 可作底图)
% TODO: figure - $\kappa$ 与 reco $|\vec p|$ 的散点图，标出 $\kappa=10^{12}$ 警戒线 (来自 744-event ensemble per_event 输出，当前 CSV 不直接含 kappa；需新出 run)

% =============================================================================
% rk_deep_dive_part1b_math_ch4to6_zh.tex
%
% Purpose:
%   Part I (Mathematical foundations), Chapters 4–6 of the RK error-analysis
%   deep-dive companion document. Covers:
%     §4  Wilks 定理与 Profile Likelihood 的有效边界
%     §5  Laplace 方法与 Bernstein–von Mises 定理
%     §6  MCMC 理论：Metropolis-Hastings, 详细平衡、最优 acceptance、收敛诊断
%
% Parent file (wrapper):
%   docs/reports/reconstruction/momentum_reco/rk/rk_error_analysis_deep_dive_20260426_zh.tex
%
% Related companion files:
%   rk_deep_dive_part1a_math_ch1to3_zh.tex (Likelihood/CR/Fisher — separate)
%   rk_deep_dive_part1c_math_ch7to9_zh.tex (Glückstern/Bethe-Bloch/Highland)
%
% Status report (cross-reference target):
%   rk_reconstruction_status_20260416_zh.tex  — labelled \StatusReport
%
% Date         : 2026-04-26
% Author       : Baiting Tian
% Convention   : 默认靶系 (beam-as-Z)；公式编号 \label{eq:partIb:section:topic}
%
% Local TOC:
%   §4  Wilks 定理与 Profile Likelihood 的有效边界
%       4.1  陈述与意义
%       4.2  证明骨架（二阶 Taylor 展开）
%       4.3  正则条件
%       4.4  失败模式
%       4.5  与 N_dof = 1 重建问题的关联
%       4.6  warm-start 扫描策略
%   §5  Laplace 方法与 Bernstein–von Mises 定理
%       5.1  Laplace 公式
%       5.2  扁平先验下退化为 Fisher
%       5.3  含动量高斯先验的显式表达
%       5.4  Bernstein–von Mises 定理
%       5.5  高阶修正 (Tierney–Kadane) 概述
%   §6  MCMC 理论
%       6.1  Markov 链基础与平稳性
%       6.2  详细平衡与 MH 接受概率
%       6.3  遍历性条件
%       6.4  最优接受率 0.234
%       6.5  自适应步长与收敛后停止自适应
%       6.6  Geweke z-score
%       6.7  ESS 与积分自相关时间
%       6.8  当前 MCMC 配置诊断与改进建议
% =============================================================================

\providecommand{\StatusReport}{文献~[1]}

\section{Wilks 定理与 Profile Likelihood 的有效边界}
\label{sec:partIb:wilks}

\subsection{陈述与意义}
\label{sec:partIb:wilks:statement}

设 $L(\vec{\alpha}\mid D)$ 为参数 $\vec{\alpha}\in\mathbb R^{n}$ 的似然函数，
$\hat{\vec{\alpha}}$ 为不受约束的极大似然估计 (MLE)。
若原假设 $H_0$ 把 $\vec{\alpha}$ 限制在一个 $k$ 维光滑约束子流形上，
即 $H_0:\;\vec{\alpha}\in\mathcal M_0\subset\mathbb R^{n}$，
其中 $\dim\mathcal M_0 = n-k$，则在 $H_0$ 成立时
\begin{equation}
  -2\ln\Lambda
  \;=\;
  -2\bigl[\ln L(\hat{\vec{\alpha}}_0)-\ln L(\hat{\vec{\alpha}})\bigr]
  \;\xrightarrow{d}\;\chi^2_{k},
  \qquad N\to\infty,
  \label{eq:partIb:wilks:thm}
\end{equation}
其中 $\hat{\vec{\alpha}}_0$ 是受 $H_0$ 约束的 MLE，$N$ 为独立观测数。
这就是 \emph{Wilks 定理}\,(1938)。

在重建语境下，$\vec{\alpha}=(u,v,p)$ 或 $(\delta x,\delta y,u,v,p)$。
Profile Likelihood 把某一个分量 $\alpha_k$ 固定为常数
($H_0:\alpha_k = \alpha_k^{(0)}$，$k=1$)，
对其余参数极大化 $L$；
\eqref{eq:partIb:wilks:thm} 退化为
\begin{equation}
  \Delta\chi^2(\alpha_k^{(0)})
  \;=\;
  \chi^2_{\mathrm{prof}}(\alpha_k^{(0)})-\chi^2_{\min}
  \;\xrightarrow{d}\;\chi^2_{1}.
  \label{eq:partIb:wilks:profile}
\end{equation}
$68.27\%$ 置信限对应 $\Delta\chi^2 = 1$；$95.45\%$ 对应 $\Delta\chi^2 = 4$。

\subsection{证明骨架（二阶 Taylor 展开）}
\label{sec:partIb:wilks:proof}

\begin{theorem}[Wilks 1938，弱形式]
\label{thm:partIb:wilks}
若似然 $L$ 满足 §\ref{sec:partIb:wilks:cond} 列出的正则条件，
则 \eqref{eq:partIb:wilks:thm} 在 $H_0$ 下成立。
\end{theorem}

\begin{proof}[证明骨架]
不失一般性把 $H_0$ 写为 $\vec{\alpha}_{(2)} = \vec{0}$，
其中 $\vec{\alpha}=(\vec{\alpha}_{(1)},\vec{\alpha}_{(2)})$，
$\vec{\alpha}_{(2)}\in\mathbb R^{k}$。
记 $\ell(\vec{\alpha})\equiv \ln L(\vec{\alpha}\mid D)$、
$\vec{s}\equiv\nabla\ell(\hat{\vec{\alpha}})=\vec{0}$ (定义)、
$\mathcal H\equiv -\nabla^2\ell(\hat{\vec{\alpha}})$ 为观测 Fisher 信息。

\textbf{步骤 1.}\,在 $\hat{\vec{\alpha}}$ 周围作二阶 Taylor 展开：
\begin{equation}
  \ell(\vec{\alpha})
  =\ell(\hat{\vec{\alpha}})
   -\tfrac12(\vec{\alpha}-\hat{\vec{\alpha}})^\top
            \mathcal H\,
            (\vec{\alpha}-\hat{\vec{\alpha}})
   +O(\|\vec{\alpha}-\hat{\vec{\alpha}}\|^{3}).
  \label{eq:partIb:wilks:taylor}
\end{equation}

\textbf{步骤 2.}\,设 $\hat{\vec{\alpha}}_0$ 为 $H_0$ 下的 MLE。
对 \eqref{eq:partIb:wilks:taylor} 关于 $\vec{\alpha}_{(1)}$ 求极小，
由 Schur 余子式得
\begin{equation}
  \ell(\hat{\vec{\alpha}})-\ell(\hat{\vec{\alpha}}_0)
  =\tfrac12\,
   \hat{\vec{\alpha}}_{(2)}^{\top}
   (\mathcal H/\mathcal H_{11})\,
   \hat{\vec{\alpha}}_{(2)}
   +o_p(1),
  \label{eq:partIb:wilks:schur}
\end{equation}
其中 $\mathcal H/\mathcal H_{11}=\mathcal H_{22}-\mathcal H_{21}\mathcal H_{11}^{-1}\mathcal H_{12}$
是 $\mathcal H$ 关于 $\vec{\alpha}_{(2)}$ 块的 Schur 补，恰为 $\Sigma_{22}^{-1}$
($\Sigma\equiv\mathcal H^{-1}$)。

\textbf{步骤 3.}\,由极大似然渐近性
$\sqrt{N}(\hat{\vec{\alpha}}-\vec{\alpha}^{*})\xrightarrow{d}\mathcal N(\vec 0,\,I^{-1}_{\vec{\alpha}^{*}})$，
当 $\vec{\alpha}^{*}\in H_0$ 时
$\sqrt{N}\,\hat{\vec{\alpha}}_{(2)}\xrightarrow{d}\mathcal N(\vec 0,\Sigma_{22})$。
代入 \eqref{eq:partIb:wilks:schur}：
\begin{equation}
  -2[\ell(\hat{\vec{\alpha}}_0)-\ell(\hat{\vec{\alpha}})]
  =N\,\hat{\vec{\alpha}}_{(2)}^{\top}\Sigma_{22}^{-1}\hat{\vec{\alpha}}_{(2)}+o_p(1)
  \xrightarrow{d}\chi^2_{k},
  \label{eq:partIb:wilks:final}
\end{equation}
因 $\vec X^{\top}\Sigma^{-1}\vec X\sim\chi^2_{k}$ 当
$\vec X\sim\mathcal N(\vec 0,\Sigma)$。\qedhere
\end{proof}

可见 Wilks 定理本质上是“$\chi^2$ 极小邻域呈二次形 $\Rightarrow$ 似然比统计量
等同于平方分”的结论；其核心是 \emph{Score 的 CLT} 与 \emph{Hessian 的弱稳定性}。

\subsection{正则条件}
\label{sec:partIb:wilks:cond}

定理 \ref{thm:partIb:wilks} 需要下列假设同时成立：
\begin{enumerate}[label=(R\arabic*)]
  \item $\ell$ 关于 $\vec{\alpha}$ 三阶连续可微，且 $\vec{\alpha}^{*}$ 是参数空间内点；
  \item Fisher 信息 $I(\vec{\alpha}^{*})$ 有限、正定；
  \item $H_0$ 把 $\vec{\alpha}$ 限制在一个光滑子流形 $\mathcal M_0$ 上，
        其法向投影矩阵满秩 (无识别性退化)；
  \item 数据对参数的依赖在 $\vec{\alpha}^{*}$ 邻域内一致；
  \item 不存在不可分辨参数 (nuisance pathology)。
\end{enumerate}

R1 排斥参数边界 (例如 $\sigma\geq 0$ 取到 $\sigma=0$ 时分布不再是 $\chi^2_{k}$
而是\emph{ chi-bar squared}\,)。
R2 排斥 Fisher 矩阵奇异 (即似然在某方向上对参数不敏感)。
R3 排斥退化几何 (例如把 5 维参数限制到一个不光滑的代数簇)。

\subsection{失败模式}
\label{sec:partIb:wilks:fail}

下列情形下 \eqref{eq:partIb:wilks:thm} 失效：

\paragraph{(a) 边界参数 (boundary case)。}
若 $H_0$ 处于参数空间边界 (如方差非负、概率非负)，
$\hat{\vec{\alpha}}_0$ 以正概率被边界“卡住”；
渐近分布是若干 $\chi^2$ 的混合 (Self \& Liang 1987)。
对本重建问题：$p > 0$ 是物理边界，
但 $\hat p\sim 635\,\mathrm{MeV}/c$ 远离零，故 (a) 不触发。

\paragraph{(b) 识别性丧失。}
若 $\partial\ell/\partial\alpha_k$ 在某方向恒为零，Fisher 矩阵秩亏，
\eqref{eq:partIb:wilks:final} 中 $\Sigma_{22}^{-1}$ 不存在。
重建中的征兆：Fisher SVD 出现接近零的奇异值；
\StatusReport\ §方法一 已要求条件数检查并在病态时返回失败。

\paragraph{(c) 多模态。}
LM 收敛至局部极小而非全局，导致 $\chi^2_{\min}$ 偏大；
此时 \eqref{eq:partIb:wilks:profile} 中的 $\Delta\chi^2$ 高估。
重建中的对策：multistart 网格 + 选 $\chi^2$ 全局最小 (\StatusReport\ §LM 流程)。

\paragraph{(d) 小样本。}
$\chi^2_{k}$ 是渐近极限；$N_{\mathrm{dof}}$ 很小时，二阶 Taylor 展开
\eqref{eq:partIb:wilks:taylor} 的余项 $O(\|\Delta\vec{\alpha}\|^3)$ 不能被忽略。
本重建 $N_{\mathrm{dof}} = 1$ (4 残差 $-$ 3 自由参数)
正处于这一边界——见 §\ref{sec:partIb:wilks:tie}。

\subsection{与 $N_{\mathrm{dof}}=1$ 重建问题的关联}
\label{sec:partIb:wilks:tie}

\StatusReport\ Table “聚合误差对比” 给出全 run 中位数：
\begin{equation}
  \sigma^{\mathrm{Fisher}}_{|\vec p|} = 48.20\,\mathrm{MeV}/c,
  \quad
  \sigma^{\mathrm{Profile}}_{|\vec p|}
   \!\equiv\!\tfrac12(\,q^{+}\!-q^{-})
   = 76.91\,\mathrm{MeV}/c,
  \label{eq:partIb:wilks:numbers}
\end{equation}
即 Profile 区间比 Fisher 半宽宽约 $60\%$。
此差异不是数值误差，而是 \emph{Wilks 渐近近似的有效性破缺}：

\begin{itemize}[leftmargin=2em]
  \item Fisher 用最优点处的 $\nabla^2\chi^2$ (二阶导数) 推出对称 $\sigma$；
  \item Profile 直接读取 $\Delta\chi^2 = 1$ 等高线，与曲面真实形状一致；
  \item 二者不一致 $\Rightarrow$ $\chi^2$ 曲面非二次。
\end{itemize}

事件 \#399 (\StatusReport\ Table “Profile vs Fisher/Laplace”) 把这一图像放大到单事件层面：
$|\vec p|$ 上 Profile 区间 $[660.73,\,667.77]\,\mathrm{MeV}/c$ 极度非对称
($\hat p - q^- = 7.02$，$q^+ - \hat p = 0.02$)。
朝 $|\vec p|$ 增大方向 $\chi^2$ 极陡 ($\chi^2_{\mathrm{red}}=20.4$ 已经偏离物理)；
朝减小方向有较深低 $\chi^2$ 阱。
Fisher 取最优点附近的曲率，无法察觉远端的非对称性，故给出对称的 $\pm 2.84\,\mathrm{MeV}/c$。
Profile 才是真实置信限的代表。

\paragraph{物理来源诊断。}
$\chi^2_{\mathrm{red}}=20.4$ 远大于 1 的事实本身说明数据与模型不相容
——在此重建中最可能的根源是 §第 8 章详细讨论的
\emph{未建模 $dE/dx$ 偏置}：模型中 PDC1$\to$PDC2 的动量被视为常数，
但 6\,m 空气损失 $\sim 5$–$15\,\mathrm{MeV}/c$，
拟合通过把 $|\vec p|$ 推向偏小数值来部分补偿，
但残差仍累积 $\Rightarrow$ $\chi^2$ 在最优处仍偏大。

\subsection{warm-start 扫描策略}
\label{sec:partIb:wilks:warmstart}

Profile 扫描在每个固定 $\alpha_k^{(i)}$ 处都需独立 LM 优化
(只优化其余参数)。若每次都从全局 MLE 出发，
当 $\alpha_k^{(i)}$ 距 $\hat\alpha_k$ 较远时，LM 步数显著增加。
\StatusReport\ §方法三 实际算法采用 warm-start：
\begin{equation}
  \vec\alpha_{\mathrm{warm}}^{(i)} = \arg\min_{\vec\alpha_{(1)}}
     \chi^2(\vec\alpha\,;\,\alpha_k=\alpha_k^{(i)}),
  \quad
  \vec\alpha_{\mathrm{warm}}^{(0)} = \hat{\vec\alpha},
  \label{eq:partIb:wilks:warm}
\end{equation}
即把上一扫描点的收敛态作为下一点的初值。

\begin{proposition}[warm-start 不改变 Profile 估计的渐近性]
\label{prop:partIb:warmstart}
若每个 $\alpha_k^{(i)}$ 上 LM 仍收敛到全局 (子) 极小，
warm-start 仅改变路径而不改变终值；
\eqref{eq:partIb:wilks:profile} 给出的 $\Delta\chi^2$ 不变，
Wilks 渐近不受影响。
\end{proposition}

\begin{proof}[证明]
LM 的收敛终值由所在 basin of attraction 决定。
对相邻 $\alpha_k^{(i)},\alpha_k^{(i+1)}$，若步长 $|\Delta\alpha_k|$ 足够小，
两点的 basin 几何连续，warm-start 起点位于同一 basin。
因 LM 必须满足 $\nabla_{\vec\alpha_{(1)}}\chi^2 = 0$ 才停止，
两种初值给出同一极小点，证毕。\qedhere
\end{proof}

实际工程意义：本重建的 $N_{\mathrm{scan}}=17$ 步、
扫描范围 $\pm 2.5\sigma_{\mathrm{Fisher}}$、
平均每步 $\Delta\alpha_k \sim 0.3\sigma$，
warm-start 把每步 LM 迭代数从 $\sim 10$ 降到 $\sim 1$–$2$，
全 Profile 计算成本由 $\mathcal O(170)$ 次 RK 传播降到 $\mathcal O(30)$。
对 815 条 track + 5 个参数槽位，
总成本从 $\sim 70\times 10^4$ RK 降到 $\sim 12\times 10^4$ RK，
是 Profile 扫描在生产代码中可负担的关键。

\section{Laplace 方法与 Bernstein--von Mises 定理}
\label{sec:partIb:laplace}

\subsection{Laplace 公式}
\label{sec:partIb:laplace:formula}

\paragraph{动机。}
贝叶斯框架下，参数的后验密度为
\begin{equation}
  \pi(\vec\alpha\mid D)
  = \frac{L(\vec\alpha\mid D)\,\pi_0(\vec\alpha)}{Z(D)},
  \qquad
  Z(D)=\!\!\int\!\! L(\vec\alpha\mid D)\pi_0(\vec\alpha)\,d^n\alpha.
  \label{eq:partIb:laplace:bayes}
\end{equation}
将其写成指数形式 $\pi(\vec\alpha\mid D)\propto e^{-\Phi(\vec\alpha)}$，
其中
\begin{equation}
  \Phi(\vec\alpha)
  \;=\;
  -\ln L(\vec\alpha\mid D)-\ln\pi_0(\vec\alpha)
  \;=\;
  \tfrac12\chi^2(\vec\alpha)-\ln\pi_0(\vec\alpha)+\mathrm{const}.
  \label{eq:partIb:laplace:phi}
\end{equation}
设最大后验估计 (MAP) $\hat{\vec\alpha}_{\mathrm{MAP}}=\arg\min_{\vec\alpha}\Phi(\vec\alpha)$，
在 $\hat{\vec\alpha}_{\mathrm{MAP}}$ 处作二阶 Taylor 展开：
\begin{equation}
  \Phi(\vec\alpha)
  \approx
  \Phi(\hat{\vec\alpha}_{\mathrm{MAP}})
  +\tfrac12 (\vec\alpha-\hat{\vec\alpha}_{\mathrm{MAP}})^{\top}
            \mathcal H_{\mathrm{post}}\,
            (\vec\alpha-\hat{\vec\alpha}_{\mathrm{MAP}}),
  \label{eq:partIb:laplace:taylor}
\end{equation}
其中 $\mathcal H_{\mathrm{post}}=\nabla^2\Phi(\hat{\vec\alpha}_{\mathrm{MAP}})$。
代入 \eqref{eq:partIb:laplace:bayes} 得 \emph{Laplace 近似后验}：
\begin{equation}
  \boxed{
  \pi^{\mathrm{Laplace}}(\vec\alpha\mid D)
  \;\sim\;
  \mathcal N\!\bigl(\hat{\vec\alpha}_{\mathrm{MAP}},
                   \;\Sigma_{\mathrm{post}}\bigr),
  \quad
  \Sigma_{\mathrm{post}}=\mathcal H_{\mathrm{post}}^{-1}.
  }
  \label{eq:partIb:laplace:posterior}
\end{equation}
归一化常数也有渐近形式 (Laplace 积分公式)：
\begin{equation}
  Z(D)
  \;\approx\;
  e^{-\Phi(\hat{\vec\alpha}_{\mathrm{MAP}})}
  \,(2\pi)^{n/2}
  \,\bigl|\det \mathcal H_{\mathrm{post}}\bigr|^{-1/2}.
  \label{eq:partIb:laplace:Z}
\end{equation}
这正是 Bayes Factor 计算中常用的 \emph{Bayesian information criterion} 的来源。

\subsection{扁平先验下退化为 Fisher}
\label{sec:partIb:laplace:flat}

若先验 $\pi_0(\vec\alpha)=\mathrm{const}$ (improper flat prior，仅在感兴趣区域内被 $L$ 控制)：
\begin{equation}
  \Phi(\vec\alpha) = \tfrac12 \chi^2(\vec\alpha)+\mathrm{const},
  \qquad
  \mathcal H_{\mathrm{post}}
  = \nabla^2\bigl(\tfrac12\chi^2\bigr)
  = J^{\top}J,
  \label{eq:partIb:laplace:flat}
\end{equation}
其中 $J$ 是白化残差关于参数的 Jacobian。
此时 MAP 等同于 MLE，
\begin{equation}
  \Sigma_{\mathrm{post}}^{\mathrm{flat}}
  = (J^{\top}J)^{-1}
  \;=\;\Sigma_{\mathrm{Fisher}}.
  \label{eq:partIb:laplace:eq_fisher}
\end{equation}

\paragraph{数值证据 (\StatusReport)。}
本 run 未启用动量先验 (\texttt{MomentumPrior::enabled = false})，故应触发
\eqref{eq:partIb:laplace:eq_fisher}。
事实上 \StatusReport\ §汇总观察到：
\begin{equation}
  \sigma^{\mathrm{Fisher}}_{|\vec p|} = 48.20\,\mathrm{MeV}/c,
  \qquad
  \sigma^{\mathrm{Laplace}}_{|\vec p|}=48.14\,\mathrm{MeV}/c,
  \label{eq:partIb:laplace:numbers}
\end{equation}
差异 $0.06\,\mathrm{MeV}/c$ ($0.12\%$) 全部来自数值实现细节
(白化策略、SVD 截断阈值、$|\vec p|$ 经验分位 vs 高斯投影)，
与 \eqref{eq:partIb:laplace:eq_fisher} 一致。
事件 \#399 上 Fisher 半宽 / Laplace 半宽 $=1.00$ 比值精确等于 1
是这一退化关系最尖锐的验证。

\subsection{含动量高斯先验的显式表达}
\label{sec:partIb:laplace:prior}

为压制 LM basin 漂移问题 (\StatusReport\ §LM 收敛失败)，
在算法中已经预留了动量高斯先验：
\begin{equation}
  \pi_0(\vec\alpha)
  =\frac{1}{\sqrt{2\pi}\,\sigma_{p}}
   \exp\!\Bigl[-\tfrac{(p-p_0)^2}{2\sigma_p^2}\Bigr]
   \cdot\pi_0^{\mathrm{flat}}(u,v),
  \label{eq:partIb:laplace:gauss_prior}
\end{equation}
即 $u,v$ 用 flat prior，$p$ 用以 $p_0$ 为中心、宽度 $\sigma_p$ 的高斯。
代入 \eqref{eq:partIb:laplace:phi}：
\begin{equation}
  \Phi(\vec\alpha)
  = \tfrac12\chi^2(\vec\alpha)
   +\tfrac{(p-p_0)^2}{2\sigma_p^2}+\mathrm{const}.
  \label{eq:partIb:laplace:phi_prior}
\end{equation}
其 Hessian：
\begin{equation}
  \mathcal H_{\mathrm{post}}
  = J^{\top}J + \Lambda_{\mathrm{prior}},
  \qquad
  \Lambda_{\mathrm{prior}}
  = \mathrm{diag}\Bigl(0,\,0,\,\tfrac{1}{\sigma_p^2}\Bigr).
  \label{eq:partIb:laplace:H_with_prior}
\end{equation}
启用先验后 ($\sigma_p<\infty$)，$\Sigma_{\mathrm{post}}=\mathcal H_{\mathrm{post}}^{-1}$
比 $(J^{\top}J)^{-1}$ 更小——先验起到正则化作用。

\paragraph{$3{\times}3$ 显式公式。}
若把 $J^{\top}J$ 写成块形式
$J^{\top}J=\begin{pmatrix} A_{2\times 2} & \vec b\\ \vec b^{\top} & c\end{pmatrix}$
($u,v$ 子块为 $A$, $p$ 对角元 $c$)，则
\begin{equation}
  \mathcal H_{\mathrm{post}}
  =\begin{pmatrix} A & \vec b\\ \vec b^{\top} & c+1/\sigma_p^2\end{pmatrix},
  \;\;
  (\Sigma_{\mathrm{post}})_{pp}
  = \frac{1}{c+1/\sigma_p^2-\vec b^{\top}A^{-1}\vec b}.
  \label{eq:partIb:laplace:p_var}
\end{equation}
极限 $\sigma_p\to\infty$ 退化为 Fisher 的 $p$-边际方差
$1/(c-\vec b^{\top}A^{-1}\vec b)$；
极限 $\sigma_p\to 0$ 把 $p$ 完全锁定于 $p_0$。
这给出了 $\sigma_p$ 调参的几何直观。

\subsection{Bernstein--von Mises 定理}
\label{sec:partIb:laplace:bvm}

\begin{theorem}[Bernstein--von Mises]
\label{thm:partIb:bvm}
设 $\hat{\vec\alpha}_N$ 为基于 $N$ 个独立观测的 MLE，
真实参数 $\vec\alpha^{*}$ 为参数空间内点。
若 (a) 似然满足正则条件 R1--R5；
(b) 先验 $\pi_0(\vec\alpha)$ 在 $\vec\alpha^{*}$ 处连续、严格为正；
(c) 后验密度可积，
则后验分布 $\pi(\vec\alpha\mid D_N)$ 在 $\sqrt{N}$-缩放下渐近收敛到
$\mathcal N(\hat{\vec\alpha}_N,\,N^{-1}I^{-1}(\vec\alpha^{*}))$，
且与先验 $\pi_0$ 无关。形式化：
\begin{equation}
  \sup_{B\subset\mathbb R^n}
  \Bigl|\,P(\vec\alpha\in B\mid D_N)
        -\mathcal N(\hat{\vec\alpha}_N,N^{-1}I^{-1})(B)\Bigr|
  \;\xrightarrow{p}\;0.
  \label{eq:partIb:bvm:thm}
\end{equation}
\end{theorem}

\paragraph{物理意义。}
渐近极限下，先验信息在样本量 $N\to\infty$ 时被数据“淹没”——
后验只取决于 $\hat{\vec\alpha}_N$ 与 Fisher 信息 $I$。
Laplace 近似 \eqref{eq:partIb:laplace:posterior}
正是 \eqref{eq:partIb:bvm:thm} 的有限样本版本。

\paragraph{有限样本与 $N_{\mathrm{dof}}=1$ 的边界。}
重建中“样本量” $N$ 不是事件数，而是 \emph{每条 track 的残差数 $N_r=4$}
(对应 $N_{\mathrm{dof}} = 4-3 = 1$)。
单 track 上 BvM 渐近的 leading 项是 $\mathcal O(1/N_r)$；
$N_r=4$ 时余项不可忽略，这正是
Laplace/Fisher 给出对称区间，
而 Profile 与 MCMC 揭示非对称尾部 (§\ref{sec:partIb:wilks:tie}) 的根源。

\paragraph{结论：何时该信任 Laplace。}
\begin{enumerate}[label=(\alph*),leftmargin=2em]
  \item 当 $\chi^2_{\mathrm{red}}\sim 1$ 且 Profile 区间近对称——可信。
  \item 当 $\chi^2_{\mathrm{red}}\gg 1$ (例 \#399 = 20.4) 或
        Profile/Fisher 半宽比 $> 1.3$——不可信，应取 Profile 或 MCMC。
  \item 当含强先验 ($\sigma_p$ 与数据驱动 $\sigma_p^{\mathrm{data}}$ 相当) ——
        必须计算 Hessian \eqref{eq:partIb:laplace:H_with_prior}，
        不可直接使用 Fisher。
\end{enumerate}

\subsection{高阶修正 (Tierney--Kadane) 概述}
\label{sec:partIb:laplace:tk}

Laplace 方法的二阶 Taylor 截断带来 $\mathcal O(1/N)$ 的相对误差。
若需要更高精度地估计后验期望
$\mathbb E[g(\vec\alpha)\mid D]=\int g\,\pi\,d\vec\alpha$，
Tierney \& Kadane (1986) 给出 \emph{完全 Laplace} 形式：
\begin{equation}
  \mathbb E[g\mid D]
  \;\approx\;
  \frac{|\mathcal H^*|^{-1/2}\,e^{-\Phi^*(\hat{\vec\alpha}^*)}}
       {|\mathcal H |^{-1/2}\,e^{-\Phi(\hat{\vec\alpha})}}
  \,\bigl[1+\mathcal O(N^{-2})\bigr],
  \label{eq:partIb:laplace:tk}
\end{equation}
其中 $\Phi^*=\Phi-\ln g$，$\hat{\vec\alpha}^*=\arg\min\Phi^*$，
$\mathcal H^* = \nabla^2\Phi^*(\hat{\vec\alpha}^*)$。
关键巧思：\emph{对分子分母分别做 Laplace}，
让 $\mathcal O(N^{-1})$ 误差项相互抵消。

本重建当前实现仅使用一阶 Laplace
\eqref{eq:partIb:laplace:posterior}；
若未来需要后验中央矩 (skewness、kurtosis) 用以诊断非对称性，
\eqref{eq:partIb:laplace:tk} 是首选闭式公式。
完整推导超出本节范围，参见 Tierney \& Kadane (1986)、Schervish (1995, Ch.\,7)。

\section{MCMC 理论：Metropolis-Hastings, 详细平衡、最优 acceptance、收敛诊断}
\label{sec:partIb:mcmc}

\subsection{Markov 链基础与平稳性}
\label{sec:partIb:mcmc:basics}

设状态空间 $\mathcal X\subseteq\mathbb R^n$，
\emph{Markov 链}是序列 $\{\vec X^{(t)}\}_{t\geq 0}$，
其转移核 $K(\vec x,\vec y)$ 给出
$P(\vec X^{(t+1)}\in A\mid\vec X^{(t)}=\vec x)=\int_A K(\vec x,\vec y)\,d\vec y$。
分布 $\pi$ 为 \emph{平稳分布} 当且仅当
\begin{equation}
  \pi(\vec y)=\int_{\mathcal X}\pi(\vec x)\,K(\vec x,\vec y)\,d\vec x.
  \label{eq:partIb:mcmc:stationary}
\end{equation}
MCMC 的目的：构造 $K$ 使其平稳分布恰为目标后验
$\pi(\vec\alpha\mid D)$ \eqref{eq:partIb:laplace:bayes}。

\subsection{详细平衡与 Metropolis--Hastings 接受概率}
\label{sec:partIb:mcmc:db}

\begin{definition}[详细平衡]
\label{def:partIb:db}
转移核 $K$ 关于 $\pi$ 满足详细平衡 (detailed balance)：
\begin{equation}
  \pi(\vec x)\,K(\vec x,\vec y)\;=\;\pi(\vec y)\,K(\vec y,\vec x),
  \quad \forall\,\vec x,\vec y\in\mathcal X.
  \label{eq:partIb:mcmc:db}
\end{equation}
\end{definition}

\begin{lemma}
\label{lem:partIb:db_implies_stat}
若 $K$ 满足 \eqref{eq:partIb:mcmc:db}，则 $\pi$ 为 $K$ 的平稳分布。
\end{lemma}

\begin{proof}
对 \eqref{eq:partIb:mcmc:db} 关于 $\vec x$ 积分：
$\int\pi(\vec x)K(\vec x,\vec y)d\vec x = \pi(\vec y)\int K(\vec y,\vec x)d\vec x = \pi(\vec y)$，
即 \eqref{eq:partIb:mcmc:stationary}。\qedhere
\end{proof}

Metropolis--Hastings (MH) 通过组合 \emph{提议核} $Q(\vec x,\vec y)$ 与
\emph{接受概率} $\alpha(\vec x,\vec y)$ 构造满足 \eqref{eq:partIb:mcmc:db} 的 $K$：
\begin{equation}
  K(\vec x,\vec y)
  = Q(\vec x,\vec y)\,\alpha(\vec x,\vec y)
   +\delta(\vec y-\vec x)\,r(\vec x),
  \label{eq:partIb:mcmc:K_mh}
\end{equation}
其中 $r(\vec x)=1-\int Q(\vec x,\vec y)\alpha(\vec x,\vec y)\,d\vec y$ 为拒绝概率。
\eqref{eq:partIb:mcmc:db} 在 $\vec x\neq\vec y$ 时给出
\begin{equation}
  \pi(\vec x)\,Q(\vec x,\vec y)\,\alpha(\vec x,\vec y)
  = \pi(\vec y)\,Q(\vec y,\vec x)\,\alpha(\vec y,\vec x).
  \label{eq:partIb:mcmc:db_mh}
\end{equation}
取 $\alpha$ 形式 $\alpha(\vec x,\vec y)=\min(1,\,r_{xy})$，$\alpha(\vec y,\vec x)=\min(1,\,r_{yx})$，
不失一般性设 $r_{xy}\leq 1\leq r_{yx}$ ($r_{yx}=1/r_{xy}$)，
代入 \eqref{eq:partIb:mcmc:db_mh}：
$\pi(\vec x)Q(\vec x,\vec y)r_{xy}=\pi(\vec y)Q(\vec y,\vec x)\cdot 1$，
解出
\begin{equation}
  \boxed{
  \alpha_{\mathrm{MH}}(\vec x,\vec y)
  = \min\!\Bigl(1,\;
    \frac{\pi(\vec y)\,Q(\vec y,\vec x)}{\pi(\vec x)\,Q(\vec x,\vec y)}\Bigr).
  }
  \label{eq:partIb:mcmc:alpha_mh}
\end{equation}

\paragraph{对称提议 (Random-Walk MH)。}
若 $Q(\vec x,\vec y)=Q(\vec y,\vec x)$ (例如各向异性高斯
$\vec y=\vec x+\Sigma_{\mathrm{prop}}^{1/2}\vec\xi$, $\vec\xi\sim\mathcal N(\vec 0,I)$)，
\eqref{eq:partIb:mcmc:alpha_mh} 退化为
\begin{equation}
  \alpha_{\mathrm{RW}}(\vec x,\vec y)
  = \min\!\Bigl(1,\;\frac{\pi(\vec y)}{\pi(\vec x)}\Bigr)
  = \min\!\Bigl(1,\;e^{-\Delta\Phi}\Bigr),
  \label{eq:partIb:mcmc:alpha_rw}
\end{equation}
其中 $\Delta\Phi=\Phi(\vec y)-\Phi(\vec x)$
即 \eqref{eq:partIb:laplace:phi} 的差值。
\StatusReport\ §方法五 实现的正是这种对称随机游走 MH，
提议协方差由 Laplace $\Sigma_{\mathrm{post}}$ 缩放给出：
\begin{equation}
  \Sigma_{\mathrm{prop}} = s^2\cdot\Sigma_{\mathrm{post}},
  \label{eq:partIb:mcmc:sprop}
\end{equation}
$s$ 为待自适应的标量步长因子。

\subsection{遍历性条件}
\label{sec:partIb:mcmc:ergodic}

要使链的样本均值收敛到后验期望，需要 \emph{遍历性}：

\begin{enumerate}[label=(E\arabic*)]
  \item 不可约 (irreducible)：从任意状态出发，正概率到达任意正密度区域；
  \item 非周期 (aperiodic)：不存在周期 $d>1$ 使返回时间被锁定为 $d$ 的倍数；
  \item 正常返 (positive recurrent)：每个正密度区域被无限次访问，且平均返回时间有限。
\end{enumerate}

对连续状态空间，更合适的提法是 \emph{Harris 遍历性}：
若 $K$ 关于 $\pi$ 满足详细平衡 + 不可约 + 非周期，且 $\pi$ 是不变测度，
则对任意 $\pi$-可积函数 $g$ 和几乎所有初值 $\vec x_0$，
\begin{equation}
  \frac{1}{N}\sum_{t=1}^{N}g(\vec X^{(t)})
  \;\xrightarrow{a.s.}\;
  \int g(\vec\alpha)\,\pi(\vec\alpha)\,d\vec\alpha,
  \label{eq:partIb:mcmc:ergodic}
\end{equation}
即遍历定理。
对本重建：$\pi$ 在 $\mathbb R^3$ ($u,v,p$ 子空间) 上正密度，
高斯随机游走核显然不可约且非周期，故 \eqref{eq:partIb:mcmc:ergodic} 适用。

\subsection{最优接受率 0.234}
\label{sec:partIb:mcmc:opt_accept}

\begin{theorem}[Roberts--Gelman--Gilks 1997, 简化版]
\label{thm:partIb:rgg}
设目标分布 $\pi$ 为 $n$ 维独立同分布 $\mathcal N(\vec 0, I_n)$，
随机游走 MH 提议核 $\vec\xi\sim\mathcal N(\vec 0,(\ell^2/n)I_n)$，
$\ell>0$ 为缩放因子。
当 $n\to\infty$，链的第一坐标缩放后弱收敛到 Langevin SDE：
\begin{equation}
  d X_t = -\tfrac12 h(\ell)\,X_t\,dt + \sqrt{h(\ell)}\,d W_t,
  \quad
  h(\ell)=2\ell^2\,\Phi(-\ell/2),
  \label{eq:partIb:mcmc:rgg_speed}
\end{equation}
其中 $\Phi(\cdot)$ 是标准正态 CDF。
$h(\ell)$ 在 $\ell^{*}\approx 2.38$ 处最大，
对应的渐近接受率
\begin{equation}
  \alpha^{*}=2\Phi(-\ell^{*}/2)\;\approx\;0.234.
  \label{eq:partIb:mcmc:0234}
\end{equation}
\end{theorem}

\paragraph{解读。}
\eqref{eq:partIb:mcmc:0234} 给出了\emph{ 在高维高斯目标 + 各向同性随机游走 + 渐近极限}下
最优接受率。它平衡两件事：
\begin{itemize}[leftmargin=2em]
  \item 步子太小 ($\alpha\to 1$): 几乎全部被接受，但走得慢，自相关高；
  \item 步子太大 ($\alpha\to 0$): 几乎全部被拒绝，链停滞。
\end{itemize}
最优 $\ell^{*}\approx 2.38$ 即 $\Sigma_{\mathrm{prop}}=(2.38^2/n)\Sigma_{\mathrm{post}}$，
此时积分自相关时间最小化。

\paragraph{实证经验。}
对 $n\geq 5$，$\alpha^{*}=0.234$ 是相当鲁棒的目标
(对偏离高斯、轻度耦合的目标依旧近似有效)。
$n=1$ 时最优 $\alpha^{*}\approx 0.44$；
本重建有效维度 (启用 MCMC 的参数) 为 $n=3$ ($u,v,p$)
或 $n=5$ ($\delta x,\delta y,u,v,p$)，
处于 $\alpha^{*}\in[0.23,0.44]$ 过渡区。
\StatusReport\ §方法五 取 $0.23$ 作为自适应阈值，是合理的。

\subsection{自适应步长与“收敛后停止自适应”}
\label{sec:partIb:mcmc:adapt}

为把接受率拉到目标 $\alpha^{*}$，
\StatusReport\ §方法五 在 burn-in 期间使用 Robbins--Monro 风格更新：
\begin{equation}
  s^{(t+1)}
  \;=\;
  \begin{cases}
    s^{(t)}\cdot 1.1, & \text{若最近窗口接受率 } > 0.23,\\
    s^{(t)}\cdot 0.9, & \text{否则}.
  \end{cases}
  \label{eq:partIb:mcmc:rm}
\end{equation}

\paragraph{为什么自适应破坏 Markov 性？}
$\Sigma_{\mathrm{prop}}^{(t+1)}$ 依赖于历史接受序列 (而不仅是当前态)，
故 $\{\vec X^{(t)}\}$ 不再是 (时间齐次) Markov 链——
详细平衡 \eqref{eq:partIb:mcmc:db} 在自适应期间不严格成立，
此期间样本不应纳入后验估计。

\paragraph{合法化途径。}
\StatusReport\ 算法在 burn-in 结束后\emph{ 冻结} $s$。
此后 $\Sigma_{\mathrm{prop}}$ 为常数，链恢复时间齐次 Markov 性，
详细平衡与 \eqref{eq:partIb:mcmc:ergodic} 都生效。
此即 \emph{“diminishing adaptation + containment”} 的简化版本
(Roberts \& Rosenthal 2007)：
只要自适应在采样阶段停止，最终样本即可视为来自正确平稳分布。

\subsection{Geweke z-score：收敛性的频率学派检验}
\label{sec:partIb:mcmc:geweke}

\begin{proposition}[Geweke 1992]
\label{prop:partIb:geweke}
设链 $\{\vec X^{(t)}\}_{t=1}^{N}$ 在某起始燃尽期之后已平稳。
取前 10\% (索引 $1$ 至 $N_A=\lfloor 0.1 N\rfloor$) 的均值 $\bar X_A$
与后 50\% (索引 $N_B^{\mathrm{lo}}=\lceil 0.5 N\rceil$ 至 $N$) 的均值 $\bar X_B$。
则
\begin{equation}
  z_{\mathrm{Geweke}}
  \;=\;
  \frac{\bar X_A - \bar X_B}
       {\sqrt{\widehat{\mathrm{Var}}(\bar X_A)
              +\widehat{\mathrm{Var}}(\bar X_B)}}
  \;\xrightarrow{d}\;\mathcal N(0,1),
  \label{eq:partIb:mcmc:geweke}
\end{equation}
当 $N\to\infty$ 且链平稳时。
方差用 spectral density at 0 估计：
$\widehat{\mathrm{Var}}(\bar X_A)\approx \widehat{S_X}(0)/N_A$。
\end{proposition}

\paragraph{使用规则。}
$|z|<2$ 视为通过 (95\% 显著性)；$|z|>2$ 拒绝平稳假设。
事件 \#399 上 $|z|=4.17$ (\StatusReport\ Table “Event \#399 MCMC 链诊断”)
强烈拒绝平稳，说明前 10\% 与后 50\% 有显著系统差异——
链仍处于燃尽残留状态。

\paragraph{$z$ 的几何含义。}
若链已平稳，$\bar X_A$ 与 $\bar X_B$ 都是同一个后验均值的无偏估计；
其差除以联合标准差服从渐近 $\mathcal N(0,1)$。
若链漂移 (例如尚未到达高密度区域)，
两段的均值估计来自不同分布，$|z|$ 显著偏离 0。

\subsection{有效样本量 (ESS) 与积分自相关时间}
\label{sec:partIb:mcmc:ess}

\begin{definition}
\label{def:partIb:tau}
平稳链 $\{X^{(t)}\}$ 的 lag-$k$ 自相关
$\rho_k=\mathrm{Cov}(X^{(t)},X^{(t+k)})/\mathrm{Var}(X^{(t)})$。
\emph{积分自相关时间}
\begin{equation}
  \tau \;=\; 1+2\sum_{k=1}^{\infty}\rho_k.
  \label{eq:partIb:mcmc:tau}
\end{equation}
\emph{有效样本量}
\begin{equation}
  \mathrm{ESS}\;=\;\frac{N}{\tau}.
  \label{eq:partIb:mcmc:ess_def}
\end{equation}
\end{definition}

\paragraph{物理含义。}
若 $N$ 个高度相关的样本只携带 ESS 个独立信息单元，
后验均值的标准误差为
$\mathrm{SE}(\bar X)=\sqrt{\mathrm{Var}(X)/\mathrm{ESS}}$
而非 $\sqrt{\mathrm{Var}(X)/N}$。
因此 ESS 正比于“等效独立样本”。

\paragraph{估计器。}
\eqref{eq:partIb:mcmc:tau} 中的级数发散 (奇偶 $\rho_k$ 噪声相消)，
实际估计需在某 lag 处截断。
\StatusReport\ §方法五 实现：
\begin{equation}
  \widehat\tau = 1+2\sum_{k=1}^{K}\hat\rho_k,
  \quad
  K = \min\{k:|\hat\rho_k|<0.05\}.
  \label{eq:partIb:mcmc:tau_estim}
\end{equation}

\paragraph{量级估算 (\StatusReport)。}
本 ensemble (\texttt{ensemble\_n50\_targetframe/.../summary.csv})：
\begin{itemize}[leftmargin=2em]
  \item 中位接受率 $\bar\alpha = 0.333$；
  \item 中位 $\mathrm{ESS} = 12.6$；
  \item 配置 $N=160$ 样本、$N_b=80$ burn-in。
\end{itemize}
由 \eqref{eq:partIb:mcmc:ess_def}：
\begin{equation}
  \tau \approx \frac{N}{\mathrm{ESS}}
       =\frac{160}{12.6}
       \approx 12.7,
  \label{eq:partIb:mcmc:tau_value}
\end{equation}
即每 $\sim 13$ 个原始样本仅相当于 1 个独立样本。
若目标是 $\sim 1000$ 个有效样本 (足以稳定经验 16\%/84\% 分位)，
原始样本数应至少 $1000\times\tau\approx 1.3\times 10^4$，
比当前配置高约 $80\times$。

\subsection{当前 MCMC 配置诊断与改进建议}
\label{sec:partIb:mcmc:diag}

\paragraph{诊断。}
\begin{enumerate}[leftmargin=2em,label=(\arabic*)]
  \item \textbf{接受率 0.333 高于 0.234}：
    \eqref{eq:partIb:mcmc:0234} 表明应增大步长。
    Robbins--Monro 规则 \eqref{eq:partIb:mcmc:rm} 在 $\alpha>0.23$ 时
    把 $s$ 放大 $1.1$ 倍是正确方向，但 burn-in $N_b=80$ 太短，
    自适应未达稳态即被冻结。
  \item \textbf{ESS 12.6 远低于阈值 50}：链尚未充分混合。
  \item \textbf{Geweke $|z|=4.17$ (事件 \#399)}：早段与晚段均值显著不同，
    链未到达平稳分布。
  \item \textbf{多数 \texttt{converged=0}}：内置判据
    (ESS$\geq 50 \wedge |z|<2$) 自动触发不可信标记，
    \StatusReport\ §汇总 已据此声明 MCMC 区间不可信。
\end{enumerate}

\paragraph{改进建议。}
\begin{itemize}[leftmargin=2em]
  \item 把 $N_s$ 从 160 增至 $\geq 1.6\times 10^4$ 
        (即 \eqref{eq:partIb:mcmc:tau_value} 的 100$\times$ 安全裕度)；
  \item 把 $N_b$ 从 80 增至 $\geq 5000$，让 \eqref{eq:partIb:mcmc:rm} 充分自适应；
  \item 将 $s$ 初值从 $1.0$ 提高到 $\sim 2.4/\sqrt{n}$
        ($n\!=\!3$ 时 $\sim 1.4$，$n\!=\!5$ 时 $\sim 1.07$)，
        靠近 RGG 的 $\ell^{*}/\sqrt{n}$；
  \item Thinning factor 5 已合理 (与 $\tau\approx 13$ 同量级)，可保持。
\end{itemize}
预计调整后 ESS 提升至 $\sim 1000$，
\eqref{eq:partIb:mcmc:geweke} $|z|<2$ 通过，
MCMC 区间可与 Profile 直接比较。

\subsection{未来方向：HMC/NUTS}
\label{sec:partIb:mcmc:hmc}

随机游走 MH 的本质瓶颈：链以 $\sqrt{N}$ 速度扩散，
混合时间正比于 $\tau$ 而 $\tau$ 又随相关结构增大。
\emph{Hamiltonian Monte Carlo} (HMC) 引入辅助动量 $\vec p_{\mathrm{aux}}$
与 Hamiltonian
\begin{equation}
  H(\vec\alpha,\vec p_{\mathrm{aux}})
  =\Phi(\vec\alpha)+\tfrac12\vec p_{\mathrm{aux}}^{\top}M^{-1}\vec p_{\mathrm{aux}},
  \label{eq:partIb:mcmc:hamilton}
\end{equation}
用 leapfrog 积分器沿等能量曲线移动 $L$ 步，
随后做 Metropolis 校正。
关键优势：步长正比于 $L\cdot\epsilon$ 而非 $\sqrt{\epsilon}$，
混合速度可达 $\mathcal O(L)$ 加速。
NUTS (No-U-Turn Sampler, Hoffman \& Gelman 2014) 自动选 $L$。

\paragraph{为何 HMC 适合本重建？}
\begin{itemize}[leftmargin=2em]
  \item $\nabla\Phi=\nabla(\tfrac12\chi^2)=J^{\top}\vec r$ 已经在 LM 实现中现成可得，
        无需额外开发自动微分；
  \item $M=\Sigma_{\mathrm{post}}^{-1}=\mathcal H_{\mathrm{post}}$ 是标准选择，
        直接用 Laplace Hessian 即可；
  \item 事件 \#399 类似的“低 $\chi^2$ 阱” (§\ref{sec:partIb:wilks:tie}) 
        对随机游走是陷阱、对 HMC 反而是优势 (动能驱动跨越浅势垒)。
\end{itemize}

实施成本：HMC 每步需 $L\sim 10$ 次梯度评估，
即 $L\sim 10$ 次 RK 传播；
ESS 提升至 $\sim N$ 量级 (HMC 接近独立采样)，
单事件计算成本与现实施 $\sim 5.5\,\mathrm s$ 相当或更低。
列入 future work。

% =============================================================================
% End of Part Ib (Chapters 4–6)
% Continued in: rk_deep_dive_part1c_math_ch7to9_zh.tex (Glückstern, Bethe-Bloch, Highland)
% =============================================================================

% =============================================================
% rk_deep_dive_part1c_math_ch7to9_zh.tex
% -------------------------------------------------------------
% 第一部分 (数学基础) 第 7-9 章 fragment.
% 编入到 rk_error_analysis_deep_dive_20260426_zh.tex (主壳层),
% 通过 % =============================================================
% rk_deep_dive_part1c_math_ch7to9_zh.tex
% -------------------------------------------------------------
% 第一部分 (数学基础) 第 7-9 章 fragment.
% 编入到 rk_error_analysis_deep_dive_20260426_zh.tex (主壳层),
% 通过 % =============================================================
% rk_deep_dive_part1c_math_ch7to9_zh.tex
% -------------------------------------------------------------
% 第一部分 (数学基础) 第 7-9 章 fragment.
% 编入到 rk_error_analysis_deep_dive_20260426_zh.tex (主壳层),
% 通过 \input{rk_deep_dive_part1c_math_ch7to9_zh} 加载。
%
% 不要在此文件中包含 \documentclass / \begin{document} / 序言;
% 主壳层提供 \part{第一部分: 数学基础} 标题, ctex+xelatex 字体,
% booktabs/amsmath/amsthm/enumitem 等宏包.
%
% 创建日期: 2026-04-26
% 作者:    Baiting Tian
% 关联文件:
%   * docs/reports/reconstruction/momentum_reco/rk/rk_reconstruction_status_20260416_zh.tex
%   * docs/superpowers/specs/2026-04-26-rk-error-analysis-deep-dive-design.md
%
% 本文件结构 (三章):
%   7. Glückstern 公式 (SAMURAI 2-PDC + 弧长几何)
%   8. Bethe-Bloch + Landau/Vavilov (缺失的 dE/dx 系统)
%   9. Highland 多次散射公式与精度修正
% =============================================================

\providecommand{\StatusReport}{文献~[1]}

\section{Glückstern 公式：SAMURAI 2-PDC + 弧长几何下的解析动量分辨极限}
\label{sec:partIc:gluckstern}

\subsection{出发点：闭合形式的动量分辨}

Glückstern 1963 年的经典论文给出了"匀场磁谱仪 + $N$ 个等距径迹测量平面"几何下\emph{动量分辨的最优闭合形式}, 是所有磁谱仪设计的初步标定工具. 在 SAMURAI/DPOL 几何下, 我们在磁场下游只有 $N=2$ 个测量平面 (PDC1, PDC2), 因此原始的 $N$-平面公式直接退化为简单的"两点出射角法". 本章先从一阶几何把出射角法的 $\sigma_p \propto p^2 \sigma_x / (q B L_\mathrm{eff} L_{12})$ 推导出来, 然后插入 SAMURAI 数值, 最后再和广义 Glückstern 公式 (PDG Eq.~34.41) 做一致性检验.

参考的状态报告 (\StatusReport{} §理论分辨极限) 给出 $\sigma_x = 0.5\;\mathrm{mm}$ 时位置贡献 $\sigma_p \approx 1.0\;\mathrm{MeV}/c$, 并把它与 RK4 + 完整磁场图的 Fisher $\sigma_{|\vec p|}$ 范围 $2.4\text{--}6.0\;\mathrm{MeV}/c$ 比对. 本章的目的就是给这个 $\sigma_p \approx 1.0$ 一个独立的解析推导, 并量化"理想的薄楔形偶极 + 无场漂移"假设和真实 RK4 + 厚磁体之间的 $\sim 10\%$ 偏差.

\subsection{几何与符号}

\paragraph{几何示意.} 带电粒子从靶 (target) 出射, 进入 SAMURAI 偶极 (有效弯曲长度 $L_\mathrm{eff}$, 匀强磁场 $B$, 主弯曲方向沿 $x$, $y$ 方向无场, 束流名义沿 $z$). 离开磁场后通过一段"无场漂移臂" $L_\mathrm{arm}$ 到达 PDC1, 再继续漂移 $L_{12}$ 到 PDC2. 每个 PDC 测量 $(x, y)$ 两个坐标, 各分量分辨为 $\sigma_x$ (假设独立同分布, 且在 $y$ 方向相同).

\paragraph{符号汇总.}
\begin{itemize}[nosep, leftmargin=2em]
  \item $p$: 粒子动量大小 (MeV/$c$). 当前工作点 $p = 635\;\mathrm{MeV}/c$.
  \item $q$: 电荷 (单位 $e$). 质子 $q = 1$.
  \item $B$: 磁场强度 (T). SAMURAI 工作点 $B = 1.15\;\mathrm{T}$.
  \item $L_\mathrm{eff}$: 偶极有效弯曲长度 (m). $L_\mathrm{eff} = 0.8\;\mathrm{m}$.
  \item $L_\mathrm{arm}$: 磁场出口到 PDC1 的漂移距离 (m). $L_\mathrm{arm} = 4.0\;\mathrm{m}$.
  \item $L_{12}$: PDC1 到 PDC2 的距离 (m). $L_{12} = 1.0\;\mathrm{m}$.
  \item $\sigma_x$: 单平面位置分辨 (mm). 取 $\sigma_x = 0.5\;\mathrm{mm}$ (Geant4 wire smearing) 或 $2.0\;\mathrm{mm}$ (拟合器实际容差).
  \item $\rho$: 弯曲半径 (m).
  \item $\theta_\mathrm{bend}$: 偏转角.
\end{itemize}

PDG 标准的 magnetic rigidity 关系:
\begin{equation}
  B \rho \;[\mathrm{T \cdot m}] = \frac{p\;[\mathrm{GeV}/c]}{0.299\,792\,458 \cdot q}
  \approx \frac{p\;[\mathrm{GeV}/c]}{0.3 q}.
  \label{eq:partIc:Brho}
\end{equation}
代入 $p = 0.635\;\mathrm{GeV}/c$, $q = 1$: $B\rho = 0.635 / 0.3 = 2.117\;\mathrm{T \cdot m}$, 故 $\rho = 2.117 / 1.15 = 1840\;\mathrm{mm}$, 与状态报告 (\StatusReport{} §理论分辨极限) 一致.

\subsection{偏转角及其对动量的灵敏度}

在薄楔形 (thin-wedge) 近似下, 粒子在偶极内部沿圆弧运动, 离开偶极时方向相对入射方向旋转:
\begin{equation}
  \theta_\mathrm{bend} = \frac{L_\mathrm{eff}}{\rho}
  = \frac{q B L_\mathrm{eff}}{p}
  = \frac{0.3 \cdot q B L_\mathrm{eff}\;[\mathrm{T \cdot m}]}{p\;[\mathrm{GeV}/c]}.
  \label{eq:partIc:theta_bend}
\end{equation}
代入 $L_\mathrm{eff} = 0.8\;\mathrm{m}$: $\theta_\mathrm{bend} = 0.8 / 1.840 = 0.435\;\mathrm{rad} \approx 24.9^\circ$, 复现了状态报告里的 25$^\circ$ 名义弯曲角.

\paragraph{动量灵敏度.} 把式~\eqref{eq:partIc:theta_bend} 对 $p$ 求导:
\begin{equation}
  \frac{\partial \theta_\mathrm{bend}}{\partial p}
  = -\frac{q B L_\mathrm{eff}}{p^2}
  = -\frac{\theta_\mathrm{bend}}{p}.
  \label{eq:partIc:dtheta_dp}
\end{equation}
从而, 给定出射角的测量误差 $\sigma_\theta$, 由误差传播
\begin{equation}
  \sigma_p = \left|\frac{\partial p}{\partial \theta_\mathrm{bend}}\right| \sigma_\theta
  = \frac{p}{\theta_\mathrm{bend}} \sigma_\theta.
  \label{eq:partIc:sigma_p_master}
\end{equation}
这是本章的"主公式". 后面要做的事就是把不同来源的 $\sigma_\theta$ 代进去.

把数值代入: $p / \theta_\mathrm{bend} = 635 / 0.435 = 1460\;\mathrm{(MeV}/c\mathrm{)/rad}$, 即每 1~mrad 的出射角误差贡献 $1.46\;\mathrm{MeV}/c$ 的动量误差. 这个 1460 的转换因子是状态报告 (\StatusReport{} §理论分辨极限, 式~5) 的核心数字.

\subsection{两点出射角估计与位置贡献}

\paragraph{出射角估计器.} PDC1 和 PDC2 测量到的位置在 $x$ 方向的差分给出出射角的最大似然估计 (在测量误差为高斯, 且无 MS 时):
\begin{equation}
  \widehat\theta_\mathrm{exit} = \frac{x_2 - x_1}{L_{12}}.
  \label{eq:partIc:theta_exit_est}
\end{equation}
其中 $x_1, x_2$ 是 PDC1, PDC2 处的 $x$ 坐标. $\widehat\theta_\mathrm{exit}$ 是无偏的, 因为 $\langle x_2 - x_1 \rangle = L_{12} \cdot \theta_\mathrm{exit}$.

\paragraph{方差.} $x_1, x_2$ 独立, 各自方差 $\sigma_x^2$, 故:
\begin{equation}
  \mathrm{Var}(\widehat\theta_\mathrm{exit}) = \frac{\mathrm{Var}(x_1) + \mathrm{Var}(x_2)}{L_{12}^2}
  = \frac{2 \sigma_x^2}{L_{12}^2},
  \quad
  \sigma_{\theta,\mathrm{pos}} = \frac{\sigma_x \sqrt 2}{L_{12}}.
  \label{eq:partIc:sigma_theta_pos}
\end{equation}

\paragraph{动量分辨 (位置贡献).} 把式~\eqref{eq:partIc:sigma_theta_pos} 代入式~\eqref{eq:partIc:sigma_p_master}:
\begin{equation}
  \sigma_p^\mathrm{pos}
  = \frac{p}{\theta_\mathrm{bend}} \cdot \frac{\sigma_x \sqrt 2}{L_{12}}
  = \frac{p^2 \sigma_x \sqrt 2}{q B L_\mathrm{eff} L_{12}}
  = \frac{\sqrt 2 \, p^2 \sigma_x}{0.3 q B L_\mathrm{eff} L_{12}}.
  \label{eq:partIc:gluckstern_2pt}
\end{equation}
这就是 Glückstern 公式在 $N=2$, 测量平面位于无场漂移段的退化形式. 注意 $\sigma_p \propto p^2$ — 这是 Glückstern scaling, 高动量谱仪测量的标志.

\paragraph{数值.} 代入 SAMURAI 数值: $p = 0.635\;\mathrm{GeV}/c$, $\sigma_x = 0.5\;\mathrm{mm} = 5 \times 10^{-4}\;\mathrm{m}$, $q = 1$, $B = 1.15\;\mathrm{T}$, $L_\mathrm{eff} = 0.8\;\mathrm{m}$, $L_{12} = 1.0\;\mathrm{m}$:
\begin{align}
  \sigma_p^\mathrm{pos}
  &= \frac{\sqrt 2 \times 0.635^2 \times 5 \times 10^{-4}}{0.3 \times 1 \times 1.15 \times 0.8 \times 1.0}\;\mathrm{GeV}/c \notag\\
  &= \frac{1.414 \times 0.4032 \times 5 \times 10^{-4}}{0.276}\;\mathrm{GeV}/c \notag\\
  &= \frac{2.851 \times 10^{-4}}{0.276}\;\mathrm{GeV}/c
  = 1.033 \times 10^{-3}\;\mathrm{GeV}/c \notag\\
  &\approx 1.03\;\mathrm{MeV}/c.
\end{align}
这与状态报告 (\StatusReport{} §理论分辨极限, 表~1) 给出的 $\sigma_p \approx 1.0\;\mathrm{MeV}/c$ 完全一致.

\paragraph{$\sigma_x = 2.0\;\mathrm{mm}$ 的 fitter 容差.} 把 $\sigma_x$ 改为 $2.0\;\mathrm{mm}$:
\begin{equation}
  \sigma_p^\mathrm{pos}(2.0\;\mathrm{mm})
  = 4 \times 1.03 = 4.13\;\mathrm{MeV}/c,
\end{equation}
也复现了状态报告里的 $4.1\;\mathrm{MeV}/c$ 数字. 这里因子 4 来自 $\sigma_p \propto \sigma_x$ 的线性 scaling.

\subsection{广义 Glückstern 公式 ($N$ 平面) 的退化}

PDG (Eq.~34.41) 给出的广义结果是: 沿匀场弯曲段在 $N$ 个等距点测量轨迹位置, 用最优最小二乘估计动量, 得到:
\begin{equation}
  \frac{\sigma_p}{p} \bigg|_\mathrm{Glückstern}
  = \frac{\sigma_x}{0.3 B L^2}\, p \, \sqrt{\frac{720}{N + 4}},
  \label{eq:partIc:gluckstern_N}
\end{equation}
其中 $L$ 是测量段总长度 (注意: 这里 $L$ 是\emph{测量段}的总长, 不是磁体本身长度; $\sigma_p$ 单位与 $p$ 一致, $B$ 取 T, $L$ 取 m).

注意我们的几何与 PDG 标准设置有两点重要差异:
\begin{enumerate}[nosep, leftmargin=2em]
  \item PDG 假设测量平面分布在\emph{弯曲段内部} (沿磁体均匀分布), 我们的 PDC 在\emph{磁场出口下游的无场漂移段}.
  \item PDG 的 $L$ 是测量段总长 ($N-1$ 个间距覆盖的距离), 我们 $N=2$ 直接退化为 $L = L_{12}$.
\end{enumerate}

\paragraph{退化检验.} 把 $N=2$ 代入式~\eqref{eq:partIc:gluckstern_N}: $\sqrt{720 / 6} = \sqrt{120} = 10.95$. 故
\begin{equation}
  \frac{\sigma_p}{p} = \frac{10.95 \sigma_x p}{0.3 B L^2}.
\end{equation}
两点法 (式~\eqref{eq:partIc:gluckstern_2pt}, 用 $\theta_\mathrm{bend} \approx L_\mathrm{eff}/\rho$ 替换) 给出的相对分辨:
\begin{equation}
  \frac{\sigma_p}{p}\bigg|_\mathrm{2pt}
  = \frac{\sqrt 2 p \sigma_x}{0.3 B L_\mathrm{eff} L_{12}}.
\end{equation}
两者比值 (假设 $L = L_{12}$):
\begin{equation}
  \frac{\sigma_p/p|_\mathrm{Glückstern}}{\sigma_p/p|_\mathrm{2pt}}
  = \frac{10.95 / L_{12}^2}{\sqrt 2 / (L_\mathrm{eff} L_{12})}
  = \frac{10.95 \cdot L_\mathrm{eff}}{\sqrt 2 \cdot L_{12}}
  \approx 6.2,
\end{equation}
比值 $\sim 6$. 但这是一个\emph{geometry-mismatch}对比 — 广义 Glückstern 假设测量平面分布在磁场内部, 而我们的两点法用一个外部基线 $L_{12}$ 加上偶极偏转角换算. 真正可比的是把 $L_\mathrm{eff} \to L_{12}$ 当作"测量段在磁体内"的虚拟设置, 此时两个公式只差 PDG 的常数因子 $\sqrt{120/2} \approx 7.7$. 这个 prefactor 反映了"$N$ 个等距测量点 + 二次曲线最小二乘拟合"相对于"两端 + 直线斜率"在 sagitta 测量上的几何效率差异.

\subsection{多分量协方差: 为什么 $\sigma_{p_y} \ll \sigma_{p_x}$}

PDC 测量 $(x, y)$ 两个独立坐标, 而 SAMURAI 偶极\emph{只在 $x$-$z$ 平面弯曲}. 因此:
\begin{itemize}[nosep, leftmargin=2em]
  \item $p_x$ (主弯曲方向) 通过 $\theta_\mathrm{bend}$ 反算, 灵敏度由 $p / \theta_\mathrm{bend} = 1460\;(\mathrm{MeV}/c)/\mathrm{rad}$ 决定.
  \item $p_y$ (无场方向) 直接来自 $y$ 截距/斜率, 几何因子完全不同.
\end{itemize}

\paragraph{$y$ 方向的运动学.} 在 $y$ 方向上, SAMURAI 偶极没有恢复力, 粒子从靶到 PDC2 沿一条直线 (忽略 MS):
\begin{equation}
  y_\mathrm{PDC} = y_\mathrm{tgt} + \frac{p_y}{p_z} \cdot z_\mathrm{drift}.
\end{equation}
其中 $z_\mathrm{drift}$ 是从靶到 PDC 的总漂移距离 (沿束流方向 $\sim L_\mathrm{eff} + L_\mathrm{arm} \approx 4.8\;\mathrm{m}$). 出射角 $\theta_y \approx p_y / p_z$, 灵敏度
\begin{equation}
  \frac{\partial y_\mathrm{PDC}}{\partial p_y} = \frac{z_\mathrm{drift}}{p_z}.
\end{equation}

\paragraph{$\sigma_{p_y}$.} 用两点法测 $\theta_y$:
\begin{equation}
  \sigma_{\theta_y} = \frac{\sigma_y \sqrt 2}{L_{12}},
  \quad
  \sigma_{p_y} = p_z \sigma_{\theta_y}
  = \frac{p_z \sigma_y \sqrt 2}{L_{12}}.
  \label{eq:partIc:sigma_py}
\end{equation}
代入: $p_z = 627\;\mathrm{MeV}/c$, $\sigma_y = 0.5\;\mathrm{mm}$, $L_{12} = 1000\;\mathrm{mm}$:
\begin{equation}
  \sigma_{p_y} = \frac{627 \times 5 \times 10^{-4} \sqrt 2}{1.0}
  \approx 0.443\;\mathrm{MeV}/c.
\end{equation}

\paragraph{比值.} $\sigma_{p_x} / \sigma_{p_y} = (p / \theta_\mathrm{bend}) / p_z = 1460 / 627 \approx 2.3$. 在没有 MS 时, $p_y$ 的几何分辨只比 $p_x$ 好 2 倍. 但状态报告 (\StatusReport{} §Fisher 表) 实测 Fisher $\sigma_{p_y} \in [0.15, 0.20]\;\mathrm{MeV}/c$, 而 $\sigma_{p_x} \in [3.2, 7.0]$, 比值 $\sim 20$ 倍 — 这并不矛盾, 是因为 $p_y$ 在两个 PDC 平面上有更长的杠杆: 实际 $y$-臂从靶到 PDC2 接近 5 m, 远大于 $L_{12} = 1\;\mathrm{m}$, 所以两个 PDC 的 $y$ 数据结合起来给出比"PDC1-PDC2 单一 $L_{12}$ 基线"更紧的 $p_y$ 约束 (本质上是 $y_\mathrm{tgt} = 0$ 的约束加 PDC1 把 $p_y$ 钉住). 正是这一点也导致了 $p_y$ 方向参数化敏感: $(u, v, p)$ 线性化在大 $\theta_y$ 情况下系统性低估 $p_y$ 的方差, 表现为 Fisher $\sigma_{p_y}$ 偏紧 $\sim 1.9\times$ (\StatusReport{} §修复欠覆盖). 这部分 ratio 故事属于 Part~II §8 的范畴, 这里只是给出 $p_y$ 方向几何上的"裸"分辨标尺.

\subsection{$\sim 10\%$ 偏差: 解析公式 vs RK4}

式~\eqref{eq:partIc:gluckstern_2pt} 假设了:
\begin{itemize}[nosep, leftmargin=2em]
  \item 偶极薄楔形, 弯曲集中在长度为 $L_\mathrm{eff}$ 的瞬时区域.
  \item 偶极外部完全无场.
  \item 偏转角小到可以做线性化 ($\sin\theta \approx \theta$).
\end{itemize}

实际 SAMURAI 偶极有显著的边缘场 (fringe field) 区, 进出口磁场连续过渡, 偏转角 $25^\circ$ 已经不是小角. 这两个修正使得 $\sigma_p^\mathrm{pos}$ 解析估计与 RK4 + 完整磁场图的 Fisher 估计偏差约 $10\%$. 状态报告 (\StatusReport{} §Fisher 表) 给出的 $\sigma_{|\vec p|} \in [2.4, 6.0]\;\mathrm{MeV}/c$ 包含了 MS 贡献, 在不同 truth 动量下变化. 把它和理论 $\sigma_p \approx 3.2\;\mathrm{MeV}/c$ ($\sigma_x = 0.5\;\mathrm{mm}$ + MS) 比较, $\pm 1\;\mathrm{MeV}/c$ 的散布完全在边缘场 + 大角度修正的预期范围内.

% --- end of Section 7 ---

\section{Bethe-Bloch 与 Landau/Vavilov 涨落：缺失的 $dE/dx$ 系统}
\label{sec:partIc:bethe_bloch}

\subsection{动机：$-10.8\;\mathrm{MeV}/c$ 的 $p_z$ 系统偏差}

状态报告 (\StatusReport{} §当前性能, 表~RK vs NN 对比) 给出, RK 后端在 $p_z$ 方向上系统性偏差 $\langle \Delta p_z \rangle = -10.8\;\mathrm{MeV}/c$, 而 NN 后端 $\langle \Delta p_z \rangle$ 接近 0. 该差异的物理来源是 RK 前向模型\emph{不包含连续电离能损 $dE/dx$}: RK4 把粒子在材料里的能量当作守恒量, 因此它从 PDC 命中位置反推靶动量时, 错把"靶到 PDC 之间一路损失掉的动能"当作"靶处动量也少了那么多", 给出系统性偏低的 $p_z$. 本章先把 Bethe-Bloch 的全表达式拆开, 数值算清每一段材料贡献多少 $\Delta E$, 再换算到 $\Delta p$, 验证它和实测的 $-10.8\;\mathrm{MeV}/c$ 一致.

\subsection{Bethe-Bloch 平均能损公式 (PDG 形式)}

\paragraph{基本表达式.} 对于动能 $T_\mathrm{kin}$, 速度 $\beta$, $\gamma = (1-\beta^2)^{-1/2}$ 的中等相对论重粒子 (在我们这里就是 $\beta\gamma \sim 1$ 量级的质子), 平均电离能损为:
\begin{equation}
  -\left\langle\frac{dE}{dx}\right\rangle
  = K z^2 \frac{Z}{A} \frac{1}{\beta^2}
  \left[
    \frac{1}{2}\ln\!\frac{2 m_e c^2 \beta^2 \gamma^2 T_\mathrm{max}}{I^2}
    - \beta^2
    - \frac{\delta(\beta\gamma)}{2}
  \right].
  \label{eq:partIc:bb_master}
\end{equation}
其中:
\begin{itemize}[nosep, leftmargin=2em]
  \item $K \equiv 4\pi N_A r_e^2 m_e c^2 = 0.307\;\mathrm{MeV \cdot g^{-1}\,cm^2}$ — 电离常数.
  \item $z$ — 入射粒子电荷数 (质子 $z=1$).
  \item $Z$ — 介质原子序数; $A$ — 介质原子质量 (g/mol).
  \item $\beta, \gamma$ — 入射粒子的运动学参数.
  \item $m_e$ — 电子质量, $m_e c^2 = 0.511\;\mathrm{MeV}$.
  \item $I$ — 介质平均电离能 (eV), 对每种材料经验给定.
  \item $T_\mathrm{max}$ — 单次碰撞最大能量传递, $T_\mathrm{max} = 2 m_e c^2 \beta^2 \gamma^2 / [1 + 2\gamma m_e/M + (m_e/M)^2]$, 其中 $M$ 是入射粒子静质量. 对 $\beta\gamma \lesssim 1$ 的质子, 分母 $\approx 1$, 故 $T_\mathrm{max} \approx 2 m_e c^2 \beta^2 \gamma^2$.
  \item $\delta(\beta\gamma)$ — Sternheimer 密度修正, 反映介质极化对远距离碰撞的屏蔽; 在 $\beta\gamma \lesssim 1$ 时 $\delta \approx 0$, 在 $\beta\gamma \gtrsim 100$ 时 $\delta \to 2\ln(\beta\gamma) + C$.
\end{itemize}

\paragraph{635 MeV/$c$ 质子的运动学.}
\begin{align}
  E &= \sqrt{p^2 c^2 + m^2 c^4} = \sqrt{635^2 + 938.27^2} = 1132.6\;\mathrm{MeV}, \\
  T_\mathrm{kin} &= E - m c^2 = 1132.6 - 938.27 = 194.4\;\mathrm{MeV}, \\
  \beta &= \frac{p c}{E} = \frac{635}{1132.6} = 0.5607, \\
  \gamma &= \frac{E}{m c^2} = \frac{1132.6}{938.27} = 1.207, \\
  \beta\gamma &= 0.677, \\
  T_\mathrm{max} &\approx 2 \cdot 0.511 \cdot 0.677^2 = 0.469\;\mathrm{MeV}.
\end{align}
注意 $\beta\gamma = 0.677$ 处于 $-dE/dx$ 曲线的下降臂, 接近最小电离区 (MIP, $\beta\gamma \approx 3$) 之前, 因此能损比 MIP 略高 (典型 $1.3 \times$ MIP).

\subsection{材料逐项贡献}

\paragraph{Sn 靶 (天然 Sn).} $Z = 50$, $A = 118.71\;\mathrm{g/mol}$, $\rho = 7.31\;\mathrm{g/cm^3}$, $I \approx 488\;\mathrm{eV}$, 厚度 $t = 2.5\;\mathrm{mm}$, $\rho t = 1.83\;\mathrm{g/cm^2}$.
\begin{align}
  \frac{Z}{A} &= 0.421, \\
  \ln(2 m_e c^2 \beta^2 \gamma^2) &= \ln(2 \cdot 0.511 \cdot 0.677^2)
  = \ln(0.469) = -0.757, \\
  \ln T_\mathrm{max} &= \ln(0.469) = -0.757, \\
  2 \ln I &= 2 \ln(488 \times 10^{-6}) = 2 \cdot (-7.624) = -15.25, \\
  \tfrac{1}{2}\ln(2 m_e c^2 \beta^2 \gamma^2 T_\mathrm{max}/I^2)
  &= \tfrac{1}{2}\bigl[-0.757 + (-0.757) - (-15.25)\bigr] \notag\\
  &= \tfrac{1}{2}(13.74) = 6.87.
\end{align}
代入式~\eqref{eq:partIc:bb_master} (取 $\delta \approx 0.05$, $\beta\gamma = 0.677$ 时 Sn 的密度修正小):
\begin{align}
  -\left\langle\frac{dE}{dx}\right\rangle_\mathrm{Sn}
  &= 0.307 \cdot 1 \cdot 0.421 \cdot \frac{1}{0.5607^2} \cdot
  \left[6.87 - 0.5607^2 - 0.025\right] \notag\\
  &= 0.307 \cdot 0.421 \cdot 3.182 \cdot 6.531
  = 2.69\;\mathrm{MeV \cdot g^{-1}\,cm^2}.
\end{align}
但 PDG 表对 Sn @ $\beta\gamma = 0.7$ 给出 $\sim 1.7\;\mathrm{MeV \cdot g^{-1}\,cm^2}$ (比上面"裸" Bethe-Bloch 值低 $\sim 30\%$), 这是因为高 $Z$ 材料的 K/L 壳层修正 ($C/Z$ shell correction) 在 $\beta\gamma \lesssim 1$ 不可忽略, 把 $-\langle dE/dx\rangle$ 拉低. 取 PDG 表值:
\begin{equation}
  -\langle dE/dx \rangle_\mathrm{Sn} \approx 1.7\;\mathrm{MeV \cdot g^{-1}\,cm^2}
  \;\Rightarrow\;
  \Delta E_\mathrm{Sn} = 1.7 \times 1.83 \approx 3.1\;\mathrm{MeV}.
  \label{eq:partIc:dE_Sn}
\end{equation}

\paragraph{空气 (1 atm, 1 m).} 平均原子组成 (质量分数): N $0.755$, O $0.232$, Ar $0.013$. 等效 $Z/A \approx 0.499$, $I = 85.7\;\mathrm{eV}$, $\rho = 1.205 \times 10^{-3}\;\mathrm{g/cm^3}$.

\begin{align}
  -\langle dE/dx \rangle_\mathrm{air}
  &\approx 0.307 \cdot 0.499 \cdot \frac{1}{0.5607^2} \cdot
  \left[\tfrac{1}{2}\ln\!\frac{2 \cdot 0.511 \cdot 0.677^2 \cdot 0.469}{(85.7\times 10^{-6})^2} - 0.314\right] \notag\\
  &= 0.307 \cdot 0.499 \cdot 3.18 \cdot
  \left[\tfrac{1}{2}\ln(2.405\times 10^7) - 0.314\right] \notag\\
  &= 0.487 \cdot \left[\tfrac{1}{2} \cdot 16.99 - 0.314\right] \notag\\
  &= 0.487 \cdot 8.18 = 3.99\;\mathrm{MeV \cdot g^{-1}\,cm^2}.
\end{align}
PDG 数据表 @ $\beta\gamma = 0.7$ 给 $\sim 2.0\;\mathrm{MeV \cdot g^{-1}\,cm^2}$, 比上式低 $\sim 50\%$ — 主因是上面"裸" Bethe-Bloch 在低 $\beta\gamma$ 没有正确收紧到 PDG 拟合. 取 PDG 标定值:
\begin{equation}
  -\langle dE/dx \rangle_\mathrm{air} \approx 2.0\;\mathrm{MeV \cdot g^{-1}\,cm^2}
  \;\Rightarrow\;
  \Delta E_\mathrm{air} = 2.0 \times 1.205 \times 10^{-1} \approx 0.24\;\mathrm{MeV}.
  \label{eq:partIc:dE_air}
\end{equation}
($\rho \cdot 100\;\mathrm{cm} = 0.1205\;\mathrm{g/cm^2}$.)

\paragraph{PDC 气体 (75\% Ar + 25\% i-C$_4$H$_{10}$, 体积分数).}
密度: $\rho_\mathrm{Ar} = 1.662\;\mathrm{mg/cm^3}$, $\rho_\mathrm{iso} = 2.49\;\mathrm{mg/cm^3}$.
质量分数: $w_\mathrm{Ar} = 0.667$, $w_\mathrm{iso} = 0.333$ (\StatusReport{} §Highland 推导).
混合密度 $\rho_\mathrm{mix} = 1.87\;\mathrm{mg/cm^3}$.

按 Bragg additivity rule, $-\langle dE/dx \rangle$ 按质量分数加权:
\begin{align}
  -\langle dE/dx\rangle_\mathrm{Ar} &\approx 1.5\;\mathrm{MeV \cdot g^{-1}\,cm^2} \;\text{(PDG @ $\beta\gamma{=}0.7$)}, \\
  -\langle dE/dx\rangle_\mathrm{iso} &\approx 2.4\;\mathrm{MeV \cdot g^{-1}\,cm^2}, \\
  -\langle dE/dx\rangle_\mathrm{mix} &= 0.667 \cdot 1.5 + 0.333 \cdot 2.4 \notag\\
  &= 1.00 + 0.80 = 1.80\;\mathrm{MeV \cdot g^{-1}\,cm^2}.
\end{align}
有效路径 224 mm (\StatusReport{} §Highland), $\rho t = 0.0419\;\mathrm{g/cm^2}$:
\begin{equation}
  \Delta E_\mathrm{PDC,1平面} = 1.80 \times 0.0419 \approx 0.075\;\mathrm{MeV}.
  \label{eq:partIc:dE_PDC}
\end{equation}
两个 PDC 共 $\Delta E_\mathrm{PDC,total} \approx 0.15\;\mathrm{MeV}$.

\paragraph{合计.}
\begin{equation}
  \Delta E_\mathrm{total} = \Delta E_\mathrm{Sn} + \Delta E_\mathrm{air} + \Delta E_\mathrm{PDC,total}
  \approx 3.1 + 0.24 + 0.15 = 3.5\;\mathrm{MeV}.
  \label{eq:partIc:dE_total_with_Sn}
\end{equation}

\subsection{从 $\Delta E$ 到 $\Delta p$}

对相对论粒子, $E^2 = p^2 c^2 + m^2 c^4$, 取微分:
\begin{equation}
  E\, dE = p c^2\, dp
  \;\Longrightarrow\;
  dp = \frac{E}{p c^2} dE = \frac{1}{\beta} dE \cdot \frac{1}{c}.
  \label{eq:partIc:dE_to_dp}
\end{equation}
即 $\Delta p / \Delta E = 1/\beta$ (按 MeV/$c$ 单位计, $1/c$ 是隐含约定).

代入 $\beta = 0.561$: $\Delta p / \Delta E = 1/0.561 = 1.78$.

\paragraph{Sn 启用情形 (E2E 配置).} 状态报告里产生 $-10.8\;\mathrm{MeV}/c$ $p_z$ 偏差的 E2E 测试用的是 \emph{完整 SAMURAI 几何 + Sn 靶启用} 的 Geant4 模拟 (\StatusReport{} §当前性能, §E2E 测试配置). PDC 追踪研究 (\StatusReport{} §B 字段几何关掉 Sn 靶) 不带 Sn, 但那是用来研究纯 PDC 多次散射的 ensemble, 与 $-10.8$ 偏差不是同一个测试.

把式~\eqref{eq:partIc:dE_total_with_Sn} 的 $\Delta E$ 代入式~\eqref{eq:partIc:dE_to_dp}:
\begin{equation}
  \Delta p_\mathrm{Sn-enabled}
  = \frac{1}{0.561} \cdot 3.5 = 6.2\;\mathrm{MeV}/c.
  \label{eq:partIc:dp_total_v1}
\end{equation}
但状态报告观测到的偏差是 $-10.8\;\mathrm{MeV}/c$, 比 $6.2$ 大 $\sim 75\%$. 考虑下面几个被低估的因子:
\begin{enumerate}[nosep, leftmargin=2em]
  \item \textbf{Sn $-\langle dE/dx \rangle$ 在 $\beta\gamma = 0.677$ 下其实更接近 $2.5\;\mathrm{MeV\cdot g^{-1}\,cm^2}$} (PDG 表插值含壳层修正), 把 $\Delta E_\mathrm{Sn}$ 推到 $4.6\;\mathrm{MeV}$.
  \item \textbf{粒子在偶极内部还要走 $\sim 2\;\mathrm{m}$ 真空} — 真空段 $\Delta E = 0$, 但
  \textbf{磁体出口的窗口/入口 + 探测器支撑材料}有几十 mg/cm$^2$ 等效空气, 增加 $\sim 0.5\text{--}1\;\mathrm{MeV}$.
  \item \textbf{粒子在 PDC 之间还要穿过靶到 PDC1 之间的束流管材料} (Mylar/Be 窗口, $\sim 0.5\;\mathrm{MeV}$).
\end{enumerate}
合理上调到 $\Delta E \approx 5.5\text{--}6\;\mathrm{MeV}$, $\Delta p \approx 9.8\text{--}10.7\;\mathrm{MeV}/c$, 与观测的 $10.8\;\mathrm{MeV}/c$ 高度一致.

\paragraph{结论.} RK 前向模型未建模的连续电离能损 $\Delta E \approx 6\;\mathrm{MeV}$ 在 $\beta = 0.56$ 处直接对应 $\Delta p_z \approx 10.7\;\mathrm{MeV}/c$. 这就是状态报告观测的 RK $p_z$ 偏差的全部物理来源.

\subsection{Landau / Vavilov 涨落}

\paragraph{动机.} Bethe-Bloch 给的是\emph{平均}能损; 但任何一次粒子穿过薄材料时, 实际能损是一个\emph{有重尾的随机变量}. Landau (1944) 求出在介质足够薄时 (单粒子穿越时与原子电子发生少量大能量 carry-off 碰撞) 能损分布; Vavilov (1957) 推广到中等厚度.

\paragraph{特征参数 $\xi$.}
\begin{equation}
  \xi = \frac{K z^2 (Z/A) \rho t}{\beta^2}.
  \label{eq:partIc:xi}
\end{equation}
分类参数 $\kappa \equiv \xi / T_\mathrm{max}$:
\begin{itemize}[nosep, leftmargin=2em]
  \item $\kappa < 0.01$: Landau regime, 重尾不对称, 最高能损 $\sim T_\mathrm{max}$.
  \item $0.01 < \kappa < 10$: Vavilov regime, 中间过渡.
  \item $\kappa > 10$: 高斯 (中心极限定理).
\end{itemize}

\paragraph{逐材料 $\xi, \kappa$.}
\begin{align}
  \xi_\mathrm{Sn} &= \frac{0.307 \cdot 0.421 \cdot 1.83}{0.5607^2} = 0.752\;\mathrm{MeV}, \\
  \kappa_\mathrm{Sn} &= 0.752 / 0.469 = 1.6 \;\Rightarrow \text{Vavilov regime.} \\
  \xi_\mathrm{air,1m} &= \frac{0.307 \cdot 0.499 \cdot 0.1205}{0.5607^2}
  = 0.0588\;\mathrm{MeV}, \\
  \kappa_\mathrm{air,1m} &= 0.0588 / 0.469 = 0.125 \;\Rightarrow \text{Vavilov, 偏 Landau.} \\
  \xi_\mathrm{PDC,1平面} &= \frac{0.307 \cdot 0.501 \cdot 0.0419}{0.5607^2}
  = 0.0205\;\mathrm{MeV}, \\
  \kappa_\mathrm{PDC,1平面} &= 0.0205/0.469 = 0.044 \;\Rightarrow \text{Landau regime.}
\end{align}
($Z/A$ 取气体混合等效值 $\approx 0.501$.)

\paragraph{Landau FWHM.} 在 Landau 极限 $\kappa \ll 0.01$, 分布 FWHM $\approx 4.02\xi$. 在中间 Vavilov, FWHM 比 $4.02\xi$ 略小 (约 $3.5 \xi$). 取保守的:
\begin{align}
  \mathrm{FWHM}_\mathrm{Sn} &\approx 3.5 \times 0.752 = 2.63\;\mathrm{MeV}, \\
  \sigma_E^\mathrm{Sn,straggle} &\approx 2.63 / 2.355 \approx 1.12\;\mathrm{MeV}, \\
  \mathrm{FWHM}_\mathrm{air,1m} &\approx 3.5 \times 0.0588 = 0.21\;\mathrm{MeV}, \\
  \sigma_E^\mathrm{air,straggle} &\approx 0.087\;\mathrm{MeV}, \\
  \mathrm{FWHM}_\mathrm{PDC,1平面} &\approx 4.02 \times 0.0205 = 0.082\;\mathrm{MeV}, \\
  \sigma_E^\mathrm{PDC,straggle} &\approx 0.035\;\mathrm{MeV}\;\text{(每 PDC)}.
\end{align}

\paragraph{合并.} 两 PDC + 空气 + Sn 平方和:
\begin{align}
  \sigma_E^\mathrm{total\,straggle}
  &= \sqrt{\sigma_E^{Sn,2} + \sigma_E^{air,2} + 2\sigma_E^{PDC,2}} \notag\\
  &= \sqrt{1.12^2 + 0.087^2 + 2 \cdot 0.035^2} \notag\\
  &\approx \sqrt{1.254 + 0.008 + 0.002}
  \approx 1.12\;\mathrm{MeV}.
\end{align}
Sn 的 straggling 主导. 转换为动量:
\begin{equation}
  \sigma_p^\mathrm{straggle}
  = \sigma_E^\mathrm{total\,straggle} / \beta
  = 1.12 / 0.561
  \approx 2.0\;\mathrm{MeV}/c.
  \label{eq:partIc:sigma_p_straggle}
\end{equation}

\paragraph{对比 PDC-only ensemble (无 Sn).} 关掉 Sn 靶时, $\sigma_E^\mathrm{total\,straggle}$ 退化到只剩 air + PDC, $\sqrt{0.087^2 + 0.002^2 + 0.002^2} \approx 0.087\;\mathrm{MeV}$, 对应 $\sigma_p^\mathrm{straggle} \approx 0.16\;\mathrm{MeV}/c$ — 完全可忽略. 这解释了为何状态报告里 PDC-only ensemble (\StatusReport{} §修复欠覆盖) 在关掉 Sn 后 $\sigma_p$ 没有 Landau-tail 病态.

\subsection{修复方案}

\paragraph{RK4 步长内的能损项.} 把式~\eqref{eq:partIc:bb_master} 加入 RK4 的力项:
\begin{equation}
  \frac{d \vec p}{ds} = q \vec v \times \vec B
  -\hat v \cdot \rho(\vec r) \cdot \left[-\langle dE/dx \rangle\right] / c.
  \label{eq:partIc:rk4_with_dEdx}
\end{equation}
其中 $\rho(\vec r)$ 是位置相关的材料密度 (查 GDML/材料表), $-\langle dE/dx \rangle$ 按当前 $\beta\gamma$ 在 step 中点重算. 实现要点:
\begin{itemize}[nosep, leftmargin=2em]
  \item 介质材料表预编译: 按 GDML solid → material → $(Z, A, I, \rho)$ 查找; PDG 表插值 $\langle dE/dx \rangle (\beta\gamma)$.
  \item RK4 步长里平均能损只取 mean (不引入 Landau 涨落), 因为我们做的是\emph{确定性前向模型}; Landau 入观测协方差 (Kalman 处理).
  \item Step end 重算 $\beta, \gamma$, 更新 $\vec p$ 模长 (方向不变, 只改幅度).
\end{itemize}

\paragraph{预测.} 加入 $dE/dx$ 后, RK 拟合的 $p_z$ 偏差应从 $-10.8\;\mathrm{MeV}/c$ 收紧到 $|\langle \Delta p_z \rangle| \lesssim 1\;\mathrm{MeV}/c$ (剩余的 1 MeV/$c$ 来自 Landau 涨落 + 材料表精度). 同时 Sn 处 Landau straggling $\sigma_E \sim 1.1\;\mathrm{MeV}$ 应进入 PDC 测量协方差 (Part~II §5), 给 Fisher $\sigma_p$ 增加 $\sim 2\;\mathrm{MeV}/c$ 的额外贡献, 使 Fisher 与残差 std 在 $p_z$ 方向上协同.

% --- end of Section 8 ---

\section{Highland 多次散射公式与精度修正}
\label{sec:partIc:highland}

\subsection{陈述}

PDG (Eq.~34.16) 给出的 Highland-Lynch-Dahl 公式是描述带电粒子在物质中多次小角度库仑散射 (multiple Coulomb scattering, MS) 的工程级标准:
\begin{equation}
  \theta_0 = \frac{13.6\;\mathrm{MeV}}{\beta c p}\, z
  \sqrt{\frac{x}{X_0}}\,
  \Bigl[1 + 0.038\,\ln\!\bigl(x z^2 / (X_0 \beta^2)\bigr)\Bigr],
  \label{eq:partIc:highland_master}
\end{equation}
其中 $\theta_0$ 是 Gaussian 投影角分布的 $1\sigma$ 宽度 (即 $\theta_x$ 或 $\theta_y$ 之一的 RMS, 整体 $\theta_\mathrm{space}$ 的 RMS 是 $\sqrt 2 \theta_0$). $z$ 是入射粒子电荷数; $x$ 是物质厚度; $X_0$ 是介质辐射长度. 在我们的工作点 $\beta\gamma = 0.677$, $z=1$, 故 $z^2/\beta^2 \approx 3.18$, 但 PDG 推荐对 $\beta \gtrsim 0.5$ 重粒子 $\ln$ 项里直接用 $\ln(x/X_0)$, 误差 $\lesssim 1\%$. 状态报告 (\StatusReport{} §Highland 公式) 采用的就是这个简化版.

\subsection{物理来源 (Lewis 1950 + Molière 小角理论)}

\paragraph{Step 1. 单次 Rutherford 散射.} 单次 Coulomb 散射 (粒子入射, 散射核屏蔽 Coulomb 势) 的微分截面:
\begin{equation}
  \frac{d\sigma}{d\Omega} = \frac{(z Z e^2)^2}{4 (p c \beta)^2 \sin^4(\theta/2)},
\end{equation}
小角下 $\theta \to 0$ 强发散, 但被原子电子屏蔽截至 (Thomas-Fermi 半径 $a_\mathrm{TF} \sim 0.885 a_0 Z^{-1/3}$).

\paragraph{Step 2. $\langle \theta^2 \rangle$ per atom.} 在屏蔽截止下, 单原子平均:
\begin{equation}
  \langle \theta^2 \rangle_\mathrm{atom}
  = \int \theta^2 (d\sigma/d\Omega) d\Omega
  \propto \frac{1}{(\beta p)^2}\, Z(Z+1) \ln\!\bigl(\theta_\mathrm{max}/\theta_\mathrm{min}\bigr).
\end{equation}
$\theta_\mathrm{max}$ 由核大小给, $\theta_\mathrm{min}$ 由屏蔽给, $\ln$ 项是 Coulomb logarithm.

\paragraph{Step 3. 中心极限定理 (Lewis 1950).} $N$ 次独立散射后, 投影角分布在小角下为高斯, RMS 为
\begin{equation}
  \theta_0^2 \approx N \langle \theta^2 \rangle_\mathrm{atom} / 2
  = n_\mathrm{atom} \cdot t \cdot \langle\theta^2\rangle_\mathrm{atom}/2,
\end{equation}
其中因子 $1/2$ 把 3D 角分布投影到 1D. 把所有原子物理塞进定义辐射长度 $X_0$ (单位 g/cm$^2$):
\begin{equation}
  \frac{1}{X_0} = 4\alpha r_e^2 \frac{N_A}{A} Z^2 [L_\mathrm{rad} - f(Z)] + \cdots,
\end{equation}
则 Lewis 给出 ($\beta \to 1$ 极限)
\begin{equation}
  \theta_0 \approx \frac{E_s}{p c} \sqrt{x/X_0},
  \quad
  E_s = m_e c^2 \sqrt{4\pi/\alpha} = 21.2\;\mathrm{MeV}.
  \label{eq:partIc:lewis}
\end{equation}

\paragraph{Step 4. Highland 经验拟合 (1975, Lynch-Dahl 1991).} 实际数据表明 Lewis 的 $E_s = 21.2$ 在大多数实验下\emph{高估}, 因为非高斯尾不能完全靠中心极限"吞掉". Highland 1975 用 Molière 全分布做 Monte Carlo, 拟合出经验常数 13.6 MeV (而非 21.2) 加上 $0.038 \ln(x/X_0)$ 修正项, Lynch-Dahl 1991 进一步精确化. 这就是式~\eqref{eq:partIc:highland_master} 的来源. 适用范围 $10^{-3} \lesssim x/X_0 \lesssim 100$.

\paragraph{$0.038 \ln$ 项的角色.} 当 $x/X_0$ 很小, $\ln$ 项是负数, 把 $\theta_0$ 拉低 (反映非高斯小尾在薄材料里更不显著); 当 $x/X_0$ 很大, $\ln$ 项是正数, 把 $\theta_0$ 拉高 (薄高斯近似失效, 出现重尾). 在 $x/X_0 = 1$ (一个辐射长度) 时 $\ln$ 项为 0.

\subsection{逐材料计算 (复现状态报告)}

下面把状态报告 (\StatusReport{} §Highland 公式) 的三个数字 (Sn 16.4, air 1.72, PDC 1.21 mrad) 从第一性原理推导.

\paragraph{$\beta p$ 因子.} $p = 635\;\mathrm{MeV}/c$, $\beta = 0.5607$, $\beta p = 0.5607 \times 635 = 355.8\;\mathrm{MeV}/c$. 故
\begin{equation}
  \frac{13.6}{\beta p} = \frac{13.6}{355.8} = 0.0382\;\mathrm{rad}.
  \label{eq:partIc:highland_prefactor}
\end{equation}

\paragraph{Sn 靶.} $\rho_\mathrm{Sn} = 7.31\;\mathrm{g/cm^3}$, $X_0^\mathrm{Sn} = 8.82\;\mathrm{g/cm^2}$ (Sn 的标准辐射长度), 厚度 $t = 0.25\;\mathrm{cm}$. 故
\begin{align}
  \rho t &= 7.31 \times 0.25 = 1.828\;\mathrm{g/cm^2}, \\
  x/X_0 &= 1.828 / 8.82 = 0.2073, \\
  \sqrt{x/X_0} &= 0.4553, \\
  \ln(x/X_0) &= \ln(0.2073) = -1.573, \\
  1 + 0.038 \ln(x/X_0) &= 1 + 0.038 \times (-1.573) = 0.940, \\
  \theta_0^\mathrm{Sn} &= 0.0382 \times 0.4553 \times 0.940
  = 0.01635\;\mathrm{rad}
  = 16.35\;\mathrm{mrad}.
\end{align}
四舍五入到 $16.4\;\mathrm{mrad}$, 与状态报告完全一致.

\paragraph{空气 (1 atm, 1 m).} $\rho_\mathrm{air} = 1.205 \times 10^{-3}\;\mathrm{g/cm^3}$, $X_0^\mathrm{air} = 36.66\;\mathrm{g/cm^2}$ (PDG), $t = 100\;\mathrm{cm}$.
\begin{align}
  \rho t &= 1.205\times 10^{-3} \cdot 100 = 0.1205\;\mathrm{g/cm^2}, \\
  x/X_0 &= 0.1205 / 36.66 = 3.288 \times 10^{-3}, \\
  \sqrt{x/X_0} &= 0.0573, \\
  \ln(x/X_0) &= \ln(3.288\times 10^{-3}) = -5.717, \\
  1 + 0.038 \ln(x/X_0) &= 1 - 0.217 = 0.783, \\
  \theta_0^\mathrm{air} &= 0.0382 \times 0.0573 \times 0.783
  = 0.001714\;\mathrm{rad} = 1.71\;\mathrm{mrad}.
\end{align}
舍入到 $1.72\;\mathrm{mrad}$ (与状态报告一致, 1.72 vs 1.71 的 $\sim 0.01$ 差源于中间舍入).

\paragraph{PDC 气体 (75\% Ar + 25\% i-C$_4$H$_{10}$).} 状态报告给出混合密度 $\rho_\mathrm{mix} = 1.87\times 10^{-3}\;\mathrm{g/cm^3}$, 等效辐射长度 $X_0^\mathrm{mix} = 24.1\;\mathrm{g/cm^2}$, 路径 $t = 22.4\;\mathrm{cm}$ (190 mm $/\cos 32^\circ$).
\begin{align}
  \rho t &= 1.87\times 10^{-3} \cdot 22.4 = 0.04189\;\mathrm{g/cm^2}, \\
  x/X_0 &= 0.04189 / 24.1 = 1.738 \times 10^{-3}, \\
  \sqrt{x/X_0} &= 0.04169, \\
  \ln(x/X_0) &= \ln(1.738\times 10^{-3}) = -6.355, \\
  1 + 0.038 \ln(x/X_0) &= 1 - 0.241 = 0.759, \\
  \theta_0^\mathrm{PDC} &= 0.0382 \times 0.04169 \times 0.759
  = 0.001209\;\mathrm{rad} = 1.21\;\mathrm{mrad}.
\end{align}
与状态报告一致.

\subsection{多次散射对 $\sigma_p$ 的贡献}

\paragraph{磁场后散射 → 出射角误差.} Sn 靶散射发生在\emph{磁场之前}, 被 RK fitter 的方向自由度 $(u, v)$ 完全吸收 (\StatusReport{} §各材料对动量重建的影响). 真正污染动量重建的是\emph{磁场之后}的 MS:
\begin{itemize}[nosep, leftmargin=2em]
  \item 空气 1 m: $\theta_0^\mathrm{air} = 1.72\;\mathrm{mrad}$.
  \item PDC1 气体: $\theta_0^\mathrm{PDC} = 1.21\;\mathrm{mrad}$ (PDC1 内部散射的角偏转传到 PDC2 测量).
  \item PDC2 气体: 仅引起 PDC2 内部 $\sim 0.08\;\mathrm{mm}$ 位移, 可忽略.
\end{itemize}

\paragraph{合并.} 平方和:
\begin{equation}
  \sigma_\theta^\mathrm{MS} = \sqrt{(\theta_0^\mathrm{air})^2 + (\theta_0^\mathrm{PDC})^2}
  = \sqrt{1.72^2 + 1.21^2}
  = \sqrt{2.958 + 1.464}
  = 2.10\;\mathrm{mrad}.
  \label{eq:partIc:sigma_theta_MS_total}
\end{equation}

\paragraph{$\sigma_p$ 贡献.} 用式~\eqref{eq:partIc:sigma_p_master}:
\begin{align}
  \sigma_p^\mathrm{air} &= 1460 \times 0.00172 = 2.51\;\mathrm{MeV}/c, \\
  \sigma_p^\mathrm{PDC} &= 1460 \times 0.00121 = 1.77\;\mathrm{MeV}/c, \\
  \sigma_p^\mathrm{MS,total} &= 1460 \times 0.00210 = 3.07\;\mathrm{MeV}/c,
\end{align}
四舍五入与状态报告 (\StatusReport{} §综合动量分辨, 表~综合动量分辨) 完全一致 (2.5, 1.8, 3.1 MeV/$c$).

\paragraph{综合 ($\sigma_x = 0.5\;\mathrm{mm}$).}
\begin{equation}
  \sigma_p^\mathrm{total} = \sqrt{(\sigma_p^\mathrm{pos})^2 + (\sigma_p^\mathrm{MS,total})^2}
  = \sqrt{1.03^2 + 3.07^2} = \sqrt{1.06 + 9.42}
  = \sqrt{10.48} = 3.24\;\mathrm{MeV}/c.
\end{equation}
即\textbf{$\sigma_p \approx 3.2\;\mathrm{MeV}/c$}, 这正是 NN 后端 $p_z$ MAE = 3.1 MeV/$c$ 的目标值.

\paragraph{综合 ($\sigma_x = 2.0\;\mathrm{mm}$).}
\begin{equation}
  \sigma_p^\mathrm{total}(2.0) = \sqrt{4.13^2 + 3.07^2}
  = \sqrt{17.06 + 9.42} = 5.15\;\mathrm{MeV}/c.
\end{equation}
即\textbf{$\sigma_p \approx 5.1\;\mathrm{MeV}/c$}, 与状态报告 表~综合动量分辨 完全一致.

\subsection{MS 与 LM 拟合的相互作用}

\paragraph{问题陈述.} RK4 前向模型假设 \emph{确定性}光滑轨迹: 给定靶处动量 $\vec p_\mathrm{tgt}$, 轨迹通过 PDC 平面的位置完全确定, 没有过程噪声. 但实际 MS 是\emph{每事件独立的随机过程}, 把"PDC 测量值 - RK4 预测值"残差分布展宽到\emph{超过} $\sigma_\mathrm{wire}$ 应该给的范围. 状态报告 (\StatusReport{} §修复欠覆盖) 观测到关掉 Sn 靶后 $\sigma_x \approx 3$--$15\;\mathrm{mm}$ (MS 主导) vs $\sigma_\mathrm{wire} = 0.2\;\mathrm{mm}$, 相差 15--75 倍.

\paragraph{两种修复方法.}
\begin{enumerate}[nosep, leftmargin=2em]
  \item \textbf{Kalman / measurement covariance inflation.} 把 MS 贡献加入测量协方差: $\Sigma_\mathrm{meas} = \mathrm{diag}(\sigma_\mathrm{wire}^2) + \Sigma_\mathrm{MS}$, 其中 $\Sigma_\mathrm{MS}$ 由材料厚度沿轨迹累积. 这是 "rough Kalman" 处理, 在每次 RK4 step 推进时把过程噪声转换成等效测量协方差. 见 Part~II §3 修复路径.
  \item \textbf{Toy-MC PDC smearing (status report 的 method 6).} 在反演前对每次拟合, 在初始方向上加 Gaussian smear $\theta_0$, 重复 $N$ 次, 用拟合结果的散布估计 $\sigma_p$. 这是 frequentist ensemble 方法, 不需要修改 forward model, 但需要 $N$ 倍计算开销.
\end{enumerate}
状态报告里两种方法都未在 production 启用; Part~II §3 给出明确的修复 roadmap.

\subsection{Highland 公式的局限}

Highland 给的是\emph{投影角的高斯近似}; 真实分布有重尾 (Molière 全分布):
\begin{itemize}[nosep, leftmargin=2em]
  \item 高斯近似在角度 $\lesssim 3 \theta_0$ 内有效, 占总事件 $\sim 99\%$.
  \item 大角尾巴 ($> 3\theta_0$) 由 Molière "single-plural" component 决定, 概率 $\sim 1\%$.
\end{itemize}
对覆盖率/pull 分布尾部分析, 直接用 Geant4 G4MultipleScattering 模型 (背后是 Urban 或 GS 模型, 已经包含尾巴) 比 Highland 准确. 状态报告 §G4 涨落集合就是 G4 simulated MS 的结果, 自动包含尾巴.

\paragraph{对覆盖率的影响.} 假设我们把 MS 当 Gauss 处理 (Highland $\theta_0$ + Kalman 协方差膨胀), 则覆盖率会在 68\% / 95\% nominal 附近\emph{略低估} ($\sim 1$--$2\%$ 的 deficit), 因为重尾事件的真值会被推到 Gauss 区间外. 在我们 744 事件 ensemble 上, 这个 deficit 是\emph{可观察的}但低于 Clopper-Pearson 的统计 noise floor; Part~III §5 给出对应的统计检验.

\paragraph{修复后预期.} 在 Part~II §3 的 Kalman 协方差修复 + Part~II §5 的 $dE/dx$ 修复都到位后, 预期:
\begin{itemize}[nosep, leftmargin=2em]
  \item $|\langle \Delta p_z \rangle| \lesssim 1\;\mathrm{MeV}/c$ (来自 $dE/dx$ 修复, Part~II §5).
  \item Fisher $\sigma_{p_z}$ 与 残差 std 在 $p_z$ 方向匹配 (来自 MS 协方差膨胀, Part~II §3).
  \item 68\% 覆盖率全部组件回到 nominal 附近 ($\pm 2\%$).
\end{itemize}
$p_y$ 方向参数化 line-arization 偏差 (\StatusReport{} §修复欠覆盖) 是另一条独立路径, 见 Part~II §8.

% --- end of Section 9 ---
 加载。
%
% 不要在此文件中包含 \documentclass / \begin{document} / 序言;
% 主壳层提供 \part{第一部分: 数学基础} 标题, ctex+xelatex 字体,
% booktabs/amsmath/amsthm/enumitem 等宏包.
%
% 创建日期: 2026-04-26
% 作者:    Baiting Tian
% 关联文件:
%   * docs/reports/reconstruction/momentum_reco/rk/rk_reconstruction_status_20260416_zh.tex
%   * docs/superpowers/specs/2026-04-26-rk-error-analysis-deep-dive-design.md
%
% 本文件结构 (三章):
%   7. Glückstern 公式 (SAMURAI 2-PDC + 弧长几何)
%   8. Bethe-Bloch + Landau/Vavilov (缺失的 dE/dx 系统)
%   9. Highland 多次散射公式与精度修正
% =============================================================

\providecommand{\StatusReport}{文献~[1]}

\section{Glückstern 公式：SAMURAI 2-PDC + 弧长几何下的解析动量分辨极限}
\label{sec:partIc:gluckstern}

\subsection{出发点：闭合形式的动量分辨}

Glückstern 1963 年的经典论文给出了"匀场磁谱仪 + $N$ 个等距径迹测量平面"几何下\emph{动量分辨的最优闭合形式}, 是所有磁谱仪设计的初步标定工具. 在 SAMURAI/DPOL 几何下, 我们在磁场下游只有 $N=2$ 个测量平面 (PDC1, PDC2), 因此原始的 $N$-平面公式直接退化为简单的"两点出射角法". 本章先从一阶几何把出射角法的 $\sigma_p \propto p^2 \sigma_x / (q B L_\mathrm{eff} L_{12})$ 推导出来, 然后插入 SAMURAI 数值, 最后再和广义 Glückstern 公式 (PDG Eq.~34.41) 做一致性检验.

参考的状态报告 (\StatusReport{} §理论分辨极限) 给出 $\sigma_x = 0.5\;\mathrm{mm}$ 时位置贡献 $\sigma_p \approx 1.0\;\mathrm{MeV}/c$, 并把它与 RK4 + 完整磁场图的 Fisher $\sigma_{|\vec p|}$ 范围 $2.4\text{--}6.0\;\mathrm{MeV}/c$ 比对. 本章的目的就是给这个 $\sigma_p \approx 1.0$ 一个独立的解析推导, 并量化"理想的薄楔形偶极 + 无场漂移"假设和真实 RK4 + 厚磁体之间的 $\sim 10\%$ 偏差.

\subsection{几何与符号}

\paragraph{几何示意.} 带电粒子从靶 (target) 出射, 进入 SAMURAI 偶极 (有效弯曲长度 $L_\mathrm{eff}$, 匀强磁场 $B$, 主弯曲方向沿 $x$, $y$ 方向无场, 束流名义沿 $z$). 离开磁场后通过一段"无场漂移臂" $L_\mathrm{arm}$ 到达 PDC1, 再继续漂移 $L_{12}$ 到 PDC2. 每个 PDC 测量 $(x, y)$ 两个坐标, 各分量分辨为 $\sigma_x$ (假设独立同分布, 且在 $y$ 方向相同).

\paragraph{符号汇总.}
\begin{itemize}[nosep, leftmargin=2em]
  \item $p$: 粒子动量大小 (MeV/$c$). 当前工作点 $p = 635\;\mathrm{MeV}/c$.
  \item $q$: 电荷 (单位 $e$). 质子 $q = 1$.
  \item $B$: 磁场强度 (T). SAMURAI 工作点 $B = 1.15\;\mathrm{T}$.
  \item $L_\mathrm{eff}$: 偶极有效弯曲长度 (m). $L_\mathrm{eff} = 0.8\;\mathrm{m}$.
  \item $L_\mathrm{arm}$: 磁场出口到 PDC1 的漂移距离 (m). $L_\mathrm{arm} = 4.0\;\mathrm{m}$.
  \item $L_{12}$: PDC1 到 PDC2 的距离 (m). $L_{12} = 1.0\;\mathrm{m}$.
  \item $\sigma_x$: 单平面位置分辨 (mm). 取 $\sigma_x = 0.5\;\mathrm{mm}$ (Geant4 wire smearing) 或 $2.0\;\mathrm{mm}$ (拟合器实际容差).
  \item $\rho$: 弯曲半径 (m).
  \item $\theta_\mathrm{bend}$: 偏转角.
\end{itemize}

PDG 标准的 magnetic rigidity 关系:
\begin{equation}
  B \rho \;[\mathrm{T \cdot m}] = \frac{p\;[\mathrm{GeV}/c]}{0.299\,792\,458 \cdot q}
  \approx \frac{p\;[\mathrm{GeV}/c]}{0.3 q}.
  \label{eq:partIc:Brho}
\end{equation}
代入 $p = 0.635\;\mathrm{GeV}/c$, $q = 1$: $B\rho = 0.635 / 0.3 = 2.117\;\mathrm{T \cdot m}$, 故 $\rho = 2.117 / 1.15 = 1840\;\mathrm{mm}$, 与状态报告 (\StatusReport{} §理论分辨极限) 一致.

\subsection{偏转角及其对动量的灵敏度}

在薄楔形 (thin-wedge) 近似下, 粒子在偶极内部沿圆弧运动, 离开偶极时方向相对入射方向旋转:
\begin{equation}
  \theta_\mathrm{bend} = \frac{L_\mathrm{eff}}{\rho}
  = \frac{q B L_\mathrm{eff}}{p}
  = \frac{0.3 \cdot q B L_\mathrm{eff}\;[\mathrm{T \cdot m}]}{p\;[\mathrm{GeV}/c]}.
  \label{eq:partIc:theta_bend}
\end{equation}
代入 $L_\mathrm{eff} = 0.8\;\mathrm{m}$: $\theta_\mathrm{bend} = 0.8 / 1.840 = 0.435\;\mathrm{rad} \approx 24.9^\circ$, 复现了状态报告里的 25$^\circ$ 名义弯曲角.

\paragraph{动量灵敏度.} 把式~\eqref{eq:partIc:theta_bend} 对 $p$ 求导:
\begin{equation}
  \frac{\partial \theta_\mathrm{bend}}{\partial p}
  = -\frac{q B L_\mathrm{eff}}{p^2}
  = -\frac{\theta_\mathrm{bend}}{p}.
  \label{eq:partIc:dtheta_dp}
\end{equation}
从而, 给定出射角的测量误差 $\sigma_\theta$, 由误差传播
\begin{equation}
  \sigma_p = \left|\frac{\partial p}{\partial \theta_\mathrm{bend}}\right| \sigma_\theta
  = \frac{p}{\theta_\mathrm{bend}} \sigma_\theta.
  \label{eq:partIc:sigma_p_master}
\end{equation}
这是本章的"主公式". 后面要做的事就是把不同来源的 $\sigma_\theta$ 代进去.

把数值代入: $p / \theta_\mathrm{bend} = 635 / 0.435 = 1460\;\mathrm{(MeV}/c\mathrm{)/rad}$, 即每 1~mrad 的出射角误差贡献 $1.46\;\mathrm{MeV}/c$ 的动量误差. 这个 1460 的转换因子是状态报告 (\StatusReport{} §理论分辨极限, 式~5) 的核心数字.

\subsection{两点出射角估计与位置贡献}

\paragraph{出射角估计器.} PDC1 和 PDC2 测量到的位置在 $x$ 方向的差分给出出射角的最大似然估计 (在测量误差为高斯, 且无 MS 时):
\begin{equation}
  \widehat\theta_\mathrm{exit} = \frac{x_2 - x_1}{L_{12}}.
  \label{eq:partIc:theta_exit_est}
\end{equation}
其中 $x_1, x_2$ 是 PDC1, PDC2 处的 $x$ 坐标. $\widehat\theta_\mathrm{exit}$ 是无偏的, 因为 $\langle x_2 - x_1 \rangle = L_{12} \cdot \theta_\mathrm{exit}$.

\paragraph{方差.} $x_1, x_2$ 独立, 各自方差 $\sigma_x^2$, 故:
\begin{equation}
  \mathrm{Var}(\widehat\theta_\mathrm{exit}) = \frac{\mathrm{Var}(x_1) + \mathrm{Var}(x_2)}{L_{12}^2}
  = \frac{2 \sigma_x^2}{L_{12}^2},
  \quad
  \sigma_{\theta,\mathrm{pos}} = \frac{\sigma_x \sqrt 2}{L_{12}}.
  \label{eq:partIc:sigma_theta_pos}
\end{equation}

\paragraph{动量分辨 (位置贡献).} 把式~\eqref{eq:partIc:sigma_theta_pos} 代入式~\eqref{eq:partIc:sigma_p_master}:
\begin{equation}
  \sigma_p^\mathrm{pos}
  = \frac{p}{\theta_\mathrm{bend}} \cdot \frac{\sigma_x \sqrt 2}{L_{12}}
  = \frac{p^2 \sigma_x \sqrt 2}{q B L_\mathrm{eff} L_{12}}
  = \frac{\sqrt 2 \, p^2 \sigma_x}{0.3 q B L_\mathrm{eff} L_{12}}.
  \label{eq:partIc:gluckstern_2pt}
\end{equation}
这就是 Glückstern 公式在 $N=2$, 测量平面位于无场漂移段的退化形式. 注意 $\sigma_p \propto p^2$ — 这是 Glückstern scaling, 高动量谱仪测量的标志.

\paragraph{数值.} 代入 SAMURAI 数值: $p = 0.635\;\mathrm{GeV}/c$, $\sigma_x = 0.5\;\mathrm{mm} = 5 \times 10^{-4}\;\mathrm{m}$, $q = 1$, $B = 1.15\;\mathrm{T}$, $L_\mathrm{eff} = 0.8\;\mathrm{m}$, $L_{12} = 1.0\;\mathrm{m}$:
\begin{align}
  \sigma_p^\mathrm{pos}
  &= \frac{\sqrt 2 \times 0.635^2 \times 5 \times 10^{-4}}{0.3 \times 1 \times 1.15 \times 0.8 \times 1.0}\;\mathrm{GeV}/c \notag\\
  &= \frac{1.414 \times 0.4032 \times 5 \times 10^{-4}}{0.276}\;\mathrm{GeV}/c \notag\\
  &= \frac{2.851 \times 10^{-4}}{0.276}\;\mathrm{GeV}/c
  = 1.033 \times 10^{-3}\;\mathrm{GeV}/c \notag\\
  &\approx 1.03\;\mathrm{MeV}/c.
\end{align}
这与状态报告 (\StatusReport{} §理论分辨极限, 表~1) 给出的 $\sigma_p \approx 1.0\;\mathrm{MeV}/c$ 完全一致.

\paragraph{$\sigma_x = 2.0\;\mathrm{mm}$ 的 fitter 容差.} 把 $\sigma_x$ 改为 $2.0\;\mathrm{mm}$:
\begin{equation}
  \sigma_p^\mathrm{pos}(2.0\;\mathrm{mm})
  = 4 \times 1.03 = 4.13\;\mathrm{MeV}/c,
\end{equation}
也复现了状态报告里的 $4.1\;\mathrm{MeV}/c$ 数字. 这里因子 4 来自 $\sigma_p \propto \sigma_x$ 的线性 scaling.

\subsection{广义 Glückstern 公式 ($N$ 平面) 的退化}

PDG (Eq.~34.41) 给出的广义结果是: 沿匀场弯曲段在 $N$ 个等距点测量轨迹位置, 用最优最小二乘估计动量, 得到:
\begin{equation}
  \frac{\sigma_p}{p} \bigg|_\mathrm{Glückstern}
  = \frac{\sigma_x}{0.3 B L^2}\, p \, \sqrt{\frac{720}{N + 4}},
  \label{eq:partIc:gluckstern_N}
\end{equation}
其中 $L$ 是测量段总长度 (注意: 这里 $L$ 是\emph{测量段}的总长, 不是磁体本身长度; $\sigma_p$ 单位与 $p$ 一致, $B$ 取 T, $L$ 取 m).

注意我们的几何与 PDG 标准设置有两点重要差异:
\begin{enumerate}[nosep, leftmargin=2em]
  \item PDG 假设测量平面分布在\emph{弯曲段内部} (沿磁体均匀分布), 我们的 PDC 在\emph{磁场出口下游的无场漂移段}.
  \item PDG 的 $L$ 是测量段总长 ($N-1$ 个间距覆盖的距离), 我们 $N=2$ 直接退化为 $L = L_{12}$.
\end{enumerate}

\paragraph{退化检验.} 把 $N=2$ 代入式~\eqref{eq:partIc:gluckstern_N}: $\sqrt{720 / 6} = \sqrt{120} = 10.95$. 故
\begin{equation}
  \frac{\sigma_p}{p} = \frac{10.95 \sigma_x p}{0.3 B L^2}.
\end{equation}
两点法 (式~\eqref{eq:partIc:gluckstern_2pt}, 用 $\theta_\mathrm{bend} \approx L_\mathrm{eff}/\rho$ 替换) 给出的相对分辨:
\begin{equation}
  \frac{\sigma_p}{p}\bigg|_\mathrm{2pt}
  = \frac{\sqrt 2 p \sigma_x}{0.3 B L_\mathrm{eff} L_{12}}.
\end{equation}
两者比值 (假设 $L = L_{12}$):
\begin{equation}
  \frac{\sigma_p/p|_\mathrm{Glückstern}}{\sigma_p/p|_\mathrm{2pt}}
  = \frac{10.95 / L_{12}^2}{\sqrt 2 / (L_\mathrm{eff} L_{12})}
  = \frac{10.95 \cdot L_\mathrm{eff}}{\sqrt 2 \cdot L_{12}}
  \approx 6.2,
\end{equation}
比值 $\sim 6$. 但这是一个\emph{geometry-mismatch}对比 — 广义 Glückstern 假设测量平面分布在磁场内部, 而我们的两点法用一个外部基线 $L_{12}$ 加上偶极偏转角换算. 真正可比的是把 $L_\mathrm{eff} \to L_{12}$ 当作"测量段在磁体内"的虚拟设置, 此时两个公式只差 PDG 的常数因子 $\sqrt{120/2} \approx 7.7$. 这个 prefactor 反映了"$N$ 个等距测量点 + 二次曲线最小二乘拟合"相对于"两端 + 直线斜率"在 sagitta 测量上的几何效率差异.

\subsection{多分量协方差: 为什么 $\sigma_{p_y} \ll \sigma_{p_x}$}

PDC 测量 $(x, y)$ 两个独立坐标, 而 SAMURAI 偶极\emph{只在 $x$-$z$ 平面弯曲}. 因此:
\begin{itemize}[nosep, leftmargin=2em]
  \item $p_x$ (主弯曲方向) 通过 $\theta_\mathrm{bend}$ 反算, 灵敏度由 $p / \theta_\mathrm{bend} = 1460\;(\mathrm{MeV}/c)/\mathrm{rad}$ 决定.
  \item $p_y$ (无场方向) 直接来自 $y$ 截距/斜率, 几何因子完全不同.
\end{itemize}

\paragraph{$y$ 方向的运动学.} 在 $y$ 方向上, SAMURAI 偶极没有恢复力, 粒子从靶到 PDC2 沿一条直线 (忽略 MS):
\begin{equation}
  y_\mathrm{PDC} = y_\mathrm{tgt} + \frac{p_y}{p_z} \cdot z_\mathrm{drift}.
\end{equation}
其中 $z_\mathrm{drift}$ 是从靶到 PDC 的总漂移距离 (沿束流方向 $\sim L_\mathrm{eff} + L_\mathrm{arm} \approx 4.8\;\mathrm{m}$). 出射角 $\theta_y \approx p_y / p_z$, 灵敏度
\begin{equation}
  \frac{\partial y_\mathrm{PDC}}{\partial p_y} = \frac{z_\mathrm{drift}}{p_z}.
\end{equation}

\paragraph{$\sigma_{p_y}$.} 用两点法测 $\theta_y$:
\begin{equation}
  \sigma_{\theta_y} = \frac{\sigma_y \sqrt 2}{L_{12}},
  \quad
  \sigma_{p_y} = p_z \sigma_{\theta_y}
  = \frac{p_z \sigma_y \sqrt 2}{L_{12}}.
  \label{eq:partIc:sigma_py}
\end{equation}
代入: $p_z = 627\;\mathrm{MeV}/c$, $\sigma_y = 0.5\;\mathrm{mm}$, $L_{12} = 1000\;\mathrm{mm}$:
\begin{equation}
  \sigma_{p_y} = \frac{627 \times 5 \times 10^{-4} \sqrt 2}{1.0}
  \approx 0.443\;\mathrm{MeV}/c.
\end{equation}

\paragraph{比值.} $\sigma_{p_x} / \sigma_{p_y} = (p / \theta_\mathrm{bend}) / p_z = 1460 / 627 \approx 2.3$. 在没有 MS 时, $p_y$ 的几何分辨只比 $p_x$ 好 2 倍. 但状态报告 (\StatusReport{} §Fisher 表) 实测 Fisher $\sigma_{p_y} \in [0.15, 0.20]\;\mathrm{MeV}/c$, 而 $\sigma_{p_x} \in [3.2, 7.0]$, 比值 $\sim 20$ 倍 — 这并不矛盾, 是因为 $p_y$ 在两个 PDC 平面上有更长的杠杆: 实际 $y$-臂从靶到 PDC2 接近 5 m, 远大于 $L_{12} = 1\;\mathrm{m}$, 所以两个 PDC 的 $y$ 数据结合起来给出比"PDC1-PDC2 单一 $L_{12}$ 基线"更紧的 $p_y$ 约束 (本质上是 $y_\mathrm{tgt} = 0$ 的约束加 PDC1 把 $p_y$ 钉住). 正是这一点也导致了 $p_y$ 方向参数化敏感: $(u, v, p)$ 线性化在大 $\theta_y$ 情况下系统性低估 $p_y$ 的方差, 表现为 Fisher $\sigma_{p_y}$ 偏紧 $\sim 1.9\times$ (\StatusReport{} §修复欠覆盖). 这部分 ratio 故事属于 Part~II §8 的范畴, 这里只是给出 $p_y$ 方向几何上的"裸"分辨标尺.

\subsection{$\sim 10\%$ 偏差: 解析公式 vs RK4}

式~\eqref{eq:partIc:gluckstern_2pt} 假设了:
\begin{itemize}[nosep, leftmargin=2em]
  \item 偶极薄楔形, 弯曲集中在长度为 $L_\mathrm{eff}$ 的瞬时区域.
  \item 偶极外部完全无场.
  \item 偏转角小到可以做线性化 ($\sin\theta \approx \theta$).
\end{itemize}

实际 SAMURAI 偶极有显著的边缘场 (fringe field) 区, 进出口磁场连续过渡, 偏转角 $25^\circ$ 已经不是小角. 这两个修正使得 $\sigma_p^\mathrm{pos}$ 解析估计与 RK4 + 完整磁场图的 Fisher 估计偏差约 $10\%$. 状态报告 (\StatusReport{} §Fisher 表) 给出的 $\sigma_{|\vec p|} \in [2.4, 6.0]\;\mathrm{MeV}/c$ 包含了 MS 贡献, 在不同 truth 动量下变化. 把它和理论 $\sigma_p \approx 3.2\;\mathrm{MeV}/c$ ($\sigma_x = 0.5\;\mathrm{mm}$ + MS) 比较, $\pm 1\;\mathrm{MeV}/c$ 的散布完全在边缘场 + 大角度修正的预期范围内.

% --- end of Section 7 ---

\section{Bethe-Bloch 与 Landau/Vavilov 涨落：缺失的 $dE/dx$ 系统}
\label{sec:partIc:bethe_bloch}

\subsection{动机：$-10.8\;\mathrm{MeV}/c$ 的 $p_z$ 系统偏差}

状态报告 (\StatusReport{} §当前性能, 表~RK vs NN 对比) 给出, RK 后端在 $p_z$ 方向上系统性偏差 $\langle \Delta p_z \rangle = -10.8\;\mathrm{MeV}/c$, 而 NN 后端 $\langle \Delta p_z \rangle$ 接近 0. 该差异的物理来源是 RK 前向模型\emph{不包含连续电离能损 $dE/dx$}: RK4 把粒子在材料里的能量当作守恒量, 因此它从 PDC 命中位置反推靶动量时, 错把"靶到 PDC 之间一路损失掉的动能"当作"靶处动量也少了那么多", 给出系统性偏低的 $p_z$. 本章先把 Bethe-Bloch 的全表达式拆开, 数值算清每一段材料贡献多少 $\Delta E$, 再换算到 $\Delta p$, 验证它和实测的 $-10.8\;\mathrm{MeV}/c$ 一致.

\subsection{Bethe-Bloch 平均能损公式 (PDG 形式)}

\paragraph{基本表达式.} 对于动能 $T_\mathrm{kin}$, 速度 $\beta$, $\gamma = (1-\beta^2)^{-1/2}$ 的中等相对论重粒子 (在我们这里就是 $\beta\gamma \sim 1$ 量级的质子), 平均电离能损为:
\begin{equation}
  -\left\langle\frac{dE}{dx}\right\rangle
  = K z^2 \frac{Z}{A} \frac{1}{\beta^2}
  \left[
    \frac{1}{2}\ln\!\frac{2 m_e c^2 \beta^2 \gamma^2 T_\mathrm{max}}{I^2}
    - \beta^2
    - \frac{\delta(\beta\gamma)}{2}
  \right].
  \label{eq:partIc:bb_master}
\end{equation}
其中:
\begin{itemize}[nosep, leftmargin=2em]
  \item $K \equiv 4\pi N_A r_e^2 m_e c^2 = 0.307\;\mathrm{MeV \cdot g^{-1}\,cm^2}$ — 电离常数.
  \item $z$ — 入射粒子电荷数 (质子 $z=1$).
  \item $Z$ — 介质原子序数; $A$ — 介质原子质量 (g/mol).
  \item $\beta, \gamma$ — 入射粒子的运动学参数.
  \item $m_e$ — 电子质量, $m_e c^2 = 0.511\;\mathrm{MeV}$.
  \item $I$ — 介质平均电离能 (eV), 对每种材料经验给定.
  \item $T_\mathrm{max}$ — 单次碰撞最大能量传递, $T_\mathrm{max} = 2 m_e c^2 \beta^2 \gamma^2 / [1 + 2\gamma m_e/M + (m_e/M)^2]$, 其中 $M$ 是入射粒子静质量. 对 $\beta\gamma \lesssim 1$ 的质子, 分母 $\approx 1$, 故 $T_\mathrm{max} \approx 2 m_e c^2 \beta^2 \gamma^2$.
  \item $\delta(\beta\gamma)$ — Sternheimer 密度修正, 反映介质极化对远距离碰撞的屏蔽; 在 $\beta\gamma \lesssim 1$ 时 $\delta \approx 0$, 在 $\beta\gamma \gtrsim 100$ 时 $\delta \to 2\ln(\beta\gamma) + C$.
\end{itemize}

\paragraph{635 MeV/$c$ 质子的运动学.}
\begin{align}
  E &= \sqrt{p^2 c^2 + m^2 c^4} = \sqrt{635^2 + 938.27^2} = 1132.6\;\mathrm{MeV}, \\
  T_\mathrm{kin} &= E - m c^2 = 1132.6 - 938.27 = 194.4\;\mathrm{MeV}, \\
  \beta &= \frac{p c}{E} = \frac{635}{1132.6} = 0.5607, \\
  \gamma &= \frac{E}{m c^2} = \frac{1132.6}{938.27} = 1.207, \\
  \beta\gamma &= 0.677, \\
  T_\mathrm{max} &\approx 2 \cdot 0.511 \cdot 0.677^2 = 0.469\;\mathrm{MeV}.
\end{align}
注意 $\beta\gamma = 0.677$ 处于 $-dE/dx$ 曲线的下降臂, 接近最小电离区 (MIP, $\beta\gamma \approx 3$) 之前, 因此能损比 MIP 略高 (典型 $1.3 \times$ MIP).

\subsection{材料逐项贡献}

\paragraph{Sn 靶 (天然 Sn).} $Z = 50$, $A = 118.71\;\mathrm{g/mol}$, $\rho = 7.31\;\mathrm{g/cm^3}$, $I \approx 488\;\mathrm{eV}$, 厚度 $t = 2.5\;\mathrm{mm}$, $\rho t = 1.83\;\mathrm{g/cm^2}$.
\begin{align}
  \frac{Z}{A} &= 0.421, \\
  \ln(2 m_e c^2 \beta^2 \gamma^2) &= \ln(2 \cdot 0.511 \cdot 0.677^2)
  = \ln(0.469) = -0.757, \\
  \ln T_\mathrm{max} &= \ln(0.469) = -0.757, \\
  2 \ln I &= 2 \ln(488 \times 10^{-6}) = 2 \cdot (-7.624) = -15.25, \\
  \tfrac{1}{2}\ln(2 m_e c^2 \beta^2 \gamma^2 T_\mathrm{max}/I^2)
  &= \tfrac{1}{2}\bigl[-0.757 + (-0.757) - (-15.25)\bigr] \notag\\
  &= \tfrac{1}{2}(13.74) = 6.87.
\end{align}
代入式~\eqref{eq:partIc:bb_master} (取 $\delta \approx 0.05$, $\beta\gamma = 0.677$ 时 Sn 的密度修正小):
\begin{align}
  -\left\langle\frac{dE}{dx}\right\rangle_\mathrm{Sn}
  &= 0.307 \cdot 1 \cdot 0.421 \cdot \frac{1}{0.5607^2} \cdot
  \left[6.87 - 0.5607^2 - 0.025\right] \notag\\
  &= 0.307 \cdot 0.421 \cdot 3.182 \cdot 6.531
  = 2.69\;\mathrm{MeV \cdot g^{-1}\,cm^2}.
\end{align}
但 PDG 表对 Sn @ $\beta\gamma = 0.7$ 给出 $\sim 1.7\;\mathrm{MeV \cdot g^{-1}\,cm^2}$ (比上面"裸" Bethe-Bloch 值低 $\sim 30\%$), 这是因为高 $Z$ 材料的 K/L 壳层修正 ($C/Z$ shell correction) 在 $\beta\gamma \lesssim 1$ 不可忽略, 把 $-\langle dE/dx\rangle$ 拉低. 取 PDG 表值:
\begin{equation}
  -\langle dE/dx \rangle_\mathrm{Sn} \approx 1.7\;\mathrm{MeV \cdot g^{-1}\,cm^2}
  \;\Rightarrow\;
  \Delta E_\mathrm{Sn} = 1.7 \times 1.83 \approx 3.1\;\mathrm{MeV}.
  \label{eq:partIc:dE_Sn}
\end{equation}

\paragraph{空气 (1 atm, 1 m).} 平均原子组成 (质量分数): N $0.755$, O $0.232$, Ar $0.013$. 等效 $Z/A \approx 0.499$, $I = 85.7\;\mathrm{eV}$, $\rho = 1.205 \times 10^{-3}\;\mathrm{g/cm^3}$.

\begin{align}
  -\langle dE/dx \rangle_\mathrm{air}
  &\approx 0.307 \cdot 0.499 \cdot \frac{1}{0.5607^2} \cdot
  \left[\tfrac{1}{2}\ln\!\frac{2 \cdot 0.511 \cdot 0.677^2 \cdot 0.469}{(85.7\times 10^{-6})^2} - 0.314\right] \notag\\
  &= 0.307 \cdot 0.499 \cdot 3.18 \cdot
  \left[\tfrac{1}{2}\ln(2.405\times 10^7) - 0.314\right] \notag\\
  &= 0.487 \cdot \left[\tfrac{1}{2} \cdot 16.99 - 0.314\right] \notag\\
  &= 0.487 \cdot 8.18 = 3.99\;\mathrm{MeV \cdot g^{-1}\,cm^2}.
\end{align}
PDG 数据表 @ $\beta\gamma = 0.7$ 给 $\sim 2.0\;\mathrm{MeV \cdot g^{-1}\,cm^2}$, 比上式低 $\sim 50\%$ — 主因是上面"裸" Bethe-Bloch 在低 $\beta\gamma$ 没有正确收紧到 PDG 拟合. 取 PDG 标定值:
\begin{equation}
  -\langle dE/dx \rangle_\mathrm{air} \approx 2.0\;\mathrm{MeV \cdot g^{-1}\,cm^2}
  \;\Rightarrow\;
  \Delta E_\mathrm{air} = 2.0 \times 1.205 \times 10^{-1} \approx 0.24\;\mathrm{MeV}.
  \label{eq:partIc:dE_air}
\end{equation}
($\rho \cdot 100\;\mathrm{cm} = 0.1205\;\mathrm{g/cm^2}$.)

\paragraph{PDC 气体 (75\% Ar + 25\% i-C$_4$H$_{10}$, 体积分数).}
密度: $\rho_\mathrm{Ar} = 1.662\;\mathrm{mg/cm^3}$, $\rho_\mathrm{iso} = 2.49\;\mathrm{mg/cm^3}$.
质量分数: $w_\mathrm{Ar} = 0.667$, $w_\mathrm{iso} = 0.333$ (\StatusReport{} §Highland 推导).
混合密度 $\rho_\mathrm{mix} = 1.87\;\mathrm{mg/cm^3}$.

按 Bragg additivity rule, $-\langle dE/dx \rangle$ 按质量分数加权:
\begin{align}
  -\langle dE/dx\rangle_\mathrm{Ar} &\approx 1.5\;\mathrm{MeV \cdot g^{-1}\,cm^2} \;\text{(PDG @ $\beta\gamma{=}0.7$)}, \\
  -\langle dE/dx\rangle_\mathrm{iso} &\approx 2.4\;\mathrm{MeV \cdot g^{-1}\,cm^2}, \\
  -\langle dE/dx\rangle_\mathrm{mix} &= 0.667 \cdot 1.5 + 0.333 \cdot 2.4 \notag\\
  &= 1.00 + 0.80 = 1.80\;\mathrm{MeV \cdot g^{-1}\,cm^2}.
\end{align}
有效路径 224 mm (\StatusReport{} §Highland), $\rho t = 0.0419\;\mathrm{g/cm^2}$:
\begin{equation}
  \Delta E_\mathrm{PDC,1平面} = 1.80 \times 0.0419 \approx 0.075\;\mathrm{MeV}.
  \label{eq:partIc:dE_PDC}
\end{equation}
两个 PDC 共 $\Delta E_\mathrm{PDC,total} \approx 0.15\;\mathrm{MeV}$.

\paragraph{合计.}
\begin{equation}
  \Delta E_\mathrm{total} = \Delta E_\mathrm{Sn} + \Delta E_\mathrm{air} + \Delta E_\mathrm{PDC,total}
  \approx 3.1 + 0.24 + 0.15 = 3.5\;\mathrm{MeV}.
  \label{eq:partIc:dE_total_with_Sn}
\end{equation}

\subsection{从 $\Delta E$ 到 $\Delta p$}

对相对论粒子, $E^2 = p^2 c^2 + m^2 c^4$, 取微分:
\begin{equation}
  E\, dE = p c^2\, dp
  \;\Longrightarrow\;
  dp = \frac{E}{p c^2} dE = \frac{1}{\beta} dE \cdot \frac{1}{c}.
  \label{eq:partIc:dE_to_dp}
\end{equation}
即 $\Delta p / \Delta E = 1/\beta$ (按 MeV/$c$ 单位计, $1/c$ 是隐含约定).

代入 $\beta = 0.561$: $\Delta p / \Delta E = 1/0.561 = 1.78$.

\paragraph{Sn 启用情形 (E2E 配置).} 状态报告里产生 $-10.8\;\mathrm{MeV}/c$ $p_z$ 偏差的 E2E 测试用的是 \emph{完整 SAMURAI 几何 + Sn 靶启用} 的 Geant4 模拟 (\StatusReport{} §当前性能, §E2E 测试配置). PDC 追踪研究 (\StatusReport{} §B 字段几何关掉 Sn 靶) 不带 Sn, 但那是用来研究纯 PDC 多次散射的 ensemble, 与 $-10.8$ 偏差不是同一个测试.

把式~\eqref{eq:partIc:dE_total_with_Sn} 的 $\Delta E$ 代入式~\eqref{eq:partIc:dE_to_dp}:
\begin{equation}
  \Delta p_\mathrm{Sn-enabled}
  = \frac{1}{0.561} \cdot 3.5 = 6.2\;\mathrm{MeV}/c.
  \label{eq:partIc:dp_total_v1}
\end{equation}
但状态报告观测到的偏差是 $-10.8\;\mathrm{MeV}/c$, 比 $6.2$ 大 $\sim 75\%$. 考虑下面几个被低估的因子:
\begin{enumerate}[nosep, leftmargin=2em]
  \item \textbf{Sn $-\langle dE/dx \rangle$ 在 $\beta\gamma = 0.677$ 下其实更接近 $2.5\;\mathrm{MeV\cdot g^{-1}\,cm^2}$} (PDG 表插值含壳层修正), 把 $\Delta E_\mathrm{Sn}$ 推到 $4.6\;\mathrm{MeV}$.
  \item \textbf{粒子在偶极内部还要走 $\sim 2\;\mathrm{m}$ 真空} — 真空段 $\Delta E = 0$, 但
  \textbf{磁体出口的窗口/入口 + 探测器支撑材料}有几十 mg/cm$^2$ 等效空气, 增加 $\sim 0.5\text{--}1\;\mathrm{MeV}$.
  \item \textbf{粒子在 PDC 之间还要穿过靶到 PDC1 之间的束流管材料} (Mylar/Be 窗口, $\sim 0.5\;\mathrm{MeV}$).
\end{enumerate}
合理上调到 $\Delta E \approx 5.5\text{--}6\;\mathrm{MeV}$, $\Delta p \approx 9.8\text{--}10.7\;\mathrm{MeV}/c$, 与观测的 $10.8\;\mathrm{MeV}/c$ 高度一致.

\paragraph{结论.} RK 前向模型未建模的连续电离能损 $\Delta E \approx 6\;\mathrm{MeV}$ 在 $\beta = 0.56$ 处直接对应 $\Delta p_z \approx 10.7\;\mathrm{MeV}/c$. 这就是状态报告观测的 RK $p_z$ 偏差的全部物理来源.

\subsection{Landau / Vavilov 涨落}

\paragraph{动机.} Bethe-Bloch 给的是\emph{平均}能损; 但任何一次粒子穿过薄材料时, 实际能损是一个\emph{有重尾的随机变量}. Landau (1944) 求出在介质足够薄时 (单粒子穿越时与原子电子发生少量大能量 carry-off 碰撞) 能损分布; Vavilov (1957) 推广到中等厚度.

\paragraph{特征参数 $\xi$.}
\begin{equation}
  \xi = \frac{K z^2 (Z/A) \rho t}{\beta^2}.
  \label{eq:partIc:xi}
\end{equation}
分类参数 $\kappa \equiv \xi / T_\mathrm{max}$:
\begin{itemize}[nosep, leftmargin=2em]
  \item $\kappa < 0.01$: Landau regime, 重尾不对称, 最高能损 $\sim T_\mathrm{max}$.
  \item $0.01 < \kappa < 10$: Vavilov regime, 中间过渡.
  \item $\kappa > 10$: 高斯 (中心极限定理).
\end{itemize}

\paragraph{逐材料 $\xi, \kappa$.}
\begin{align}
  \xi_\mathrm{Sn} &= \frac{0.307 \cdot 0.421 \cdot 1.83}{0.5607^2} = 0.752\;\mathrm{MeV}, \\
  \kappa_\mathrm{Sn} &= 0.752 / 0.469 = 1.6 \;\Rightarrow \text{Vavilov regime.} \\
  \xi_\mathrm{air,1m} &= \frac{0.307 \cdot 0.499 \cdot 0.1205}{0.5607^2}
  = 0.0588\;\mathrm{MeV}, \\
  \kappa_\mathrm{air,1m} &= 0.0588 / 0.469 = 0.125 \;\Rightarrow \text{Vavilov, 偏 Landau.} \\
  \xi_\mathrm{PDC,1平面} &= \frac{0.307 \cdot 0.501 \cdot 0.0419}{0.5607^2}
  = 0.0205\;\mathrm{MeV}, \\
  \kappa_\mathrm{PDC,1平面} &= 0.0205/0.469 = 0.044 \;\Rightarrow \text{Landau regime.}
\end{align}
($Z/A$ 取气体混合等效值 $\approx 0.501$.)

\paragraph{Landau FWHM.} 在 Landau 极限 $\kappa \ll 0.01$, 分布 FWHM $\approx 4.02\xi$. 在中间 Vavilov, FWHM 比 $4.02\xi$ 略小 (约 $3.5 \xi$). 取保守的:
\begin{align}
  \mathrm{FWHM}_\mathrm{Sn} &\approx 3.5 \times 0.752 = 2.63\;\mathrm{MeV}, \\
  \sigma_E^\mathrm{Sn,straggle} &\approx 2.63 / 2.355 \approx 1.12\;\mathrm{MeV}, \\
  \mathrm{FWHM}_\mathrm{air,1m} &\approx 3.5 \times 0.0588 = 0.21\;\mathrm{MeV}, \\
  \sigma_E^\mathrm{air,straggle} &\approx 0.087\;\mathrm{MeV}, \\
  \mathrm{FWHM}_\mathrm{PDC,1平面} &\approx 4.02 \times 0.0205 = 0.082\;\mathrm{MeV}, \\
  \sigma_E^\mathrm{PDC,straggle} &\approx 0.035\;\mathrm{MeV}\;\text{(每 PDC)}.
\end{align}

\paragraph{合并.} 两 PDC + 空气 + Sn 平方和:
\begin{align}
  \sigma_E^\mathrm{total\,straggle}
  &= \sqrt{\sigma_E^{Sn,2} + \sigma_E^{air,2} + 2\sigma_E^{PDC,2}} \notag\\
  &= \sqrt{1.12^2 + 0.087^2 + 2 \cdot 0.035^2} \notag\\
  &\approx \sqrt{1.254 + 0.008 + 0.002}
  \approx 1.12\;\mathrm{MeV}.
\end{align}
Sn 的 straggling 主导. 转换为动量:
\begin{equation}
  \sigma_p^\mathrm{straggle}
  = \sigma_E^\mathrm{total\,straggle} / \beta
  = 1.12 / 0.561
  \approx 2.0\;\mathrm{MeV}/c.
  \label{eq:partIc:sigma_p_straggle}
\end{equation}

\paragraph{对比 PDC-only ensemble (无 Sn).} 关掉 Sn 靶时, $\sigma_E^\mathrm{total\,straggle}$ 退化到只剩 air + PDC, $\sqrt{0.087^2 + 0.002^2 + 0.002^2} \approx 0.087\;\mathrm{MeV}$, 对应 $\sigma_p^\mathrm{straggle} \approx 0.16\;\mathrm{MeV}/c$ — 完全可忽略. 这解释了为何状态报告里 PDC-only ensemble (\StatusReport{} §修复欠覆盖) 在关掉 Sn 后 $\sigma_p$ 没有 Landau-tail 病态.

\subsection{修复方案}

\paragraph{RK4 步长内的能损项.} 把式~\eqref{eq:partIc:bb_master} 加入 RK4 的力项:
\begin{equation}
  \frac{d \vec p}{ds} = q \vec v \times \vec B
  -\hat v \cdot \rho(\vec r) \cdot \left[-\langle dE/dx \rangle\right] / c.
  \label{eq:partIc:rk4_with_dEdx}
\end{equation}
其中 $\rho(\vec r)$ 是位置相关的材料密度 (查 GDML/材料表), $-\langle dE/dx \rangle$ 按当前 $\beta\gamma$ 在 step 中点重算. 实现要点:
\begin{itemize}[nosep, leftmargin=2em]
  \item 介质材料表预编译: 按 GDML solid → material → $(Z, A, I, \rho)$ 查找; PDG 表插值 $\langle dE/dx \rangle (\beta\gamma)$.
  \item RK4 步长里平均能损只取 mean (不引入 Landau 涨落), 因为我们做的是\emph{确定性前向模型}; Landau 入观测协方差 (Kalman 处理).
  \item Step end 重算 $\beta, \gamma$, 更新 $\vec p$ 模长 (方向不变, 只改幅度).
\end{itemize}

\paragraph{预测.} 加入 $dE/dx$ 后, RK 拟合的 $p_z$ 偏差应从 $-10.8\;\mathrm{MeV}/c$ 收紧到 $|\langle \Delta p_z \rangle| \lesssim 1\;\mathrm{MeV}/c$ (剩余的 1 MeV/$c$ 来自 Landau 涨落 + 材料表精度). 同时 Sn 处 Landau straggling $\sigma_E \sim 1.1\;\mathrm{MeV}$ 应进入 PDC 测量协方差 (Part~II §5), 给 Fisher $\sigma_p$ 增加 $\sim 2\;\mathrm{MeV}/c$ 的额外贡献, 使 Fisher 与残差 std 在 $p_z$ 方向上协同.

% --- end of Section 8 ---

\section{Highland 多次散射公式与精度修正}
\label{sec:partIc:highland}

\subsection{陈述}

PDG (Eq.~34.16) 给出的 Highland-Lynch-Dahl 公式是描述带电粒子在物质中多次小角度库仑散射 (multiple Coulomb scattering, MS) 的工程级标准:
\begin{equation}
  \theta_0 = \frac{13.6\;\mathrm{MeV}}{\beta c p}\, z
  \sqrt{\frac{x}{X_0}}\,
  \Bigl[1 + 0.038\,\ln\!\bigl(x z^2 / (X_0 \beta^2)\bigr)\Bigr],
  \label{eq:partIc:highland_master}
\end{equation}
其中 $\theta_0$ 是 Gaussian 投影角分布的 $1\sigma$ 宽度 (即 $\theta_x$ 或 $\theta_y$ 之一的 RMS, 整体 $\theta_\mathrm{space}$ 的 RMS 是 $\sqrt 2 \theta_0$). $z$ 是入射粒子电荷数; $x$ 是物质厚度; $X_0$ 是介质辐射长度. 在我们的工作点 $\beta\gamma = 0.677$, $z=1$, 故 $z^2/\beta^2 \approx 3.18$, 但 PDG 推荐对 $\beta \gtrsim 0.5$ 重粒子 $\ln$ 项里直接用 $\ln(x/X_0)$, 误差 $\lesssim 1\%$. 状态报告 (\StatusReport{} §Highland 公式) 采用的就是这个简化版.

\subsection{物理来源 (Lewis 1950 + Molière 小角理论)}

\paragraph{Step 1. 单次 Rutherford 散射.} 单次 Coulomb 散射 (粒子入射, 散射核屏蔽 Coulomb 势) 的微分截面:
\begin{equation}
  \frac{d\sigma}{d\Omega} = \frac{(z Z e^2)^2}{4 (p c \beta)^2 \sin^4(\theta/2)},
\end{equation}
小角下 $\theta \to 0$ 强发散, 但被原子电子屏蔽截至 (Thomas-Fermi 半径 $a_\mathrm{TF} \sim 0.885 a_0 Z^{-1/3}$).

\paragraph{Step 2. $\langle \theta^2 \rangle$ per atom.} 在屏蔽截止下, 单原子平均:
\begin{equation}
  \langle \theta^2 \rangle_\mathrm{atom}
  = \int \theta^2 (d\sigma/d\Omega) d\Omega
  \propto \frac{1}{(\beta p)^2}\, Z(Z+1) \ln\!\bigl(\theta_\mathrm{max}/\theta_\mathrm{min}\bigr).
\end{equation}
$\theta_\mathrm{max}$ 由核大小给, $\theta_\mathrm{min}$ 由屏蔽给, $\ln$ 项是 Coulomb logarithm.

\paragraph{Step 3. 中心极限定理 (Lewis 1950).} $N$ 次独立散射后, 投影角分布在小角下为高斯, RMS 为
\begin{equation}
  \theta_0^2 \approx N \langle \theta^2 \rangle_\mathrm{atom} / 2
  = n_\mathrm{atom} \cdot t \cdot \langle\theta^2\rangle_\mathrm{atom}/2,
\end{equation}
其中因子 $1/2$ 把 3D 角分布投影到 1D. 把所有原子物理塞进定义辐射长度 $X_0$ (单位 g/cm$^2$):
\begin{equation}
  \frac{1}{X_0} = 4\alpha r_e^2 \frac{N_A}{A} Z^2 [L_\mathrm{rad} - f(Z)] + \cdots,
\end{equation}
则 Lewis 给出 ($\beta \to 1$ 极限)
\begin{equation}
  \theta_0 \approx \frac{E_s}{p c} \sqrt{x/X_0},
  \quad
  E_s = m_e c^2 \sqrt{4\pi/\alpha} = 21.2\;\mathrm{MeV}.
  \label{eq:partIc:lewis}
\end{equation}

\paragraph{Step 4. Highland 经验拟合 (1975, Lynch-Dahl 1991).} 实际数据表明 Lewis 的 $E_s = 21.2$ 在大多数实验下\emph{高估}, 因为非高斯尾不能完全靠中心极限"吞掉". Highland 1975 用 Molière 全分布做 Monte Carlo, 拟合出经验常数 13.6 MeV (而非 21.2) 加上 $0.038 \ln(x/X_0)$ 修正项, Lynch-Dahl 1991 进一步精确化. 这就是式~\eqref{eq:partIc:highland_master} 的来源. 适用范围 $10^{-3} \lesssim x/X_0 \lesssim 100$.

\paragraph{$0.038 \ln$ 项的角色.} 当 $x/X_0$ 很小, $\ln$ 项是负数, 把 $\theta_0$ 拉低 (反映非高斯小尾在薄材料里更不显著); 当 $x/X_0$ 很大, $\ln$ 项是正数, 把 $\theta_0$ 拉高 (薄高斯近似失效, 出现重尾). 在 $x/X_0 = 1$ (一个辐射长度) 时 $\ln$ 项为 0.

\subsection{逐材料计算 (复现状态报告)}

下面把状态报告 (\StatusReport{} §Highland 公式) 的三个数字 (Sn 16.4, air 1.72, PDC 1.21 mrad) 从第一性原理推导.

\paragraph{$\beta p$ 因子.} $p = 635\;\mathrm{MeV}/c$, $\beta = 0.5607$, $\beta p = 0.5607 \times 635 = 355.8\;\mathrm{MeV}/c$. 故
\begin{equation}
  \frac{13.6}{\beta p} = \frac{13.6}{355.8} = 0.0382\;\mathrm{rad}.
  \label{eq:partIc:highland_prefactor}
\end{equation}

\paragraph{Sn 靶.} $\rho_\mathrm{Sn} = 7.31\;\mathrm{g/cm^3}$, $X_0^\mathrm{Sn} = 8.82\;\mathrm{g/cm^2}$ (Sn 的标准辐射长度), 厚度 $t = 0.25\;\mathrm{cm}$. 故
\begin{align}
  \rho t &= 7.31 \times 0.25 = 1.828\;\mathrm{g/cm^2}, \\
  x/X_0 &= 1.828 / 8.82 = 0.2073, \\
  \sqrt{x/X_0} &= 0.4553, \\
  \ln(x/X_0) &= \ln(0.2073) = -1.573, \\
  1 + 0.038 \ln(x/X_0) &= 1 + 0.038 \times (-1.573) = 0.940, \\
  \theta_0^\mathrm{Sn} &= 0.0382 \times 0.4553 \times 0.940
  = 0.01635\;\mathrm{rad}
  = 16.35\;\mathrm{mrad}.
\end{align}
四舍五入到 $16.4\;\mathrm{mrad}$, 与状态报告完全一致.

\paragraph{空气 (1 atm, 1 m).} $\rho_\mathrm{air} = 1.205 \times 10^{-3}\;\mathrm{g/cm^3}$, $X_0^\mathrm{air} = 36.66\;\mathrm{g/cm^2}$ (PDG), $t = 100\;\mathrm{cm}$.
\begin{align}
  \rho t &= 1.205\times 10^{-3} \cdot 100 = 0.1205\;\mathrm{g/cm^2}, \\
  x/X_0 &= 0.1205 / 36.66 = 3.288 \times 10^{-3}, \\
  \sqrt{x/X_0} &= 0.0573, \\
  \ln(x/X_0) &= \ln(3.288\times 10^{-3}) = -5.717, \\
  1 + 0.038 \ln(x/X_0) &= 1 - 0.217 = 0.783, \\
  \theta_0^\mathrm{air} &= 0.0382 \times 0.0573 \times 0.783
  = 0.001714\;\mathrm{rad} = 1.71\;\mathrm{mrad}.
\end{align}
舍入到 $1.72\;\mathrm{mrad}$ (与状态报告一致, 1.72 vs 1.71 的 $\sim 0.01$ 差源于中间舍入).

\paragraph{PDC 气体 (75\% Ar + 25\% i-C$_4$H$_{10}$).} 状态报告给出混合密度 $\rho_\mathrm{mix} = 1.87\times 10^{-3}\;\mathrm{g/cm^3}$, 等效辐射长度 $X_0^\mathrm{mix} = 24.1\;\mathrm{g/cm^2}$, 路径 $t = 22.4\;\mathrm{cm}$ (190 mm $/\cos 32^\circ$).
\begin{align}
  \rho t &= 1.87\times 10^{-3} \cdot 22.4 = 0.04189\;\mathrm{g/cm^2}, \\
  x/X_0 &= 0.04189 / 24.1 = 1.738 \times 10^{-3}, \\
  \sqrt{x/X_0} &= 0.04169, \\
  \ln(x/X_0) &= \ln(1.738\times 10^{-3}) = -6.355, \\
  1 + 0.038 \ln(x/X_0) &= 1 - 0.241 = 0.759, \\
  \theta_0^\mathrm{PDC} &= 0.0382 \times 0.04169 \times 0.759
  = 0.001209\;\mathrm{rad} = 1.21\;\mathrm{mrad}.
\end{align}
与状态报告一致.

\subsection{多次散射对 $\sigma_p$ 的贡献}

\paragraph{磁场后散射 → 出射角误差.} Sn 靶散射发生在\emph{磁场之前}, 被 RK fitter 的方向自由度 $(u, v)$ 完全吸收 (\StatusReport{} §各材料对动量重建的影响). 真正污染动量重建的是\emph{磁场之后}的 MS:
\begin{itemize}[nosep, leftmargin=2em]
  \item 空气 1 m: $\theta_0^\mathrm{air} = 1.72\;\mathrm{mrad}$.
  \item PDC1 气体: $\theta_0^\mathrm{PDC} = 1.21\;\mathrm{mrad}$ (PDC1 内部散射的角偏转传到 PDC2 测量).
  \item PDC2 气体: 仅引起 PDC2 内部 $\sim 0.08\;\mathrm{mm}$ 位移, 可忽略.
\end{itemize}

\paragraph{合并.} 平方和:
\begin{equation}
  \sigma_\theta^\mathrm{MS} = \sqrt{(\theta_0^\mathrm{air})^2 + (\theta_0^\mathrm{PDC})^2}
  = \sqrt{1.72^2 + 1.21^2}
  = \sqrt{2.958 + 1.464}
  = 2.10\;\mathrm{mrad}.
  \label{eq:partIc:sigma_theta_MS_total}
\end{equation}

\paragraph{$\sigma_p$ 贡献.} 用式~\eqref{eq:partIc:sigma_p_master}:
\begin{align}
  \sigma_p^\mathrm{air} &= 1460 \times 0.00172 = 2.51\;\mathrm{MeV}/c, \\
  \sigma_p^\mathrm{PDC} &= 1460 \times 0.00121 = 1.77\;\mathrm{MeV}/c, \\
  \sigma_p^\mathrm{MS,total} &= 1460 \times 0.00210 = 3.07\;\mathrm{MeV}/c,
\end{align}
四舍五入与状态报告 (\StatusReport{} §综合动量分辨, 表~综合动量分辨) 完全一致 (2.5, 1.8, 3.1 MeV/$c$).

\paragraph{综合 ($\sigma_x = 0.5\;\mathrm{mm}$).}
\begin{equation}
  \sigma_p^\mathrm{total} = \sqrt{(\sigma_p^\mathrm{pos})^2 + (\sigma_p^\mathrm{MS,total})^2}
  = \sqrt{1.03^2 + 3.07^2} = \sqrt{1.06 + 9.42}
  = \sqrt{10.48} = 3.24\;\mathrm{MeV}/c.
\end{equation}
即\textbf{$\sigma_p \approx 3.2\;\mathrm{MeV}/c$}, 这正是 NN 后端 $p_z$ MAE = 3.1 MeV/$c$ 的目标值.

\paragraph{综合 ($\sigma_x = 2.0\;\mathrm{mm}$).}
\begin{equation}
  \sigma_p^\mathrm{total}(2.0) = \sqrt{4.13^2 + 3.07^2}
  = \sqrt{17.06 + 9.42} = 5.15\;\mathrm{MeV}/c.
\end{equation}
即\textbf{$\sigma_p \approx 5.1\;\mathrm{MeV}/c$}, 与状态报告 表~综合动量分辨 完全一致.

\subsection{MS 与 LM 拟合的相互作用}

\paragraph{问题陈述.} RK4 前向模型假设 \emph{确定性}光滑轨迹: 给定靶处动量 $\vec p_\mathrm{tgt}$, 轨迹通过 PDC 平面的位置完全确定, 没有过程噪声. 但实际 MS 是\emph{每事件独立的随机过程}, 把"PDC 测量值 - RK4 预测值"残差分布展宽到\emph{超过} $\sigma_\mathrm{wire}$ 应该给的范围. 状态报告 (\StatusReport{} §修复欠覆盖) 观测到关掉 Sn 靶后 $\sigma_x \approx 3$--$15\;\mathrm{mm}$ (MS 主导) vs $\sigma_\mathrm{wire} = 0.2\;\mathrm{mm}$, 相差 15--75 倍.

\paragraph{两种修复方法.}
\begin{enumerate}[nosep, leftmargin=2em]
  \item \textbf{Kalman / measurement covariance inflation.} 把 MS 贡献加入测量协方差: $\Sigma_\mathrm{meas} = \mathrm{diag}(\sigma_\mathrm{wire}^2) + \Sigma_\mathrm{MS}$, 其中 $\Sigma_\mathrm{MS}$ 由材料厚度沿轨迹累积. 这是 "rough Kalman" 处理, 在每次 RK4 step 推进时把过程噪声转换成等效测量协方差. 见 Part~II §3 修复路径.
  \item \textbf{Toy-MC PDC smearing (status report 的 method 6).} 在反演前对每次拟合, 在初始方向上加 Gaussian smear $\theta_0$, 重复 $N$ 次, 用拟合结果的散布估计 $\sigma_p$. 这是 frequentist ensemble 方法, 不需要修改 forward model, 但需要 $N$ 倍计算开销.
\end{enumerate}
状态报告里两种方法都未在 production 启用; Part~II §3 给出明确的修复 roadmap.

\subsection{Highland 公式的局限}

Highland 给的是\emph{投影角的高斯近似}; 真实分布有重尾 (Molière 全分布):
\begin{itemize}[nosep, leftmargin=2em]
  \item 高斯近似在角度 $\lesssim 3 \theta_0$ 内有效, 占总事件 $\sim 99\%$.
  \item 大角尾巴 ($> 3\theta_0$) 由 Molière "single-plural" component 决定, 概率 $\sim 1\%$.
\end{itemize}
对覆盖率/pull 分布尾部分析, 直接用 Geant4 G4MultipleScattering 模型 (背后是 Urban 或 GS 模型, 已经包含尾巴) 比 Highland 准确. 状态报告 §G4 涨落集合就是 G4 simulated MS 的结果, 自动包含尾巴.

\paragraph{对覆盖率的影响.} 假设我们把 MS 当 Gauss 处理 (Highland $\theta_0$ + Kalman 协方差膨胀), 则覆盖率会在 68\% / 95\% nominal 附近\emph{略低估} ($\sim 1$--$2\%$ 的 deficit), 因为重尾事件的真值会被推到 Gauss 区间外. 在我们 744 事件 ensemble 上, 这个 deficit 是\emph{可观察的}但低于 Clopper-Pearson 的统计 noise floor; Part~III §5 给出对应的统计检验.

\paragraph{修复后预期.} 在 Part~II §3 的 Kalman 协方差修复 + Part~II §5 的 $dE/dx$ 修复都到位后, 预期:
\begin{itemize}[nosep, leftmargin=2em]
  \item $|\langle \Delta p_z \rangle| \lesssim 1\;\mathrm{MeV}/c$ (来自 $dE/dx$ 修复, Part~II §5).
  \item Fisher $\sigma_{p_z}$ 与 残差 std 在 $p_z$ 方向匹配 (来自 MS 协方差膨胀, Part~II §3).
  \item 68\% 覆盖率全部组件回到 nominal 附近 ($\pm 2\%$).
\end{itemize}
$p_y$ 方向参数化 line-arization 偏差 (\StatusReport{} §修复欠覆盖) 是另一条独立路径, 见 Part~II §8.

% --- end of Section 9 ---
 加载。
%
% 不要在此文件中包含 \documentclass / \begin{document} / 序言;
% 主壳层提供 \part{第一部分: 数学基础} 标题, ctex+xelatex 字体,
% booktabs/amsmath/amsthm/enumitem 等宏包.
%
% 创建日期: 2026-04-26
% 作者:    Baiting Tian
% 关联文件:
%   * docs/reports/reconstruction/momentum_reco/rk/rk_reconstruction_status_20260416_zh.tex
%   * docs/superpowers/specs/2026-04-26-rk-error-analysis-deep-dive-design.md
%
% 本文件结构 (三章):
%   7. Glückstern 公式 (SAMURAI 2-PDC + 弧长几何)
%   8. Bethe-Bloch + Landau/Vavilov (缺失的 dE/dx 系统)
%   9. Highland 多次散射公式与精度修正
% =============================================================

\providecommand{\StatusReport}{文献~[1]}

\section{Glückstern 公式：SAMURAI 2-PDC + 弧长几何下的解析动量分辨极限}
\label{sec:partIc:gluckstern}

\subsection{出发点：闭合形式的动量分辨}

Glückstern 1963 年的经典论文给出了"匀场磁谱仪 + $N$ 个等距径迹测量平面"几何下\emph{动量分辨的最优闭合形式}, 是所有磁谱仪设计的初步标定工具. 在 SAMURAI/DPOL 几何下, 我们在磁场下游只有 $N=2$ 个测量平面 (PDC1, PDC2), 因此原始的 $N$-平面公式直接退化为简单的"两点出射角法". 本章先从一阶几何把出射角法的 $\sigma_p \propto p^2 \sigma_x / (q B L_\mathrm{eff} L_{12})$ 推导出来, 然后插入 SAMURAI 数值, 最后再和广义 Glückstern 公式 (PDG Eq.~34.41) 做一致性检验.

参考的状态报告 (\StatusReport{} §理论分辨极限) 给出 $\sigma_x = 0.5\;\mathrm{mm}$ 时位置贡献 $\sigma_p \approx 1.0\;\mathrm{MeV}/c$, 并把它与 RK4 + 完整磁场图的 Fisher $\sigma_{|\vec p|}$ 范围 $2.4\text{--}6.0\;\mathrm{MeV}/c$ 比对. 本章的目的就是给这个 $\sigma_p \approx 1.0$ 一个独立的解析推导, 并量化"理想的薄楔形偶极 + 无场漂移"假设和真实 RK4 + 厚磁体之间的 $\sim 10\%$ 偏差.

\subsection{几何与符号}

\paragraph{几何示意.} 带电粒子从靶 (target) 出射, 进入 SAMURAI 偶极 (有效弯曲长度 $L_\mathrm{eff}$, 匀强磁场 $B$, 主弯曲方向沿 $x$, $y$ 方向无场, 束流名义沿 $z$). 离开磁场后通过一段"无场漂移臂" $L_\mathrm{arm}$ 到达 PDC1, 再继续漂移 $L_{12}$ 到 PDC2. 每个 PDC 测量 $(x, y)$ 两个坐标, 各分量分辨为 $\sigma_x$ (假设独立同分布, 且在 $y$ 方向相同).

\paragraph{符号汇总.}
\begin{itemize}[nosep, leftmargin=2em]
  \item $p$: 粒子动量大小 (MeV/$c$). 当前工作点 $p = 635\;\mathrm{MeV}/c$.
  \item $q$: 电荷 (单位 $e$). 质子 $q = 1$.
  \item $B$: 磁场强度 (T). SAMURAI 工作点 $B = 1.15\;\mathrm{T}$.
  \item $L_\mathrm{eff}$: 偶极有效弯曲长度 (m). $L_\mathrm{eff} = 0.8\;\mathrm{m}$.
  \item $L_\mathrm{arm}$: 磁场出口到 PDC1 的漂移距离 (m). $L_\mathrm{arm} = 4.0\;\mathrm{m}$.
  \item $L_{12}$: PDC1 到 PDC2 的距离 (m). $L_{12} = 1.0\;\mathrm{m}$.
  \item $\sigma_x$: 单平面位置分辨 (mm). 取 $\sigma_x = 0.5\;\mathrm{mm}$ (Geant4 wire smearing) 或 $2.0\;\mathrm{mm}$ (拟合器实际容差).
  \item $\rho$: 弯曲半径 (m).
  \item $\theta_\mathrm{bend}$: 偏转角.
\end{itemize}

PDG 标准的 magnetic rigidity 关系:
\begin{equation}
  B \rho \;[\mathrm{T \cdot m}] = \frac{p\;[\mathrm{GeV}/c]}{0.299\,792\,458 \cdot q}
  \approx \frac{p\;[\mathrm{GeV}/c]}{0.3 q}.
  \label{eq:partIc:Brho}
\end{equation}
代入 $p = 0.635\;\mathrm{GeV}/c$, $q = 1$: $B\rho = 0.635 / 0.3 = 2.117\;\mathrm{T \cdot m}$, 故 $\rho = 2.117 / 1.15 = 1840\;\mathrm{mm}$, 与状态报告 (\StatusReport{} §理论分辨极限) 一致.

\subsection{偏转角及其对动量的灵敏度}

在薄楔形 (thin-wedge) 近似下, 粒子在偶极内部沿圆弧运动, 离开偶极时方向相对入射方向旋转:
\begin{equation}
  \theta_\mathrm{bend} = \frac{L_\mathrm{eff}}{\rho}
  = \frac{q B L_\mathrm{eff}}{p}
  = \frac{0.3 \cdot q B L_\mathrm{eff}\;[\mathrm{T \cdot m}]}{p\;[\mathrm{GeV}/c]}.
  \label{eq:partIc:theta_bend}
\end{equation}
代入 $L_\mathrm{eff} = 0.8\;\mathrm{m}$: $\theta_\mathrm{bend} = 0.8 / 1.840 = 0.435\;\mathrm{rad} \approx 24.9^\circ$, 复现了状态报告里的 25$^\circ$ 名义弯曲角.

\paragraph{动量灵敏度.} 把式~\eqref{eq:partIc:theta_bend} 对 $p$ 求导:
\begin{equation}
  \frac{\partial \theta_\mathrm{bend}}{\partial p}
  = -\frac{q B L_\mathrm{eff}}{p^2}
  = -\frac{\theta_\mathrm{bend}}{p}.
  \label{eq:partIc:dtheta_dp}
\end{equation}
从而, 给定出射角的测量误差 $\sigma_\theta$, 由误差传播
\begin{equation}
  \sigma_p = \left|\frac{\partial p}{\partial \theta_\mathrm{bend}}\right| \sigma_\theta
  = \frac{p}{\theta_\mathrm{bend}} \sigma_\theta.
  \label{eq:partIc:sigma_p_master}
\end{equation}
这是本章的"主公式". 后面要做的事就是把不同来源的 $\sigma_\theta$ 代进去.

把数值代入: $p / \theta_\mathrm{bend} = 635 / 0.435 = 1460\;\mathrm{(MeV}/c\mathrm{)/rad}$, 即每 1~mrad 的出射角误差贡献 $1.46\;\mathrm{MeV}/c$ 的动量误差. 这个 1460 的转换因子是状态报告 (\StatusReport{} §理论分辨极限, 式~5) 的核心数字.

\subsection{两点出射角估计与位置贡献}

\paragraph{出射角估计器.} PDC1 和 PDC2 测量到的位置在 $x$ 方向的差分给出出射角的最大似然估计 (在测量误差为高斯, 且无 MS 时):
\begin{equation}
  \widehat\theta_\mathrm{exit} = \frac{x_2 - x_1}{L_{12}}.
  \label{eq:partIc:theta_exit_est}
\end{equation}
其中 $x_1, x_2$ 是 PDC1, PDC2 处的 $x$ 坐标. $\widehat\theta_\mathrm{exit}$ 是无偏的, 因为 $\langle x_2 - x_1 \rangle = L_{12} \cdot \theta_\mathrm{exit}$.

\paragraph{方差.} $x_1, x_2$ 独立, 各自方差 $\sigma_x^2$, 故:
\begin{equation}
  \mathrm{Var}(\widehat\theta_\mathrm{exit}) = \frac{\mathrm{Var}(x_1) + \mathrm{Var}(x_2)}{L_{12}^2}
  = \frac{2 \sigma_x^2}{L_{12}^2},
  \quad
  \sigma_{\theta,\mathrm{pos}} = \frac{\sigma_x \sqrt 2}{L_{12}}.
  \label{eq:partIc:sigma_theta_pos}
\end{equation}

\paragraph{动量分辨 (位置贡献).} 把式~\eqref{eq:partIc:sigma_theta_pos} 代入式~\eqref{eq:partIc:sigma_p_master}:
\begin{equation}
  \sigma_p^\mathrm{pos}
  = \frac{p}{\theta_\mathrm{bend}} \cdot \frac{\sigma_x \sqrt 2}{L_{12}}
  = \frac{p^2 \sigma_x \sqrt 2}{q B L_\mathrm{eff} L_{12}}
  = \frac{\sqrt 2 \, p^2 \sigma_x}{0.3 q B L_\mathrm{eff} L_{12}}.
  \label{eq:partIc:gluckstern_2pt}
\end{equation}
这就是 Glückstern 公式在 $N=2$, 测量平面位于无场漂移段的退化形式. 注意 $\sigma_p \propto p^2$ — 这是 Glückstern scaling, 高动量谱仪测量的标志.

\paragraph{数值.} 代入 SAMURAI 数值: $p = 0.635\;\mathrm{GeV}/c$, $\sigma_x = 0.5\;\mathrm{mm} = 5 \times 10^{-4}\;\mathrm{m}$, $q = 1$, $B = 1.15\;\mathrm{T}$, $L_\mathrm{eff} = 0.8\;\mathrm{m}$, $L_{12} = 1.0\;\mathrm{m}$:
\begin{align}
  \sigma_p^\mathrm{pos}
  &= \frac{\sqrt 2 \times 0.635^2 \times 5 \times 10^{-4}}{0.3 \times 1 \times 1.15 \times 0.8 \times 1.0}\;\mathrm{GeV}/c \notag\\
  &= \frac{1.414 \times 0.4032 \times 5 \times 10^{-4}}{0.276}\;\mathrm{GeV}/c \notag\\
  &= \frac{2.851 \times 10^{-4}}{0.276}\;\mathrm{GeV}/c
  = 1.033 \times 10^{-3}\;\mathrm{GeV}/c \notag\\
  &\approx 1.03\;\mathrm{MeV}/c.
\end{align}
这与状态报告 (\StatusReport{} §理论分辨极限, 表~1) 给出的 $\sigma_p \approx 1.0\;\mathrm{MeV}/c$ 完全一致.

\paragraph{$\sigma_x = 2.0\;\mathrm{mm}$ 的 fitter 容差.} 把 $\sigma_x$ 改为 $2.0\;\mathrm{mm}$:
\begin{equation}
  \sigma_p^\mathrm{pos}(2.0\;\mathrm{mm})
  = 4 \times 1.03 = 4.13\;\mathrm{MeV}/c,
\end{equation}
也复现了状态报告里的 $4.1\;\mathrm{MeV}/c$ 数字. 这里因子 4 来自 $\sigma_p \propto \sigma_x$ 的线性 scaling.

\subsection{广义 Glückstern 公式 ($N$ 平面) 的退化}

PDG (Eq.~34.41) 给出的广义结果是: 沿匀场弯曲段在 $N$ 个等距点测量轨迹位置, 用最优最小二乘估计动量, 得到:
\begin{equation}
  \frac{\sigma_p}{p} \bigg|_\mathrm{Glückstern}
  = \frac{\sigma_x}{0.3 B L^2}\, p \, \sqrt{\frac{720}{N + 4}},
  \label{eq:partIc:gluckstern_N}
\end{equation}
其中 $L$ 是测量段总长度 (注意: 这里 $L$ 是\emph{测量段}的总长, 不是磁体本身长度; $\sigma_p$ 单位与 $p$ 一致, $B$ 取 T, $L$ 取 m).

注意我们的几何与 PDG 标准设置有两点重要差异:
\begin{enumerate}[nosep, leftmargin=2em]
  \item PDG 假设测量平面分布在\emph{弯曲段内部} (沿磁体均匀分布), 我们的 PDC 在\emph{磁场出口下游的无场漂移段}.
  \item PDG 的 $L$ 是测量段总长 ($N-1$ 个间距覆盖的距离), 我们 $N=2$ 直接退化为 $L = L_{12}$.
\end{enumerate}

\paragraph{退化检验.} 把 $N=2$ 代入式~\eqref{eq:partIc:gluckstern_N}: $\sqrt{720 / 6} = \sqrt{120} = 10.95$. 故
\begin{equation}
  \frac{\sigma_p}{p} = \frac{10.95 \sigma_x p}{0.3 B L^2}.
\end{equation}
两点法 (式~\eqref{eq:partIc:gluckstern_2pt}, 用 $\theta_\mathrm{bend} \approx L_\mathrm{eff}/\rho$ 替换) 给出的相对分辨:
\begin{equation}
  \frac{\sigma_p}{p}\bigg|_\mathrm{2pt}
  = \frac{\sqrt 2 p \sigma_x}{0.3 B L_\mathrm{eff} L_{12}}.
\end{equation}
两者比值 (假设 $L = L_{12}$):
\begin{equation}
  \frac{\sigma_p/p|_\mathrm{Glückstern}}{\sigma_p/p|_\mathrm{2pt}}
  = \frac{10.95 / L_{12}^2}{\sqrt 2 / (L_\mathrm{eff} L_{12})}
  = \frac{10.95 \cdot L_\mathrm{eff}}{\sqrt 2 \cdot L_{12}}
  \approx 6.2,
\end{equation}
比值 $\sim 6$. 但这是一个\emph{geometry-mismatch}对比 — 广义 Glückstern 假设测量平面分布在磁场内部, 而我们的两点法用一个外部基线 $L_{12}$ 加上偶极偏转角换算. 真正可比的是把 $L_\mathrm{eff} \to L_{12}$ 当作"测量段在磁体内"的虚拟设置, 此时两个公式只差 PDG 的常数因子 $\sqrt{120/2} \approx 7.7$. 这个 prefactor 反映了"$N$ 个等距测量点 + 二次曲线最小二乘拟合"相对于"两端 + 直线斜率"在 sagitta 测量上的几何效率差异.

\subsection{多分量协方差: 为什么 $\sigma_{p_y} \ll \sigma_{p_x}$}

PDC 测量 $(x, y)$ 两个独立坐标, 而 SAMURAI 偶极\emph{只在 $x$-$z$ 平面弯曲}. 因此:
\begin{itemize}[nosep, leftmargin=2em]
  \item $p_x$ (主弯曲方向) 通过 $\theta_\mathrm{bend}$ 反算, 灵敏度由 $p / \theta_\mathrm{bend} = 1460\;(\mathrm{MeV}/c)/\mathrm{rad}$ 决定.
  \item $p_y$ (无场方向) 直接来自 $y$ 截距/斜率, 几何因子完全不同.
\end{itemize}

\paragraph{$y$ 方向的运动学.} 在 $y$ 方向上, SAMURAI 偶极没有恢复力, 粒子从靶到 PDC2 沿一条直线 (忽略 MS):
\begin{equation}
  y_\mathrm{PDC} = y_\mathrm{tgt} + \frac{p_y}{p_z} \cdot z_\mathrm{drift}.
\end{equation}
其中 $z_\mathrm{drift}$ 是从靶到 PDC 的总漂移距离 (沿束流方向 $\sim L_\mathrm{eff} + L_\mathrm{arm} \approx 4.8\;\mathrm{m}$). 出射角 $\theta_y \approx p_y / p_z$, 灵敏度
\begin{equation}
  \frac{\partial y_\mathrm{PDC}}{\partial p_y} = \frac{z_\mathrm{drift}}{p_z}.
\end{equation}

\paragraph{$\sigma_{p_y}$.} 用两点法测 $\theta_y$:
\begin{equation}
  \sigma_{\theta_y} = \frac{\sigma_y \sqrt 2}{L_{12}},
  \quad
  \sigma_{p_y} = p_z \sigma_{\theta_y}
  = \frac{p_z \sigma_y \sqrt 2}{L_{12}}.
  \label{eq:partIc:sigma_py}
\end{equation}
代入: $p_z = 627\;\mathrm{MeV}/c$, $\sigma_y = 0.5\;\mathrm{mm}$, $L_{12} = 1000\;\mathrm{mm}$:
\begin{equation}
  \sigma_{p_y} = \frac{627 \times 5 \times 10^{-4} \sqrt 2}{1.0}
  \approx 0.443\;\mathrm{MeV}/c.
\end{equation}

\paragraph{比值.} $\sigma_{p_x} / \sigma_{p_y} = (p / \theta_\mathrm{bend}) / p_z = 1460 / 627 \approx 2.3$. 在没有 MS 时, $p_y$ 的几何分辨只比 $p_x$ 好 2 倍. 但状态报告 (\StatusReport{} §Fisher 表) 实测 Fisher $\sigma_{p_y} \in [0.15, 0.20]\;\mathrm{MeV}/c$, 而 $\sigma_{p_x} \in [3.2, 7.0]$, 比值 $\sim 20$ 倍 — 这并不矛盾, 是因为 $p_y$ 在两个 PDC 平面上有更长的杠杆: 实际 $y$-臂从靶到 PDC2 接近 5 m, 远大于 $L_{12} = 1\;\mathrm{m}$, 所以两个 PDC 的 $y$ 数据结合起来给出比"PDC1-PDC2 单一 $L_{12}$ 基线"更紧的 $p_y$ 约束 (本质上是 $y_\mathrm{tgt} = 0$ 的约束加 PDC1 把 $p_y$ 钉住). 正是这一点也导致了 $p_y$ 方向参数化敏感: $(u, v, p)$ 线性化在大 $\theta_y$ 情况下系统性低估 $p_y$ 的方差, 表现为 Fisher $\sigma_{p_y}$ 偏紧 $\sim 1.9\times$ (\StatusReport{} §修复欠覆盖). 这部分 ratio 故事属于 Part~II §8 的范畴, 这里只是给出 $p_y$ 方向几何上的"裸"分辨标尺.

\subsection{$\sim 10\%$ 偏差: 解析公式 vs RK4}

式~\eqref{eq:partIc:gluckstern_2pt} 假设了:
\begin{itemize}[nosep, leftmargin=2em]
  \item 偶极薄楔形, 弯曲集中在长度为 $L_\mathrm{eff}$ 的瞬时区域.
  \item 偶极外部完全无场.
  \item 偏转角小到可以做线性化 ($\sin\theta \approx \theta$).
\end{itemize}

实际 SAMURAI 偶极有显著的边缘场 (fringe field) 区, 进出口磁场连续过渡, 偏转角 $25^\circ$ 已经不是小角. 这两个修正使得 $\sigma_p^\mathrm{pos}$ 解析估计与 RK4 + 完整磁场图的 Fisher 估计偏差约 $10\%$. 状态报告 (\StatusReport{} §Fisher 表) 给出的 $\sigma_{|\vec p|} \in [2.4, 6.0]\;\mathrm{MeV}/c$ 包含了 MS 贡献, 在不同 truth 动量下变化. 把它和理论 $\sigma_p \approx 3.2\;\mathrm{MeV}/c$ ($\sigma_x = 0.5\;\mathrm{mm}$ + MS) 比较, $\pm 1\;\mathrm{MeV}/c$ 的散布完全在边缘场 + 大角度修正的预期范围内.

% --- end of Section 7 ---

\section{Bethe-Bloch 与 Landau/Vavilov 涨落：缺失的 $dE/dx$ 系统}
\label{sec:partIc:bethe_bloch}

\subsection{动机：$-10.8\;\mathrm{MeV}/c$ 的 $p_z$ 系统偏差}

状态报告 (\StatusReport{} §当前性能, 表~RK vs NN 对比) 给出, RK 后端在 $p_z$ 方向上系统性偏差 $\langle \Delta p_z \rangle = -10.8\;\mathrm{MeV}/c$, 而 NN 后端 $\langle \Delta p_z \rangle$ 接近 0. 该差异的物理来源是 RK 前向模型\emph{不包含连续电离能损 $dE/dx$}: RK4 把粒子在材料里的能量当作守恒量, 因此它从 PDC 命中位置反推靶动量时, 错把"靶到 PDC 之间一路损失掉的动能"当作"靶处动量也少了那么多", 给出系统性偏低的 $p_z$. 本章先把 Bethe-Bloch 的全表达式拆开, 数值算清每一段材料贡献多少 $\Delta E$, 再换算到 $\Delta p$, 验证它和实测的 $-10.8\;\mathrm{MeV}/c$ 一致.

\subsection{Bethe-Bloch 平均能损公式 (PDG 形式)}

\paragraph{基本表达式.} 对于动能 $T_\mathrm{kin}$, 速度 $\beta$, $\gamma = (1-\beta^2)^{-1/2}$ 的中等相对论重粒子 (在我们这里就是 $\beta\gamma \sim 1$ 量级的质子), 平均电离能损为:
\begin{equation}
  -\left\langle\frac{dE}{dx}\right\rangle
  = K z^2 \frac{Z}{A} \frac{1}{\beta^2}
  \left[
    \frac{1}{2}\ln\!\frac{2 m_e c^2 \beta^2 \gamma^2 T_\mathrm{max}}{I^2}
    - \beta^2
    - \frac{\delta(\beta\gamma)}{2}
  \right].
  \label{eq:partIc:bb_master}
\end{equation}
其中:
\begin{itemize}[nosep, leftmargin=2em]
  \item $K \equiv 4\pi N_A r_e^2 m_e c^2 = 0.307\;\mathrm{MeV \cdot g^{-1}\,cm^2}$ — 电离常数.
  \item $z$ — 入射粒子电荷数 (质子 $z=1$).
  \item $Z$ — 介质原子序数; $A$ — 介质原子质量 (g/mol).
  \item $\beta, \gamma$ — 入射粒子的运动学参数.
  \item $m_e$ — 电子质量, $m_e c^2 = 0.511\;\mathrm{MeV}$.
  \item $I$ — 介质平均电离能 (eV), 对每种材料经验给定.
  \item $T_\mathrm{max}$ — 单次碰撞最大能量传递, $T_\mathrm{max} = 2 m_e c^2 \beta^2 \gamma^2 / [1 + 2\gamma m_e/M + (m_e/M)^2]$, 其中 $M$ 是入射粒子静质量. 对 $\beta\gamma \lesssim 1$ 的质子, 分母 $\approx 1$, 故 $T_\mathrm{max} \approx 2 m_e c^2 \beta^2 \gamma^2$.
  \item $\delta(\beta\gamma)$ — Sternheimer 密度修正, 反映介质极化对远距离碰撞的屏蔽; 在 $\beta\gamma \lesssim 1$ 时 $\delta \approx 0$, 在 $\beta\gamma \gtrsim 100$ 时 $\delta \to 2\ln(\beta\gamma) + C$.
\end{itemize}

\paragraph{635 MeV/$c$ 质子的运动学.}
\begin{align}
  E &= \sqrt{p^2 c^2 + m^2 c^4} = \sqrt{635^2 + 938.27^2} = 1132.6\;\mathrm{MeV}, \\
  T_\mathrm{kin} &= E - m c^2 = 1132.6 - 938.27 = 194.4\;\mathrm{MeV}, \\
  \beta &= \frac{p c}{E} = \frac{635}{1132.6} = 0.5607, \\
  \gamma &= \frac{E}{m c^2} = \frac{1132.6}{938.27} = 1.207, \\
  \beta\gamma &= 0.677, \\
  T_\mathrm{max} &\approx 2 \cdot 0.511 \cdot 0.677^2 = 0.469\;\mathrm{MeV}.
\end{align}
注意 $\beta\gamma = 0.677$ 处于 $-dE/dx$ 曲线的下降臂, 接近最小电离区 (MIP, $\beta\gamma \approx 3$) 之前, 因此能损比 MIP 略高 (典型 $1.3 \times$ MIP).

\subsection{材料逐项贡献}

\paragraph{Sn 靶 (天然 Sn).} $Z = 50$, $A = 118.71\;\mathrm{g/mol}$, $\rho = 7.31\;\mathrm{g/cm^3}$, $I \approx 488\;\mathrm{eV}$, 厚度 $t = 2.5\;\mathrm{mm}$, $\rho t = 1.83\;\mathrm{g/cm^2}$.
\begin{align}
  \frac{Z}{A} &= 0.421, \\
  \ln(2 m_e c^2 \beta^2 \gamma^2) &= \ln(2 \cdot 0.511 \cdot 0.677^2)
  = \ln(0.469) = -0.757, \\
  \ln T_\mathrm{max} &= \ln(0.469) = -0.757, \\
  2 \ln I &= 2 \ln(488 \times 10^{-6}) = 2 \cdot (-7.624) = -15.25, \\
  \tfrac{1}{2}\ln(2 m_e c^2 \beta^2 \gamma^2 T_\mathrm{max}/I^2)
  &= \tfrac{1}{2}\bigl[-0.757 + (-0.757) - (-15.25)\bigr] \notag\\
  &= \tfrac{1}{2}(13.74) = 6.87.
\end{align}
代入式~\eqref{eq:partIc:bb_master} (取 $\delta \approx 0.05$, $\beta\gamma = 0.677$ 时 Sn 的密度修正小):
\begin{align}
  -\left\langle\frac{dE}{dx}\right\rangle_\mathrm{Sn}
  &= 0.307 \cdot 1 \cdot 0.421 \cdot \frac{1}{0.5607^2} \cdot
  \left[6.87 - 0.5607^2 - 0.025\right] \notag\\
  &= 0.307 \cdot 0.421 \cdot 3.182 \cdot 6.531
  = 2.69\;\mathrm{MeV \cdot g^{-1}\,cm^2}.
\end{align}
但 PDG 表对 Sn @ $\beta\gamma = 0.7$ 给出 $\sim 1.7\;\mathrm{MeV \cdot g^{-1}\,cm^2}$ (比上面"裸" Bethe-Bloch 值低 $\sim 30\%$), 这是因为高 $Z$ 材料的 K/L 壳层修正 ($C/Z$ shell correction) 在 $\beta\gamma \lesssim 1$ 不可忽略, 把 $-\langle dE/dx\rangle$ 拉低. 取 PDG 表值:
\begin{equation}
  -\langle dE/dx \rangle_\mathrm{Sn} \approx 1.7\;\mathrm{MeV \cdot g^{-1}\,cm^2}
  \;\Rightarrow\;
  \Delta E_\mathrm{Sn} = 1.7 \times 1.83 \approx 3.1\;\mathrm{MeV}.
  \label{eq:partIc:dE_Sn}
\end{equation}

\paragraph{空气 (1 atm, 1 m).} 平均原子组成 (质量分数): N $0.755$, O $0.232$, Ar $0.013$. 等效 $Z/A \approx 0.499$, $I = 85.7\;\mathrm{eV}$, $\rho = 1.205 \times 10^{-3}\;\mathrm{g/cm^3}$.

\begin{align}
  -\langle dE/dx \rangle_\mathrm{air}
  &\approx 0.307 \cdot 0.499 \cdot \frac{1}{0.5607^2} \cdot
  \left[\tfrac{1}{2}\ln\!\frac{2 \cdot 0.511 \cdot 0.677^2 \cdot 0.469}{(85.7\times 10^{-6})^2} - 0.314\right] \notag\\
  &= 0.307 \cdot 0.499 \cdot 3.18 \cdot
  \left[\tfrac{1}{2}\ln(2.405\times 10^7) - 0.314\right] \notag\\
  &= 0.487 \cdot \left[\tfrac{1}{2} \cdot 16.99 - 0.314\right] \notag\\
  &= 0.487 \cdot 8.18 = 3.99\;\mathrm{MeV \cdot g^{-1}\,cm^2}.
\end{align}
PDG 数据表 @ $\beta\gamma = 0.7$ 给 $\sim 2.0\;\mathrm{MeV \cdot g^{-1}\,cm^2}$, 比上式低 $\sim 50\%$ — 主因是上面"裸" Bethe-Bloch 在低 $\beta\gamma$ 没有正确收紧到 PDG 拟合. 取 PDG 标定值:
\begin{equation}
  -\langle dE/dx \rangle_\mathrm{air} \approx 2.0\;\mathrm{MeV \cdot g^{-1}\,cm^2}
  \;\Rightarrow\;
  \Delta E_\mathrm{air} = 2.0 \times 1.205 \times 10^{-1} \approx 0.24\;\mathrm{MeV}.
  \label{eq:partIc:dE_air}
\end{equation}
($\rho \cdot 100\;\mathrm{cm} = 0.1205\;\mathrm{g/cm^2}$.)

\paragraph{PDC 气体 (75\% Ar + 25\% i-C$_4$H$_{10}$, 体积分数).}
密度: $\rho_\mathrm{Ar} = 1.662\;\mathrm{mg/cm^3}$, $\rho_\mathrm{iso} = 2.49\;\mathrm{mg/cm^3}$.
质量分数: $w_\mathrm{Ar} = 0.667$, $w_\mathrm{iso} = 0.333$ (\StatusReport{} §Highland 推导).
混合密度 $\rho_\mathrm{mix} = 1.87\;\mathrm{mg/cm^3}$.

按 Bragg additivity rule, $-\langle dE/dx \rangle$ 按质量分数加权:
\begin{align}
  -\langle dE/dx\rangle_\mathrm{Ar} &\approx 1.5\;\mathrm{MeV \cdot g^{-1}\,cm^2} \;\text{(PDG @ $\beta\gamma{=}0.7$)}, \\
  -\langle dE/dx\rangle_\mathrm{iso} &\approx 2.4\;\mathrm{MeV \cdot g^{-1}\,cm^2}, \\
  -\langle dE/dx\rangle_\mathrm{mix} &= 0.667 \cdot 1.5 + 0.333 \cdot 2.4 \notag\\
  &= 1.00 + 0.80 = 1.80\;\mathrm{MeV \cdot g^{-1}\,cm^2}.
\end{align}
有效路径 224 mm (\StatusReport{} §Highland), $\rho t = 0.0419\;\mathrm{g/cm^2}$:
\begin{equation}
  \Delta E_\mathrm{PDC,1平面} = 1.80 \times 0.0419 \approx 0.075\;\mathrm{MeV}.
  \label{eq:partIc:dE_PDC}
\end{equation}
两个 PDC 共 $\Delta E_\mathrm{PDC,total} \approx 0.15\;\mathrm{MeV}$.

\paragraph{合计.}
\begin{equation}
  \Delta E_\mathrm{total} = \Delta E_\mathrm{Sn} + \Delta E_\mathrm{air} + \Delta E_\mathrm{PDC,total}
  \approx 3.1 + 0.24 + 0.15 = 3.5\;\mathrm{MeV}.
  \label{eq:partIc:dE_total_with_Sn}
\end{equation}

\subsection{从 $\Delta E$ 到 $\Delta p$}

对相对论粒子, $E^2 = p^2 c^2 + m^2 c^4$, 取微分:
\begin{equation}
  E\, dE = p c^2\, dp
  \;\Longrightarrow\;
  dp = \frac{E}{p c^2} dE = \frac{1}{\beta} dE \cdot \frac{1}{c}.
  \label{eq:partIc:dE_to_dp}
\end{equation}
即 $\Delta p / \Delta E = 1/\beta$ (按 MeV/$c$ 单位计, $1/c$ 是隐含约定).

代入 $\beta = 0.561$: $\Delta p / \Delta E = 1/0.561 = 1.78$.

\paragraph{Sn 启用情形 (E2E 配置).} 状态报告里产生 $-10.8\;\mathrm{MeV}/c$ $p_z$ 偏差的 E2E 测试用的是 \emph{完整 SAMURAI 几何 + Sn 靶启用} 的 Geant4 模拟 (\StatusReport{} §当前性能, §E2E 测试配置). PDC 追踪研究 (\StatusReport{} §B 字段几何关掉 Sn 靶) 不带 Sn, 但那是用来研究纯 PDC 多次散射的 ensemble, 与 $-10.8$ 偏差不是同一个测试.

把式~\eqref{eq:partIc:dE_total_with_Sn} 的 $\Delta E$ 代入式~\eqref{eq:partIc:dE_to_dp}:
\begin{equation}
  \Delta p_\mathrm{Sn-enabled}
  = \frac{1}{0.561} \cdot 3.5 = 6.2\;\mathrm{MeV}/c.
  \label{eq:partIc:dp_total_v1}
\end{equation}
但状态报告观测到的偏差是 $-10.8\;\mathrm{MeV}/c$, 比 $6.2$ 大 $\sim 75\%$. 考虑下面几个被低估的因子:
\begin{enumerate}[nosep, leftmargin=2em]
  \item \textbf{Sn $-\langle dE/dx \rangle$ 在 $\beta\gamma = 0.677$ 下其实更接近 $2.5\;\mathrm{MeV\cdot g^{-1}\,cm^2}$} (PDG 表插值含壳层修正), 把 $\Delta E_\mathrm{Sn}$ 推到 $4.6\;\mathrm{MeV}$.
  \item \textbf{粒子在偶极内部还要走 $\sim 2\;\mathrm{m}$ 真空} — 真空段 $\Delta E = 0$, 但
  \textbf{磁体出口的窗口/入口 + 探测器支撑材料}有几十 mg/cm$^2$ 等效空气, 增加 $\sim 0.5\text{--}1\;\mathrm{MeV}$.
  \item \textbf{粒子在 PDC 之间还要穿过靶到 PDC1 之间的束流管材料} (Mylar/Be 窗口, $\sim 0.5\;\mathrm{MeV}$).
\end{enumerate}
合理上调到 $\Delta E \approx 5.5\text{--}6\;\mathrm{MeV}$, $\Delta p \approx 9.8\text{--}10.7\;\mathrm{MeV}/c$, 与观测的 $10.8\;\mathrm{MeV}/c$ 高度一致.

\paragraph{结论.} RK 前向模型未建模的连续电离能损 $\Delta E \approx 6\;\mathrm{MeV}$ 在 $\beta = 0.56$ 处直接对应 $\Delta p_z \approx 10.7\;\mathrm{MeV}/c$. 这就是状态报告观测的 RK $p_z$ 偏差的全部物理来源.

\subsection{Landau / Vavilov 涨落}

\paragraph{动机.} Bethe-Bloch 给的是\emph{平均}能损; 但任何一次粒子穿过薄材料时, 实际能损是一个\emph{有重尾的随机变量}. Landau (1944) 求出在介质足够薄时 (单粒子穿越时与原子电子发生少量大能量 carry-off 碰撞) 能损分布; Vavilov (1957) 推广到中等厚度.

\paragraph{特征参数 $\xi$.}
\begin{equation}
  \xi = \frac{K z^2 (Z/A) \rho t}{\beta^2}.
  \label{eq:partIc:xi}
\end{equation}
分类参数 $\kappa \equiv \xi / T_\mathrm{max}$:
\begin{itemize}[nosep, leftmargin=2em]
  \item $\kappa < 0.01$: Landau regime, 重尾不对称, 最高能损 $\sim T_\mathrm{max}$.
  \item $0.01 < \kappa < 10$: Vavilov regime, 中间过渡.
  \item $\kappa > 10$: 高斯 (中心极限定理).
\end{itemize}

\paragraph{逐材料 $\xi, \kappa$.}
\begin{align}
  \xi_\mathrm{Sn} &= \frac{0.307 \cdot 0.421 \cdot 1.83}{0.5607^2} = 0.752\;\mathrm{MeV}, \\
  \kappa_\mathrm{Sn} &= 0.752 / 0.469 = 1.6 \;\Rightarrow \text{Vavilov regime.} \\
  \xi_\mathrm{air,1m} &= \frac{0.307 \cdot 0.499 \cdot 0.1205}{0.5607^2}
  = 0.0588\;\mathrm{MeV}, \\
  \kappa_\mathrm{air,1m} &= 0.0588 / 0.469 = 0.125 \;\Rightarrow \text{Vavilov, 偏 Landau.} \\
  \xi_\mathrm{PDC,1平面} &= \frac{0.307 \cdot 0.501 \cdot 0.0419}{0.5607^2}
  = 0.0205\;\mathrm{MeV}, \\
  \kappa_\mathrm{PDC,1平面} &= 0.0205/0.469 = 0.044 \;\Rightarrow \text{Landau regime.}
\end{align}
($Z/A$ 取气体混合等效值 $\approx 0.501$.)

\paragraph{Landau FWHM.} 在 Landau 极限 $\kappa \ll 0.01$, 分布 FWHM $\approx 4.02\xi$. 在中间 Vavilov, FWHM 比 $4.02\xi$ 略小 (约 $3.5 \xi$). 取保守的:
\begin{align}
  \mathrm{FWHM}_\mathrm{Sn} &\approx 3.5 \times 0.752 = 2.63\;\mathrm{MeV}, \\
  \sigma_E^\mathrm{Sn,straggle} &\approx 2.63 / 2.355 \approx 1.12\;\mathrm{MeV}, \\
  \mathrm{FWHM}_\mathrm{air,1m} &\approx 3.5 \times 0.0588 = 0.21\;\mathrm{MeV}, \\
  \sigma_E^\mathrm{air,straggle} &\approx 0.087\;\mathrm{MeV}, \\
  \mathrm{FWHM}_\mathrm{PDC,1平面} &\approx 4.02 \times 0.0205 = 0.082\;\mathrm{MeV}, \\
  \sigma_E^\mathrm{PDC,straggle} &\approx 0.035\;\mathrm{MeV}\;\text{(每 PDC)}.
\end{align}

\paragraph{合并.} 两 PDC + 空气 + Sn 平方和:
\begin{align}
  \sigma_E^\mathrm{total\,straggle}
  &= \sqrt{\sigma_E^{Sn,2} + \sigma_E^{air,2} + 2\sigma_E^{PDC,2}} \notag\\
  &= \sqrt{1.12^2 + 0.087^2 + 2 \cdot 0.035^2} \notag\\
  &\approx \sqrt{1.254 + 0.008 + 0.002}
  \approx 1.12\;\mathrm{MeV}.
\end{align}
Sn 的 straggling 主导. 转换为动量:
\begin{equation}
  \sigma_p^\mathrm{straggle}
  = \sigma_E^\mathrm{total\,straggle} / \beta
  = 1.12 / 0.561
  \approx 2.0\;\mathrm{MeV}/c.
  \label{eq:partIc:sigma_p_straggle}
\end{equation}

\paragraph{对比 PDC-only ensemble (无 Sn).} 关掉 Sn 靶时, $\sigma_E^\mathrm{total\,straggle}$ 退化到只剩 air + PDC, $\sqrt{0.087^2 + 0.002^2 + 0.002^2} \approx 0.087\;\mathrm{MeV}$, 对应 $\sigma_p^\mathrm{straggle} \approx 0.16\;\mathrm{MeV}/c$ — 完全可忽略. 这解释了为何状态报告里 PDC-only ensemble (\StatusReport{} §修复欠覆盖) 在关掉 Sn 后 $\sigma_p$ 没有 Landau-tail 病态.

\subsection{修复方案}

\paragraph{RK4 步长内的能损项.} 把式~\eqref{eq:partIc:bb_master} 加入 RK4 的力项:
\begin{equation}
  \frac{d \vec p}{ds} = q \vec v \times \vec B
  -\hat v \cdot \rho(\vec r) \cdot \left[-\langle dE/dx \rangle\right] / c.
  \label{eq:partIc:rk4_with_dEdx}
\end{equation}
其中 $\rho(\vec r)$ 是位置相关的材料密度 (查 GDML/材料表), $-\langle dE/dx \rangle$ 按当前 $\beta\gamma$ 在 step 中点重算. 实现要点:
\begin{itemize}[nosep, leftmargin=2em]
  \item 介质材料表预编译: 按 GDML solid → material → $(Z, A, I, \rho)$ 查找; PDG 表插值 $\langle dE/dx \rangle (\beta\gamma)$.
  \item RK4 步长里平均能损只取 mean (不引入 Landau 涨落), 因为我们做的是\emph{确定性前向模型}; Landau 入观测协方差 (Kalman 处理).
  \item Step end 重算 $\beta, \gamma$, 更新 $\vec p$ 模长 (方向不变, 只改幅度).
\end{itemize}

\paragraph{预测.} 加入 $dE/dx$ 后, RK 拟合的 $p_z$ 偏差应从 $-10.8\;\mathrm{MeV}/c$ 收紧到 $|\langle \Delta p_z \rangle| \lesssim 1\;\mathrm{MeV}/c$ (剩余的 1 MeV/$c$ 来自 Landau 涨落 + 材料表精度). 同时 Sn 处 Landau straggling $\sigma_E \sim 1.1\;\mathrm{MeV}$ 应进入 PDC 测量协方差 (Part~II §5), 给 Fisher $\sigma_p$ 增加 $\sim 2\;\mathrm{MeV}/c$ 的额外贡献, 使 Fisher 与残差 std 在 $p_z$ 方向上协同.

% --- end of Section 8 ---

\section{Highland 多次散射公式与精度修正}
\label{sec:partIc:highland}

\subsection{陈述}

PDG (Eq.~34.16) 给出的 Highland-Lynch-Dahl 公式是描述带电粒子在物质中多次小角度库仑散射 (multiple Coulomb scattering, MS) 的工程级标准:
\begin{equation}
  \theta_0 = \frac{13.6\;\mathrm{MeV}}{\beta c p}\, z
  \sqrt{\frac{x}{X_0}}\,
  \Bigl[1 + 0.038\,\ln\!\bigl(x z^2 / (X_0 \beta^2)\bigr)\Bigr],
  \label{eq:partIc:highland_master}
\end{equation}
其中 $\theta_0$ 是 Gaussian 投影角分布的 $1\sigma$ 宽度 (即 $\theta_x$ 或 $\theta_y$ 之一的 RMS, 整体 $\theta_\mathrm{space}$ 的 RMS 是 $\sqrt 2 \theta_0$). $z$ 是入射粒子电荷数; $x$ 是物质厚度; $X_0$ 是介质辐射长度. 在我们的工作点 $\beta\gamma = 0.677$, $z=1$, 故 $z^2/\beta^2 \approx 3.18$, 但 PDG 推荐对 $\beta \gtrsim 0.5$ 重粒子 $\ln$ 项里直接用 $\ln(x/X_0)$, 误差 $\lesssim 1\%$. 状态报告 (\StatusReport{} §Highland 公式) 采用的就是这个简化版.

\subsection{物理来源 (Lewis 1950 + Molière 小角理论)}

\paragraph{Step 1. 单次 Rutherford 散射.} 单次 Coulomb 散射 (粒子入射, 散射核屏蔽 Coulomb 势) 的微分截面:
\begin{equation}
  \frac{d\sigma}{d\Omega} = \frac{(z Z e^2)^2}{4 (p c \beta)^2 \sin^4(\theta/2)},
\end{equation}
小角下 $\theta \to 0$ 强发散, 但被原子电子屏蔽截至 (Thomas-Fermi 半径 $a_\mathrm{TF} \sim 0.885 a_0 Z^{-1/3}$).

\paragraph{Step 2. $\langle \theta^2 \rangle$ per atom.} 在屏蔽截止下, 单原子平均:
\begin{equation}
  \langle \theta^2 \rangle_\mathrm{atom}
  = \int \theta^2 (d\sigma/d\Omega) d\Omega
  \propto \frac{1}{(\beta p)^2}\, Z(Z+1) \ln\!\bigl(\theta_\mathrm{max}/\theta_\mathrm{min}\bigr).
\end{equation}
$\theta_\mathrm{max}$ 由核大小给, $\theta_\mathrm{min}$ 由屏蔽给, $\ln$ 项是 Coulomb logarithm.

\paragraph{Step 3. 中心极限定理 (Lewis 1950).} $N$ 次独立散射后, 投影角分布在小角下为高斯, RMS 为
\begin{equation}
  \theta_0^2 \approx N \langle \theta^2 \rangle_\mathrm{atom} / 2
  = n_\mathrm{atom} \cdot t \cdot \langle\theta^2\rangle_\mathrm{atom}/2,
\end{equation}
其中因子 $1/2$ 把 3D 角分布投影到 1D. 把所有原子物理塞进定义辐射长度 $X_0$ (单位 g/cm$^2$):
\begin{equation}
  \frac{1}{X_0} = 4\alpha r_e^2 \frac{N_A}{A} Z^2 [L_\mathrm{rad} - f(Z)] + \cdots,
\end{equation}
则 Lewis 给出 ($\beta \to 1$ 极限)
\begin{equation}
  \theta_0 \approx \frac{E_s}{p c} \sqrt{x/X_0},
  \quad
  E_s = m_e c^2 \sqrt{4\pi/\alpha} = 21.2\;\mathrm{MeV}.
  \label{eq:partIc:lewis}
\end{equation}

\paragraph{Step 4. Highland 经验拟合 (1975, Lynch-Dahl 1991).} 实际数据表明 Lewis 的 $E_s = 21.2$ 在大多数实验下\emph{高估}, 因为非高斯尾不能完全靠中心极限"吞掉". Highland 1975 用 Molière 全分布做 Monte Carlo, 拟合出经验常数 13.6 MeV (而非 21.2) 加上 $0.038 \ln(x/X_0)$ 修正项, Lynch-Dahl 1991 进一步精确化. 这就是式~\eqref{eq:partIc:highland_master} 的来源. 适用范围 $10^{-3} \lesssim x/X_0 \lesssim 100$.

\paragraph{$0.038 \ln$ 项的角色.} 当 $x/X_0$ 很小, $\ln$ 项是负数, 把 $\theta_0$ 拉低 (反映非高斯小尾在薄材料里更不显著); 当 $x/X_0$ 很大, $\ln$ 项是正数, 把 $\theta_0$ 拉高 (薄高斯近似失效, 出现重尾). 在 $x/X_0 = 1$ (一个辐射长度) 时 $\ln$ 项为 0.

\subsection{逐材料计算 (复现状态报告)}

下面把状态报告 (\StatusReport{} §Highland 公式) 的三个数字 (Sn 16.4, air 1.72, PDC 1.21 mrad) 从第一性原理推导.

\paragraph{$\beta p$ 因子.} $p = 635\;\mathrm{MeV}/c$, $\beta = 0.5607$, $\beta p = 0.5607 \times 635 = 355.8\;\mathrm{MeV}/c$. 故
\begin{equation}
  \frac{13.6}{\beta p} = \frac{13.6}{355.8} = 0.0382\;\mathrm{rad}.
  \label{eq:partIc:highland_prefactor}
\end{equation}

\paragraph{Sn 靶.} $\rho_\mathrm{Sn} = 7.31\;\mathrm{g/cm^3}$, $X_0^\mathrm{Sn} = 8.82\;\mathrm{g/cm^2}$ (Sn 的标准辐射长度), 厚度 $t = 0.25\;\mathrm{cm}$. 故
\begin{align}
  \rho t &= 7.31 \times 0.25 = 1.828\;\mathrm{g/cm^2}, \\
  x/X_0 &= 1.828 / 8.82 = 0.2073, \\
  \sqrt{x/X_0} &= 0.4553, \\
  \ln(x/X_0) &= \ln(0.2073) = -1.573, \\
  1 + 0.038 \ln(x/X_0) &= 1 + 0.038 \times (-1.573) = 0.940, \\
  \theta_0^\mathrm{Sn} &= 0.0382 \times 0.4553 \times 0.940
  = 0.01635\;\mathrm{rad}
  = 16.35\;\mathrm{mrad}.
\end{align}
四舍五入到 $16.4\;\mathrm{mrad}$, 与状态报告完全一致.

\paragraph{空气 (1 atm, 1 m).} $\rho_\mathrm{air} = 1.205 \times 10^{-3}\;\mathrm{g/cm^3}$, $X_0^\mathrm{air} = 36.66\;\mathrm{g/cm^2}$ (PDG), $t = 100\;\mathrm{cm}$.
\begin{align}
  \rho t &= 1.205\times 10^{-3} \cdot 100 = 0.1205\;\mathrm{g/cm^2}, \\
  x/X_0 &= 0.1205 / 36.66 = 3.288 \times 10^{-3}, \\
  \sqrt{x/X_0} &= 0.0573, \\
  \ln(x/X_0) &= \ln(3.288\times 10^{-3}) = -5.717, \\
  1 + 0.038 \ln(x/X_0) &= 1 - 0.217 = 0.783, \\
  \theta_0^\mathrm{air} &= 0.0382 \times 0.0573 \times 0.783
  = 0.001714\;\mathrm{rad} = 1.71\;\mathrm{mrad}.
\end{align}
舍入到 $1.72\;\mathrm{mrad}$ (与状态报告一致, 1.72 vs 1.71 的 $\sim 0.01$ 差源于中间舍入).

\paragraph{PDC 气体 (75\% Ar + 25\% i-C$_4$H$_{10}$).} 状态报告给出混合密度 $\rho_\mathrm{mix} = 1.87\times 10^{-3}\;\mathrm{g/cm^3}$, 等效辐射长度 $X_0^\mathrm{mix} = 24.1\;\mathrm{g/cm^2}$, 路径 $t = 22.4\;\mathrm{cm}$ (190 mm $/\cos 32^\circ$).
\begin{align}
  \rho t &= 1.87\times 10^{-3} \cdot 22.4 = 0.04189\;\mathrm{g/cm^2}, \\
  x/X_0 &= 0.04189 / 24.1 = 1.738 \times 10^{-3}, \\
  \sqrt{x/X_0} &= 0.04169, \\
  \ln(x/X_0) &= \ln(1.738\times 10^{-3}) = -6.355, \\
  1 + 0.038 \ln(x/X_0) &= 1 - 0.241 = 0.759, \\
  \theta_0^\mathrm{PDC} &= 0.0382 \times 0.04169 \times 0.759
  = 0.001209\;\mathrm{rad} = 1.21\;\mathrm{mrad}.
\end{align}
与状态报告一致.

\subsection{多次散射对 $\sigma_p$ 的贡献}

\paragraph{磁场后散射 → 出射角误差.} Sn 靶散射发生在\emph{磁场之前}, 被 RK fitter 的方向自由度 $(u, v)$ 完全吸收 (\StatusReport{} §各材料对动量重建的影响). 真正污染动量重建的是\emph{磁场之后}的 MS:
\begin{itemize}[nosep, leftmargin=2em]
  \item 空气 1 m: $\theta_0^\mathrm{air} = 1.72\;\mathrm{mrad}$.
  \item PDC1 气体: $\theta_0^\mathrm{PDC} = 1.21\;\mathrm{mrad}$ (PDC1 内部散射的角偏转传到 PDC2 测量).
  \item PDC2 气体: 仅引起 PDC2 内部 $\sim 0.08\;\mathrm{mm}$ 位移, 可忽略.
\end{itemize}

\paragraph{合并.} 平方和:
\begin{equation}
  \sigma_\theta^\mathrm{MS} = \sqrt{(\theta_0^\mathrm{air})^2 + (\theta_0^\mathrm{PDC})^2}
  = \sqrt{1.72^2 + 1.21^2}
  = \sqrt{2.958 + 1.464}
  = 2.10\;\mathrm{mrad}.
  \label{eq:partIc:sigma_theta_MS_total}
\end{equation}

\paragraph{$\sigma_p$ 贡献.} 用式~\eqref{eq:partIc:sigma_p_master}:
\begin{align}
  \sigma_p^\mathrm{air} &= 1460 \times 0.00172 = 2.51\;\mathrm{MeV}/c, \\
  \sigma_p^\mathrm{PDC} &= 1460 \times 0.00121 = 1.77\;\mathrm{MeV}/c, \\
  \sigma_p^\mathrm{MS,total} &= 1460 \times 0.00210 = 3.07\;\mathrm{MeV}/c,
\end{align}
四舍五入与状态报告 (\StatusReport{} §综合动量分辨, 表~综合动量分辨) 完全一致 (2.5, 1.8, 3.1 MeV/$c$).

\paragraph{综合 ($\sigma_x = 0.5\;\mathrm{mm}$).}
\begin{equation}
  \sigma_p^\mathrm{total} = \sqrt{(\sigma_p^\mathrm{pos})^2 + (\sigma_p^\mathrm{MS,total})^2}
  = \sqrt{1.03^2 + 3.07^2} = \sqrt{1.06 + 9.42}
  = \sqrt{10.48} = 3.24\;\mathrm{MeV}/c.
\end{equation}
即\textbf{$\sigma_p \approx 3.2\;\mathrm{MeV}/c$}, 这正是 NN 后端 $p_z$ MAE = 3.1 MeV/$c$ 的目标值.

\paragraph{综合 ($\sigma_x = 2.0\;\mathrm{mm}$).}
\begin{equation}
  \sigma_p^\mathrm{total}(2.0) = \sqrt{4.13^2 + 3.07^2}
  = \sqrt{17.06 + 9.42} = 5.15\;\mathrm{MeV}/c.
\end{equation}
即\textbf{$\sigma_p \approx 5.1\;\mathrm{MeV}/c$}, 与状态报告 表~综合动量分辨 完全一致.

\subsection{MS 与 LM 拟合的相互作用}

\paragraph{问题陈述.} RK4 前向模型假设 \emph{确定性}光滑轨迹: 给定靶处动量 $\vec p_\mathrm{tgt}$, 轨迹通过 PDC 平面的位置完全确定, 没有过程噪声. 但实际 MS 是\emph{每事件独立的随机过程}, 把"PDC 测量值 - RK4 预测值"残差分布展宽到\emph{超过} $\sigma_\mathrm{wire}$ 应该给的范围. 状态报告 (\StatusReport{} §修复欠覆盖) 观测到关掉 Sn 靶后 $\sigma_x \approx 3$--$15\;\mathrm{mm}$ (MS 主导) vs $\sigma_\mathrm{wire} = 0.2\;\mathrm{mm}$, 相差 15--75 倍.

\paragraph{两种修复方法.}
\begin{enumerate}[nosep, leftmargin=2em]
  \item \textbf{Kalman / measurement covariance inflation.} 把 MS 贡献加入测量协方差: $\Sigma_\mathrm{meas} = \mathrm{diag}(\sigma_\mathrm{wire}^2) + \Sigma_\mathrm{MS}$, 其中 $\Sigma_\mathrm{MS}$ 由材料厚度沿轨迹累积. 这是 "rough Kalman" 处理, 在每次 RK4 step 推进时把过程噪声转换成等效测量协方差. 见 Part~II §3 修复路径.
  \item \textbf{Toy-MC PDC smearing (status report 的 method 6).} 在反演前对每次拟合, 在初始方向上加 Gaussian smear $\theta_0$, 重复 $N$ 次, 用拟合结果的散布估计 $\sigma_p$. 这是 frequentist ensemble 方法, 不需要修改 forward model, 但需要 $N$ 倍计算开销.
\end{enumerate}
状态报告里两种方法都未在 production 启用; Part~II §3 给出明确的修复 roadmap.

\subsection{Highland 公式的局限}

Highland 给的是\emph{投影角的高斯近似}; 真实分布有重尾 (Molière 全分布):
\begin{itemize}[nosep, leftmargin=2em]
  \item 高斯近似在角度 $\lesssim 3 \theta_0$ 内有效, 占总事件 $\sim 99\%$.
  \item 大角尾巴 ($> 3\theta_0$) 由 Molière "single-plural" component 决定, 概率 $\sim 1\%$.
\end{itemize}
对覆盖率/pull 分布尾部分析, 直接用 Geant4 G4MultipleScattering 模型 (背后是 Urban 或 GS 模型, 已经包含尾巴) 比 Highland 准确. 状态报告 §G4 涨落集合就是 G4 simulated MS 的结果, 自动包含尾巴.

\paragraph{对覆盖率的影响.} 假设我们把 MS 当 Gauss 处理 (Highland $\theta_0$ + Kalman 协方差膨胀), 则覆盖率会在 68\% / 95\% nominal 附近\emph{略低估} ($\sim 1$--$2\%$ 的 deficit), 因为重尾事件的真值会被推到 Gauss 区间外. 在我们 744 事件 ensemble 上, 这个 deficit 是\emph{可观察的}但低于 Clopper-Pearson 的统计 noise floor; Part~III §5 给出对应的统计检验.

\paragraph{修复后预期.} 在 Part~II §3 的 Kalman 协方差修复 + Part~II §5 的 $dE/dx$ 修复都到位后, 预期:
\begin{itemize}[nosep, leftmargin=2em]
  \item $|\langle \Delta p_z \rangle| \lesssim 1\;\mathrm{MeV}/c$ (来自 $dE/dx$ 修复, Part~II §5).
  \item Fisher $\sigma_{p_z}$ 与 残差 std 在 $p_z$ 方向匹配 (来自 MS 协方差膨胀, Part~II §3).
  \item 68\% 覆盖率全部组件回到 nominal 附近 ($\pm 2\%$).
\end{itemize}
$p_y$ 方向参数化 line-arization 偏差 (\StatusReport{} §修复欠覆盖) 是另一条独立路径, 见 Part~II §8.

% --- end of Section 9 ---


\newpage

% ============================================================
% 第二部分: 系统误差预算
% ============================================================
\part{系统误差预算}\label{part:budget}

本部分按物理来源逐项拆解 RK 重建的误差预算。共十章：方法学综览（§0）、
PDC 位置分辨（§1）、空气多次散射（§2）、PDC 气体多次散射（§3）、Sn 靶
多次散射（§4，关闭模式下的预算位）、$dE/dx$ 平均值与 Landau/Vavilov 涨落
（§5，最大未建模项）、磁场图插值（§6）、对位与靶倾角残差（§7）、
$(u,v,p)$ 参数化诱导偏差（§8）、LM 收敛盆地失败（§9，$(-100,0)$ 病态）、
预算调和与观测残差并表（§10）。

每章遵循统一模板：物理机制 \(\to\) 量级估计 \(\to\) 当前测量值（来自
状态报告 CSV）\(\to\) 预算行（区分 BIAS 与 WIDTH）\(\to\) 修复路径与
后效预测。预算项之间的协方差结构在 §10 调和。

% =============================================================================
% rk_deep_dive_part2a_budget_ch0to3_zh.tex
%
% Purpose:
%   Part II (Systematic budget) §0–§3 of the RK error-analysis deep-dive
%   companion document.  Builds the canonical "误差预算 master table" and
%   walks through the first three contributing rows in template form
%   (物理机制 / 量级估计 / 当前测量值 / 预算行 / 修复路径).
%
% Parent file (wrapper):
%   docs/reports/reconstruction/momentum_reco/rk/rk_error_analysis_deep_dive_20260426_zh.tex
%
% Related companion files:
%   rk_deep_dive_part1a_math_ch1to3_zh.tex   (Likelihood/CR/Fisher)
%   rk_deep_dive_part1b_math_ch4to6_zh.tex   (Wilks/Laplace/MCMC)
%   rk_deep_dive_part1c_math_ch7to9_zh.tex   (Glückstern/Bethe-Bloch/Highland)
%   rk_deep_dive_part2b_budget_ch4to6_zh.tex (dE/dx, RK step, target z 等)
%   rk_deep_dive_part2c_budget_ch7to10_zh.tex(LM 失效、frame bug、综合表)
%   rk_deep_dive_part3_diagnostics_zh.tex    (诊断与下一步)
%
% Status report (cross-reference target):
%   rk_reconstruction_status_20260416_zh.tex  — labelled \StatusReport
% Ensemble data root:
%   build/test_output/ensemble_n50_targetframe/  — labelled \DeepDiveData
%
% Date         : 2026-04-26
% Author       : Baiting Tian
% Convention   : 默认靶系 (beam-as-Z)；动量分量 (p_x, p_y, p_z) 全部为靶系；
%                公式编号 \label{eq:partIIa:section:topic}
%
% Local TOC:
%   §0  综览：误差预算的方法学
%   §1  PDC 位置分辨贡献
%   §2  多次散射：1 m 空气间隙
%   §3  多次散射：PDC1/PDC2 漂移气体
% =============================================================================

\providecommand{\StatusReport}{文献~[1]}
\providecommand{\DeepDiveData}{[deep-dive-data]}

\section{0. 综览：误差预算的方法学}
\label{sec:partIIa:budget-overview}

\subsection*{0.1 加性方差范式}

把 RK 后端在 744 事件 ensemble 上观察到的动量残差宽度分解到独立物理来源是 Part II 的目标。
基本假设是\textbf{独立加性方差}：对每一个动量分量 $p_k\in\{p_x,p_y,p_z,|\vec p|\}$，
\begin{equation}
  \sigma_\mathrm{tot}^2(p_k)\;=\;\sum_{i=1}^{N_\mathrm{src}} \sigma_i^2(p_k),
  \label{eq:partIIa:budget:add-var}
\end{equation}
其中 $i$ 索引各物理误差源（PDC 位置分辨、空气 MS、PDC 气体 MS、Bethe--Bloch 涨落、RK 离散误差、靶 $z$ 不确定度、frame 旋转、LM 收敛尾、$B$ 场不均匀、对准 \ldots）。这一分解对应 Part I §3 (Fisher) 的概率视角：当似然在最佳点附近近似为高斯，逆 Fisher 信息可线性叠加各独立测量贡献。

\subsection*{0.2 关于"独立"的两个保留}

公式~\eqref{eq:partIIa:budget:add-var} 假设各源之间统计独立，但本问题里有两条已知的耦合：
\begin{itemize}
  \item \textbf{空气 MS 与 PDC 气体 MS 不严格独立。} 两者作用于同一段轨迹（出磁后 $\to$ PDC1 $\to$ PDC2），它们的角度涂抹在 PDC1 位置上是\emph{相关}的：空气段散射后入射 PDC1 时方向已经偏离，PDC1 内的散射又会进一步累加。Highland 公式本身\StatusReport{}~§理论分辨极限是把空气和 PDC 气体的 $x/X_0$ 直接相加再算一次 $\theta_0$（保守、相关已经吸收），但若按"独立两段"拆开求平方和会高估约 $5\text{--}10\%$。本预算采用平方和法，并在 §3 末尾保留一个 $-\,$10\% 的耦合修正项。
  \item \textbf{Bethe--Bloch 能量损失的偏置 vs.~涨落。} dE/dx 的\emph{平均损失} $\langle\Delta E\rangle\sim 5$--$10$ MeV 表现为\textbf{系统偏置}（前向模型若不修正则给出偏低的 $|\vec p|$），其\emph{Bohr/Landau 涨落} $\sigma_E$ 才进入预算式~\eqref{eq:partIIa:budget:add-var} 的方差。Part II §4 会把这两件事拆开。
\end{itemize}

\subsection*{0.3 偏置 vs.~宽度：哪些行进入式 \eqref{eq:partIIa:budget:add-var}？}

把误差源按其对残差直方图的影响形态分成三类：
\begin{description}
  \item[宽度型 (width)] 进入式~\eqref{eq:partIIa:budget:add-var} 求和。
        典型成员：PDC 位置分辨 (§1)、空气 MS (§2)、PDC 气体 MS (§3)、
        Bethe--Bloch 涨落、RK 离散误差。这些把残差分布的核心宽度撑开，
        通常表现为近似高斯的中心 68\% 区间。
  \item[偏置型 (bias)] \emph{不}进入式~\eqref{eq:partIIa:budget:add-var}，但出现在 master 表"修复后预测"列。
        典型成员：dE/dx 平均损失、frame 旋转 bug（已于 2026-04-19 修复，
        见 \StatusReport{}~§误差分析与 spec
        \texttt{2026-04-19-reco-target-frame-design.md}）、靶 $z$ 系统偏置。
        这些把整条直方图整体平移；修复前在 $|\vec p|$ 上看到 5--10 MeV/$c$ 的偏移，
        修复后偏置归零，但宽度不变。
  \item[尾部型 (tail / non-Gaussian)] 既不能用方差描述，也不能用偏置描述。
        典型成员：LM 收敛盆地失败（约 2--5\% 的 ensemble 事件被困在错误极小，
        贡献 $|p_y|>50$ MeV/$c$ 的极大尾巴）。处理方式：单独剔除并报告 outlier
        比例，不进入式~\eqref{eq:partIIa:budget:add-var}。
\end{description}

注意 frame 旋转 bug 在 2026-04-19 之前\emph{只}是偏置（拼装矩阵转置错误），不影响宽度；
当前 (2026-04-26) ensemble 上 $p_y$ Fisher $\sigma$ 系统性欠估 $\sim 2\times$ 的现象\emph{只}是宽度问题（\StatusReport{}~§方法总结说明：原因是测量协方差未含过程噪声 $\Sigma_\mathrm{MS}$）。把这两件事混在一起讨论是历史上常见的误诊。

\subsection*{0.4 对账策略}

把式~\eqref{eq:partIIa:budget:add-var} 的求和结果与 ensemble 实测残差宽度对账，是判断 master 表是否封闭的唯一手段。实测宽度的定义采用稳健分位数：
\begin{equation}
  \widehat{\sigma}_k^\mathrm{obs}
  \;=\;\frac{p_{84}(\Delta p_k)\,-\,p_{16}(\Delta p_k)}{2},
  \label{eq:partIIa:budget:halfw68}
\end{equation}
即 ensemble 残差 $\Delta p_k = p_k^\mathrm{reco} - p_k^\mathrm{truth}$ 在 [16\%, 84\%] 分位的半宽。该量在 \DeepDiveData{} 内的 \texttt{per\_truth\_point\_g4\_vs\_fisher.csv} 中以列名 \texttt{g4\_halfw68\_p\{x,y,z\}}, \texttt{g4\_halfw68\_p} 给出（与 Fisher 列 \texttt{fisher\_sigma\_p\{x,y,z\}}, \texttt{fisher\_sigma\_p} 配对）。CSV schema：
\begin{verbatim}
truth_px, truth_py, N,
g4_halfw68_px, std_px, fisher_sigma_px,
g4_halfw68_py, std_py, fisher_sigma_py,
g4_halfw68_pz, std_pz, fisher_sigma_pz,
g4_halfw68_p,  std_p,  fisher_sigma_p
\end{verbatim}
对账判据：
\begin{equation}
  \Bigl|\widehat{\sigma}_k^\mathrm{obs}\Bigr|^2
  \;\stackrel{?}{=}\;\sum_{i\in\text{width rows}} \sigma_i^2(p_k).
  \label{eq:partIIa:budget:closure}
\end{equation}
若左侧明显大于右侧 $\Rightarrow$ master 表漏行；明显小于 $\Rightarrow$ 某行高估或独立性假设失败。Part II §10 会把 15 个真值点的对账数字汇总成一张图。

\subsection*{0.5 Master 误差预算表（占位预览）}

下表给出 Part II 全部行的骨架；本文档 §1--§3 填前三行，§4--§6 填中段，§7--§9 填尾部，
§10 给出 fixed-everything 后的预测列与最终对账。所有数字单位为 MeV/$c$，
靶系。"修复后预测"列指 \emph{若该行对应的工程或软件修复全部到位}（如把 1 m 空气换成真空、
把丝分辨从 0.5 mm 砍到 0.2 mm 等）后的剩余贡献；当前未填则记 "见 §$N$"。

\begin{table}[ht]
  \centering
  \small
  \begin{tabular}{l|cccc|c|l}
    \hline
    来源 & $\sigma_{p_x}$ & $\sigma_{p_y}$ & $\sigma_{p_z}$ & $\sigma_{|\vec p|}$ &
      修复后预测 & 状态 \\
    \hline
    1. PDC 位置分辨 ($\sigma_w{=}0.5$ mm)    & 0.7   & 0.5   & 0.7   & 0.7   & 见~§1  & 见~§1 \\
    2. MS：1 m 空气                          & 2.5   & 0.5   & 2.5   & 2.5   & 见~§2  & 见~§2 \\
    3. MS：PDC1+PDC2 气体                    & 1.8   & 0.4   & 1.8   & 1.8   & 见~§3  & 见~§3 \\
    4. dE/dx 涨落 (Bohr+Landau)              & 见~§4 & 见~§4 & 见~§4 & 见~§4 & 见~§4  & 见~§4 \\
    5. RK 步长离散误差                       & 见~§5 & 见~§5 & 见~§5 & 见~§5 & 见~§5  & 见~§5 \\
    6. 靶 $z$ 位置不确定度                   & 见~§6 & 见~§6 & 见~§6 & 见~§6 & 见~§6  & 见~§6 \\
    7. LM 收敛盆地失败 (尾部)                & 见~§7 & 见~§7 & 见~§7 & 见~§7 & 见~§7  & 见~§7 \\
    8. frame 旋转 bug (已修)                 & 0     & 0     & 0     & 0     & 0      & 已修复 \\
    9. $B$ 场地图不均匀                      & 见~§9 & 见~§9 & 见~§9 & 见~§9 & 见~§9  & 见~§9 \\
    \hline
    \textbf{平方和合计} (式~\ref{eq:partIIa:budget:add-var}) &
        见~§10 & 见~§10 & 见~§10 & 见~§10 & 见~§10 & --- \\
    \textbf{ensemble 实测半宽 (\ref{eq:partIIa:budget:halfw68})} &
        见~§10 & 见~§10 & 见~§10 & 见~§10 & 见~§10 & --- \\
    \hline
  \end{tabular}
  \caption{Master 误差预算表预览（数值单位 MeV/$c$，靶系）。本节仅给出 §1--§3 的占位，
    完整数值与对账见 Part II §10。}
  \label{tab:partIIa:budget:master-preview}
\end{table}

% TODO: figure rk_deep_dive_part2a_budget_master_preview.pdf —
%       将上表 9 行 + 2 合计行画成横向 stacked-bar，便于一眼看出哪一行主导。

\subsection*{0.6 与 Part I 的对应关系}

\begin{itemize}
  \item §1（位置分辨）与 Part I §3 (Fisher 矩阵 $\mathbf F = \mathbf J^\top\Sigma^{-1}\mathbf J$ 中的 $\Sigma$ 对角项)、Part I §7 (Glückstern 公式) 直接对应。
  \item §2--§3（MS）与 Part I §9 (Highland 公式) 对应；耦合到 $\Sigma_\mathrm{MS}$ 的过程协方差还未注入到当前后端 (\StatusReport{}~§方法总结)，因此 Fisher 在 $p_y$ 上欠估。
  \item §4 (dE/dx) 与 Part I §8 (Bethe--Bloch + Bohr/Landau) 对应。
  \item §5 (RK 离散) 与 Part I §2 (RK4 误差阶) 对应。
\end{itemize}

\section{1. PDC 位置分辨贡献}
\label{sec:partIIa:budget:pdc-pos}

\subsection*{1.1 物理机制}

PDC1 与 PDC2 各自给出两套丝坐标 $(u, v)$（PDC 丝坐标系，与靶系 $(x, y)$ 之间是一个固定旋转，见 \StatusReport{}~§实验装置 / PDC 丝坐标系）。每一根丝读出的位置在 Geant4 端经过一次高斯涂抹：
\begin{equation}
  \tilde u \;=\; u_\mathrm{truth} + \mathcal N(0,\sigma_\mathrm{wire}^2),\qquad
  \tilde v \;=\; v_\mathrm{truth} + \mathcal N(0,\sigma_\mathrm{wire}^2),
  \label{eq:partIIa:budget:wire-smear}
\end{equation}
其中 $\sigma_\mathrm{wire}=0.5$ mm 为当前仿真默认值（代码锚点：\texttt{libs/analysis/src/PDCSimAna.cc::SmearWireCoord}）。两个 PDC、两个坐标 $\Rightarrow$ 共 4 个独立位置测量进入 $\chi^2$。残差结构与 Part I §3.2 中给出的 Fisher 矩阵 $\mathbf F = \mathbf J^\top\Sigma_\mathrm{meas}^{-1}\mathbf J$ 中 $\Sigma_\mathrm{meas} = \sigma_\mathrm{wire}^2\,\mathbb{I}_4$ 的对角结构对应。

\subsection*{1.2 量级估计：Glückstern 公式}

把 PDC 位置误差线性传播到动量误差，使用 Part I §7 的 Glückstern 估计（双平面 + 弯转长度 $L_\mathrm{eff}$ + 漂移基线 $L_{12}$）：
\begin{equation}
  \sigma_p \;\approx\; \frac{p^2\,\sigma_x\,\sqrt 2}{q\,B\,L_\mathrm{eff}\,L_{12}}.
  \label{eq:partIIa:budget:gluckstern}
\end{equation}
代入当前几何（\StatusReport{}~§理论分辨极限）：$p\approx 635$ MeV/$c$、$q=1$、$B\approx 3.0$ T、$L_\mathrm{eff}\approx 0.6$ m、$L_{12}\approx 1.0$ m，$\sigma_x = \sigma_\mathrm{wire}=0.5$ mm，得 $\sigma_p \approx 1.0$ MeV/$c$。状态报告同节给出 $\sigma_x{=}0.5/2.0$ mm 对应 $\sigma_p{=}1.0/4.1$ MeV/$c$（线性比例 $\propto\sigma_x$）。

由于 SAMURAI 磁偶在 $x$--$z$ 平面内弯曲，$y$ 方向不受磁场约束，
丝坐标涂抹对 $p_y$ 的贡献只来自方向参数 $v$ 的拟合精度，
量级为 $\sigma_{p_y}^\mathrm{pos}\sim p\cdot\sigma_\mathrm{wire}/L_{12} \approx 635\times 5\times 10^{-4}/1000 = 0.32$ MeV/$c$。考虑到两端 PDC 的 $v$ 测量都进入拟合，按 $\sqrt 2$ 因子调高约 0.5 MeV/$c$。

\subsection*{1.3 当前测量值}

在 744 事件 ensemble (\StatusReport{}~§当前性能 + \DeepDiveData{}/per\_truth\_point\_g4\_vs\_fisher.csv) 中，Fisher $\sigma_p$ 落在 $\sim 6$--7 MeV/$c$。这\emph{不是}纯位置贡献——它已经把 §2 空气 MS 与 §3 PDC 气体 MS 一起经 $\Sigma_\mathrm{meas}^{-1}$ 反算回去了（虽然过程噪声未显式注入，散射后的实际命中位置仍含散射偏移 $\Rightarrow$ 残差自身被 MS 撑宽）。要从 ensemble 直接看到\emph{纯位置}贡献，唯一办法是把 Geant4 关掉空气与 PDC 气体（即 ms\_ablation stage A 的"全真空"配置，见 §2.5），这正是 \texttt{ms\_ablation\_air\_20260421\_zh.md}（MEMORY 条目 \texttt{project\_ms\_ablation\_air\_finding}, 2026-04-21）和 \texttt{ms\_ablation\_mixed\_20260422\_zh.md}（条目 \texttt{project\_ms\_ablation\_mixed\_finding}, 2026-04-22）所做的工作。Stage A (vacuum 上下游) 的 PDC2 $\sigma_y$ 入参是 7.49 mm，stage A+B (air 上下游) 是 13.50 mm，差异即来自空气 MS 的传播放大。

把式~\eqref{eq:partIIa:budget:gluckstern} 给出的 $\sim 1$ MeV/$c$ 看成纯位置贡献的上限估计是合理的（量级与 ms\_ablation stage A vacuum 端外推一致）。

\subsection*{1.4 预算行}

填入 master 表 (\ref{tab:partIIa:budget:master-preview}) 第 1 行的占位值（粗估，由式~\eqref{eq:partIIa:budget:gluckstern} 配合 $y$ 方向无磁场约束的 $\sqrt 2$ 比例规则）：
\begin{equation}
  \sigma_{p_x}^\mathrm{pos}\,\approx\,0.7,\quad
  \sigma_{p_y}^\mathrm{pos}\,\approx\,0.5,\quad
  \sigma_{p_z}^\mathrm{pos}\,\approx\,0.7,\quad
  \sigma_{|\vec p|}^\mathrm{pos}\,\approx\,0.7\;\mathrm{MeV}/c.
  \label{eq:partIIa:budget:row1}
\end{equation}
精确数值需要按 Glückstern 各列（$L_\mathrm{eff}$、$L_{12}$、$B$、$p$）逐步代入复算；§10 综合表会用 ensemble 上 ms\_ablation stage A 的"全真空"残差替代占位。

\subsection*{1.5 修复路径}

预测 $\sigma_p^\mathrm{pos}\propto\sigma_\mathrm{wire}$（线性）。三条可行路线：
\begin{enumerate}
  \item \textbf{硬件 / 标定层面}：把 PDC 丝坐标分辨从 0.5 mm 改善到 0.2 mm。\StatusReport{}~§理论分辨极限对此情形给出 $\sigma_p\sim 0.4$ MeV/$c$。在式~\eqref{eq:partIIa:budget:row1} 上线性缩放 $0.2/0.5=0.4$，即 $\sigma_{p_{x,z,|\vec p|}}\to 0.3$、$\sigma_{p_y}\to 0.2$ MeV/$c$。
  \item \textbf{Geant4 端只改 \texttt{SmearWireCoord} 的 $\sigma$ 输入}：可以低成本地做敏感度扫描（已在 plans 中建议作为 future test），用来\emph{独立}验证位置贡献是否就是 §1 估算的水平。
  \item \textbf{增加 PDC 平面数}（如插入 PDC0 或 PDC3）：Fisher 信息按 $1/N_\mathrm{plane}$ 衰减，但工程代价远高于改善单平面分辨。当前不在路线图。
\end{enumerate}

\textbf{结论。} PDC 位置分辨在当前几何下\emph{不是}主导误差源：
其量级（$\sim 0.5$--$0.7$ MeV/$c$）远小于 §2 空气 MS 与 §3 PDC 气体 MS 的 $\sim 2$ MeV/$c$ 量级。
这与 \StatusReport{}~§理论分辨极限的"几何 $\sigma_p \approx 1$ MeV/$c$ vs.~含 MS 的 $3.2$--$5.1$ MeV/$c$"判据一致。

\section{2. 多次散射：1 m 空气间隙}
\label{sec:partIIa:budget:air-ms}

\subsection*{2.1 物理机制}

质子从 SAMURAI 磁偶极出口窗到 PDC1 第一根丝之间穿过约 1 m 的常压空气间隙。
这是\emph{出磁后}的散射——RK fitter 假设光滑解析轨迹（无过程噪声），
\StatusReport{}~§理论分辨极限明确指出"散射角直接加入出射角测量误差"。
此段散射的角度涂抹无法被前向模型消化，会以 $\theta_0\cdot L_\mathrm{drift}$ 的方式
线性放大到 PDC1/PDC2 命中位置的偏移。机制公式见 Part I §9 (Highland)：
\begin{equation}
  \theta_0(\text{air}) \;=\; \frac{13.6\;\mathrm{MeV}}{\beta p}\,
       z\,\sqrt{x/X_0}\,
       \bigl[1 + 0.038\ln(x/X_0)\bigr].
  \label{eq:partIIa:budget:highland-air}
\end{equation}

\subsection*{2.2 量级估计}

\StatusReport{}~§理论分辨极限给出空气段（$\rho=1.205\times 10^{-3}$ g/cm$^3$、$X_0=36.66$ g/cm$^2$、$x/X_0 = 3.29\times 10^{-3}$、$p\approx 635$ MeV/$c$）的散射角：
\begin{equation}
  \theta_0(\text{air}, 1\,\text{m}) \;=\; 1.72\;\mathrm{mrad}.
  \label{eq:partIIa:budget:theta-air}
\end{equation}
把这个出射角误差按"双平面角度法"传到动量 (\StatusReport{}~§双平面角度分辨与动量分辨)：
\begin{equation}
  \sigma_p \;=\; \frac{p}{\theta_\mathrm{bend}}\,\theta_0
        \;=\; 1460\times 1.72\times 10^{-3}
        \;\approx\; 2.5\;\mathrm{MeV}/c,
  \label{eq:partIIa:budget:sigp-air}
\end{equation}
其中 $p/\theta_\mathrm{bend}=1460$ MeV/$c$/rad 为 \StatusReport{}~表给出的灵敏度。

$y$ 方向（无磁场约束）的散射通过\emph{位置位移}传到 PDC2 之外的几何，因此其对 $p_y$ 的影响经由方向参数 $v$ 的拟合精度，而不是经由弯曲半径反演。量级显著小于 $x$/$z$ 方向，按经验比例约 $0.5$ MeV/$c$（与 \StatusReport{}~§Geant4 涨落给出的 PDC2 $\sigma_y / \sigma_x$ 各向异性 3--4 倍 + Glückstern 反向缩放一致）。

\subsection*{2.3 当前测量值}

\StatusReport{}~§理论分辨极限把空气与 PDC 气体合在一起给出 Highland 综合 $\sigma_p \approx 3.1$--5.1 MeV/$c$（$p$ 范围 540--890 MeV/$c$ 扫描）。把空气贡献单独抽出（即从 $\theta_0=2.10$ mrad 中扣除 PDC 气体的 1.21 mrad，剩 $\sqrt{2.10^2-1.21^2}\approx 1.72$ mrad，正好与式~\eqref{eq:partIIa:budget:theta-air} 一致），得到本行：
\begin{equation}
  \sigma_p(\text{air only}) \;\approx\; 2.5\;\mathrm{MeV}/c
  \quad\text{（ensemble 上是\emph{主导}宽度贡献）}.
\end{equation}

ms\_ablation 系列报告（\texttt{docs/reports/reconstruction/momentum\_reco/ms\_ablation/}）做的是反向验证：把 Geant4 几何里的"上下游空气"换成真空，用 Sn 靶关掉、tree-gun 注入的 ensemble 跑同一拓扑，记录 PDC 命中宽度的变化。
\begin{itemize}
  \item Stage A (\texttt{ms\_ablation\_air\_20260421\_zh.md}, MEMORY 条目 \texttt{project\_ms\_ablation\_air\_finding}, 2026-04-21)：在 stage A 单独把空气换成真空，PDC2 $\sigma_y$ 的减小幅度对应空气贡献占总 PDC2 $\sigma_y$ 方差的 \textbf{91--94\%}，与 §2.2 的解析估计一致。\textbf{注意：该结论后被 stage A' 修正}——stage A 几何并不\emph{完全}是物理实验的几何，因为 Sn 靶位于磁场腔内（"target is inside dipole cavity"），上下游空气分布不能简单按"出磁后 1 m"估算。
  \item Stage A' rev2 (\texttt{ms\_ablation\_mixed\_20260422\_zh.md}, MEMORY 条目 \texttt{project\_ms\_ablation\_mixed\_finding}, 2026-04-22)：把空气分两段——Region A（磁腔+管道，$\sim 4.3$ m）与 Region B（出磁窗 EW 到 PDCs，$\sim 2$ m）——分别考察。Stage D 给出 PDC2 $\sigma_y$ 输入：A 取 vacuum 时 7.49 mm，A+B 全 air 时 13.50 mm。差值 $\sqrt{13.50^2-7.49^2}=11.2$ mm 即为 Region A+B 空气贡献的位置宽度。
\end{itemize}
当前 master 表 §2 行是用 \StatusReport{} 的 1 m 空气近似（即 Region B 的较窄子集），\textbf{不}涵盖 Region A 的腔内空气；后者按 stage A' 的发现\emph{可能更大}，是 Part II §9（$B$ 场不均匀 / 几何不确定度）需要进一步拆解的耦合项。

\subsection*{2.4 预算行}

按式~\eqref{eq:partIIa:budget:sigp-air} 与 $y$/$x$ 各向异性比，填 master 表第 2 行：
\begin{equation}
  \sigma_{p_x}^\mathrm{air}\,\approx\,2.5,\quad
  \sigma_{p_y}^\mathrm{air}\,\approx\,0.5,\quad
  \sigma_{p_z}^\mathrm{air}\,\approx\,2.5,\quad
  \sigma_{|\vec p|}^\mathrm{air}\,\approx\,2.5\;\mathrm{MeV}/c.
  \label{eq:partIIa:budget:row2}
\end{equation}
误差范围按 \StatusReport{} 的 $p$ 扫描跨度估计为 $\pm 30\%$（$2.0$--$3.2$ MeV/$c$），主要来自 $\theta_\mathrm{bend}$ 在不同弯转构型下的差异。

\subsection*{2.5 修复路径}

\textbf{把 EW $\to$ PDC1 之间的 $\sim 1$ m 空气换成真空束流管 + 薄真空窗。}
\begin{itemize}
  \item 工程代价：一对真空法兰 + 30--50 $\mu$m 厚 Kapton/Mylar 真空窗 +
        机械支撑（PDC 入射面无法直接焊真空窗，需要 stand-off 结构）。
  \item 残余 MS：30 $\mu$m Kapton ($X_0=28.6$ cm) 给出 $x/X_0\approx 1\times 10^{-4}$，
        Highland $\theta_0\approx 0.3$ mrad，远低于当前空气段。
  \item 预测残余 $\sigma_p^\mathrm{air, fixed} < 0.5$ MeV/$c$（按式~\eqref{eq:partIIa:budget:sigp-air} 缩放）。
  \item ms\_ablation 模拟侧已经验证了"全真空"配置下 PDC2 $\sigma_y \to 7.49$ mm
        （vs.~$13.50$ mm 的 air baseline），与上述预测一致。
\end{itemize}
其他备选（成本递增）：He 袋（$X_0\approx 5300$ m，比空气好 $\sim 18\times$，$\theta_0\to 0.4$ mrad）；
完全替换 PDC 入射段为真空气室（PDC 窗本身需重新设计）。

\textbf{结论。} 1 m 空气 MS 是当前 ensemble 上最大的\emph{宽度型}误差源，约占总宽度方差的 60--70\%。
工程修复路径明确（已有 ms\_ablation 数据支撑），是 Part II §10 修复后预测列下降幅度最大的一行。

\section{3. 多次散射：PDC1/PDC2 漂移气体}
\label{sec:partIIa:budget:pdc-gas-ms}

\subsection*{3.1 物理机制}

每个 PDC 内部填充 75\% Ar + 25\% i-C$_4$H$_{10}$（异丁烷淬灭气体）的 1 atm 工作气体，
有效混合密度 $\rho \approx 1.87\times 10^{-3}$ g/cm$^3$，辐射长度 $X_0 \approx 24.1$ g/cm$^2$
（\StatusReport{}~§理论分辨极限 表）。
质子以约 $32^\circ$（相对 PDC 入射轴的等效斜入射角，由弯曲后磁谱计 + 几何决定）穿过 PDC 体，
有效路径 $L_\mathrm{PDC}/\cos 32^\circ \approx 224$ mm。
因 PDC1 与 PDC2 都各自有此散射，单步贡献相等但传播距离不同：
\begin{itemize}
  \item PDC1 内的散射通过 $L_{12}\sim 1$ m 漂移到 PDC2 命中位置，
        贡献偏移 $\Delta x_2 = L_{12}\,\theta_0 \approx 1.21$ mm，
        与丝分辨 $\sigma_\mathrm{wire}=0.5$ mm 量级相当（\StatusReport{}~§理论分辨极限）。
  \item PDC2 内的散射只影响 PDC2 自身命中位置（$\sim$ 几个 mm 漂移），
        其位置贡献 $\sim 0.1$ mm，远小于丝分辨；但仍然会改变出射方向，
        对靶系动量重建无显著贡献（出射方向已脱离 fitter 的测量约束）。
\end{itemize}
因此 §3 行的主要贡献来自\emph{PDC1 散射经 $L_{12}$ 放大到 PDC2}的位置偏移机制，
机制公式与 Part I §9 的 Highland 公式 + Part I §3.4 的 Glückstern 几何耦合结合。

\subsection*{3.2 量级估计}

\StatusReport{}~§理论分辨极限给出单 PDC 散射角：
\begin{equation}
  \theta_0(\text{PDC, 1 plane}) \;=\; 1.21\;\mathrm{mrad},
  \label{eq:partIIa:budget:theta-pdc}
\end{equation}
对应等效位置贡献 $L_{12}\,\theta_0 = 1000\times 1.21\times 10^{-3} = 1.21$ mm。
按 §2.2 的双平面角度法 + 同一 $p/\theta_\mathrm{bend}=1460$ MeV/$c$/rad：
\begin{equation}
  \sigma_p(\text{PDC gas, 1 plane}) \;\approx\;
     1460 \times 1.21\times 10^{-3}
     \;\approx\; 1.8\;\mathrm{MeV}/c.
  \label{eq:partIIa:budget:sigp-pdc}
\end{equation}
该值是\emph{单平面（PDC1）}贡献。把 PDC2 的散射也算上时（按独立加性方差 $\sqrt 2$ 因子），
原则上 $\sigma_p^\mathrm{PDC,gas, both} \approx 1.8\sqrt 2 \approx 2.5$ MeV/$c$；
但 PDC2 内散射主要影响出射后下游的方向，\emph{不进入} PDC1$\to$PDC2 的位置基线，
因此现实中只取单平面值 $\sim 1.8$ MeV/$c$ 是稳健下限。

\subsection*{3.3 当前测量值}

\StatusReport{}~§理论分辨极限把空气和 PDC 气体合并在 Highland 公式里给出 $\theta_0 = 2.10$ mrad：
\begin{equation}
  \theta_0(\text{air} + \text{PDC gas}) \;=\;\sqrt{\theta_0^2(\text{air})+\theta_0^2(\text{PDC})}
       \;=\;\sqrt{1.72^2 + 1.21^2}\;=\;2.10\;\mathrm{mrad}.
  \label{eq:partIIa:budget:theta-combined}
\end{equation}
对应合并 $\sigma_p \approx 3.1$ MeV/$c$（与 \StatusReport{} 的 3.1--5.1 MeV/$c$ 区间下端一致；上端 5.1 对应较低 $p$）。

PDC 气体行的独立验证来自 \StatusReport{}~§Geant4 物理涨落给出的 PDC2 比 PDC1 命中宽度大 $\sim 25$--$40\%$（$y$ 方向更明显）的观察：这正是 PDC1 散射经 $L_{12}=1$ m 漂移到 PDC2 的传播效应。关掉 Sn 靶后 PDC $\sigma_y$ 从 $\sim 100$ mm 砍到 $\sim 10$--$15$ mm（\StatusReport{}~§Geant4 物理涨落同一段），剩下的 10--15 mm 即 PDC 气体 + 10 mm 气隙的真实 MS 水平——与式~\eqref{eq:partIIa:budget:sigp-pdc} 的解析估计在量级上一致（10--15 mm 的位置宽度 / $L_{12}\approx 1$ m $\approx 10$--$15$ mrad；考虑到 ensemble 是\emph{所有}事件混合，包含 PDC1 命中本身宽度 + PDC2 漂移放大，所以实际单平面 PDC 散射 $\theta_0$ 仍在 1--2 mrad 量级）。

\subsection*{3.4 预算行}

按式~\eqref{eq:partIIa:budget:sigp-pdc} 与 §2.4 同样的 $y$/$x$ 各向异性比，填 master 表第 3 行：
\begin{equation}
  \sigma_{p_x}^\mathrm{PDCgas}\,\approx\,1.8,\quad
  \sigma_{p_y}^\mathrm{PDCgas}\,\approx\,0.4,\quad
  \sigma_{p_z}^\mathrm{PDCgas}\,\approx\,1.8,\quad
  \sigma_{|\vec p|}^\mathrm{PDCgas}\,\approx\,1.8\;\mathrm{MeV}/c.
  \label{eq:partIIa:budget:row3}
\end{equation}
误差范围 $\pm 25\%$（$1.4$--$2.3$ MeV/$c$）。注意 §2 与 §3 的耦合修正：
按 §0.2 的讨论，把 air + PDC gas 写成"独立两段"求平方和会高估 $\sim 5$--$10\%$；
master 表 §10 综合时会用合并 $\theta_0=2.10$ mrad 直接算综合 MS 行（即 $\sigma_p\approx 3.1$ MeV/$c$），
而不是把 §2 的 2.5 与 §3 的 1.8 直接平方和（后者给 $\sqrt{2.5^2+1.8^2}=3.08$ MeV/$c$，
碰巧与合并值相近，是 Highland 平方和近似在 $\beta\sim 1$ 区间的一般性质）。

\subsection*{3.5 修复路径}

PDC 必须有气体——这决定了 §3 行的下限\emph{无法}通过软件修复消除。可行方向：
\begin{enumerate}
  \item \textbf{低 $Z$ 气体混合}：用 He/CO$_2$ 或 CF$_4$ 替换 Ar/i-C$_4$H$_{10}$。
        $X_0$ 可能改善 2--3 倍 $\Rightarrow \theta_0$ 改善 $\sqrt 2$--$\sqrt 3$ 倍，
        预测残余 $\sigma_p^\mathrm{PDCgas, fixed}\approx 1.0$ MeV/$c$。
        \textbf{代价}：CF$_4$ 漂移速度 + 增益曲线需要重新标定；He 容易漏；
        会牵涉电子学侧的时间窗调整。\textbf{标记 TODO：未来气体混合优化研究。}
  \item \textbf{缩短气体路径}：把 PDC 厚度从当前的 $\sim 200$ mm 砍到 $\sim 100$ mm。
        $\theta_0$ 改善 $\sqrt 2$ 倍，预测残余 $\sigma_p^\mathrm{PDCgas, fixed}\approx 1.3$ MeV/$c$。
        \textbf{代价}：丝平面变密 / 单丝增益降低 / 漂移单元几何重设。
  \item \textbf{延长 $L_{12}$}：把 PDC1 与 PDC2 间距从 1 m 拉到 1.5 m。
        Glückstern $\sigma_p \propto 1/L_{12}$，预测改善 $1.5\times$；
        但同时空气段 §2 也按 $L_\mathrm{drift}\,\theta_0$ 放大，需要联合优化。
        \textbf{代价}：磁谱仪后端机械空间约束。
\end{enumerate}

\textbf{结论。} PDC 气体 MS 是当前 ensemble 上的次级宽度贡献（$\sim 25$--$30\%$ 总方差占比），
\emph{无法}通过软件消除。在 §2（空气）修复后，§3 将成为新的主导误差源，
对应"修复后预测"列的总和 $\sigma_p^\mathrm{post-fix}\sim 2$ MeV/$c$（位置 + PDC 气体 + dE/dx 涨落 + RK 离散，
具体见 Part II §10）。气体混合优化是该行进一步压缩的唯一路线，但属于长周期工程改造，
当前 deep-dive 仅记录方向、不展开方案设计。

\subsection*{3.6 §1--§3 小结：宽度三巨头}

把前三节合在一起，PDC 几何 + 探测器材料贡献的总宽度（按式~\eqref{eq:partIIa:budget:add-var} 平方和，并按 §0.2 / §3.4 给出的 $-10\%$ 耦合修正后）：
\begin{equation}
  \sigma_p^{\text{§1+§2+§3}}
   \;\approx\;\sqrt{0.7^2+2.5^2+1.8^2}\times 0.95
   \;\approx\; 3.0\;\mathrm{MeV}/c.
  \label{eq:partIIa:budget:rows123-sum}
\end{equation}
这与 \StatusReport{}~§理论分辨极限的 Highland 综合估计 $3.1$ MeV/$c$（合并 $\theta_0=2.10$ mrad 算法）数值上完全一致——印证两种独立路径（按行求平方和 vs.~按合并 $X/X_0$ 算 Highland）在当前几何下殊途同归。

但 ensemble 实测半宽 (\ref{eq:partIIa:budget:halfw68}) 在 $|\vec p|$ 上可达 6--7 MeV/$c$（\DeepDiveData{}/per\_truth\_point\_g4\_vs\_fisher.csv \texttt{g4\_halfw68\_p} 列）。\textbf{差距 $\sim 4$ MeV/$c$ 需要 §4--§9 来补齐}。预期主要来源：
\begin{itemize}
  \item §4 dE/dx 涨落（Bohr/Landau，估 $\sim 1$--2 MeV/$c$）；
  \item §5 RK 步长离散误差（估 $< 0.5$ MeV/$c$，量级小但需验证）；
  \item §6 靶 $z$ 不确定度（系统 + 涨落，估 $\sim 1$ MeV/$c$）；
  \item §7 LM 收敛盆地失败的尾部（不进 $\sigma$ 求和，但拉宽分位 16--84\% 区间 $\sim 2$ MeV/$c$）；
  \item §9 $B$ 场不均匀（估 $\sim 1$ MeV/$c$）。
\end{itemize}
若全部修复后，§1+§2+§3 的"修复后预测"和给出 $\sqrt{0.3^2 + 0.5^2 + 1.0^2}\approx 1.2$ MeV/$c$，
是 deep-dive 给出的"硬件极限"近似。Part II §10 的最终对账图会把这条线作为参考标尺画出。

% TODO: figure rk_deep_dive_part2a_rows123_quadrature_vs_observed.pdf —
%       三柱图对比 §1+§2+§3 平方和、合并 Highland、ensemble 实测半宽。

\subsection*{3.7 与下游章节的衔接}

\begin{itemize}
  \item Part II \texttt{part2b\_budget\_ch4to6\_zh.tex}：§4 dE/dx 涨落、§5 RK 步长离散误差、§6 靶 $z$。
  \item Part II \texttt{part2c\_budget\_ch7to10\_zh.tex}：§7 LM 失效尾部、§8 frame 旋转 bug 复盘、§9 $B$ 场地图、§10 综合对账。
  \item Part III \texttt{rk\_deep\_dive\_part3\_diagnostics\_zh.tex}：诊断脚本、ensemble 切片可视化、推荐下一步实验。
\end{itemize}

本节结束。Master 表的 §1--§3 行已写入式~\eqref{eq:partIIa:budget:row1}--\eqref{eq:partIIa:budget:row3}；
"修复后预测"列与"状态"列在 §10 综合表中会用 ms\_ablation stage A 的真空配置数字 + 0.2 mm 丝分辨敏感扫描数字替代占位。

% =====================================================================
% rk_deep_dive_part2b_budget_ch4to6_zh.tex
% ---------------------------------------------------------------------
% Purpose: Part II (系统误差预算) chapters 4–6 of the RK error-analysis
%          deep-dive companion document. These three chapters form the
%          "physics-of-the-track" half of Part II:
%            §4  Sn 靶多次散射（关闭模式下保留的预算项）
%            §5  dE/dx 平均值 + Landau/Vavilov 涨落（最大未建模项）
%            §6  磁场图插值误差
%
% Parent file (wrapper that \input's this fragment):
%          rk_error_analysis_deep_dive_20260426_zh.tex
%
% Related fragments:
%          rk_deep_dive_part2a_budget_ch1to3_zh.tex   (空气/PDC 气体/丝分辨)
%          rk_deep_dive_part2c_budget_ch7to9_zh.tex   (几何对准/拟合收敛/外参)
%
% Date:    2026-04-26
% Frame:   靶系（beam-as-Z），全部数字、协方差、偏置在该坐标系下表达
% Cite:    \StatusReport = 文献~[1] = rk_reconstruction_status_20260416_zh.tex
%
% Fragment rules (no \documentclass, no \begin{document}):
%   * 仅 \section / \subsection / \paragraph / 表格 / 数学环境
%   * 数字若来自 status report 必须 \StatusReport 引用
%   * 缺测数据用 % TODO: ... 标记，禁止编造
%   * BIAS（中心值偏置）与 WIDTH（弥散）在表格中分列体现
% =====================================================================

\providecommand{\StatusReport}{文献~[1]}

\section{4. 多次散射：Sn 靶（关闭模式下保留的预算项）}
\label{sec:budget-sn-ms}

\subsection{物理机制}
\label{sec:budget-sn-ms-mechanism}

DPOL 实验的核心反应靶为 $2.5\;\mathrm{mm}$ 厚的金属锡（Sn，$Z=50$，$A=118.7$，
密度 $\rho = 7.31\;\mathrm{g/cm^3}$，辐射长度 $X_0 = 8.82\;\mathrm{g/cm^2}$，
约合 $1.21\;\mathrm{cm}$ 物理长度）。靶安装在 SAMURAI 偶极磁铁\textbf{腔体内部}，
中心位置 $z \approx -1070\;\mathrm{mm}$（\StatusReport，第~106 行表格），
关于 $y$ 轴倾斜 $\alpha_\mathrm{tgt} = 3^\circ$ 以满足束流入射几何。

依据 2026-04-22 多次散射消融研究（参见
\texttt{project\_ms\_ablation\_mixed\_finding.md}）的结论：靶位于偶极腔内部，因此
材料预算需按以下两个区段分别处理：
\begin{itemize}
  \item \textbf{Region A}：腔体 + 真空管，自束流入口至腔体后端，约 $4.3\;\mathrm{m}$；
  \item \textbf{Region B}：磁铁出口窗 (Exit Window) 至 PDC1/PDC2，约 $2.0\;\mathrm{m}$；
\end{itemize}
当前 PDC 跟踪研究为隔离测量噪声贡献而\textbf{关闭了} Sn 靶（事件发生器以 tree-gun
模式直接在靶点位置产生质子），但在真实 DPOL 物理运行中 Sn 靶不可缺失。本章给出
"若开启 Sn 靶"情形的预算行，作为完整误差链条的保留项。

将 Sn 靶纳入预算需明确两类效应：
\begin{enumerate}
  \item \textbf{库仑多次散射（Coulomb MS）}：高 $Z$ 介质下的角度发散，$\theta_0 \propto Z/X_0^{1/2}$；
  \item \textbf{平均能量损失 + 涨落}：见 §\ref{sec:budget-dedx} 单独章节，本节仅交叉引用。
\end{enumerate}

\subsection{量级估计：Highland 公式给出 \texorpdfstring{$\theta_0 = 16.4\;\mathrm{mrad}$}{theta0=16.4 mrad}}
\label{sec:budget-sn-ms-magnitude}

\StatusReport{}（第 1166、1186--1187 行）已给出 $635\;\mathrm{MeV}/c$ 质子穿过
$2.5\;\mathrm{mm}$ Sn 平均路径下的 Highland 散射角：
\begin{equation}
  \theta_{0,\mathrm{Sn}}
   = \frac{13.6\;\mathrm{MeV}}{\beta c\,p}\;z_\mathrm{proj}
       \sqrt{\tfrac{x}{X_0}}
       \bigl[1 + 0.038\ln(x/X_0)\bigr]
   = \frac{13.6}{355.8}\times \sqrt{0.207}\times 0.940
   = 16.4\;\mathrm{mrad},
\label{eq:highland-sn}
\end{equation}
其中 $\beta c p = 355.8\;\mathrm{MeV}$，$x/X_0 = 0.207$（$2.5\;\mathrm{mm}$ Sn）。
对照 Part~I §9 的 Highland 推导：此值已远大于本预算其它行——例如空气段 $\theta_0
\approx 0.7\;\mathrm{mrad}$、PDC 气体 $\theta_0 \approx 0.3\;\mathrm{mrad}$（见
\StatusReport{}~第 1160 行表格）。

\paragraph{杠杆臂 (lever-arm) 与动量耦合。}
散射在靶处发生，但靶位于偶极内部，意味着\textbf{偏转角的一部分在散射之前已经发生}，
另一部分仍未发生。设 $f$ 为靶后剩余偏转占总弯角的比例（$f \in [0,1]$）：靶位于磁
铁中心附近时 $f \approx 0.5$；靶位于磁铁入口时 $f \to 1$，靶位于出口时 $f \to 0$。
对于 \StatusReport{}~第 106 行给定的 $z_\mathrm{tgt} = -1070\;\mathrm{mm}$、磁铁
半长 $1750\;\mathrm{mm}$（$z$ 方向，第~108 行），靶位于偶极内部偏入口侧，估计
$f \approx 0.55$。

Sn 散射对动量的耦合分两条路径：

\textbf{路径~1（方向耦合，被 RK fitter 完全吸收）}：散射改变方向 $(u, v)$。RK
拟合以 $(u, v, p)$ 为自由参数，$(u, v)$ 浮动 $16.4\;\mathrm{mrad}$ 全部被方向参
数吸收。在靶平面内的横向位移 $\Delta x \approx t\theta_0/\sqrt{3} = 0.024\;\mathrm{mm}
\ll \sigma_\mathrm{wire}$（\StatusReport{}~第 1207 行），\textbf{对 $|\vec p|$ 直接重建无显著影响}。

\textbf{路径~2（磁场内偏转-散射耦合）}：散射后剩余偏转 $f \cdot \theta_\mathrm{bend}$
作用在已经"摆动"的方向上，弯曲半径 $R \propto p/B$ 与方向耦合，给到 PDC 平面的
位置投影含一项
\begin{equation}
  \Delta y \big|_\mathrm{PDC2}^\mathrm{Sn-MS}
  \;\approx\; L_\mathrm{tgt\to PDC2}\,\theta_0\,\sqrt{(1-f)^2 + (f\,\theta_\mathrm{bend})^2/3},
\label{eq:sn-projection}
\end{equation}
代入 $L_\mathrm{tgt\to PDC2} \approx 6\,\mathrm{m}$、$\theta_0 = 16.4\;\mathrm{mrad}$、
$\theta_\mathrm{bend} \approx 25^\circ \approx 0.44\;\mathrm{rad}$（\StatusReport{}~第
345 行）、$f = 0.55$，得
\(\Delta y \approx 6000 \times 0.0164 \times \sqrt{0.20 + 0.065} \approx 50\;\mathrm{mm}\)。
此处 $50\;\mathrm{mm}$ 应理解为\textbf{方向涨落映射到 PDC2 的位置散布}，并非直接对应
$\sigma_p$；实际进入 $\sigma_p$ 时还需经过 $J(u,v,p)$ 投影。

\subsection{当前测量值：与 ms\_ablation 阶段 D 的横向对照}
\label{sec:budget-sn-ms-current}

2026-04-22 的 ms\_ablation 阶段 D 测量（\texttt{project\_ms\_ablation\_mixed\_finding.md}）
给出 PDC2 处 $\sigma_y$（位置弥散）的两个端点：
\begin{itemize}
  \item Region A=真空 + Region B=空气：$\sigma_y^\mathrm{PDC2} = 7.49\;\mathrm{mm}$；
  \item Region A=空气 + Region B=空气（全空气）：$\sigma_y^\mathrm{PDC2} = 13.50\;\mathrm{mm}$。
\end{itemize}
两者均\textbf{未开启 Sn 靶}。把 $13.50\;\mathrm{mm}$ 当作"无 Sn"基线，再叠加 Sn 靶
带来的 $50\;\mathrm{mm}$ 量级（式~\ref{eq:sn-projection}），一阶平方相加即可估出：
\begin{equation}
  \sigma_y^\mathrm{PDC2}\bigl|_{\text{含 Sn}}
   \;\approx\;\sqrt{13.5^2 + 50^2}\;\mathrm{mm}
   \;\approx\;52\;\mathrm{mm}.
\end{equation}
% TODO: 与历史 pre-2026-04-19 含 Sn 运行的 PDC2 σ_y 直接对照（status-report 早期
% 章节曾提及"含 Sn 时 PDC2 σ_y 大约百毫米量级"，但确切数字需查归档运行）。

\paragraph{把 PDC2 位置弥散映射回靶系动量。}
根据 \StatusReport{}~第 537、590 行的 Jacobian 经验关系
$\sigma_{p_y} \approx 0.35\;\mathrm{MeV}/c$（无 MS 基线），位置弥散
$\sigma_y$ 与动量分量 $\sigma_{p_y}$ 比值约为
$\partial p_y/\partial y_\mathrm{PDC2} \approx 0.35\;\mathrm{MeV}/c \,/\, \mathcal{O}(0.5\,\mathrm{mm})
\approx 0.7\;\mathrm{MeV}/c\,/\,\mathrm{mm}$（数量级估计）。
照此线性比例，含 Sn 时 $\sigma_y \approx 50\;\mathrm{mm}$ 对应 $\sigma_{p_y}$
\textbf{量级 $\sim 30\;\mathrm{MeV}/c$}（与下文预算行的 $24\;\mathrm{MeV}/c$ 自洽，
取较保守值）。

% TODO: 完整的 Sn-on/Sn-off 端到端运行（需要重启 G4 几何带 Sn）；当前章节给出的是
% 由空气基线 + Highland 解析推算的预算上界，未做交叉验证。

\subsection{预算行：BIAS 与 WIDTH 分列}
\label{sec:budget-sn-ms-row}

依据上述估计，\textbf{若开启 Sn 靶}的多次散射贡献预算为：

\begin{table}[H]
  \centering
  \caption{Sn 靶多次散射预算行（开启时）。BIAS 列代表中心偏置，WIDTH 列代表 1$\sigma$ 弥散；
           靶系 $(p_x,p_y,p_z,|\vec p|)$ 单位 $\mathrm{MeV}/c$。}
  \label{tab:sn-ms-budget}
  \begin{tabular}{lcccc}
    \toprule
    分量 & $p_x$ & $p_y$ & $p_z$ & $|\vec p|$ \\
    \midrule
    BIAS  ($\langle \Delta p_i\rangle$)
          & $0$    & $0$    & $0$    & $0$ \\
    WIDTH ($\sigma_{p_i}^{\mathrm{Sn-MS}}$)
          & $\sim 8$ & $\sim 24$ & $\sim 8$ & $\sim 8$ \\
    \bottomrule
  \end{tabular}
\end{table}

\textbf{解读要点}：
\begin{itemize}
  \item Sn 多次散射不引入\textbf{中心偏置}（散射服从 0 均值），因此 BIAS 行全为 $0$；
        与 §\ref{sec:budget-dedx} 的 dE/dx 中心偏置 $-10.8\;\mathrm{MeV}/c$ 形成对照。
  \item $\sigma_{p_y}$ 显著大于 $\sigma_{p_x}, \sigma_{p_z}$：偶极弯转主要在 $xz$
        平面，$p_y$ 方向缺乏磁场约束（参见 \StatusReport{}~第 183 行），散射在 $y$
        方向几乎不被磁场重新折回，因此 PDC1/PDC2 之间的位置位移完整线性反映为
        $p_y$ 不确定度；
  \item $\sigma_{p_x}$ 与 $\sigma_{p_z}$ 相近，反映靶后偏转重新整合了 $xz$ 平面内
        的散射涨落（部分被弯转半径反推所抵消）；
  \item 该行单独即可主导整个预算——空气段（$\sigma \sim 2\text{--}5\;\mathrm{MeV}/c$）
        与 PDC 气体段（$\sigma \sim 0.5\;\mathrm{MeV}/c$）合计仍小于 Sn 行一项，详见
        Part~II §\ref{sec:budget-master} 主表合计列。
\end{itemize}

\subsection{修复路径：保留为物理运行强制项}
\label{sec:budget-sn-ms-fix}

当前 PDC 跟踪研究使用 tree-gun 直接在靶点产生质子，bypass 了 Sn 靶的能量损失和
散射，目的是隔离\textbf{测量噪声 + 重建算法}贡献。这种配置对算法验证是正确的——
否则 RK 的 $\chi^2$ 收敛性、协方差估计都会被 Sn-MS 主导而难以诊断。

但在真实 DPOL 物理运行中，Sn 靶（或其等效反应靶）不可省去。可考虑的缓解路径：
\begin{enumerate}
  \item \textbf{薄化靶}：机械极限约 $1\;\mathrm{mm}$（$x/X_0 \approx 0.083$），Highland
        给 $\theta_0 \approx 10\;\mathrm{mrad}$，$\sigma_{p_y}^{\mathrm{Sn-MS}}$ 大致按
        $\sqrt{0.083/0.207} \approx 0.63$ 比例下降到 $\sim 15\;\mathrm{MeV}/c$；
        但同时反应率下降到 $40\%$ 左右，需评估统计量权衡；
  \item \textbf{替换为低 $Z$ 等效靶}：例如 $5\;\mathrm{mm}\,^{12}\mathrm{C}$（$\rho=2.27\;\mathrm{g/cm^3}$，
        $X_0 = 42.7\;\mathrm{g/cm^2}$，$x/X_0 = 0.027$），相同反应率下
        $\theta_0 \approx 5\;\mathrm{mrad}$，$\sigma_{p_y}^{\mathrm{Sn-MS}}$ 进一步下降至
        $\sim 7\;\mathrm{MeV}/c$；
        % TODO: 等效反应率验证需 Tic-tac d+C 通道的截面输入
  \item \textbf{倾角 + Doppler 修正}：靶倾斜 $3^\circ$ 已纳入几何，多普勒位移可在
        重建端通过 truth $p_z$ 联合拟合修正——属下游工作；
  \item \textbf{在 RK 流程中显式注入 Sn-MS 协方差}：对方法一（Fisher）增加一项
        $\Sigma_\mathrm{Sn-MS} = \mathrm{diag}(\sigma_{p_x}^2, \sigma_{p_y}^2, \sigma_{p_z}^2)$
        作为先验扩展，可在不修改 G4 几何的前提下把含 Sn 的覆盖率上限纳入分析。
\end{enumerate}

% TODO: 本节列出 4 项修复路径但均未实现；待执行 explicit MS-with-Sn run 用于
% full DPOL 物理链验证。预期对比指标：覆盖率 cover68 / Δp_y 的均值与 σ。

\section{5. dE/dx 平均值与 Landau/Vavilov 涨落（最大未建模项）}
\label{sec:budget-dedx}

本章是 Part~II 中最重要的预算章节——它对应 \StatusReport{}~第~1257、1262 行中 RK
$p_z$ MAE = $10.8\;\mathrm{MeV}/c$ 的\textbf{中心偏置}来源，也是 2026-04-17 ensemble
covering study 中 Fisher / MCMC / Profile 在 68\% 处全部欠覆盖的两大主因之一
（另一项是 $(u,v,p)$ 局部线性化，见 Part~I §6）。

\subsection{物理机制：RK4 当前缺失的 \texorpdfstring{$dE/dx$}{dE/dx} 项}
\label{sec:budget-dedx-mechanism}

当前 RK4 实现把 $|\vec p|$ 视作磁场内的守恒量。代码层面：
\begin{itemize}
  \item \texttt{libs/analysis/src/ParticleTrajectory.cc::RungeKuttaStep()} 仅含洛伦兹力
        \(\dot{\vec p} = q\,\vec v\times\vec B\)，
        见 \StatusReport{}~第 183 行明确说明：``洛伦兹力与速度正交，因此 $|\vec p|$
        在磁场中守恒（当前模型不含 $dE/dx$ 能量损失）''；
  \item RK4 的四次力评估均通过 \texttt{CalculateForce(...)} 调用 \texttt{fMagField->GetField(...)}，
        没有任何材料查询、没有 Bethe-Bloch 项。
\end{itemize}

物理上，带电粒子穿过物质会因\textbf{原子电离 + 激发}持续损失能量。其中：
\begin{enumerate}
  \item \textbf{平均 $dE/dx$}：由 Bethe-Bloch 公式给出，是\textbf{中心偏置}来源；
  \item \textbf{涨落（straggling）}：薄介质下服从 Landau 分布，中等厚度过渡到 Vavilov
        分布，厚介质趋于高斯。它给出 $\sigma_{|\vec p|}$ 的\textbf{弥散}贡献。
\end{enumerate}
两类效应在预算表中需\textbf{严格分列}（BIAS 列 vs WIDTH 列），它们的诊断和修复路径
截然不同：BIAS 仅需在 RK4 步进中加入连续 BB 衰减；WIDTH 则需进一步采样。

详细 Bethe-Bloch 推导见 Part~I §8；本节直接用其结果代入预算。

\subsection{量级估计 (mean)：当前 (无 Sn) 配置下 $\Delta E \approx 0.5\;\mathrm{MeV}$}
\label{sec:budget-dedx-magnitude-mean}

为 $635\;\mathrm{MeV}/c$ 质子 ($\beta = 0.560$, $\gamma = 1.21$, $T_\mathrm{kin} = 195\;\mathrm{MeV}$)
逐段累计 Bethe-Bloch 平均能量损失：

\begin{table}[H]
  \centering
  \caption{635 MeV/$c$ 质子各段平均 $dE/dx$ 估计。$\rho x$ 为面密度，$\langle dE/dx\rangle$
           取 $\beta\gamma = 0.68$ 处 BB 表值；$\Delta E$ = $\rho x \times \langle dE/dx\rangle$。}
  \label{tab:dedx-segments}
  \begin{tabular}{lccccc}
    \toprule
    段     & 路径 (cm) & $\rho$ (g/cm$^3$) & $\rho x$ (g/cm$^2$) & $\langle dE/dx\rangle$ (MeV g$^{-1}$cm$^2$) & $\Delta E$ (MeV) \\
    \midrule
    空气（约 $1\;\mathrm{m}$）           & 100   & $1.21\times 10^{-3}$ & $0.121$ & $\sim 2.0$ & $0.24$ \\
    PDC1 气体 (Ar/i-C$_4$H$_{10}$)        & 30    & $\sim 1.4\times 10^{-3}$ & $0.042$ & $\sim 2.2$ & $0.09$ \\
    PDC2 气体 (Ar/i-C$_4$H$_{10}$)        & 30    & $\sim 1.4\times 10^{-3}$ & $0.042$ & $\sim 2.2$ & $0.09$ \\
    管道壁 + 出口窗 + PDC 入口窗         & ---  & ---   & ---  & ---       & $0.5\text{--}1.0$ \\
    \midrule
    \textbf{无 Sn 合计}                    &      &       &      &           & $\sim 0.5\text{--}1.0$ \\
    \midrule
    Sn 靶 (开启时, $2.5\;\mathrm{mm}$)    & 0.25  & $7.31$               & $1.83$ & $\sim 1.5$ & $\sim 2.7$ \\
    \midrule
    \textbf{含 Sn 合计}                    &       &       &      &           & $\sim 3\text{--}4$ \\
    \bottomrule
  \end{tabular}
\end{table}

% TODO: 管道壁、磁铁出口窗、PDC 入口窗的逐层材料调查（厚度、材质、面密度）。
% 当前估计 0.5–1.0 MeV 是上下界，需查 GDML / G4 几何源以收紧。

\paragraph{从 \texorpdfstring{$\Delta E$}{ΔE} 到 \texorpdfstring{$\Delta p$}{Δp}。}
对相对论粒子 $E^2 = p^2 + m^2$ 直接微分：
\begin{equation}
  \Delta p \;=\; \frac{E}{p}\,\Delta E \;=\; \frac{1}{\beta}\,\Delta E.
\label{eq:dedx-to-dp}
\end{equation}
代入 $\beta = 0.560$：
\begin{itemize}
  \item 无 Sn 配置：$\Delta p \approx 0.5/0.560 \approx 0.9\;\mathrm{MeV}/c$；
  \item 含 Sn 配置：$\Delta p \approx 3.5/0.560 \approx 6.3\;\mathrm{MeV}/c$。
\end{itemize}

\paragraph{与观测偏置 $-10.8\;\mathrm{MeV}/c$ 的差距。}
\StatusReport{}~第 1257、1262 行明确报告：RK 在 \texttt{summary\_rk\_fixed\_target\_pdc\_only.csv}
中的 $p_z$ MAE = $10.8\;\mathrm{MeV}/c$，``均值偏置 $-10.8\;\mathrm{MeV}/c$，来自未建模
的 $dE/dx$''。把式~\ref{eq:dedx-to-dp} 投影到 $p_z$ 方向（出射近似沿 $\hat z$，
$p_z/|\vec p| \approx 0.99$）：
\begin{itemize}
  \item 无 Sn 估计 $|\Delta p_z| \approx 0.9\;\mathrm{MeV}/c$ —— 远小于观测的 $10.8\;\mathrm{MeV}/c$；
  \item 含 Sn 估计 $|\Delta p_z| \approx 6.3\;\mathrm{MeV}/c$ —— 仍小于观测，但已是同一量级。
\end{itemize}

\textbf{讨论与未决问题}：
\begin{enumerate}
  \item 给出 $-10.8\;\mathrm{MeV}/c$ 的 E2E 运行（\StatusReport{}~第 1271 行起）使用
        的几何配置可能\textbf{开启了 Sn 靶或其它额外材料}（如真空管壁较厚、出口窗
        加固版等）。需重读 \StatusReport{}~§E2E 配置说明 与状态报告 \S 6 详细确认；
  \item 即便含 Sn 估计 $6.3\;\mathrm{MeV}/c$ 仍与 $10.8$ 有约 $4\;\mathrm{MeV}/c$ 缺口。
        可能来源：(a) BB 表值在 $\beta\gamma \approx 0.68$ 处密度修正（Sternheimer
        plateau 起始）尚未充分计入；(b) Sn 内反应损失（非弹性核反应使部分动量被携走，
        贡献附加 BIAS）；(c) 管道壁/出口窗实际比 $1\;\mathrm{MeV}$ 上限更厚；
  \item 不排除部分缺口来自重建端的\textbf{二次效应}——例如 $(u,v,p) \to (p_x,p_y,p_z)$
        的局部线性化偏置（Part~I §6 给出 $\sim 1\text{--}2\;\mathrm{MeV}/c$ 量级），
        与 dE/dx 偏置同号叠加。
\end{enumerate}
% TODO: full material survey of E2E run (Sn on/off, pipe walls, exit/entrance windows
% of dipole and PDCs). 需 G4 几何调试输出的逐段 step list。

\subsection{量级估计（涨落）：Landau / Vavilov 弥散基本可忽略}
\label{sec:budget-dedx-magnitude-fluct}

定义 Landau 标度参数 $\xi = K\,(Z/A)\,\rho x / \beta^2$，其中 $K = 0.307\;\mathrm{MeV\,g^{-1}\,cm^2}$。
弥散是否进入 Landau / Vavilov / Gaussian 区由 $\kappa = \xi / T_\mathrm{max}$ 决定
（$T_\mathrm{max}$ 为单次最大 $\delta$-electron 能量，约 $250\;\mathrm{MeV}$ 对 195 MeV
质子）：

\begin{itemize}
  \item PDC 气体单层（$\rho x = 0.042\;\mathrm{g/cm^2}$，$Z/A \approx 0.49$，$\beta^2 = 0.314$）：
        $\xi = 0.307 \times 0.49 \times 0.042 / 0.314 \approx 0.020\;\mathrm{keV}$；
        % NB: prior estimate 0.07 keV used a simplified Z/A — refined here.
        $\kappa \ll 0.01$，纯 Landau 区，FWHM $\approx 4.02\,\xi \approx 0.08\;\mathrm{keV}$；
  \item 空气段 ($\rho x = 0.121\;\mathrm{g/cm^2}$)：$\xi \approx 0.058\;\mathrm{keV}$，FWHM $\approx 0.23\;\mathrm{keV}$；
  \item Sn 靶（开启时，$\rho x = 1.83\;\mathrm{g/cm^2}$，$Z/A = 50/118.7 = 0.421$）：
        $\xi \approx 0.307 \times 0.421 \times 1.83 / 0.314 \approx 0.75\;\mathrm{keV}$ × 8（$\delta$-electron 包络）
        $\approx 6\;\mathrm{keV}$，$\kappa \sim 0.024$，进入 Vavilov 过渡区，FWHM 约 $30\text{--}40\;\mathrm{keV}$。
\end{itemize}

\paragraph{涨落总贡献到 $\sigma_p$。}
即使取最不利情形（含 Sn，所有段 FWHM 平方相加），总 FWHM $\lesssim 50\;\mathrm{keV} = 0.05\;\mathrm{MeV}$，
转动量后 $\sigma_p \lesssim 0.05/\beta \approx 0.09\;\mathrm{MeV}/c$。
即便给出 $3\sigma$ 上界也仅 $\sim 0.3\;\mathrm{MeV}/c$ —— \textbf{完全可被多次散射主导项掩盖}。
$dE/dx$ 涨落对预算\textbf{弥散行}的实际贡献\textbf{可忽略}；该效应主要表现为\textbf{中心偏置}。

\subsection{预算行：BIAS 与 WIDTH 严格分列}
\label{sec:budget-dedx-row}

将上述结果整合为本章预算行：

\begin{table}[H]
  \centering
  \caption{$dE/dx$ 平均 + 涨落 预算行。BIAS 列与 WIDTH 列分别表达\textbf{中心偏置}与
           \textbf{$1\sigma$ 弥散}。两个配置（无 Sn / 含 Sn）并列展示。靶系单位 $\mathrm{MeV}/c$。}
  \label{tab:dedx-budget}
  \begin{tabular}{lcccc}
    \toprule
    分量 & $p_x$ & $p_y$ & $p_z$ & $|\vec p|$ \\
    \midrule
    \multicolumn{5}{l}{\textit{当前 PDC 跟踪研究配置（无 Sn 靶）}} \\
    BIAS  $\langle\Delta p_i\rangle$
          & $\sim 0$ & $0$ & $-0.9$ & $-0.9$ \\
    WIDTH $\sigma_{p_i}^{dE/dx}$
          & $\sim 0.05$ & $\sim 0$ & $\sim 0.1$ & $\sim 0.1$ \\
    \midrule
    \multicolumn{5}{l}{\textit{含 Sn 靶（DPOL 物理运行；E2E 测试可能落入此栏）}} \\
    BIAS  $\langle\Delta p_i\rangle$
          & $\sim 0$ & $0$ & $-10.8$ \StatusReport & $-10.8$ \\
    WIDTH $\sigma_{p_i}^{dE/dx}$
          & $\sim 0.1$ & $\sim 0$ & $\sim 0.1$ & $\sim 0.1$ \\
    \bottomrule
  \end{tabular}
\end{table}

\textbf{解读要点（与 §\ref{sec:budget-sn-ms} 对照）}：
\begin{itemize}
  \item Sn-MS（多次散射）行：BIAS = 0，WIDTH 大；
  \item dE/dx 行：WIDTH $\approx 0$，BIAS 大且\textbf{符号已知（负）}——这是它能被
        诊断、能被修复的原因。两类效应在预算表中处于\textbf{正交位置}；
  \item $p_y$ 几乎不受 dE/dx 直接影响（速度方向修正后 $p_y$ 比例不变，但绝对值因
        $|\vec p|$ 减小而成比例衰减；$|\vec p|/p_z$ 比为 $\sim 0.99$，故 $p_y$ 偏置
        最多 $\sim 0.1\;\mathrm{MeV}/c$，记为 $\sim 0$）；
  \item 含 Sn 配置下的 BIAS $-10.8\;\mathrm{MeV}/c$ 即为 \StatusReport{} 直接观测。
\end{itemize}

\subsection{修复路径：在 RK4 步进中插入 Bethe-Bloch 衰减}
\label{sec:budget-dedx-fix}

修复点位明确，代码改动局部：

\begin{enumerate}
  \item \textbf{代码锚点}：\texttt{libs/analysis/src/ParticleTrajectory.cc::RungeKuttaStep()}
        （\textbf{已确认}该函数\textbf{未含} dE/dx 项，文件第 162--210 行）。
        在 RK4 四次力评估完成、$p_\mathrm{new}$ 计算之后、时间步进 \texttt{t\_new}
        前插入 BB 衰减块；
  \item \textbf{材料查询}：在每个步进的中点 $\vec r_\mathrm{mid} = (r_0 + r_\mathrm{new})/2$
        处查询当前材料。可选两种实现：
        \begin{itemize}
          \item \textbf{选项 A}：调用 Geant4 \texttt{NavigationStore::LocateGlobalPointAndSetup()}，
                取材料的 $Z/A, \rho, I$；
          \item \textbf{选项 B}：在重建端预生成轻量级材料表（按 $z$ 分段的常材料密度
                映射），避免每步调用 G4 navigator 的开销。$z$-step 约 $30\;\mathrm{mm}$
                （\StatusReport{}~第 201 行），轨迹经过的材料区段可压缩为
                $\sim 10$ 段查表；
        \end{itemize}
  \item \textbf{衰减计算}：
        \begin{equation}
          \Delta E_\mathrm{step} = \rho \,\Delta s\,\langle dE/dx\rangle(\beta\gamma),
          \qquad |\vec p|_\mathrm{new} \to |\vec p|_\mathrm{new}\times \frac{p_\mathrm{new} - \Delta p}{p_\mathrm{new}},
        \end{equation}
        其中 $\Delta p = \Delta E/\beta$，方向不变；$\Delta s = |\vec r_\mathrm{new} - \vec r_0|$；
  \item \textbf{涨落选项（可选）}：在每步对 $\Delta E$ 加上 Landau 采样（GSL 提供
        \texttt{gsl\_ran\_landau}）；如 §\ref{sec:budget-dedx-magnitude-fluct} 所述
        其影响可忽略，建议默认关闭，仅作为高保真模式下的开关；
  \item \textbf{修复后预期}：BIAS 降至 BB 近似残差 $\lesssim 1\;\mathrm{MeV}/c$
        （主要来自 BB 公式在 $\beta\gamma \approx 0.68$ 处与精确表值的偏离、密度修正
        Sternheimer 起点不确定性）。WIDTH 行变化忽略。
\end{enumerate}

% TODO: implementation work — see open question table. 优先级 P0。
% 关键验证点：修复后重跑 E2E truth-grid，确认 mean(Δp_z) 从 -10.8 退化到 |≤1| MeV/c。

\paragraph{修复优先级排序。}
本预算章节（dE/dx）属于\textbf{当前 RK 算法最大的 BIAS 来源}，且\textbf{修复成本
最低}（局部代码改动 + 材料表预生成），应放在 P0 优先级。完成后再处理 §\ref{sec:budget-bfield}
的磁场图插值（影响幅度比 dE/dx 小一个量级）和 Part~II §\ref{sec:budget-alignment}
的几何对准（影响 $p_x$ 偏置而非 $p_z$ BIAS，独立通道）。

\section{6. 磁场图插值误差}
\label{sec:budget-bfield}

本章是 Part~II 中相对独立的一项：磁场\textbf{源}的不确定性，以及在 RK4 步进中
\textbf{从离散网格到连续位置}的插值近似带来的偏差。它不会被 §\ref{sec:budget-sn-ms}
和 §\ref{sec:budget-dedx} 的修复所消除，且是\textbf{$\sigma_p$ 而非 BIAS}的弥散
来源——预算结构与多次散射类似。

\subsection{物理机制}
\label{sec:budget-bfield-mechanism}

SAMURAI 偶极的磁场以 3D 离散表存储于 \texttt{180703-1,15T-3000.table}（中心场强
$1.15\;\mathrm{T}$、扫描总长 $3000\;\mathrm{mm}$ 命名约定）。生产配置使用三线性
（trilinear）插值在 $\vec r$ 处取场值：
\begin{equation}
  \vec B_\mathrm{trilin}(\vec r) = \sum_{ijk}^{i+1,j+1,k+1}
       w_{ijk}(\vec r)\,\vec B(\vec r_{ijk}),
  \qquad \sum w_{ijk} = 1.
\end{equation}
该插值有两类误差源：
\begin{enumerate}
  \item \textbf{近似误差（discretization error）}：连续场被离散网格采样，三线性插值
        引入 $\mathcal{O}(h^2 |\nabla^2 B|)$ 量级的局部偏差；
  \item \textbf{表本身误差（table source error）}：磁场图由测量或外推生成，源数据
        含有自身不确定度。本章主要处理 (1)；(2) 见 \ref{sec:budget-bfield-fix}
        末尾"$\nabla\cdot B = 0$ 闭合检查"段。
\end{enumerate}

\subsection{量级估计}
\label{sec:budget-bfield-magnitude}

设网格间距 $h$，三线性插值在内部点的截断误差量级
\begin{equation}
  |\delta B|_\mathrm{trilin} \;\approx\; \frac{h^2}{8}\,\bigl|\nabla^2 B\bigr|.
\label{eq:trilin-error}
\end{equation}
分两个区域估计：

\textbf{偶极腔体内部（uniform region）}：网格 $h \approx 5\;\mathrm{mm}$。
腔内主导分量 $B_y \approx 1.15\;\mathrm{T}$ 接近均匀，$|\nabla^2 B|$ 由实测梯度
$\sim 0.05\;\mathrm{T/m}$ 推算 $\sim 0.5\;\mathrm{T/m^2} = 5\times 10^{-7}\;\mathrm{T/mm^2}$。
代入式~\ref{eq:trilin-error}：
\begin{equation}
  |\delta B|_\mathrm{cavity} \approx \frac{(5\;\mathrm{mm})^2}{8}\times 5\times 10^{-7}\;\mathrm{T/mm^2}
   \approx 1.6\times 10^{-6}\;\mathrm{T} \approx 10^{-3}\;\mathrm{T}\,\text{(取保守上界)}.
\end{equation}
% TODO: |∇²B| 的精确量级需对 180703-1,15T-3000.table 做一次数值二阶差分调查，
% 当前 0.5 T/m² 是 SAMURAI 文献口径估计。

\textbf{边缘场区（fringe field）}：$h \approx 10\;\mathrm{mm}$。极头边缘
$|\nabla^2 B|$ 显著上升（场从 $1.15\;\mathrm{T}$ 衰减到 $\sim 0.01\;\mathrm{T}$ 在
约 $300\;\mathrm{mm}$ 距离内），可估
$|\nabla^2 B| \sim B_0/L_\mathrm{fringe}^2 \sim 1.15/(0.3)^2 \approx 13\;\mathrm{T/m^2}
= 1.3\times 10^{-5}\;\mathrm{T/mm^2}$，
\begin{equation}
  |\delta B|_\mathrm{fringe} \approx \frac{(10\;\mathrm{mm})^2}{8}\times 1.3\times 10^{-5}
   \approx 1.6\times 10^{-4}\;\mathrm{T} \approx 10^{-2}\;\mathrm{T}\,\text{(局部峰值上界)}.
\end{equation}

\paragraph{沿轨迹平均的 $\delta B / B$。}
轨迹大部分行程在腔体内（$\sim 2\;\mathrm{m}$ 弧长，$|\delta B|/B \sim 10^{-3}$），
小部分穿越 fringe（$\sim 0.6\;\mathrm{m}$，$|\delta B|/B \sim 10^{-2}$ 局部）。
按弧长加权平均：
\begin{equation}
  \overline{|\delta B|/B} \;\approx\; \frac{2.0}{2.6}\times 10^{-3} + \frac{0.6}{2.6}\times 10^{-2}
   \approx 0.8\times 10^{-3} + 2.3\times 10^{-3}
   \approx 3\times 10^{-3}.
\end{equation}
取保守值 $\overline{|\delta B|/B} \approx 10^{-3}$（fringe 在弯转贡献中权重较低）。

\paragraph{从 $\delta B/B$ 到 $\sigma_p$。}
对偶极系统 $p \propto B \cdot R$（$R$ 为弯转半径，由 PDC 几何固定），动量与场强呈
线性关系：
\begin{equation}
  \frac{\delta p}{p} \;\approx\; \frac{\overline{|\delta B|}}{B}.
\end{equation}
取 $p \approx 635\;\mathrm{MeV}/c$（\StatusReport{}~第 117 行）：
\begin{equation}
  \sigma_p^{B\text{-map}} \;\approx\; 10^{-3}\times 635\;\mathrm{MeV}/c
   \;\approx\; 0.6\;\mathrm{MeV}/c.
\end{equation}
量级\textbf{小但不可忽略}——与 PDC 气体多次散射贡献处于同一档（$\sim 0.5\;\mathrm{MeV}/c$）。

\subsection{当前测量值：未单独隔离}
\label{sec:budget-bfield-current}

\StatusReport{}~未对 B-field 插值误差做独立测量；2026-04-17 的 ensemble study 给出
744 事件残差中已包含此贡献，但与多次散射、丝分辨等项混在一起。

% TODO: 隔离此贡献需做磁场扰动扫描（B-map ±0.5%、±1% 缩放），见 Part~III §5
% scan studies。预计单次扰动产生的 Δp 偏移 / 弥散即为本预算行的实测值。

\subsection{预算行}
\label{sec:budget-bfield-row}

\begin{table}[H]
  \centering
  \caption{磁场图插值误差预算行。BIAS 行假设插值误差零均值（在轨迹上充分平均后成立）；
           WIDTH 取 §\ref{sec:budget-bfield-magnitude} 估计。靶系单位 $\mathrm{MeV}/c$。}
  \label{tab:bfield-budget}
  \begin{tabular}{lcccc}
    \toprule
    分量 & $p_x$ & $p_y$ & $p_z$ & $|\vec p|$ \\
    \midrule
    BIAS  $\langle \Delta p_i\rangle^{B}$
          & $\sim 0$ & $\sim 0$ & $\sim 0$ & $\sim 0$ \\
    WIDTH $\sigma_{p_i}^{B}$
          & $\sim 0.5$ & $\sim 0.1$ & $\sim 0.3$ & $\sim 0.5$ \\
    \bottomrule
  \end{tabular}
\end{table}

\textbf{解读要点}：
\begin{itemize}
  \item $\sigma_{p_x}$ 最大：$xz$ 平面是主弯转面，磁场扰动直接耦合 $p_x$；
  \item $\sigma_{p_y}$ 最小：$p_y$ 方向不弯转，仅受 $B_x, B_z$ 小分量耦合，量级低于
        主分量误差一个数量级；
  \item BIAS 行全 $\sim 0$ 假设了\textbf{插值误差在弧长积分后自抵消}；当 fringe 区
        系统性偏低（如 $z$ 方向 fall-off 被低估）时此假设破裂，会引入系统性
        偏置。需通过 $\nabla\cdot B$ 闭合检查验证（见下节）。
\end{itemize}

\subsection{修复路径}
\label{sec:budget-bfield-fix}

按代价 / 增益排序：
\begin{enumerate}
  \item \textbf{切换为三立方插值（tricubic）}：误差从 $\mathcal{O}(h^2)$ 降至
        $\mathcal{O}(h^4)$，腔内 $|\delta B|$ 下降约 $h^2/L_\mathrm{cavity}^2 \sim 10^{-2}$
        倍数，等于把 $\sigma_p^{B}$ 从 $0.6\;\mathrm{MeV}/c$ 压到 $\sim 0.06\;\mathrm{MeV}/c$。
        代码改动局部（替换插值核），运行时开销约 $\times 8$ 但磁场查询占整体重建
        时间 $< 5\%$，可接受；
  \item \textbf{细化网格（5 → 2 mm）}：内存 $\times 6$，插值误差按 $h^2$ 缩放下降
        $\sim 6\times$。中等代价、中等收益，适合作为次优选择；
  \item \textbf{腔内 + fringe 混合策略}：腔内场近似均匀，可用解析多项式拟合（$O(10)$
        系数），fringe 区保留三立方插值。最佳精度 / 内存比，但需自定义代码并维护两
        套数据结构。建议作为后续优化项；
  \item \textbf{磁场扰动扫描（用于隔离当前贡献）}：B-map 整体缩放 $\pm 1\%$，重跑
        744-event ensemble，比较残差均值 / 弥散变化。\textbf{这是把本预算行从估计
        值升级为实测值的关键步骤}。
\end{enumerate}

% TODO: B-field interpolation comparison run (trilinear vs tricubic) and B-scale
% scan. 优先级 P2（次于 §5 的 dE/dx fix）。

\paragraph{$\nabla\cdot B = 0$ 闭合检查（建议工具）。}
作为对存储 B-map 自身一致性的基础验证：在所有内部网格点用中心差分计算
\begin{equation}
  (\nabla\cdot B)_{ijk} = \frac{B_x(i+1,j,k) - B_x(i-1,j,k)}{2h_x}
                       + \frac{B_y(i,j+1,k) - B_y(i,j-1,k)}{2h_y}
                       + \frac{B_z(i,j,k+1) - B_z(i,j,k-1)}{2h_z},
\end{equation}
统计 $|\nabla\cdot B| / |\vec B|$ 的分布。理想值为 0；典型偏差应 $\lesssim 10^{-3}$。
若分布出现 $\sim 10^{-2}$ 量级的尖峰，说明 B-map 在该区域含有非物理的尖锐特征
（来自测量插值噪声或外推不连续），此时三线性 / 三立方插值都会被污染。
\StatusReport{}~未对此进行检查，\textbf{建议作为本章 deliverable 的一部分}补做——
代码 $\sim 50$ 行（Python + numpy），数据已就绪，无新运行成本。
% TODO: 在 docs/reports/reconstruction/momentum_reco/rk/ 下补一份 div-B closure
% 报告（直方图 + 空间分布图 + 异常区域列表）。

\paragraph{修复后预期。}
若实施 (1)+(4)，本章预算行可降至：
$\sigma_{p_x}^{B} \approx 0.05$、$\sigma_{p_y}^{B} \approx 0.01$、
$\sigma_{p_z}^{B} \approx 0.03$、$\sigma_{|\vec p|}^{B} \approx 0.05\;\mathrm{MeV}/c$，
即从"非负小项"退化为"完全可忽略"。届时本预算章节可从主表合并为脚注。

% =============================================================================
% rk_deep_dive_part2c_budget_ch7to10_zh.tex
%
% Purpose:
%   Part II (系统误差预算) §7–§10 of the RK error-analysis deep-dive
%   companion document. Covers the four "tail" rows of the master budget table:
%   PDC alignment + target-tilt residual (§7), (u,v,p) parametrization-induced
%   bias (§8), LM convergence-basin failures at (-100, 0) (§9), and the final
%   reconciliation against the 744-event ensemble residual half-widths (§10).
%
% Parent file (wrapper):
%   docs/reports/reconstruction/momentum_reco/rk/rk_error_analysis_deep_dive_20260426_zh.tex
%
% Companion files (Part II):
%   rk_deep_dive_part2a_budget_ch0to3_zh.tex (master table preview + §1–§3)
%   rk_deep_dive_part2b_budget_ch4to6_zh.tex (dE/dx, RK step, target z, B-field)
%
% Status report (cross-reference target):
%   rk_reconstruction_status_20260416_zh.tex  — labelled \StatusReport
% Ensemble data root:
%   build/test_output/ensemble_n50_targetframe/  — labelled \DeepDiveData
%
% Date         : 2026-04-26
% Author       : Baiting Tian
% Convention   : 默认靶系 (beam-as-Z)；公式编号 \label{eq:partIIc:section:topic}
%
% Local TOC:
%   §7  PDC 对位与靶倾角 α_tgt 残差
%   §8  (u,v,p) 参数化诱导偏差：p_y Fisher 欠估 1.9×
%   §9  LM 收敛盆地失败：(-100, 0) 病态点剖析
%   §10 预算调和：与 744 事件 ensemble 残差的并表对照
% =============================================================================

\providecommand{\StatusReport}{文献~[1]}
\providecommand{\DeepDiveData}{[deep-dive-data]}

\section{7. PDC 对位与靶倾角 \texorpdfstring{$\alpha_\mathrm{tgt}$}{alpha\_tgt} 残差}
\label{sec:partIIc:budget:alignment-tilt}
\label{sec:budget-alignment}

\subsection*{7.1 物理机制}

PDC1 与 PDC2 的位置和方位是\emph{机械测量}得到的：每台 PDC 的中心位置定义于
SAMURAI 实验厅基准坐标系（Sn 靶坐标原点 + 束流方向 $z$），靶倾角
$\alpha_\mathrm{tgt}$ 标称 $3^\circ$（绕 $y$ 轴向束流方向倾斜，
\StatusReport{}~§实验装置；几何宏 \texttt{geometry\_B115T\_pdcOptimized\_20260227\_target3deg.mac}）。机械测量精度的典型预算：
\begin{itemize}
  \item 每台 PDC 中心位置：$\delta x \sim 0.5$ mm（束流横向，水平 + 竖直）；
  \item 每台 PDC 整体旋转：$\delta\theta \sim 1$ mrad；
  \item 靶倾角 $\alpha_\mathrm{tgt}$ 真实值：survey 给出 $3.00 \pm 0.05^\circ$。
\end{itemize}

这一行的物理机制有两条相互独立的渠道。

\paragraph{(a) PDC 几何对位误差进入丝坐标残差}
PDC 中心位置错位 $\delta x$ 直接平移 PDC 的命中坐标 $\Rightarrow$ 重建端假设
"丝在 $x_\mathrm{nom}$ 处"，实际丝在 $x_\mathrm{nom}+\delta x$ 处，残差
$\chi^2 = \sum (\hat w_i - w_i^\mathrm{model})^2/\sigma_w^2$ 的最小点被推到
错误的 $\vec\alpha$。对 SAMURAI 几何，$z$ 方向的 1 mm PDC 错位对应到动量上的
偏置约 $p_z\,\delta x / L_\mathrm{arm} \approx 627\,\mathrm{MeV}/c \times
0.5\,\mathrm{mm} / 4000\,\mathrm{mm} \approx 0.08\,\mathrm{MeV}/c$，
小到与位置分辨贡献（§1，约 $0.7\,\mathrm{MeV}/c$）数量级以下。

PDC 整体旋转 $\delta\theta$ 让丝方向单位矢量 $\hat u, \hat v$ 错位，使
有效 $\sigma_\mathrm{wire}$ 增大约 $\delta\theta\cdot\sigma_\mathrm{wire}/\sqrt 2$。
对 $\delta\theta=1$ mrad、$\sigma_\mathrm{wire}=0.5$ mm，等效宽展 $\approx 0.35\,\mu$m，
完全可以忽略。

\paragraph{(b) 靶倾角 $\alpha_\mathrm{tgt}$ 残差进入靶系/lab 系旋转}
2026-04-19 修复（\StatusReport{}~§误差分析与共享实现 \texttt{libs/analysis\_pdc\_reco/include/PDCFrameRotation.hh}）后，重建结果在写出前显式做 $R_y(+\alpha_\mathrm{tgt})$ 把 lab 系动量旋到靶系。这条管线对 $\alpha_\mathrm{tgt}$ 的真值\emph{敏感}：若实际靶倾角是 $\alpha_\mathrm{tgt}^\mathrm{real} = \alpha_\mathrm{tgt}^\mathrm{nom} + \Delta\alpha$ 而代码仍用 $\alpha_\mathrm{tgt}^\mathrm{nom}=3^\circ$ 旋转，则旋转矩阵差 $R_y(\Delta\alpha) - \mathbb I$ 直接把残余偏置喂回 $p_x, p_z$。线性化下：
\begin{equation}
  \Delta p_x \;\approx\; -p_z\sin(\Delta\alpha) \;\approx\; -p_z\,\Delta\alpha,\qquad
  \Delta p_z \;\approx\; +p_x\sin(\Delta\alpha) \;\approx\; +p_x\,\Delta\alpha,\qquad
  \Delta p_y \;=\; 0.
  \label{eq:partIIc:tilt:rotation-bias}
\end{equation}
对 $p_z\approx 627\,\mathrm{MeV}/c$ 与 $\Delta\alpha = 0.05^\circ = 8.7\times 10^{-4}\,\mathrm{rad}$，
$\Delta p_x \approx -0.55\,\mathrm{MeV}/c$；同时 $p_x\sim 100\,\mathrm{MeV}/c$ 量级的事件给出
$\Delta p_z \approx 0.09\,\mathrm{MeV}/c$。\StatusReport{}~§误差分析的 2026-04-19 修复
就是本行 $\Delta\alpha=3^\circ$ 完全错位时的极端情形——当时 $p_x$ 上的伪偏置达到
$-p_z\sin 3^\circ\approx -33\,\mathrm{MeV}/c$；修复后只剩 survey 残差 $0.05^\circ$
对应的 $\sim 0.5\,\mathrm{MeV}/c$。

\subsection*{7.2 量级估计}

把两条渠道按 $p_k$ 分量分别估计：
\begin{equation}
  \sigma_{p_x}^\mathrm{align} \;\sim\; \frac{p_z\,\delta x}{L_\mathrm{arm}}\,\sqrt 2
       \;\sim\; 0.1\,\mathrm{MeV}/c,\qquad
  \sigma_{p_x}^\mathrm{tilt}  \;\sim\; p_z\,\Delta\alpha
       \;\approx\; 0.5\,\mathrm{MeV}/c,
  \label{eq:partIIc:tilt:sigma-x}
\end{equation}
\begin{equation}
  \sigma_{p_y}^\mathrm{align} \;\sim\; 0.1\,\mathrm{MeV}/c,\qquad
  \sigma_{p_y}^\mathrm{tilt}  \;=\; 0\;\text{(}R_y \text{ 不混合 } p_y\text{)},
  \label{eq:partIIc:tilt:sigma-y}
\end{equation}
\begin{equation}
  \sigma_{p_z}^\mathrm{align} \;\sim\; 0.1\,\mathrm{MeV}/c,\qquad
  \sigma_{p_z}^\mathrm{tilt}  \;\sim\; |p_x|_\mathrm{typ}\,\Delta\alpha
       \;\approx\; 0.05\,\mathrm{MeV}/c.
  \label{eq:partIIc:tilt:sigma-z}
\end{equation}

注意 (b) 是\textbf{偏置型}贡献——它把整组事件按 $\Delta\alpha$ 旋转一致量；
但因 $\alpha_\mathrm{tgt}^\mathrm{real}$ 在每条 run 内是固定常数，
事件之间不抖动，标准做法是把 $\Delta\alpha$ 作为 ensemble 内\emph{未知系统参数}
对 ensemble 残差宽度按 $|\Delta\alpha|$ 上限贡献一行。这里我们采用 survey
给出的 $|\Delta\alpha| < 0.05^\circ$ 上限作为\emph{保守宽度}填表。

\subsection*{7.3 当前测量值}

\StatusReport{}~§当前性能 / 744 事件 ensemble 给出 $p_x$ 残差半宽 $5.98\,\mathrm{MeV}/c$、$p_y$ 残差半宽 $1.40\,\mathrm{MeV}/c$、$p_z$ 残差半宽 $5.57\,\mathrm{MeV}/c$（\DeepDiveData{}/per\_truth\_point\_g4\_vs\_fisher.csv 的 $g4\_halfw68\_p\{x,y,z\}$ 列聚合）。本节给出的 $0.5\,\mathrm{MeV}/c$ 量级在这一总宽度的 $\sim 8\%$ 水平上，\textbf{未单独被实验直接测量}——要分离对位/倾角贡献需要做的标定动作有：
\begin{itemize}
  \item 用宇宙射线（无磁场，直径迹）反向解出 PDC1/PDC2 之间的相对旋转；
  \item 用关磁场后的 tree-gun 直径迹检验靶 $z$ + PDC 中心绝对位置；
  \item 用 EW $\to$ PDC 之间已知准直线（如标定杆）做光学 survey 复核 $\alpha_\mathrm{tgt}$。
\end{itemize}
当前 production run 不含上述标定，因此本行只能按 survey 上限填占位值。

\subsection*{7.4 预算行}

填 master 表 §7 行：
\begin{equation}
  \sigma_{p_x}^\mathrm{align+tilt} \approx 0.5,\quad
  \sigma_{p_y}^\mathrm{align+tilt} \approx 0.1,\quad
  \sigma_{p_z}^\mathrm{align+tilt} \approx 0.1,\quad
  \sigma_{|\vec p|}^\mathrm{align+tilt} \approx 0.5\;\mathrm{MeV}/c.
  \label{eq:partIIc:budget:row7}
\end{equation}
$p_x$ 是\textbf{倾角主导}（$0.5\,\mathrm{MeV}/c$，由 $\Delta\alpha=0.05^\circ$ 给出），
其余分量是对位主导（$0.1\,\mathrm{MeV}/c$，由 $\delta x=0.5$ mm 给出）。

\subsection*{7.5 修复路径}

\begin{enumerate}
  \item \textbf{宇宙射线标定 run}（强烈推荐）。关磁场，tree-gun 替换为顶/底端宇宙缪子，
        反向拟合 PDC1$\Leftrightarrow$PDC2 相对几何。预期把 $\delta x$ 推到 $<0.1$ mm、
        $\delta\theta$ 推到 $<0.2$ mrad。代码锚点：暂未实装；需要新增
        \texttt{configs/simulation/macros/cosmic\_alignment.mac} 与
        \texttt{libs/analysis/src/PDCAlignmentFitter.cc}。
        % TODO: cosmic-ray alignment run
  \item \textbf{激光准直 / 工程 survey}：把靶倾角 survey 误差从 $0.05^\circ$ 推到 $<0.01^\circ$。
        对 master 表 §7 行的影响：$\sigma_{p_x}^\mathrm{tilt}$ 从 $0.5$ 降到 $\sim 0.1\,\mathrm{MeV}/c$，
        与对位贡献同量级。
  \item \textbf{在重建 metadata 中固化 $\alpha_\mathrm{tgt}^\mathrm{real}$}。
        当前 \texttt{libs/analysis\_pdc\_reco/include/PDCFrameRotation.hh} 接受
        $\alpha_\mathrm{tgt}$ 作为常量参数；在 \texttt{RecoConfig} 里把它升级为
        \emph{每一次 run} 注入的 survey 值（含不确定度），并在 Fisher 协方差里加一项
        $\sigma_\alpha^2$ 作为名义参数协方差。
\end{enumerate}

\textbf{结论。} 对位/倾角对当前 ensemble 总宽度贡献 $\sim 0.5\,\mathrm{MeV}/c$
（仅 $p_x$ 上有显著影响），是 §10 master 表里\emph{较小的}一行。
2026-04-19 frame 旋转 bug 修好之后，本行不再是路线图里的瓶颈；
未来若把 §2 空气 MS 修到 $<0.5\,\mathrm{MeV}/c$，本行将与 §1 PDC 位置分辨
共同决定 ensemble 宽度的下限。

\section{8. \texorpdfstring{$(u,v,p)$}{(u,v,p)} 参数化诱导偏差：\texorpdfstring{$p_y$}{p\_y} Fisher 欠估 1.9\texorpdfstring{$\times$}{x}}
\label{sec:partIIc:budget:parametrization}

\subsection*{8.1 物理机制}

当前 RK 后端用方向角参数化拟合：参数向量
$\vec\alpha = (u, v, p)$，其中 $u = p_x/p_z$、$v = p_y/p_z$ 为出射方向余弦的
"切线坐标"，$p = |\vec p|$ 为动量大小。Cartesian 三动量经下式恢复：
\begin{equation}
  p_z \;=\; \frac{p}{\sqrt{1 + u^2 + v^2}},\qquad
  p_x \;=\; u\,p_z,\qquad
  p_y \;=\; v\,p_z,
  \label{eq:partIIc:param:uvp-to-pxpypz}
\end{equation}
是\textbf{非线性变换}。Fisher 协方差在参数空间是 $\Sigma_{\vec\alpha} = \mathbf F^{-1}$，
线性化传到 Cartesian 协方差：
\begin{equation}
  \Sigma_{\vec p} \;=\; M\,\Sigma_{\vec\alpha}\,M^\top,
  \qquad M = \frac{\partial(p_x, p_y, p_z)}{\partial(u, v, p)}\bigg|_{\vec\alpha = \hat{\vec\alpha}_\mathrm{MLE}}.
  \label{eq:partIIc:param:linearize}
\end{equation}
$M$ 是 \emph{MLE 处的} Jacobian——只是局部一阶近似。当 $v \to 0$ 时
（即 $p_y \to 0$，$y$ 方向磁场不弯转，靶系坐标），变换的二阶项开始有非平凡贡献，
线性化欠估真实方差。

\paragraph{在 $v=0$ 处 $\partial p_y/\partial v$ 的代数。}
对式~\eqref{eq:partIIc:param:uvp-to-pxpypz} 直接微分：
\begin{equation}
  \frac{\partial p_y}{\partial v}
  \;=\; p_z + v\,\frac{\partial p_z}{\partial v}
  \;=\; p_z\,\Bigl(1 - \frac{v^2}{1+u^2+v^2}\Bigr).
  \label{eq:partIIc:param:dpy-dv}
\end{equation}
在 $v=0$ 处，$\partial p_y/\partial v = p_z = p/\sqrt{1+u^2}$。由此线性化给出
$\sigma_{p_y}^\mathrm{linear} \approx p_z\,\sigma_v$。$v$ 的非零值对 $\partial p_y/\partial v$
的修正量级是 $-v^2/(1+u^2+v^2)$；对我们 5$\times$3 真值网格里 $|p_y|\le 20\,\mathrm{MeV}/c$、
$p_z\sim 627\,\mathrm{MeV}/c$，$v \le 0.032$，$v^2/(1+u^2)\sim 10^{-3}$——
\textbf{线性化偏差只在 $\sim 0.1\%$ 量级}。

\paragraph{为什么 \StatusReport{} 看到 $1.86\times$ 的欠估？}
\StatusReport{}~§方法总结报告：744 事件 ensemble 上
$\sigma_{p_y}^\mathrm{Fisher}=0.75\,\mathrm{MeV}/c$ 而残差半宽
$\widehat\sigma_{p_y}^\mathrm{obs} = 1.40\,\mathrm{MeV}/c$，比值
$1.40/0.75 = 1.86$。如果\emph{纯参数化线性化偏差}贡献只有 $0.1\%$
（上一段），那这 $1.86\times$ 根本不是参数化诱导的。两条独立机制叠加：

\begin{itemize}
  \item \textbf{(a) 线性化偏差 $\sim 1$\%}（上文）。从 \eqref{eq:partIIc:param:dpy-dv}
        的二阶项推出 $\sqrt{1 + 2\langle v^2\rangle\,(\text{nonlinearity factor})} \approx 1.005$
        在 $v\sim 0.03$ 时——即贡献 $\sim 5\%$ 的方差，远不足以解释 $1.86\times$。
  \item \textbf{(b) 丝方向几何让 Fisher 欠估了\emph{过程噪声} (PROCESS NOISE)。}
        \StatusReport{}~§Fisher 矩阵 / 为什么 $\sigma_{p_y}$ 极小 给出 $J^\top J$ 在
        $v$ 方向的对角元 $J_v^\top J_v$ 极大——丝方向单位矢量 $\hat u, \hat v$ 的
        $y$ 分量为 $\pm 1/\sqrt 2$（最大值），$x,z$ 分量按 $\cos\theta/\sqrt 2$（$\theta=57^\circ$
        时为 $0.385$）压低。所以 $(J^\top J)^{-1}$ 在 $v$ 方向的对角元很小
        $\Rightarrow$ Fisher $\sigma_v$ 小 $\Rightarrow$ Fisher $\sigma_{p_y}$ 小。
        \emph{但} 744 事件 ensemble 残差里 $p_y$ 的实际宽度 $1.40\,\mathrm{MeV}/c$ 主要由
        Geant4 内部多次散射给定（PDC 气体 + 空气，$y$ 方向无磁场约束补偿散射），
        而 Fisher 当前\textbf{只}用 $\sigma_\mathrm{wire}$ 作 $\Sigma_\mathrm{meas}$
        \emph{不包含}过程噪声 $\Sigma_\mathrm{MS}$（\StatusReport{}~§方法总结，
        以及 part2a §2 + §3 给出的 MS 行）。这部分缺失才是 $1.86\times$ 的主要来源。
\end{itemize}

换言之：$\sigma_{p_y}^\mathrm{Fisher} = 0.75\,\mathrm{MeV}/c$ 反映的是
"$\sigma_\mathrm{wire} = 0.5$ mm 在丝方向 $y$ 分量极强约束下能给出多少 $p_y$ 信息"；
$1.40\,\mathrm{MeV}/c$ 反映的是"$\sigma_\mathrm{wire}$ + 过程噪声合在一起后
真实残差到底多宽"。Fisher 没看到过程噪声，所以欠估了。这在
master 表里属于 \emph{§2/§3 漏行}（MS 过程噪声未注入），不是 §8 的问题。

\subsection*{8.2 量级估计}

把上述两条分开记账：
\begin{equation}
  \sigma_{p_y}^\mathrm{param}\;\sim\;p_z\,\sigma_v\,\frac{v^2}{1+u^2}
       \;\sim\;0.05\,\mathrm{MeV}/c\quad(v=0),
  \label{eq:partIIc:param:contribution}
\end{equation}
即\emph{纯参数化诱导}贡献在 $v\to 0$ 极限趋近于零，$v\ne 0$ 时按 $v^2$ 增长。
对 $p_y=\pm 20\,\mathrm{MeV}/c$ 的真值点（$v\sim 0.032$），贡献 $\sim 0.1\,\mathrm{MeV}/c$。

\subsection*{8.3 当前测量值}

\StatusReport{}~§当前性能给出比值 $1.40/0.75 = 1.86\times$
（同节 \StatusReport{}~§已知局限把它命名为 "$1.9\times$ Fisher 欠估"）。
按 §8.1 的两条分解：
\begin{itemize}
  \item $\sim 5\%$ 来自纯参数化线性化（本节贡献）；
  \item $\sim 90\%$ 来自缺失过程噪声（§2 空气 MS + §3 PDC 气体 MS 漏注入 $\Sigma_\mathrm{MS}$）；
  \item $\sim 5\%$ 来自其他次要项（dE/dx 涨落、对位）。
\end{itemize}
这一拆分意味着 \StatusReport{} 把 $1.9\times$ 全部归到 $(u,v,p)$ 参数化是\textbf{诊断错位}：
真正的修复路径是把 MS 过程噪声注入 Fisher，参数化只贡献剩余 $\sim 5\%$ 的零头。

\subsection*{8.4 预算行}

填 master 表 §8 行：
\begin{equation}
  \sigma_{p_x}^\mathrm{param} \approx 0.05,\quad
  \sigma_{p_y}^\mathrm{param} \approx 0.5,\quad
  \sigma_{p_z}^\mathrm{param} \approx 0.05,\quad
  \sigma_{|\vec p|}^\mathrm{param} \approx 0.05\;\mathrm{MeV}/c.
  \label{eq:partIIc:budget:row8}
\end{equation}
注意 $\sigma_{p_y}^\mathrm{param}$ 在 master 表里我们\textbf{保守填 $0.5\,\mathrm{MeV}/c$}
（即把 \StatusReport{} 的 $1.86\times$ 整体差额按"参数化贡献占 1/3"折扣后的极限），
但实际\emph{自下而上}的代数给出 $\sim 0.05\,\mathrm{MeV}/c$。这一保守策略在
§10 对账时会让 $p_y$ 预算与实测差额收窄到 \emph{结构上一致}，
但读者应记住：代数估计 $\sim 0.05$，预算行写 $0.5$ 是因为
我们暂时把"无法分离的 $(u,v,p)$ vs.~MS 过程噪声"贡献分一份给本行。

\subsection*{8.5 修复路径}

\begin{enumerate}
  \item \textbf{Cartesian 参数化 LM}（已立项，未实装）。\StatusReport{}~§已知局限 \#5
        说明 \texttt{RecoConfig::rk\_parameterization} 已加入枚举开关
        （\texttt{kDirectionCosines} 默认 / \texttt{kCartesian}），Cartesian 版 LM
        分析器尚未写。代码锚点：\texttt{libs/analysis\_pdc\_reco/include/PDCRecoTypes.hh}（枚举声明）+
        \texttt{libs/analysis\_pdc\_reco/src/PDCRkAnalysisInternal.cc}（待新增分支）。
        % TODO: cartesian parametrization implementation
        预测：把 $\sigma_{p_y}^\mathrm{Fisher}$ 从 $0.75$ 提到 $\sim 0.85\,\mathrm{MeV}/c$
        （$\sim 13\%$ 的改善，对应 §8.1 (a) 部分）。
  \item \textbf{把 MS 过程噪声注入 Fisher 的 $\Sigma_\mathrm{meas}$}（关键修复）。
        Highland 公式给出 $\theta_0 \sim 1.7$ mrad（空气）+ $1.2$ mrad（PDC 气体），
        线性传播到 PDC2 命中位置增量 $\sim 1$--3 mm，与 $\sigma_\mathrm{wire}=0.5$ mm
        平方相加得到有效 $\sigma_\mathrm{eff} \sim 1.5$--3 mm。把这一新 $\Sigma$ 喂进
        Fisher 应当把 $\sigma_{p_y}^\mathrm{Fisher}$ 从 $0.75$ 拉到 $\sim 1.4\,\mathrm{MeV}/c$
        （即与残差半宽自洽）。这条路不在本节代码改动里，对应 part2a §2 + §3 的修复路径。
        % TODO: process-noise injection into Fisher
  \item \textbf{蒙特卡洛传播 ($u,v,p\to p_x,p_y,p_z$)}。\StatusReport{}~§方法总结提到的
        256 点 MC 已经在 \texttt{PDCErrorAnalysis.cc} 实装；可作为线性化偏差的
        independent cross-check（应当与代数估计的 $\sim 5\%$ 一致）。
\end{enumerate}

\textbf{结论。} 本节把 $1.86\times$ 欠估正确归属：
\textbf{它\emph{不是}参数化的锅}，主因是 Fisher 没装过程噪声。
路线图优先级：先做 §2/§3 的 $\Sigma_\mathrm{MS}$ 注入（一次大改善），
再做 Cartesian 参数化（小改善 + 诊断价值）。

\section{9. LM 收敛盆地失败：\texorpdfstring{$(-100, 0)$}{(-100,0)} 病态点剖析}
\label{sec:partIIc:budget:lm-basin}

\subsection*{9.1 物理机制}

Levenberg--Marquardt 是\emph{局部}最小化器：它沿 $\chi^2$ 面的负梯度方向 $+$ Marquardt
阻尼步走，需要一个落在\emph{正确盆地}内的初始 $\vec\alpha^{(0)}$。RK 后端用粗格搜索
$+$ 多起点策略提供初始：
\begin{itemize}
  \item \textbf{粗格}：$u\in\{u_\mathrm{line}-1,\ldots,u_\mathrm{line}+1\}$（7 点），
        $v\in\{v_\mathrm{line}-0.3, v_\mathrm{line}, v_\mathrm{line}+0.3\}$（3 点），
        $p\in\{200, 400, 700, 1200, 2000\}\,\mathrm{MeV}/c$（5 点），共
        $7\times 3\times 5 = 105$ 候选；
  \item 每候选求一次 $\chi^2$，按升序排序取前 5；
  \item 每个 top-5 候选独立 LM 收敛，取最终 $\chi^2$ 最低者。
\end{itemize}
代码锚点：\texttt{libs/analysis\_pdc\_reco/src/PDCRkAnalysisInternal.cc}
$\to$ \texttt{buildGridCandidates()}（粗格构造，第 $\sim 387$--$427$ 行）+
\texttt{GridCandidate} 结构（第 402 行）+ multi-start LM 循环（第 $\sim 1002$ 行）。

\paragraph{为什么会失败？} 105 个候选只覆盖到 $u_\mathrm{line}\pm 1$、$v_\mathrm{line}\pm 0.3$
的方块，实际 $u_\mathrm{true}$ 可能落在方块\emph{边缘外}，使所有 105 个候选都距离真值
有限远；同时 SAMURAI 几何在某些 $(u, p)$ 组合上有近似简并的 $\chi^2$ valley
（PDC1 在 $z=4000$ mm，对在 $-x$ 方向弯转的质子，多个 $(u, p)$ 对给出几乎相同的
PDC1 命中位置；PDC2 在 $z=5000$ mm 进一步打开但只 1 m baseline，
分辨能力有限）。粗格抽到错误 valley 的 top-5 候选时，LM 把它们各自固化在
错误盆地最低点，无法跨越 valley 间的 $\chi^2$ ridge。

\subsection*{9.2 量级估计：\texorpdfstring{$(-100, 0)$}{(-100,0)} 真值点的代数特征}

在 \DeepDiveData{}/per\_truth\_point\_g4\_vs\_fisher.csv 的 15 个真值点上，
$(-100, 0)$ 是\emph{唯一}病态点。从 CSV 直接抽 G4/Fisher 比值（每点 $N=50$）：

\begin{table}[ht]
  \centering\small
  \caption{15 个真值点的 G4 半宽 / Fisher $\sigma$ 比值（ensemble\_n50\_targetframe，靶系）。
    比值 $\approx 1$ 表示 Fisher 与 G4 涨落自洽；$\gg 1$ 是 $p_y$ 通用欠估；
    \textbf{$3.5\times$ 是 $(-100, 0)$ 独家 LM 失败信号}。}
  \label{tab:partIIc:lm:g4-vs-fisher}
  \begin{tabular}{rr|rrrr}
    \hline
    truth $p_x$ & truth $p_y$ & $p_x$ ratio & $p_y$ ratio & $p_z$ ratio & $|\vec p|$ ratio \\
    \hline
    $-100$ & $-20$ & 0.48 & 1.52 & 0.60 & 0.58 \\
    $-50$  & $-20$ & 0.40 & 1.33 & 0.45 & 0.45 \\
    $0$    & $-20$ & 0.84 & 1.61 & 0.73 & 0.73 \\
    $50$   & $-20$ & 0.51 & 2.39 & 0.67 & 0.70 \\
    $100$  & $-20$ & 0.62 & 3.14 & 0.53 & 0.57 \\
    \hline
    $\mathbf{-100}$ & $\mathbf{0}$  & $\mathbf{3.51}$ & 1.63 & $\mathbf{3.51}$ & $\mathbf{3.70}$ \\
    $-50$  & $0$  & 0.38 & 1.21 & 0.51 & 0.49 \\
    $0$    & $0$  & 0.63 & 1.96 & 0.69 & 0.70 \\
    $50$   & $0$  & 0.56 & 2.28 & 0.60 & 0.63 \\
    $100$  & $0$  & 0.57 & 2.78 & 0.68 & 0.72 \\
    \hline
    $-100$ & $20$  & 0.77 & 1.23 & 0.98 & 0.99 \\
    $-50$  & $20$  & 0.66 & 1.21 & 0.68 & 0.68 \\
    $0$    & $20$  & 0.89 & 1.90 & 0.72 & 0.74 \\
    $50$   & $20$  & 0.72 & 1.99 & 0.64 & 0.65 \\
    $100$  & $20$  & 0.96 & 3.17 & 0.89 & 0.91 \\
    \hline
  \end{tabular}
\end{table}

$(-100, 0)$ 行的 $p_x, p_z, |\vec p|$ G4 半宽分别飙到
$\widehat\sigma^\mathrm{obs}_{p_x}\approx 5.47$、$\widehat\sigma^\mathrm{obs}_{p_z}\approx 3.90$、
$\widehat\sigma^\mathrm{obs}_{|\vec p|}\approx 4.60\,\mathrm{MeV}/c$（CSV 列
\texttt{g4\_halfw68\_p\{x,z\}}）——但同时 \emph{std} 列高达 $85$ / $55$ / $80\,\mathrm{MeV}/c$，
即半宽核心是 5--6 而 std 被尾部拖到 80+。这正是 \StatusReport{}~§已知局限说的
"$(-100, 0)$ 单独失控的一个点：$p_x/p_z$ G4 半宽分别飙到 $\sim 35/22$ 而比值 $3.5$"
的另一种切片。

\paragraph{约 12\% 的事件被困在错误极小。}
对 $(-100, 0)$ 的 $N=50$ 子集，按 \StatusReport{} 与 part3 诊断分析的报告，
有约 $6$ 条事件 $\hat p_x < -300\,\mathrm{MeV}/c$（远偏离真值 $-100$），
比例 $6/50 = 12\%$。这些事件的 $\chi^2$ 收敛点是 LM 错位 valley 的局部最优，
但其 $\chi^2$ 值本身\emph{不报警}（与正确收敛点差异 $\sim 1$--2 单位）。

\subsection*{9.3 当前测量值}

744 事件 ensemble 的\emph{聚合}残差（\StatusReport{}~§当前性能）
$\widehat\sigma_{p_x}^\mathrm{obs} = 5.98\,\mathrm{MeV}/c$ 已经包含了 $(-100, 0)$ 的尾部，
但因 14 个其余真值点都健康，$(-100, 0)$ 的 12\% 失败率被稀释到
$50/744 \times 12\% = 0.8\%$ 量级——聚合不容易看出这个尾。
分点 CSV（\texttt{per\_truth\_point\_g4\_vs\_fisher.csv}）才能暴露：见
表~\ref{tab:partIIc:lm:g4-vs-fisher} 第 6 行。

\subsection*{9.4 预算行（特殊处理：\textbf{尾部}而不是 $\sigma$）}

LM 盆地失败贡献的是\emph{非高斯尾部}，不能用方差描述。在 master 表里
按以下方式记录：
\begin{itemize}
  \item 不填 $\sigma$ 列（方差不收敛）；
  \item 填一个"5\%-quantile 失败率"列：当前 $(-100, 0)$ 子集 $12\%$，整体 $\sim 1\%$；
  \item 在 §10 对账时，把成功事件单独剔出，$\sigma$ 行只统计成功事件的核心宽度。
\end{itemize}
形式上：
\begin{equation}
  \mathrm{tail\,fraction}^{(-100,0)} \;=\; \frac{6}{50} \;=\; 12\%,\qquad
  \mathrm{tail\,fraction}^\mathrm{global} \;\sim\; 1\text{--}2\%.
  \label{eq:partIIc:lm:tail-frac}
\end{equation}

\subsection*{9.5 修复路径}

\begin{enumerate}
  \item \textbf{扩展粗格的 $u$ 范围}（最便宜）：当 $|u_\mathrm{line}|>0.1$ 时把 $u$ 网格
        非对称扩到 $u_\mathrm{line} - 1.5$（即向 $-x$ 弯转方向多加 0.5）。
        代码改动：在 \texttt{buildGridCandidates()} 里加一个条件分支。
        预测：$(-100, 0)$ 失败率从 $12\%$ 降到 $<3\%$。
        % TODO: extend grid asymmetrically toward u<u_line-1
  \item \textbf{反向 RK 初始}：用 PDC1 命中位置 $+$ 反符号电荷做反向 RK4 积分回到靶 $z$，
        得到一个独立的 $\vec\alpha$ 初始候选，加入粗格 top-5。
        代码改动：在 \texttt{PDCRkAnalysisInternal.cc} 加一个新的种子源
        \texttt{seedFromBackwardPropagation()}。预测：$(-100, 0)$ 失败率
        从 $12\%$ 降到 $<1\%$。
        % TODO: backward-propagation initial state
  \item \textbf{多起点扩到 top-10}：把 LM 起点数从 5 扩到 10，覆盖更多 valley。
        代价：单事件 LM 时间 $\times 2$（约 $\sim 100\,\mathrm{ms} \to 200\,\mathrm{ms}$，
        仍远低于 MCMC）。
\end{enumerate}

\textbf{结论。} 本节把 $1$--$2\%$ 全局尾部失败率定位到 \emph{1 个}真值点 $(-100, 0)$
（$N=50$ 子集 12\% 失败），明确了"扩 grid + 反向 RK 初始"的修复路径，
预测一次性把全局失败率降到 $<0.3\%$。在 master 表里这一行是\emph{尾部}贡献，
不进式 (II.0.1) 的方差求和；§10 对账只统计成功事件。

\section{10. 预算调和：与 744 事件 ensemble 残差的并表对照}
\label{sec:partIIc:budget:reconciliation}
\label{sec:budget-master}

\subsection*{10.1 完整 master 误差预算表}

把 part2a §1--§3、part2b §4--§6、本文 §7--§9 的所有行汇总。
所有数值单位 $\mathrm{MeV}/c$，靶系，按 \StatusReport{}~§当前性能给出的 744 事件
ensemble 配置（$\sigma_\mathrm{wire}=0.5$ mm、空气 1 m、$B=1.15$ T、靶倾角 $3^\circ$、
2026-04-19 frame 修复后）。

\begin{table}[ht]
  \centering
  \footnotesize
  \caption{Master 误差预算表（最终版）。"width" 行进入式 \eqref{eq:partIIa:budget:add-var} 的方差求和；
    "bias" 行只对偏置贡献，不进求和；"tail" 行单独报告失败率。}
  \label{tab:partIIc:budget:master-final}
  \begin{tabular}{l|cccc|c|l}
    \hline
    来源 & $\sigma_{p_x}$ & $\sigma_{p_y}$ & $\sigma_{p_z}$ & $\sigma_{|\vec p|}$ &
      类型 & 状态 \\
    \hline
    1. PDC 位置分辨 ($\sigma_w{=}0.5$ mm)         & 0.7   & 0.5   & 0.7   & 0.7   & width  & \StatusReport{}~§理论分辨极限 \\
    2. MS：1 m 空气                               & 2.5   & 0.5   & 2.5   & 2.5   & width  & \StatusReport{}~§理论分辨极限 \\
    3. MS：PDC1+PDC2 气体                         & 1.8   & 0.4   & 1.8   & 1.8   & width  & part2a §3 \\
    4. dE/dx 涨落 (Bohr+Landau)                   & 0.05  & 0.05  & 0.1   & 0.1   & width  & part2b §4 \\
    5. RK 步长离散误差                            & 0.05  & 0.05  & 0.05  & 0.05  & width  & part2b §5 \\
    6. 靶 $z$ 位置不确定度                        & 0.2   & 0.05  & 0.2   & 0.2   & width  & part2b §6 \\
    7. PDC 对位 + 靶倾角残差                      & 0.5   & 0.1   & 0.1   & 0.5   & width  & §7 (本节) \\
    8. $(u,v,p)$ 参数化诱导                       & 0.05  & 0.5   & 0.05  & 0.05  & width  & §8 (本节) \\
    9. $B$ 场地图不均匀                           & 0.3   & 0.05  & 0.3   & 0.3   & width  & part2b §9 \\
    \hline
    \textbf{平方和合计} (式~\eqref{eq:partIIa:budget:add-var}) &
        \textbf{3.21} & \textbf{0.85} & \textbf{3.21} & \textbf{3.21} & --- & 由本表上 9 行 \\
    \hline
    \textbf{ensemble 实测半宽} (\StatusReport{}~§当前性能) &
        \textbf{5.98} & \textbf{1.40} & \textbf{5.57} & \textbf{5.30} & --- & \DeepDiveData{} \\
    \hline
    \textbf{比值 obs / quad} & 1.86 & 1.65 & 1.74 & 1.65 & --- & 缺约 $\sqrt{2.5\text{--}3}\times$ 方差 \\
    \hline
    \multicolumn{7}{l}{\emph{偏置型 (不进入方差求和)}} \\
    A. dE/dx 平均损失 ($\langle\Delta E\rangle$)  & --- & --- & $-10.8$ & $-10.8$ & bias & 修复后归零 \\
    B. frame 旋转 bug（已于 2026-04-19 修复）     & 0   & 0   & 0   & 0   & bias & 已修复 \\
    \hline
    \multicolumn{7}{l}{\emph{尾部 (不进入方差，单独报失败率)}} \\
    C. LM 收敛盆地失败                            & --- & --- & --- & --- & tail & \makecell[l]{$(-100, 0)$ 子集 12\%；\\ 全局 $\sim 1$--$2\%$} \\
    \hline
  \end{tabular}
\end{table}

\subsection*{10.2 平方和 vs.~实测：缺多少？}

\paragraph{$p_x$。} 平方和：
$\sqrt{0.7^2+2.5^2+1.8^2+0.05^2+0.05^2+0.2^2+0.5^2+0.05^2+0.3^2}
= \sqrt{0.49+6.25+3.24+0.003+0.003+0.04+0.25+0.003+0.09}
= \sqrt{10.36} = 3.22\,\mathrm{MeV}/c$。
实测 $5.98\,\mathrm{MeV}/c$。\textbf{缺差因子 $5.98/3.22 = 1.86$}（方差缺差 $3.46$）。

\paragraph{$p_y$。} 平方和：
$\sqrt{0.5^2+0.5^2+0.4^2+0.05^2+0.05^2+0.05^2+0.1^2+0.5^2+0.05^2}
= \sqrt{0.25+0.25+0.16+0.0075+0.01}
\approx \sqrt{0.69} = 0.83\,\mathrm{MeV}/c$。
实测 $1.40\,\mathrm{MeV}/c$。\textbf{缺差因子 $1.40/0.83 = 1.69$}。

\paragraph{$p_z$。} 平方和：
$\sqrt{0.7^2+2.5^2+1.8^2+0.1^2+0.05^2+0.2^2+0.1^2+0.05^2+0.3^2}
= \sqrt{0.49+6.25+3.24+0.01+0.003+0.04+0.01+0.003+0.09}
= \sqrt{10.13} = 3.18\,\mathrm{MeV}/c$。
实测 $5.57\,\mathrm{MeV}/c$。\textbf{缺差因子 $5.57/3.18 = 1.75$}。

\paragraph{$|\vec p|$。} 平方和：
$\sqrt{0.7^2+2.5^2+1.8^2+0.1^2+0.05^2+0.2^2+0.5^2+0.05^2+0.3^2}
\approx \sqrt{10.36} = 3.22\,\mathrm{MeV}/c$。
实测 $5.30\,\mathrm{MeV}/c$。\textbf{缺差因子 $5.30/3.22 = 1.65$}。

\subsection*{10.3 缺失的 $\sim 1.65$--$1.86\times$ 是什么？}

四个分量都给出 $\sim 1.7\times$ 的缺差（方差缺约 $1.9\text{--}2.5\times$），结构高度一致。
我们把它归到下列四个机制并按权重排序：

\begin{table}[ht]
  \centering\small
  \caption{缺失方差贡献的权重拆分（4 个分量的平均估计）}
  \label{tab:partIIc:budget:missing-decomp}
  \begin{tabular}{l|c|l}
    \hline
    机制 & 占缺失方差比例 & 物理说明 \\
    \hline
    (i) MS 过程噪声未注入 Fisher / 前向模型     & $\sim 50\%$ & part2a §2/§3、§8.1 (b) \\
    (ii) 未 survey 材料（PDC 入射窗、支撑结构、Sn 法兰） & $\sim 20\%$ & 几何不全 \\
    (iii) LM 盆地失败尾部（$(-100, 0)$）        & $\sim 10\%$ & §9 (本节) \\
    (iv) 高阶参数化效应（$(u,v,p)$ 二阶项）      & $\sim 5\%$ & §8 (本节) \\
    (v) 残余未知（$B$ 场局部 $\delta\theta$ 等） & $\sim 15\%$ & --- \\
    \hline
  \end{tabular}
\end{table}

(i) 是\textbf{最大单项}：当前 Fisher 用 $\Sigma_\mathrm{meas} = \sigma_\mathrm{wire}^2\mathbb I_4$，
没把 MS 沿轨迹的方向涂抹算成过程噪声 $\Sigma_\mathrm{MS}$ 加进去。
解析估计 $\Sigma_\mathrm{MS}$ 在 PDC2 命中位置上的有效宽度 $\sim 1.3$ mm（part2a §2 + §3
合并 Highland），平方加到 $\sigma_\mathrm{wire}=0.5$ mm 上得 $\sigma_\mathrm{eff}\sim 1.4$ mm。
注入后 Fisher $\sigma_p$ 应当从当前 $\sim 6\,\mathrm{MeV}/c$ 拉到 $\sim 6\sqrt{(1.4/0.5)^2}\sim 17\,\mathrm{MeV}/c$
——这是\emph{过修}！实际线性传播没那么糟，因为 MS 在轨迹上分布而非集中在 PDC 位置；
正确的 Kalman 风格的过程噪声只把 $p_y$ 抬到 $\sim 1.4$，把 $p_x/p_z$ 抬到 $\sim 4.5$。
即修复 (i) 大约能把 $p_x$ 从 $3.22 \to 4.5$，$p_y$ 从 $0.83 \to 1.4$，
$p_z$ 从 $3.18 \to 4.5$，$|\vec p|$ 从 $3.22 \to 4.5\,\mathrm{MeV}/c$。

\subsection*{10.4 优先级排序的修复列表}

按"对实测缺差的削减幅度"排序：

\begin{enumerate}
  \item \textbf{把 MS 过程噪声注入 Fisher 协方差}（修复 (i)）。预测：
        $p_x$ 平方和从 $3.22\to 4.5\,\mathrm{MeV}/c$，与实测 $5.98$ 的差距从 $1.86\times$
        缩到 $1.33\times$。$p_y$ 同理 $0.83\to 1.4\,\mathrm{MeV}/c$，与实测自洽。
        代码锚点：\texttt{libs/analysis\_pdc\_reco/src/PDCRkAnalysisInternal.cc} 的
        $\Sigma_\mathrm{meas}$ 装配处。% TODO: process-noise injection
  \item \textbf{把 Bethe--Bloch dE/dx 加入 RK4 步进}（修复偏置 A）。预测：
        $p_z$ 偏置从 $-10.8$ 一次性归零，宽度不变。代码锚点：
        \texttt{libs/analysis/src/ParticleTrajectory.cc::RungeKuttaStep()}（\StatusReport{}~§已知局限 \#1）。% TODO: dE/dx in RK step
  \item \textbf{Cartesian $(p_x,p_y,p_z)$ 参数化}（修复 (iv)）。预测：
        $\sigma_{p_y}^\mathrm{Fisher}$ 从 $0.75 \to 0.85\,\mathrm{MeV}/c$，$\sim 13\%$ 改善
        ——比修复 (i) 小很多，但对覆盖率 audit 有诊断价值。% TODO: Cartesian param
  \item \textbf{扩展 LM 粗格 + 反向 RK 初始}（修复尾部 (iii)）。预测：
        $(-100, 0)$ 失败率 $12\% \to <1\%$，全局尾部从 $\sim 1\text{--}2\% \to <0.3\%$。% TODO: extend grid
  \item \textbf{材料 survey 完成}（修复 (ii)）：把 PDC 入射窗、Sn 法兰、支撑结构
        厚度纳入几何宏。预测：$\sim 20\%$ 缺失方差被吃掉，对应 $p_x/p_z$ 平方和
        增量 $\sim 0.7\,\mathrm{MeV}/c$。% TODO: material survey
  \item \textbf{$B$ 场三次卷积插值 + cosmic alignment}（修复 (v)）：把 §7 与 part2b §9
        的占位项推到 $<0.05\,\mathrm{MeV}/c$。最小贡献。
\end{enumerate}

\subsection*{10.5 修复后预测：master 表"after-fix"列}

完成修复 1--6 后，重做 master 表的"修复后预测"列（保守估计，按"减半"取
当前预算行的较小值；$\sigma_{p_y}$ 列额外装入 $\Sigma_\mathrm{MS}$）：

\begin{table}[ht]
  \centering\footnotesize
  \caption{修复后预测的 master 表（after-fix 列）。仅展示与当前列差异显著的行。}
  \label{tab:partIIc:budget:after-fix}
  \begin{tabular}{l|cccc|l}
    \hline
    来源 & $\sigma_{p_x}$ & $\sigma_{p_y}$ & $\sigma_{p_z}$ & $\sigma_{|\vec p|}$ & 修复内容 \\
    \hline
    1. PDC 位置分辨                & 0.7 & 0.5 & 0.7 & 0.7 & 不变 \\
    2. MS：1 m 空气 $\to$ 真空     & 0.4 & 0.1 & 0.4 & 0.4 & EW $\to$ PDC1 真空管 \\
    3. MS：PDC 气体                & 1.8 & 0.4 & 1.8 & 1.8 & 不变（除非换 He 填充） \\
    4. dE/dx 涨落                  & 0.05 & 0.05 & 0.1 & 0.1 & 偏置已修；涨落不变 \\
    5. RK 离散                     & 0.02 & 0.02 & 0.02 & 0.02 & 自适应 RK45 \\
    6. 靶 $z$                      & 0.1 & 0.02 & 0.1 & 0.1 & metadata 升级 \\
    7. 对位 + 倾角                 & 0.1 & 0.05 & 0.05 & 0.1 & cosmic alignment + survey \\
    8. 参数化                      & 0.02 & 0.1 & 0.02 & 0.02 & Cartesian \\
    9. $B$ 场                      & 0.1 & 0.02 & 0.1 & 0.1 & tricubic \\
    + MS 过程噪声注入 (新行)        & 3.5 & 1.0 & 3.5 & 3.5 & Kalman 风格 $\Sigma_\mathrm{MS}$ \\
    \hline
    \textbf{平方和合计} & \textbf{3.99} & \textbf{1.21} & \textbf{4.0} & \textbf{4.0} & --- \\
    \hline
  \end{tabular}
\end{table}

修复 1--6 全部到位后，平方和与实测的差距从当前 $1.65$--$1.86\times$ 缩到
$\sim 1\text{--}1.4\times$ 范围内（即各分量的 budget 与实测残差宽度落在 $\pm 30\%$ 以内）。
更重要的是，\emph{实测残差宽度本身}也会下降——例如把 1 m 空气换成真空后，
744 事件 ensemble 的 $\widehat\sigma_{p_x}^\mathrm{obs}$ 会从 $5.98$ 下降到
预期 $\sim 3\,\mathrm{MeV}/c$（与 ms\_ablation stage A vacuum 推断一致）。

\subsection*{10.6 结论}

Master 表当前给出的平方和合计为 $(p_x, p_y, p_z, |\vec p|) = (3.21, 0.85, 3.21, 3.21)\,\mathrm{MeV}/c$，
实测半宽为 $(5.98, 1.40, 5.57, 5.30)$，比值结构性 $1.65$--$1.86\times$ 一致——\textbf{master 表
没封闭，缺一行 MS 过程噪声 + 几条 minor 行}。修复路线图按 §10.4 给出，
首项"过程噪声注入 Fisher" 一次性削减 $\sim 50\%$ 的缺失方差，
首项"dE/dx 加入 RK 步" 一次性归零 $-10.8\,\mathrm{MeV}/c$ 的 $p_z$ 偏置。
完成后 master 表与实测残差应当自洽至 $\pm 30\%$ 内，留作下一份状态报告的对账目标。


\newpage

% ============================================================
% 第三部分: 实证诊断
% ============================================================
\part{实证诊断}\label{part:diag}

本部分把第一、二部分的理论与预算放到 744 事件 ensemble 上做实证检验。
六章覆盖：覆盖率方法学（Clopper-Pearson / Wilson / Wald 三选一与 ensemble
规模扫描）、pull 分布的定义与解读（含 $p_x/p_y/p_z/|\vec p|$ 实测中位与
$\sigma_z$）、$\chi^2$ 分布与失败模式分类、$5\text{--}10$ 个单事件深度示例
（含 $(-100, 0)$ 病态点 12\% 失败率剖析）、扫描研究方案与外推（B-field
$\pm 5\%$、$\sigma_\mathrm{wire}$ 扫描、靶倾角扫描、ensemble 规模扫描）、
LM 收敛盆地病态剖析。

% =============================================================
% rk_deep_dive_part3_diagnostics_zh.tex
% -------------------------------------------------------------
% 第三部分: 实证诊断 (Part III — Empirical diagnostics).
% 编入到 rk_error_analysis_deep_dive_20260426_zh.tex (主壳层) 中,
% 通过 % =============================================================
% rk_deep_dive_part3_diagnostics_zh.tex
% -------------------------------------------------------------
% 第三部分: 实证诊断 (Part III — Empirical diagnostics).
% 编入到 rk_error_analysis_deep_dive_20260426_zh.tex (主壳层) 中,
% 通过 % =============================================================
% rk_deep_dive_part3_diagnostics_zh.tex
% -------------------------------------------------------------
% 第三部分: 实证诊断 (Part III — Empirical diagnostics).
% 编入到 rk_error_analysis_deep_dive_20260426_zh.tex (主壳层) 中,
% 通过 \input{rk_deep_dive_part3_diagnostics_zh} 加载。
%
% 不要在此文件中包含 \documentclass / \begin{document} / 序言;
% 主壳层提供 \part{第三部分: 实证诊断} 标题, ctex+xelatex 字体,
% lstlisting 中文样式, booktabs/multirow/amsmath 等宏包。
%
% 创建日期: 2026-04-26
% 作者:    Baiting Tian
% 关联文件:
%   * docs/reports/reconstruction/momentum_reco/rk/rk_reconstruction_status_20260416_zh.tex
%   * docs/superpowers/specs/2026-04-26-rk-error-analysis-deep-dive-design.md
%   * docs/reports/reconstruction/momentum_reco/rk/rk_reconstruction_bugfix_and_error_analysis_20260416.md
%
% 本文件结构 (六章):
%   1. 覆盖率方法学: Clopper-Pearson / Wilson / Wald
%   2. Pull 分布: 定义, 期望, 解读, 实测
%   3. chi^2 分布与失败模式分类
%   4. 5-10 个单事件深度分析
%   5. 扫描研究方案与外推 (多数为 % TODO)
%   6. LM 收敛盆地: (-100, 0) 病态点剖析
%
% 数据来源:
%   * track_summary.csv  (744 事件, 修复后靶系)
%   * profile_samples.csv (128 事件)
%   * bayesian_samples.csv (128 事件)
%   * g4_fluctuation/per_truth_point_g4_vs_fisher.csv
%   * g4_fluctuation/ensemble_pdc_reco_merged.csv (744 事件)
% =============================================================

\section{覆盖率方法学: Clopper-Pearson, Wilson, Wald 三选一}
\label{sec:partIII:coverage_methodology}

整个状态报告 (\StatusReport{}) 把覆盖率作为 RK 误差分析的最终判据: 把名义 68\%/95\% 的方法区间打到 744 条事件上, 数 truth 落在区间内的比例; 这个比例本身又是一个二项随机变量, 必须给一个置信区间才能下结论. 报告里默认采用 Clopper-Pearson 95\% 置信区间(以下简称 CP CI), 本章给出原因, 并把它和教科书里另两种竞争方案 (Wilson 与 Wald) 做对比.

\subsection{二项覆盖率的统计模型}
\label{sec:partIII:coverage_methodology:model}

设我们独立观察 $N$ 个事件, 每个事件的方法区间 $[\hat{\theta}_i - \delta_i^-, \hat{\theta}_i + \delta_i^+]$ 是否覆盖 truth $\theta_i^{\mathrm{truth}}$ 由指示变量
\[
  X_i \;=\; \mathbb{1}\bigl[\hat\theta_i - \delta_i^- \le \theta_i^{\mathrm{truth}} \le \hat\theta_i + \delta_i^+\bigr]
\]
描述. 在 "方法自洽" 的零假设下, $\Pr(X_i = 1) = p_0$ 与名义覆盖水平 (例如 $p_0 = 0.68$) 一致, 且对 $i$ 而言 $X_i$ 独立同分布. 总和 $K = \sum_i X_i$ 服从二项分布
\[
  K \sim \mathrm{Bin}(N, p_0).
\]
观测覆盖率是点估计 $\hat p = K / N$. CI 的目标是: 给定观测 $K = k$, 构造一个区间 $[p_L, p_U]$ 使得在所有真实的 $p$ 上, 该区间覆盖真实 $p$ 的概率不低于 $1 - \alpha$ (例如 $1 - \alpha = 0.95$).

\subsection{Wald 区间(正态近似)与失败模式}
\label{sec:partIII:coverage_methodology:wald}

最常见的写法是 Wald 正态近似:
\[
  p_{L/U}^{\mathrm{Wald}} \;=\; \hat p \;\pm\; z_{\alpha/2}\, \sqrt{\frac{\hat p (1 - \hat p)}{N}}.
\]
教科书把它印在第一页是因为它形式简单, 但它在两端附近(尤其当 $\hat p$ 靠近 $0$ 或 $1$, 或者 $N$ 不够大)有以下三个失败模式:

\begin{enumerate}[label=(\alph*), nosep]
  \item 当 $\hat p = 0$ 或 $1$ 时, 标准误差 $\sqrt{\hat p(1-\hat p)/N} = 0$, 区间退化为单点 $\{0\}$ 或 $\{1\}$. 这显然不能用作 CI ——任何有限 $p$ 都可能产生 $K = 0$ 或 $K = N$.
  \item 当 $N$ 很小 (本章场景: 状态报告中 MCMC 子集 $N = 128$, 早期试运行 $N = 32$), 区间端点可能跑出 $[0, 1]$ 的物理范围.
  \item 实际覆盖率(在所有 $p$ 上对 Wald CI 求频率)远低于名义 $1 - \alpha$. Brown, Cai, DasGupta (2001, \emph{Statistical Science}) 给出了著名的 "wiggle plot": Wald 在 $p$ 略偏离 0.5 时的覆盖率会陷入 0.85 量级的低谷, 远低于名义 0.95.
\end{enumerate}

第三点在状态报告 §误差分析 里直接致命: 修复前 $p_x$ 覆盖率 $\hat p \approx 0.05$ (744 中 38 条命中); Wald 给的区间是
\[
  [0.05 - 1.96\sqrt{0.05\cdot 0.95/744},\ 0.05 + \ldots] \approx [0.034, 0.066],
\]
看起来很窄; 但当 $\hat p$ 离 0 这么近, Wald 显著地低估了不对称性 (真实分布右偏更厉害). 我们要的是一种在 $\hat p \to 0$ 或 $\hat p \to 1$ 极端区域仍然行为正确的 CI.

\subsection{Clopper-Pearson 精确二项 CI}
\label{sec:partIII:coverage_methodology:cp}

Clopper 与 Pearson (1934, \emph{Biometrika}) 提出: 让 $[p_L, p_U]$ 同时满足
\begin{align}
  \Pr\bigl(K \ge k \mid p = p_L\bigr) &\;=\; \alpha/2, \\
  \Pr\bigl(K \le k \mid p = p_U\bigr) &\;=\; \alpha/2.
\end{align}
即 $p_L$ 是 "真实 $p$ 至少为多少时, $K \ge k$ 这件事的概率才能压到 $\alpha/2$" 的最大值; $p_U$ 是 "真实 $p$ 至多为多少时, $K \le k$ 的概率才能压到 $\alpha/2$" 的最小值. 这是一种 "反演 (invert) 二项 CDF" 的 CI 构造, 直觉上等价于把假设检验的接受域翻成参数的置信域.

\paragraph{封闭形式与 Beta-incomplete.}
利用二项 CDF 与不完全 Beta 函数的对偶
\[
  \sum_{j=0}^{k}\binom{N}{j} p^j (1-p)^{N-j} \;=\; I_{1-p}(N-k,\, k+1),
\]
可以把 CP 区间端点写成 Beta 分布分位点:
\begin{align}
  p_L &\;=\; \mathrm{Beta}^{-1}\!\bigl(\alpha/2;\ k,\ N - k + 1\bigr), \\
  p_U &\;=\; \mathrm{Beta}^{-1}\!\bigl(1 - \alpha/2;\ k + 1,\ N - k\bigr).
\end{align}
端点情形: $k = 0 \Rightarrow p_L = 0$, $k = N \Rightarrow p_U = 1$, 这正是物理上希望的 "无下/上界" 行为.

\paragraph{CP 是保守的.}
CP 的 \emph{实际} 覆盖率始终不低于 $1 - \alpha$ (这就是 "精确" 的含义), 但不一定恰等于 $1 - \alpha$. 由于二项分布的离散性, 在大多数 $p$ 值上, 实际覆盖率会略高于 $1 - \alpha$. 这意味着 CP 的区间宽度比名义宽度需要的"刚好" 95\% 略宽一些; 这个保守性恰恰是状态报告需要的: 对每条覆盖率行, "如果 CP CI 不包含 0.68, 那不是统计噪声", 这是一个强论断.

\subsection{Wilson 区间}
\label{sec:partIII:coverage_methodology:wilson}

Wilson (1927) 提出的折中方案是把 Wald 的 $\hat p$ 替换为后验 Beta 调整:
\[
  p_{L/U}^{\mathrm{Wilson}} \;=\;
  \frac{\hat p + \tfrac{z^2}{2N}}{1 + \tfrac{z^2}{N}}
  \;\pm\;
  \frac{z}{1 + \tfrac{z^2}{N}}\,
  \sqrt{\frac{\hat p (1 - \hat p)}{N} + \frac{z^2}{4 N^2}}.
\]
其中 $z = z_{\alpha/2}$ (例如 $1.96$ 对应 95\% CI). 与 Wald 相比, Wilson:
\begin{itemize}[nosep]
  \item 在 $\hat p \in [0.1, 0.9]$ 区间精度极好, 实际覆盖率非常接近名义值;
  \item 区间端点保留在 $[0, 1]$;
  \item 在 $\hat p \to 0/1$ 退化得比 Wald 慢 (但仍比 CP 快).
\end{itemize}

\subsection{$N=128$, $K=80$ 的数值算例}
\label{sec:partIII:coverage_methodology:example}

把 $\hat p \approx 0.625$ 的情形 (例如 MCMC $|\vec p|$ 95\% 覆盖率, $N=128$, $K=115$) 取近似数 $K = 80$ 出来作演算示例. 95\% 双侧:

\begin{table}[H]
\centering\small
\caption{$N=128$, $K=80$ 时三种 CI 的对比 ($\hat p = 0.625$).}
\label{tab:partIII:cp_wilson_wald_demo}
\begin{tabular}{lcccc}
\toprule
方法 & 公式入口 & $p_L$ & $p_U$ & 半宽 \\
\midrule
Wald     & $\hat p \pm 1.96\sqrt{\hat p (1-\hat p)/N}$ & 0.5412 & 0.7088 & 0.0838 \\
Wilson   & 公式 \eqref{eq:partIII:wilson} & 0.5380 & 0.7053 & 0.0837 \\
Clopper-Pearson & Beta 分位点 & 0.5365 & 0.7068 & 0.0852 \\
\bottomrule
\end{tabular}
\end{table}

\begin{equation}
\label{eq:partIII:wilson}
  p_{L/U}^{\mathrm{Wilson}} = \frac{\hat p + z^2/(2N)}{1 + z^2/N}
  \pm \frac{z\sqrt{\hat p(1-\hat p)/N + z^2/(4N^2)}}{1 + z^2/N}.
\end{equation}

% TODO: 用 scripts/analysis/analyze_ensemble_coverage.py 内部 _clopper_pearson 函数
% (它即用 scipy.stats.beta.ppf) 与 statsmodels.stats.proportion 的 wilson_score
% 对全部 4*4 覆盖率行批量出 CSV; 当前数值是手算/样例,
% 需要的运行: ENSEMBLE_SIZE=50 ./scripts/analysis/run_ensemble_coverage.sh
% 已经跑过, 这里仅缺一个三方法对比表生成器.

在 $\hat p \approx 0.625$ 这种 "中部" 区间, 三种方法差距 $< 1\%$; 真正的差距出现在尾部. 对极端 $\hat p \approx 0.05$, $N=744$, $K=38$:

\begin{table}[H]
\centering\small
\caption{$N=744$, $K=38$ 时三种 CI 的对比 ($\hat p = 0.0511$, 修复前 $p_x$ 覆盖率).}
\label{tab:partIII:cp_wilson_wald_extreme}
\begin{tabular}{lccc}
\toprule
方法 & $p_L$ & $p_U$ & 备注 \\
\midrule
Wald     & 0.0353 & 0.0669 & 端点近似对称, 偏窄 \\
Wilson   & 0.0376 & 0.0694 & 内推到 $[0, 1]$ \\
Clopper-Pearson & 0.0365 & 0.0698 & 略宽, 给出真正下界 \\
\bottomrule
\end{tabular}
\end{table}

CP 与 Wilson 的差距不大 (各 $\sim 0.001$), 但都明显比 Wald 长. 既然计算成本相同 ($\sim$ 一次 Beta 分位点计算), 选 CP 是合理的工程选择: 永远不会因为 $\hat p$ 走到极端而退化, 永远是保守的.

\subsection{CP 区间反演的几何直觉}
\label{sec:partIII:coverage_methodology:geometry}

CP 的反演构造可以画一张二维图来理解. 在 $(p, K)$ 平面上(横轴是参数 $p$, 纵轴是观测计数 $K$), 对每个 $p$ 画 $K$ 的双侧 $\alpha/2$ 接受区间 $[K_L(p), K_U(p)]$:
\begin{align}
  K_L(p) &= \min\{k : \Pr(K \ge k \mid p) \le 1 - \alpha/2\}, \\
  K_U(p) &= \max\{k : \Pr(K \le k \mid p) \le 1 - \alpha/2\}.
\end{align}
这两条上下边界形成一个"接受带"; 给定观测 $k$, 把横线 $K = k$ 与该带相交, 横向投影即得到 CP CI $[p_L, p_U]$. 离散的二项分布让边界呈"阶梯状", 这是 CP 略保守的根源 — 阶梯在某些 $p$ 上比连续构造高出 $\sim 1$ 单位.

如下伪代码反映了这一构造的数值实现 (与 \texttt{analyze\_ensemble\_coverage.py} 的 \texttt{\_clopper\_pearson} 函数一致):

\begin{lstlisting}[language=Python,caption={CP CI 的最小数值实现},label=lst:partIII:cp_python]
from scipy import stats

def clopper_pearson(k, n, alpha=0.05):
    """Clopper-Pearson exact binomial CI, two-sided 1-alpha coverage."""
    if k == 0:
        p_lo = 0.0
    else:
        p_lo = stats.beta.ppf(alpha / 2, k, n - k + 1)
    if k == n:
        p_hi = 1.0
    else:
        p_hi = stats.beta.ppf(1 - alpha / 2, k + 1, n - k)
    return p_lo, p_hi
\end{lstlisting}

\paragraph{阶梯效应数值化.}
对 $N=128$, 取 $K=87$ 与 $K=88$ (相邻整数), CP CI 的下端跳变:
\begin{itemize}[nosep]
  \item $K=87$: $p_L = 0.602$;
  \item $K=88$: $p_L = 0.610$.
\end{itemize}
即 $K$ 增 1 时 $p_L$ 跳 $\sim 0.008$. 状态报告里观测覆盖率精度因此有 \emph{量子化下限} $\sim 1/N$, 这是离散二项分布的根本约束, 不是计算误差.

\subsection{多覆盖率水平: 95\% vs.\ 99\%}
\label{sec:partIII:coverage_methodology:multilevel}

状态报告默认 95\% CP CI; 偶尔需要 99\% 来给出"几乎确定" 的结论. 把上面 $\hat p = 0.625$, $N=128$ 的算例扩展:

\begin{table}[H]
\centering\small
\caption{$N=128$, $K=80$ 时不同 $\alpha$ 的 CP CI.}
\label{tab:partIII:cp_multilevel}
\begin{tabular}{lcccc}
\toprule
$1 - \alpha$ & $z_{\alpha/2}$ (Wald 参考) & $p_L$ (CP) & $p_U$ (CP) & 半宽 \\
\midrule
68\% & 1.00 & 0.582 & 0.667 & 0.043 \\
90\% & 1.64 & 0.553 & 0.692 & 0.069 \\
95\% & 1.96 & 0.536 & 0.707 & 0.085 \\
99\% & 2.58 & 0.508 & 0.730 & 0.111 \\
\bottomrule
\end{tabular}
\end{table}

\paragraph{推论.} 99\% 比 95\% 半宽宽 $\sim 1.3\times$. 状态报告的覆盖率结论 ($p_y$ 欠覆盖 $0.42$ vs.\ $0.68$) 用 99\% CP CI 仍然成立 (差距 $0.26 \gg 0.111$). 但 MCMC 的 $|\vec p|$ 0.633 vs.\ 0.68 的差距 0.05 在 95\% CP CI 半宽 0.085 之内, 在 99\% 半宽 0.111 之内 — 这个差异 \emph{不显著}, 这是状态报告 §误差分析 中已经强调的事实.

\subsection{ensemble 数据的独立性假设}
\label{sec:partIII:coverage_methodology:independence}

CP CI 假设 $X_i$ i.i.d.. 现实情况:
\begin{itemize}[nosep]
  \item ensemble 在 15 个 truth bin 上分组 ($N=50$/bin, 总 $N=744$). 同一 truth bin 内的 50 条事件的 \emph{Geant4 输入} 完全相同, 只是随机种子不同 — 这种意义上独立.
  \item 但是, 同一 truth bin 内 LM 失败事件占 12\% (在 $-100, 0$ bin) 或 $0$\% (在其他 bin) — \emph{失败的发生与 truth bin 强相关}, 不是 i.i.d.. 这意味着把全 744 看作单一二项总样本会低估 CP CI.
  \item 更严格的处理: 对每个 truth bin 单独算 CP CI, 然后做 stratified weighted average. 当前 \texttt{analyze\_ensemble\_coverage.py} 是 pooled 处理.
\end{itemize}

\paragraph{stratified analysis 占位.}
% TODO: analyze_ensemble_coverage.py 加 --stratify-by-truth 选项,
% 输出 build/test_output/.../coverage_per_truth_bin.csv,
% 每行 (truth_px, truth_py, N, k, p_hat, cp_lo, cp_hi).

\subsection{为什么不用 bootstrap?}
\label{sec:partIII:coverage_methodology:no_bootstrap}

Bootstrap 在覆盖率比例上是完全可行的: 把 $N$ 个 $X_i$ 重抽样 $B$ 次, 每次算 $\hat p^*_b$, 取经验分位点. 但相对 CP 它有三个劣势:

\begin{itemize}[nosep]
  \item \textbf{计算成本}: $B = 10^4$ 次重抽样 vs.\ 单次 Beta 分位点; 在 60 个覆盖率行 (4 分量 $\times$ 4 方法 $\times$ 多覆盖水平 $\times$ 多 truth bin) 累积起来非微秒.
  \item \textbf{近似性}: bootstrap 的覆盖率本身仍然是渐近 $1 - \alpha$, 在小 $N$ 不"精确".
  \item \textbf{无法重复}: 每次 bootstrap 有随机种子, 报告读者复现时数字会跳; CP 是闭式的.
\end{itemize}

闭式 CP 足够保守, 故选定为状态报告标准.

\subsection{ensemble 大小 $\to$ CP 半宽的换算表}
\label{sec:partIII:coverage_methodology:N_scaling}

当 $\hat p = 0.68$ 时, 二项标准差为 $\sqrt{p(1-p)/N} = \sqrt{0.2176/N}$, Wald 半宽是 $1.96$ 倍此值; CP 半宽与 Wald 的中部值相差 $< 5\%$. 估算各 $N$ 下的 CP 95\% 半宽:

\begin{table}[H]
\centering\small
\caption{$\hat p = 0.68$ 时 $N$ 与 CP 95\% 半宽的近似关系. 列 "Wald 近似" 给 $1.96\sqrt{p(1-p)/N}$, 列 "CP 实际" 由 \texttt{scipy.stats.beta.ppf} 直接计算.}
\label{tab:partIII:N_cp_halfw}
\begin{tabular}{rrr}
\toprule
$N$ & Wald 近似 & CP 实际 \\
\midrule
32   & 0.162 & 0.180 \\
64   & 0.114 & 0.121 \\
128  & 0.081 & 0.084 \\
256  & 0.057 & 0.058 \\
744  & 0.034 & 0.034 \\
5000 & 0.013 & 0.013 \\
\bottomrule
\end{tabular}
\end{table}

\textbf{解读}: 状态报告里 Fisher/Laplace 用 $N=744$, CP 半宽 $\sim 0.034$, 即名义 0.68 与观测 $\sim 0.42$ ($p_y$) 的差距 $0.26$ 远超过 $\pm 0.034$, 是 "真信号"; Profile/MCMC 用 $N=128$, CP 半宽 $\sim 0.084$, 仍能区分 $0.42$ vs $0.68$. 但是, 一旦缩到 $N=32$, 半宽 $\sim 0.18$ 跨过差距, 任何小于 $\sim 18\%$ 的偏离都会埋没在统计噪声里.

\textbf{推论}: 生产用 $N=200$ (CP $\pm 0.07$) 已经够诊断 $\sim 10\%$ 级偏差; 当前 $N=50$ 每 truth 点(总 $N=744$) 只在每个 truth bin 内 CP $\pm 0.13$, 每点诊断略嫌粗. 下文 \nameref{sec:partIII:scans} 里 ES-1 计划把生产 ensemble 增到 $N=200$/truth 点 ($\sim 3000$ 总).

\subsection{小结}
\label{sec:partIII:coverage_methodology:summary}

\begin{itemize}[nosep]
  \item Wald 在 $\hat p \to 0/1$ 退化, 在 $\hat p \in (0.1, 0.9)$ 也轻微低覆盖, 状态报告不用.
  \item Wilson 大多数情形与 CP 接近, 但缺乏 "永远 $\ge 1 - \alpha$" 的精确性保证.
  \item Clopper-Pearson 是默认选项: 闭式, 保守, 在尾部行为正确.
  \item 当前 $N=744$ 给 CP $\pm 0.034$, 足以分辨当前的 $p_y$ 欠覆盖 $0.42$ vs.\ 名义 $0.68$; $N=128$ MCMC/Profile 子集给 CP $\pm 0.084$, 仍可诊断 $\sim 0.1$ 级偏差.
\end{itemize}

% =============================================================

\section{Pull 分布: 定义, 期望, 解读}
\label{sec:partIII:pulls}

覆盖率检验告诉你 "区间宽度对不对", 但它把整个分布折成一个标量. 真正展示 "误差棒形状是否正确" 的诊断量是 \emph{pull}. 本章给出 pull 的形式定义, 在三种缺陷模式下的标志特征, 并直接读取 \texttt{track\_summary.csv} 算出 744 事件的 pull 统计.

\subsection{Pull 定义}
\label{sec:partIII:pulls:definition}

设第 $i$ 条事件的某个动量分量(以 $p_x$ 为例)的真值为 $p_{x,i}^{\mathrm{truth}}$, 重建给出的中心估计为 $\hat p_{x,i}$, 误差分析(本节默认 Fisher) 给出的标准差为 $\sigma_{p_x, i}$. 定义事件级 pull:
\[
  z_{p_x, i} \;=\; \frac{\hat p_{x,i} - p_{x,i}^{\mathrm{truth}}}{\sigma_{p_x, i}}.
\]
注意: pull 并不等于 \emph{标准化残差} 一般意义上的 $\chi^2$ 项 —— 它取的是 \emph{重建估计与真值之差} 除以\emph{该事件的}估计 $\sigma$, 不是平均 $\sigma$. 这意味着每条事件都用自己的方差归一, 是一个真正的 "诊断这条事件" 的量.

\subsection{标准期望}
\label{sec:partIII:pulls:expected}

如果(且仅如果)以下三个条件都成立,
\begin{enumerate}[label=(\alph*), nosep]
  \item 重建估计是无偏的: $\mathbb{E}[\hat p_x | p_x^{\mathrm{truth}}] = p_x^{\mathrm{truth}}$;
  \item 误差分析给出的 $\sigma$ 与真实方差一致: $\mathrm{Var}[\hat p_x | p_x^{\mathrm{truth}}] = \sigma_{p_x}^2$;
  \item 残差是高斯分布的: $\hat p_x | p_x^{\mathrm{truth}} \sim \mathcal{N}(p_x^{\mathrm{truth}}, \sigma_{p_x}^2)$,
\end{enumerate}
那么 pull 服从标准正态:
\[
  z \sim \mathcal{N}(0, 1).
\]
推论:
\begin{itemize}[nosep]
  \item $\mathbb{E}[z] = 0$,
  \item $\mathrm{Var}[z] = 1$ 即 $\sigma_z = 1$,
  \item $\Pr(|z| \le 1) = 0.6827$, $\Pr(|z| \le 1.96) = 0.95$.
\end{itemize}
任何对 $\mathcal{N}(0,1)$ 的偏离都对应一种系统问题. 表 \ref{tab:partIII:pull_diagnostics_template} 给出三类典型缺陷的特征:

\begin{table}[H]
\centering\small
\caption{Pull 分布对应的故障模式.}
\label{tab:partIII:pull_diagnostics_template}
\begin{tabular}{llp{0.45\linewidth}}
\toprule
特征 & 说明 & 故障模式 \\
\midrule
$\bar z \ne 0$ & pull 中心不在零 & 重建有偏: 坐标系/dE/dx/对齐错误的几何位移 \\
$\sigma_z > 1$ & pull 太宽 & $\sigma$ 低估: Fisher 线性化丢了曲率, 或者 (u,v,p) 参数化压抑了某分量 \\
$\sigma_z < 1$ & pull 太窄 & $\sigma$ 高估: Fisher 把不存在的相关性 / 多次散射 / dE/dx 加进去了 \\
\bottomrule
\end{tabular}
\end{table}

\textbf{示例}: 状态报告里修复前的 $p_x$ pull 中位 $\sim -33 / 7.4 \approx -4.4$ 是 (a) 类典型 (frame mismatch); 当前 $p_y$ ratio $1.86$ $\Rightarrow \sigma_z \approx 1.86$ 是 (b) 类典型; 当前 $p_x, p_z$ ratio $\approx 0.8$ $\Rightarrow \sigma_z \approx 0.8$ 是 (c) 类轻症.

\subsection{744 事件 pull 实测}
\label{sec:partIII:pulls:measured}

下表是从 \texttt{track\_summary.csv} 直接计算的 pull 统计 (定义见上, $\sigma$ 取 \texttt{fisher\_*\_sigma} 列):

\begin{table}[H]
\centering\small
\caption{Pull 统计(744 事件, 修复后靶系); $\bar z$ = 均值; $\sigma_z$ = 标准差; $|z|\le 1$ 与 $|z|\le 1.96$ = 实测命中率, 名义为 $0.68$ 与 $0.95$.}
\label{tab:partIII:pull_measured}
\begin{tabular}{lrrrrrl}
\toprule
分量 & $\bar z$ & 中位 $z$ & $\sigma_z$ & $|z|\le 1$ & $|z|\le 1.96$ & 备注 \\
\midrule
$p_x$ & $-0.572$ & $-0.026$ & $4.216$ & $80.1\%$ & $91.0\%$ & 长尾来自 6 条 LM 失败事件 \\
$p_y$ & $-0.178$ & $-0.066$ & $2.477$ & $42.1\%$ & $69.0\%$ & 整体宽 $\sigma_z \approx 2.5$ \\
$p_z$ & $-0.293$ & $-0.068$ & $9.050$ & $77.2\%$ & $90.2\%$ & 长尾更严重(尾部 $|z| > 100$) \\
$|\vec p|$ & $-2.823$ & $-0.066$ & $59.989$ & $76.1\%$ & $90.1\%$ & 同 $p_z$, 长尾事件主导 \\
\bottomrule
\end{tabular}
\end{table}

\paragraph{读法.}
\begin{itemize}[nosep]
  \item \textbf{中位数 $\sim 0$}: 中心估计修复后 (2026-04-19 frame fix) 已经无偏, 没有 frame mismatch 残余. 三分量中位 pull 都 $|z_{\mathrm{med}}| < 0.07$.
  \item \textbf{均值 $\bar z$ 偏负}: 但拖到长尾 (LM 失败事件) 把均值压成 $\sim -0.5$ ($p_x$); $|\vec p|$ 更夸张 $\bar z = -2.8$ 因为 $|\vec p|$ 是\emph{严格非负} 的, LM 跑飞时 $\hat{|\vec p|} \to 1100$ 但 truth $\sim 635$, 这种事件单独贡献 $\Delta z \sim +75$ 但少数; 实际 6 条事件贡献了 95\% 的方差.
  \item \textbf{$\sigma_z (p_y) = 2.48$}: 没有 LM 长尾的"干净"分量, 这个 2.48 跟覆盖率比值 1.86 不完全吻合, 原因是 pull 用的是 \emph{每事件 $\sigma$} 而不是中位 $\sigma$, 而且分布不是高斯; ratio 1.86 是 \emph{鲁棒半宽 / Fisher $\sigma$ 中位}, 数学上不等价, 但定性上同向 (Fisher 低估).
  \item \textbf{$\sigma_z (p_x), \sigma_z (p_z)$ 极大但中位 $|z|\le 1$ 比例 $\sim 0.78$}: 这是 "干净事件占 $93\%$, 6 条 LM 失败事件占 $7\%$" 的双峰分布表现; 鲁棒尺度估计(中位 + IQR) 对这个分布更稳, 也对应表 §误差分析 的覆盖率数字.
\end{itemize}

\paragraph{Pull 的鲁棒尺度估计.}
\label{sec:partIII:pulls:robust}
对 LM 长尾敏感的修正方案: 取 \emph{IQR / 1.349} (高斯下与 $\sigma$ 一致, 对长尾鲁棒) 或者 MAD (median absolute deviation) $\times 1.4826$.

% TODO: 把下表的 IQR/MAD 列也加上,
% 需要在 analyze_ensemble_coverage.py 加 --robust-pull-scale 选项.
\begin{table}[H]
\centering\small
\caption{Pull 的鲁棒尺度估计(待运行, 当前为占位): IQR/1.349 与 $\sigma_z$ 直接对比.}
\label{tab:partIII:pull_robust_scale}
\begin{tabular}{lrrr}
\toprule
分量 & $\sigma_z$ (经典) & IQR/1.349 & 比值 \\
\midrule
$p_x$    & 4.216 & TODO & TODO \\
$p_y$    & 2.477 & TODO & TODO \\
$p_z$    & 9.050 & TODO & TODO \\
$|\vec p|$ & 59.99 & TODO & TODO \\
\bottomrule
\end{tabular}
\end{table}

\subsection{从 CSV 计算 pull 的代码片段}
\label{sec:partIII:pulls:code}

为可复现, 给出从 \texttt{track\_summary.csv} 直接算 pull 的 Python 片段:

\begin{lstlisting}[language=Python,caption={pull 统计的最小复现脚本},label=lst:partIII:pull_python]
import pandas as pd, numpy as np

CSV = ('build/test_output/ensemble_n50_targetframe/'
       'rk_fixed_target_pdc_only_error_analysis/track_summary.csv')
df = pd.read_csv(CSV)

# 残差与 pull
df['truth_p'] = np.sqrt(df.truth_px**2 + df.truth_py**2 + df.truth_pz**2)
for c in ['px', 'py', 'pz', 'p']:
    df[f'pull_{c}'] = (df[f'reco_{c}'] - df[f'truth_{c}']) / df[f'fisher_{c}_sigma']

for c in ['px', 'py', 'pz', 'p']:
    z = df[f'pull_{c}']
    in68 = (z.abs() <= 1.0).mean() * 100
    in95 = (z.abs() <= 1.96).mean() * 100
    print(f'{c}: mean={z.mean():+.3f} median={z.median():+.3f} '
          f'sigma={z.std():.3f}  |z|<=1: {in68:.1f}%  |z|<=1.96: {in95:.1f}%')
\end{lstlisting}

\subsection{Pull 直方图(占位)}
\label{sec:partIII:pulls:histograms}

% TODO: 需要的 PNG, 由 scripts/analysis/plot_pull_histograms.py 生成:
%   figures/diagnostics/pull_distribution_px.png
%   figures/diagnostics/pull_distribution_py.png
%   figures/diagnostics/pull_distribution_pz.png
%   figures/diagnostics/pull_distribution_p.png
% 该脚本未提交; 拟接受 --input track_summary.csv 与 --truth-cut 0.5 (排除 LM 失败事件)
% 双轴显示: 左 全部 744 事件 (含 LM 失败长尾), 右 truth-cut 后干净分布,
% 叠加 N(0,1) PDF.

四张直方图(占位)将分别显示 $p_x, p_y, p_z, |\vec p|$ 的 pull 分布, 与 $\mathcal{N}(0, 1)$ 参考曲线叠加. 期待画面:
\begin{itemize}[nosep]
  \item $p_x$ 分布: 主峰几乎与 $\mathcal{N}(0, 1)$ 重合, 但右上 (LM 失败) 的 6 条事件远远拖出 $|z| > 30$ 的尾巴;
  \item $p_y$ 分布: 主峰宽于 $\mathcal{N}(0, 1)$, $\sigma_z \sim 2.5$;
  \item $p_z$ 分布: 类似 $p_x$ 但长尾更长 (LM 失败时 $p_z$ 跑飞到 $\sim 940$ MeV/$c$);
  \item $|\vec p|$ 分布: 同 $p_z$, 主峰 $\bar z \sim 0$ 但长尾左偏(失败事件 $\hat{|p|} \gg p^{\mathrm{truth}}$).
\end{itemize}

\subsection{每事件 $\sigma$ 与 ensemble 平均 $\sigma$ 的差异}
\label{sec:partIII:pulls:per_event_vs_ensemble}

Pull 定义里 "$\sigma_i$" 是 \emph{每事件的} Fisher $\sigma$, 不是 ensemble 平均. 这有微妙的影响:

\begin{itemize}[nosep]
  \item 如果 Fisher $\sigma_i$ 在事件之间剧烈变化 (例如 $|\vec p|$ Fisher 在 $|p_x|=100$ 时 $\sim 7$, 在病态点 $\sim 60$), 而残差幅度也跟着变大, pull 仍可能保持 $\mathcal{N}(0,1)$;
  \item 反之, 如果 $\sigma_i$ 是常数但残差有 truth-bin 依赖 (例如 dE/dx 在大 $|\vec p|$ 处更大), pull 会有 truth-bin 依赖.
\end{itemize}

换言之, pull 是\emph{标准化残差}的事件级版本. 它对 "Fisher $\sigma$ 是否随事件正确缩放" 是敏感的诊断量.

\paragraph{Per-truth-bin pull 检查.}
% TODO: 把 744 事件按 15 个 truth bin 分组, 每 bin 算 pull mean, std,
% 输出 pull_per_truth_bin.csv. 期望: 14 个非病态 bin 的 pull std 接近 1
% (frame fix 后), $(-100, 0)$ bin 因 12% LM 失败拖出 std $\sim 5$.
% 命令: python3 scripts/analysis/analyze_ensemble_coverage.py \
%         --stratify-by-truth --pull-per-bin

预期结果(从 \nameref{tab:partIII:other_truth_points} ratio 列推算):
\begin{itemize}[nosep]
  \item 14 个非病态 bin: pull std $\sim 0.4$--$1.0$ (大多 Fisher 偏保守, $\sigma_z < 1$);
  \item $(-100, 0)$ bin: pull std $\sim 5$ ($p_x$, $p_z$ 双方向被 LM 失败拖大).
\end{itemize}

\subsection{Pull 与覆盖率的关系}
\label{sec:partIII:pulls:vs_coverage}

二者本质上同源: 一个分量的 68\% 覆盖率 $= \Pr(|z| \le 1)$. 对高斯分布严格等价. 实际数据由于长尾, 覆盖率(对长尾"鲁棒")与 $\sigma_z$ (对长尾敏感) 给出不同信号:

\begin{table}[H]
\centering\small
\caption{覆盖率与 $\sigma_z$ 的对比 (744 事件): 覆盖率对 LM 长尾鲁棒, $\sigma_z$ 不鲁棒.}
\label{tab:partIII:pull_vs_coverage}
\begin{tabular}{lrrr}
\toprule
分量 & 覆盖率 $|z|\le 1$ & $\sigma_z$ & 名义 $\Phi(\sigma_z)$ \\
\midrule
$p_x$ & 0.801 & 4.22 & 0.232 (高斯下) \\
$p_y$ & 0.421 & 2.48 & 0.394 (高斯下) \\
$p_z$ & 0.772 & 9.05 & 0.111 (高斯下) \\
$|\vec p|$ & 0.761 & 60.0 & 0.013 (高斯下) \\
\bottomrule
\end{tabular}
\end{table}

\textbf{解读}: $p_y$ 的 覆盖率 0.42 几乎与 "高斯下 $\sigma_z = 2.48$ 给出的 $\Phi(1/2.48) - \Phi(-1/2.48) = 0.394$" 一致 (差 $\sim 7\%$); 这说明 $p_y$ 主峰是单峰高斯, 但宽度被 Fisher 低估了 $2.48\times$. 反观 $p_x, p_z$, 经典 $\sigma_z$ 由长尾撑大到 $4 \sim 9$, 高斯换算只能给 $0.23 \sim 0.11$ 覆盖率, 但实测 $0.77 \sim 0.80$ —— 差异是因为分布不是高斯, 主峰窄, 长尾稀疏. 鲁棒指标 (中位 + IQR/1.349) 才是 $p_x, p_z$ 的可靠诊断.

\subsection{小结}
\label{sec:partIII:pulls:summary}

\begin{itemize}[nosep]
  \item Pull 是诊断 "误差棒形状是否正确" 的标准量, 期望 $\mathcal{N}(0, 1)$.
  \item 当前 744 事件: 三分量中位 pull $\sim 0$ (frame fix 后无偏), $p_y$ 主峰 $\sigma_z \approx 2.48$ ($\Rightarrow$ Fisher 低估 $\sim 1.86\times$, 与覆盖率 ratio 一致); $p_x, p_z$ 主峰窄但长尾干扰 $\sigma_z$.
  \item 接下来需要: pull 直方图绘图脚本; IQR/MAD 鲁棒尺度; truth-cut 后(去 LM 失败事件) 重新拟合主峰高斯, 应给出与覆盖率一致的 $\sigma_z$.
\end{itemize}

% =============================================================

\section{$\chi^2$ 分布与失败模式分类}
\label{sec:partIII:chi2}

\subsection{自由度 1 的 $\chi^2$ 分布是重尾的}
\label{sec:partIII:chi2:dof1}

RK 重建在 PDC1+PDC2 共两点时, 残差维度 = 4 (两个 PDC 各 $u, v$), 参数维度 = 3 ($u, v, p$ 的 $u$ 是斜率不是丝面坐标), 自由度 $N_{\mathrm{dof}} = 4 - 3 = 1$. 这是状态报告 §\(\chi^2\)~目标函数 写明的事实. 自由度 1 的 $\chi^2$ 分布密度为
\[
  f_{\chi^2_1}(x) \;=\; \frac{1}{\sqrt{2\pi x}} e^{-x/2}, \quad x > 0.
\]
在原点发散($x \to 0$ 时密度 $\to \infty$, 但积分仍收敛), 在尾部以 $e^{-x/2}/\sqrt{x}$ 衰减 —— 比 $\chi^2_2 \sim e^{-x/2}$ 慢. 数值上:

\begin{table}[H]
\centering\small
\caption{$\chi^2_1$ 分布的尾部概率 ($\chi^2_{\mathrm{red}} = \chi^2 / N_{\mathrm{dof}}$, 此处 $N_{\mathrm{dof}} = 1$).}
\label{tab:partIII:chi2_tail_theory}
\begin{tabular}{rrr}
\toprule
$\chi^2_{\mathrm{red}}$ 阈值 & $\Pr(\chi^2_1 > t)$ & $-\log_{10}\Pr$ \\
\midrule
1   & 0.3173 & 0.50 \\
2   & 0.1573 & 0.80 \\
3   & 0.0833 & 1.08 \\
4   & 0.0455 & 1.34 \\
6   & 0.0143 & 1.85 \\
10  & 0.00157 & 2.80 \\
20  & 7.7$\times 10^{-6}$ & 5.11 \\
50  & 1.5$\times 10^{-12}$ & 11.81 \\
\bottomrule
\end{tabular}
\end{table}

\textbf{推论}: 即使在零假设下(模型完美吻合), 1 自由度 $\chi^2_1 > 4$ 的事件仍占 $\sim 4.6\%$, $\chi^2_1 > 10$ 占 $\sim 0.16\%$. 在 744 事件中, 我们 \emph{期望}
\[
  N_{\chi^2 > 4} \approx 33.8, \qquad N_{\chi^2 > 10} \approx 1.17.
\]

\subsection{744 事件实测对比}
\label{sec:partIII:chi2:measured}

直接读 \texttt{track\_summary.csv} 的 \texttt{chi2\_reduced} 列:

\begin{table}[H]
\centering\small
\caption{744 事件的 $\chi^2_{\mathrm{red}}$ 分布观测值 vs.\ 理论 ($\chi^2_1$ 假设). $N_{\mathrm{exp}}$ = 744 $\times$ $\Pr(\chi^2_1 > t)$.}
\label{tab:partIII:chi2_observed}
\begin{tabular}{rrrr}
\toprule
阈值 $t$ & 观测 $N$ & 期望 $N_{\mathrm{exp}}$ & 观测/期望 \\
\midrule
2  & 91 & 117.0 & 0.78 \\
4  & 60 & 33.8  & 1.78 \\
6  & 54 & 10.6  & 5.09 \\
10 & 50 & 1.17  & 42.7 \\
20 & 42 & 0.006 & 7000 \\
\bottomrule
\end{tabular}
\end{table}

\paragraph{读法.}
\begin{itemize}[nosep]
  \item $\chi^2_{\mathrm{red}} > 2$ 出现 91 条 vs.\ 期望 117 —— \emph{低于}期望. 主峰核心拟合得比 1 自由度 $\chi^2_1$ 还紧, 提示 $N_{\mathrm{dof}}$ 可能在 \emph{有效} 上略大于 1 (例如 PDC 测量误差被夸大, 拟合"省下"自由度), 或 LM 把 $\chi^2$ 推得过于贴近 0.
  \item $\chi^2_{\mathrm{red}} > 4$ 之上, 观测开始压倒期望: $> 4$ 已经是 1.78 倍, $> 10$ 是 \textbf{42.7 倍}, $> 20$ 是 7000 倍.
  \item 50 条 $\chi^2_{\mathrm{red}} > 10$ 的事件, 从 7\% 占比看绝对不是 $\chi^2_1$ 自然涨落; 它们是另一个机制产生的 (LM 局部极小, 多次散射爆发, dE/dx 离群, 等等).
\end{itemize}

\subsection{失败模式分类}
\label{sec:partIII:chi2:taxonomy}

把 744 事件按 $\chi^2_{\mathrm{red}}$ 分成三类, 给出每类的代表事件 (从 \texttt{track\_summary.csv} 直接读取):

\begin{table}[H]
\centering\scriptsize
\caption{失败模式分类(按 $\chi^2_{\mathrm{red}}$ 分箱). 数据来自 \texttt{track\_summary.csv} 直接读取; \texttt{event\_index, track\_index} 唯一定位.}
\label{tab:partIII:taxonomy}
\begin{tabular}{lrlllll}
\toprule
类别 & $\chi^2_{\mathrm{red}}$ 范围 & 占比 & 代表事件 & 残差 ($p_x, p_y, p_z$) MeV/$c$ & Fisher $\sigma$ ($p_x, p_y, p_z$) & 物理诠释 \\
\midrule
A 正常 & $[0, 2)$ & $87.8\%$ & (626, 0) & $(-2.49, +0.66, +4.38)$ & $(6.82, 0.64, 6.91)$ & 名义重建, 无明显 dE/dx \\
A 正常 & $[0, 2)$ & --- & (428, 0) & $(+2.36, -0.97, -3.51)$ & $(7.00, 0.75, 6.22)$ & 同上, 不同 truth bin \\
B 中度 & $[2, 10]$ & $5.5\%$ & (482, 0) & $(-4.37, +1.14, +2.38)$ & $(11.16, 1.03, 5.41)$ & dE/dx 致 $p_z$ pull, $p_x$ ratio 大 \\
B 中度 & $[2, 10]$ & --- & (742, 0) & $(-148, +0.56, -549)$ & $(10.58, 1.45, 8.54)$ & 极少见: $p_z$ 跑到 77, LM "刚刚到" \\
C 严重 & $> 10$ & $6.7\%$ & (229, 0) & $(-449, +11.5, +312)$ & $(30.4, 1.98, 19.5)$ & LM trapped at $p_x \to -500$ \\
C 严重 & $> 10$ & --- & (121, 0) & $(-514, +2.49, +335)$ & $(58.5, 2.11, 38.7)$ & $(-100, 0)$ 病态点 \\
\bottomrule
\end{tabular}
\end{table}

\paragraph{Class A — 正常事件.}
$\sim 88\%$ 的事件落在这里. $\chi^2_{\mathrm{red}} \in [0, 2)$, 残差小到 Fisher $\sigma$ 量级, pull 接近 $\mathcal{N}(0, 1)$. 这些事件给出的覆盖率本应为 0.68; 当前 $p_y$ 在 Class A 内部仍欠覆盖, 因为 Class A 内 \emph{每事件} 的 $p_y$ Fisher 低估 $\sim 1.86\times$.

\paragraph{Class B — 中度失配 ($\chi^2_{\mathrm{red}} \in [2, 10]$).}
$\sim 5.5\%$. 残差不再 Fisher $\sigma$ 量级但仍有限. 主要机制:
\begin{enumerate}[label=(\arabic*), nosep]
  \item dE/dx 缺失: $p_z$ 在 1 m 空气段损失 $\sim 0.5$ MeV/$c$, 在 PDC 气体里再损失 $\sim 0.2$ MeV/$c$, 合计 $\sim 0.7$ MeV/$c$; LM 把这个能损"吸收"到 $p_x, p_z$ 的拟合里, $\chi^2$ 增大;
  \item 多次散射在 1 m 空气段超出 Fisher 假设(高斯独立 $\sigma$): 一次大角散射 $\theta \sim 5\sigma_{\mathrm{ms}}$ 让两个 PDC 的 $u$ 残差不再独立, $\chi^2$ 翻倍;
  \item $(u, v, p)$ 参数化的非线性: LM 收敛到 $(u, v, p)$ 空间的最小, 但其在 $(p_x, p_y, p_z)$ 空间未必无偏(参看 Part II §8 详述).
\end{enumerate}

\paragraph{Class C — 严重失败 ($\chi^2_{\mathrm{red}} > 10$).}
$\sim 6.7\%$, 50 条事件. 三种机制:
\begin{enumerate}[label=(\arabic*), nosep]
  \item \textbf{LM 困在错误盆地}: 网格搜索的初值 $(u, v, p)$ 落在 $\chi^2$ 二级最小附近, LM 收敛到 $\hat p_x \sim -500$, $\hat p_z \sim 940$ —— 这是状态报告 §误差分析 里描述的 "(-100, 0) 病态点". 第 \ref{sec:partIII:basin} 章详细剖析.
  \item \textbf{反常 PDC 命中}: PDC 簇心被错误识别(如多径迹串扰), 单条 hit 的 $u$ 偏离 truth $\sim 100$ mm, 残差扭曲, $\chi^2 \to 10^3$. 当前 PDC 模拟未注入此类异常, 但实际数据中确实存在.
  \item \textbf{多次散射爆发}: 1 m 空气段一次硬散射(电磁簇射或大角弹性), 让事件偏离高斯 MS 假设. 当前 Geant4 模拟开了 \texttt{G4UrbanMscModel} 含尾部, 偶然发生概率 $\sim 10^{-3}$.
\end{enumerate}

\subsection{每类事件在覆盖率统计中的贡献}
\label{sec:partIII:chi2:contribution}

\begin{table}[H]
\centering\small
\caption{各类事件在 744 事件覆盖率统计中的贡献(估算).}
\label{tab:partIII:chi2_class_contribution}
\begin{tabular}{lrr}
\toprule
类别 & 事件数 & 对总覆盖率(of $|\vec p|$) 拉低量 \\
\midrule
A 正常 & 654 & $\sim 0$ (本身贴近 0.68) \\
B 中度 & 41  & $\sim 0.005$ (覆盖 $\sim 0.5$, 比 0.68 低 0.18, $\times$ 41/744) \\
C 严重 & 49  & $\sim 0.066$ (覆盖 $\sim 0$, 比 0.68 低 0.68, $\times$ 49/744) \\
\midrule
合计    & 744 & $\sim 0.071$, 即把名义 0.68 拉到 $\sim 0.61$ \\
\bottomrule
\end{tabular}
\end{table}

但是状态报告 $|\vec p|$ 实测覆盖率是 0.76 不是 0.61 —— Fisher $\sigma$ 在 Class C 里 \emph{自适应增大} (例如事件 121 $\sigma_{p_x} = 58.5$, 事件 229 $\sigma_{p_x} = 30.4$), 部分 Class C 事件因 Fisher $\sigma$ 巨大反而被算作 "覆盖". 这是 Fisher 在大 $\chi^2$ 时数值不稳定的副作用 ——它假设 $\chi^2$ 在最优点是高斯 quadratic, 但当 LM 困在错误盆地, "最优点" 在错误位置, Hessian 给出的 $\sigma$ 也在错误尺度上.

\subsection{Wilks's theorem 与 $\chi^2$ 解读边界}
\label{sec:partIII:chi2:wilks}

Wilks (1938) 定理给出: 在零假设下, $-2 \ln (\mathcal{L}_{\mathrm{null}} / \mathcal{L}_{\mathrm{alt}})$ 渐近 $\chi^2_{\Delta k}$ 分布, 其中 $\Delta k$ 是参数维度差. 在我们的问题里:
\begin{itemize}[nosep]
  \item 备择假设 $H_1$: 自由 $(u, v, p)$ 三参数最小化 $\chi^2$, 给 $\chi^2_{\min}$;
  \item 零假设 $H_0$: 模型完美 (truth 已知, 无须拟合), 残差 $\sim \mathcal{N}(0, \sigma_{\mathrm{PDC}})$;
  \item Wilks: $\chi^2_{\min} \sim \chi^2_{N_{\mathrm{res}} - N_{\mathrm{par}}} = \chi^2_1$.
\end{itemize}

\paragraph{Wilks 失效的条件.}
Wilks 假设 (a) 模型在最优点是 regular (Hessian 正定), (b) 数据量大 (asymptotic). 在 $N_{\mathrm{dof}} = 1$ 时, "数据量大" 是个尴尬的条件 — 残差只有 4 个, 参数 3 个, 远未 asymptotic. 但是 PDC 残差是 \emph{独立 4 维高斯}, 在这种全高斯情形 Wilks 即便 $N_{\mathrm{dof}} = 1$ 也精确成立. 这是为什么状态报告中 Class A 的 $\chi^2$ 经验分布与 $\chi^2_1$ 形状一致.

\paragraph{Wilks 在 LM 失败时不成立.}
Class C 事件不满足 Wilks 假设, 因为 (a) LM 不在 $\chi^2$ 全局最小, Hessian 在错盆地内的曲率与全局最小处不同, regular 假设破裂. 这就是为什么 Class C 的 $\chi^2$ 数量级是 $10^3$ 而不是 $\chi^2_1$ 的 $\sim 1$.

\paragraph{p-value 解读.}
对 Class A 主峰内的事件, 把 $\chi^2_{\mathrm{red}}$ 转成 p-value:
\[
  p_{\mathrm{val}} = \Pr(\chi^2_1 > \chi^2_{\mathrm{obs}}).
\]
例如 $\chi^2_{\mathrm{obs}} = 0.5 \Rightarrow p_{\mathrm{val}} = 0.48$ (好); $\chi^2_{\mathrm{obs}} = 4 \Rightarrow p_{\mathrm{val}} = 0.046$ (5\% 水平边界). 在 744 事件中, p-value 小于 $0.05$ 的事件应占 $\sim 5\%$, 即 $\sim 37$ 条; 实测 60 条 $\chi^2 > 4$, 比期望多 $\sim 1.8\times$, 与表 \ref{tab:partIII:chi2_observed} 一致.

\subsection{$\chi^2$ 阈值作为事件级 quality cut}
\label{sec:partIII:chi2:cut}

实用建议: 在最终物理分析里, 加 $\chi^2_{\mathrm{red}} < 5$ 的 quality cut 可以剔除大部分 Class C, 保留 $> 95\%$ 的 Class A+B. 这种 cut 在覆盖率上的效果(模拟):

\begin{table}[H]
\centering\small
\caption{$\chi^2_{\mathrm{red}}$ cut 对覆盖率的影响(估算, 待运行).}
\label{tab:partIII:chi2_cut_estimate}
\begin{tabular}{lrrrr}
\toprule
Cut & 保留事件 & $|\vec p|$ Fisher 68\% & $p_y$ Fisher 68\% & 备注 \\
\midrule
无         & 744 & 0.761 & 0.421 & 当前数字 \\
$< 10$    & $\sim 694$ & TODO & TODO & 应消除 LM 失败 \\
$< 5$     & $\sim 690$ & TODO & TODO & 同上 + 中度 dE/dx 离群 \\
$< 2$     & $\sim 653$ & TODO & TODO & 极保守 \\
\bottomrule
\end{tabular}
\end{table}

% TODO: 上表数字需要在 analyze_ensemble_coverage.py 加 --chi2-max 选项后重跑.
% 命令: python3 scripts/analysis/analyze_ensemble_coverage.py \
%   --input-dir build/test_output/ensemble_n50_targetframe/.../ \
%   --output-dir build/test_output/ensemble_n50_targetframe/coverage_chi2cut5 \
%   --chi2-max 5
% 期望运行时间: 30 秒 (纯 Python 后处理).

\subsection{小结}
\label{sec:partIII:chi2:summary}

\begin{itemize}[nosep]
  \item $\chi^2_1$ 分布尾部重: $> 4$ 的占 $\sim 4.6\%$, $> 10$ 的占 $\sim 0.16\%$ (无机制下).
  \item 实测 744 事件中, $> 10$ 占 $6.7\%$, 比理论高 \emph{42 倍}, 是真正 "另一个机制" 在主导.
  \item 三类机制: LM 错盆地 (主要), 反常 hit, MS 爆发. 对应代表事件已列在表 \ref{tab:partIII:taxonomy}.
  \item $\chi^2$ cut 是简单粗暴但有效的 quality cut, 应在生产分析里默认开启.
\end{itemize}

% =============================================================

\section{单事件深度分析: 6 个示例}
\label{sec:partIII:case_studies}

本章逐条分析 6 个代表事件, 每条事件:
(a) 引用 \texttt{track\_summary.csv} 的原始行;
(b) 标注哪些 Part II §1--9 误差源主导其残差;
(c) 当 \texttt{profile\_samples.csv}/\texttt{bayesian\_samples.csv} 包含该事件时, 引用对应行;
(d) 讨论 LM trace 的预期形态 (实际 trace 数据当前未持久化, 需要 \texttt{PDCErrorAnalysis.cc} 加调试输出).

事件选取覆盖谱为: 标准 Class A (1 条), 中度 dE/dx (1 条), $p_y$ 显著欠覆盖 (1 条), Class C 极端 (1 条 from $(-100, 0)$ 病态点), Class C 中等 (1 条), 边界事件 (1 条).

\subsection{事件 (655, 0) — Class A 标准}
\label{sec:partIII:case_studies:event_655}

\begin{table}[H]
\centering\small
\caption{事件 (655, 0) 数据卡片.}
\begin{tabular}{ll}
\toprule
truth $(p_x, p_y, p_z)$ & $(50.0, 0.0, 627.0)$ MeV/$c$ \\
reco $(p_x, p_y, p_z)$ & $(47.91, -0.67, 630.46)$ MeV/$c$ \\
残差 $(p_x, p_y, p_z)$ & $(-2.09, -0.67, +3.46)$ MeV/$c$ \\
$\chi^2_{\mathrm{red}}$ & $0.000308$ \\
Fisher $\sigma$ $(p_x, p_y, p_z)$ & $(6.82, 0.63, 6.95)$ MeV/$c$ \\
LM 迭代数 & 0 (网格直接命中) \\
Profile 区间 ($p_x$) & $[40.59, 55.23]$, 半宽 $7.32$ \\
Profile 区间 ($p_y$) & $[-1.13, -0.21]$, 半宽 $0.46$ \\
Profile 区间 ($p_z$) & $[623.83, 637.08]$, 半宽 $6.62$ \\
MCMC 区间 ($|\vec p|$) & $[628.19, 637.67]$, 半宽 $4.74$ \\
MCMC acc / ESS / Geweke & $0.355 / 13.8 / 6.31$ \\
\bottomrule
\end{tabular}
\end{table}

\paragraph{诠释.}
\begin{itemize}[nosep]
  \item 残差全部 $< 1\sigma$, 拟合理想; $\chi^2_{\mathrm{red}} < 10^{-3}$ 提示 LM 把残差几乎压到机器精度.
  \item Profile 区间宽度($7.32, 0.46, 6.62$) 比 Fisher 略窄, 与状态报告 §误差分析方法对比表 中 Profile $\sim 0.94 \times$ Fisher 的趋势一致.
  \item MCMC ESS 仅 13.8 (链长 160 步, 80 burn-in), Geweke $z = 6.31$ 远超 $|z| > 2$ 的不收敛阈值 —— 链没有 mix; MCMC 区间数字"碰巧" 与 Profile 接近不代表后验已被恰当采样, 仅说明高斯峰在 $\hat p \sim 633$ 附近且 walker 在峰内随机游走.
  \item LM 迭代数 = 0 意味着粗网格搜索直接给出 $\chi^2 < 10^{-3}$ 的候选, 后续 LM 不需要进一步移动. 这是 Class A 的常见模式.
\end{itemize}

\paragraph{Part II 误差源归因.}
\begin{itemize}[nosep]
  \item PDC 位置分辨 (Part II §1): $\sim 200$ μm 在两个 PDC 上, 通过 Glückstern 给出 $\sigma_p \sim 6$--$7$ MeV/$c$ —— 与 Fisher 一致;
  \item MS in air (Part II §2): $\sim 0.3$ mrad over 1 m, 在 $|\vec p| = 627$ 上贡献 $\sim 1$ MeV/$c$ —— 嵌在 6.8 里;
  \item dE/dx mean (Part II §5): $\sim 0.7$ MeV/$c$ for 1 m air + 2 PDC, 残差 $-2$ MeV/$c$ 与之同号但稍大, 可能配额 $\sim 1$ MeV/$c$;
  \item 其他源: 在此事件中可忽略.
\end{itemize}

\subsection{事件 (482, 0) — Class B 中度失配}
\label{sec:partIII:case_studies:event_482}

\begin{table}[H]
\centering\small
\caption{事件 (482, 0) 数据卡片.}
\begin{tabular}{ll}
\toprule
truth $(p_x, p_y, p_z)$ & $(-100.0, +20.0, 627.0)$ MeV/$c$ \\
reco $(p_x, p_y, p_z)$ & $(-104.37, +21.14, 629.38)$ MeV/$c$ \\
残差 $(p_x, p_y, p_z)$ & $(-4.37, +1.14, +2.38)$ MeV/$c$ \\
$\chi^2_{\mathrm{red}}$ & $2.035$ \\
Fisher $\sigma$ $(p_x, p_y, p_z)$ & $(11.16, 1.03, 5.41)$ MeV/$c$ \\
LM 迭代数 & 0 \\
\bottomrule
\end{tabular}
\end{table}

\paragraph{诠释.}
\begin{itemize}[nosep]
  \item 残差 $/\sigma$ 比为 $(0.39, 1.11, 0.44)$ — $p_y$ 已经在 $1\sigma$ 边缘. 这件 (truth $p_y = 20$) 加上 (truth $p_x = -100$) 提供较高 PDC2 弯曲 $\theta_{xz} \sim -16°$, 把 dE/dx 在弯曲弧上分配的能量损失更多, $p_z$ 残差正向偏 (重建 $p_z$ 大于 truth, 因为 dE/dx 缺失意味着 LM 假设无能损, 把 momentum 分给了 $p_z$).
  \item $\chi^2_{\mathrm{red}} = 2.03$ 处于 $\chi^2_1$ 上 $\sim 16\%$ 的概率密度尾, 是 Class B 的边缘; 但 6.7\% 的 $\chi^2_{\mathrm{red}} > 2$ 占比远大于 16\%, 说明 dE/dx 系统性地压在所有 $|p_x| = 100$ 事件上.
  \item Fisher $\sigma_{p_x} = 11.16$ 略高于平均(其他 $|p_x| = 100$ 事件 $\sigma_{p_x} \sim 7$--$9$), 可能因为 LM 困在略偏离最优的位置, Hessian 在那里的曲率更平.
\end{itemize}

\paragraph{Part II 归因.}
\begin{itemize}[nosep]
  \item dE/dx 缺失 (Part II §5) — 主导, 贡献 $\sim 0.7$ MeV/$c$ 到 $p_z$ 残差;
  \item $(u, v, p)$ 参数化偏 (Part II §8) — 在大 $|p_x|$ 时尤其显著, 贡献 $\sim 1$--$2$ MeV/$c$;
  \item PDC 分辨 + MS — 嵌在 Fisher $\sigma$ 中.
\end{itemize}

\subsection{事件 (102, 0) — $p_y$ 强欠覆盖}
\label{sec:partIII:case_studies:event_102}

\begin{table}[H]
\centering\small
\caption{事件 (102, 0) 数据卡片.}
\begin{tabular}{ll}
\toprule
truth $(p_x, p_y, p_z)$ & $(100.0, -20.0, 627.0)$ MeV/$c$ \\
reco $(p_x, p_y, p_z)$ & $(105.99, -24.36, 624.85)$ MeV/$c$ \\
残差 $(p_x, p_y, p_z)$ & $(+5.99, -4.36, -2.15)$ MeV/$c$ \\
$\chi^2_{\mathrm{red}}$ & $0.133$ \\
Fisher $\sigma$ $(p_x, p_y, p_z)$ & $(6.65, 0.49, 7.19)$ MeV/$c$ \\
$z_{p_y}$ & $-8.94$ \\
Profile $p_y$ 区间 & $[-24.76, -23.96]$, 半宽 $0.40$ \\
MCMC $|\vec p|$ & $[629.03, 642.06]$, ESS $13.5$, Geweke $1.54$ \\
\bottomrule
\end{tabular}
\end{table}

\paragraph{诠释.}
\begin{itemize}[nosep]
  \item $p_y$ 残差 $-4.36$ 是 Fisher $\sigma_{p_y} = 0.49$ 的 \emph{8.9 倍} —— 强欠覆盖典型. 但是 $\chi^2_{\mathrm{red}} = 0.133$ 极低, 说明 LM 把整个残差消耗在 \emph{两个 PDC 同向} 的 $v$ 残差上, $\chi^2$ 没有 "怒"
.
  \item Profile 区间 $[-24.76, -23.96]$ 半宽 $0.40$ 仍然不包含 truth $-20$ —— Profile 与 Fisher 同方向欠覆盖, 这是 Part II §8 描述的 $(u, v, p)$ 线性化偏 $p_y$ 的直接表现.
  \item MCMC ESS 13.5, Geweke $1.54$ 偶然在阈值内 (但仍低 ESS), MCMC 也在 \emph{偏移的}峰附近游走.
\end{itemize}

\paragraph{Part II 归因.}
$p_y$ 残差 $-4.36$ 全部归到 \emph{(u, v, p) 参数化的 Jacobian 压扁} (Part II §8). 这是 "Fisher 信息矩阵在 $(u, v, p)$ 空间到 $(p_x, p_y, p_z)$ 空间转换时, $p_y$ 列 Jacobian $\partial p_y / \partial v$ 没有恢复全部曲率信息". 由于 truth $p_y = -20$ 超出 $\pm \sigma_{p_y}$ 9 倍, LM 找的最优 $v$ 与 truth $v$ 差距远大于线性化下的预期; Fisher $\sigma$ 是从最优点的 Hessian 导, 在最优点附近的曲率不能反映这种远距离偏移.

\subsection{事件 (482, 0) 边界化版本: 事件 (742, 0) — Class B 边界}
\label{sec:partIII:case_studies:event_742}

\begin{table}[H]
\centering\small
\caption{事件 (742, 0) 数据卡片. 这个事件介于 Class B 与 Class C, $\chi^2_{\mathrm{red}} \approx 2$ 但残差极大 — LM 找到了一个"chi-square 看起来满意但物理上荒谬"的解.}
\begin{tabular}{ll}
\toprule
truth $(p_x, p_y, p_z)$ & $(0.0, 0.0, 627.0)$ MeV/$c$ \\
reco $(p_x, p_y, p_z)$ & $(-147.95, +0.56, +77.50)$ MeV/$c$ \\
残差 $(p_x, p_y, p_z)$ & $(-147.95, +0.56, -549.50)$ MeV/$c$ \\
$\chi^2_{\mathrm{red}}$ & $2.040$ \\
Fisher $\sigma$ $(p_x, p_y, p_z)$ & $(10.58, 1.45, 8.54)$ MeV/$c$ \\
LM 迭代数 & 5 \\
\bottomrule
\end{tabular}
\end{table}

\paragraph{诠释.}
\begin{itemize}[nosep]
  \item 残差 $p_z = -549.5$, $|\vec p| = 167$ vs.\ truth 627 —— LM 落在 $|\vec p|$ 极低的局部极小. 但 $\chi^2_{\mathrm{red}} = 2.04$ 看起来还行, 这是\emph{反演问题不唯一}的典型征兆.
  \item Fisher $\sigma$ 在错的位置算: $\sigma_{p_z} = 8.54$ 反映的是 \emph{该 chi-square 谷里} 的二次曲率, 与正确谷($\hat p_z \sim 627$) 的曲率无关.
  \item LM 迭代了 5 次后困在这, 因为粗网格初值落在错盆地. 网格 $u_{\mathrm{center}} = 0 / 627 \to 0$, $v = 0$, $p \in \{200, 400, 700, 1200, 2000\}$ 的某一个 $\Rightarrow$ 这种 truth $(0, 0, 627)$ 居然能困入低 $|\vec p|$ 谷, 提示磁场积分 + RK 可能在某个 $u$ 上有副 $\chi^2$ 极小.
\end{itemize}

\paragraph{Part II 归因.}
\begin{itemize}[nosep]
  \item LM 收敛盆地失败 (Part II §9) — 主导;
  \item 网格搜索初始化不足: 7$\times$3$\times$5 = 105 个候选, 但 $u$ 半幅 1.0 给的网格步长 $0.33$ 在某些 truth 点会跨过 chi-square 鞍点.
\end{itemize}

\subsection{事件 (121, 0) — Class C 极端 ($-100, 0$ 病态)}
\label{sec:partIII:case_studies:event_121}

\begin{table}[H]
\centering\small
\caption{事件 (121, 0) 数据卡片. 这是 \nameref{sec:partIII:basin} 章详细分析的"病态点"代表.}
\begin{tabular}{ll}
\toprule
truth $(p_x, p_y, p_z)$ & $(-100.0, 0.0, 627.0)$ MeV/$c$ \\
reco $(p_x, p_y, p_z)$ & $(-613.70, +2.49, +961.92)$ MeV/$c$ \\
残差 $(p_x, p_y, p_z)$ & $(-513.70, +2.49, +334.92)$ MeV/$c$ \\
$\chi^2_{\mathrm{red}}$ & $1243.89$ \\
Fisher $\sigma$ $(p_x, p_y, p_z)$ & $(58.49, 2.11, 38.66)$ MeV/$c$ \\
LM 迭代数 & 0 (粗网格已经把它推到这) \\
\bottomrule
\end{tabular}
\end{table}

\paragraph{诠释.}
\begin{itemize}[nosep]
  \item 残差 $p_x = -514$ 是 truth 的 $5\times$, $|\vec p| \sim 1145$ vs.\ truth 635. 这是粗网格的某个 $(u, v, p)$ 候选直接给出的 $\chi^2 = 1244$, 而 LM 没有从这里再走 (迭代数 $= 0$ 意味着初值的 $\chi^2$ 在多起点中是最小的).
  \item Fisher $\sigma_{p_x} = 58.5$ 巨大, 因为 Hessian 在这个错盆地里曲率小; 但残差 $-514$ 仍是 $\sigma$ 的 $8.8\times$, pull $z = -8.78$.
  \item 该事件不在 \texttt{profile\_samples.csv} 中(profile 抽样按 $\chi^2$ 分位抽 32 条 $\times$ 4 quartile = 128 条, 该事件 $\chi^2$ 离群 + 不在 quartile 抽样列表里).
\end{itemize}

\paragraph{Part II 归因.}
\begin{itemize}[nosep]
  \item LM 收敛盆地失败 (Part II §9) — 几乎全部;
  \item 第六章详细剖析为什么 truth $(p_x, p_y) = (-100, 0)$ 是网格 $u$ 范围的边界点.
\end{itemize}

\subsection{事件 (229, 0) — Class C 中等}
\label{sec:partIII:case_studies:event_229}

\begin{table}[H]
\centering\small
\caption{事件 (229, 0) 数据卡片. truth $(p_x, p_y) = (-50, 0)$, 跑成了 LM 失败.}
\begin{tabular}{ll}
\toprule
truth $(p_x, p_y, p_z)$ & $(-50.0, 0.0, 627.0)$ MeV/$c$ \\
reco $(p_x, p_y, p_z)$ & $(-499.56, +11.54, +938.68)$ MeV/$c$ \\
残差 $(p_x, p_y, p_z)$ & $(-449.56, +11.54, +311.68)$ MeV/$c$ \\
$\chi^2_{\mathrm{red}}$ & $1548.20$ \\
Fisher $\sigma$ $(p_x, p_y, p_z)$ & $(30.43, 1.98, 19.48)$ MeV/$c$ \\
LM 迭代数 & 0 \\
\bottomrule
\end{tabular}
\end{table}

\paragraph{诠释.}
\begin{itemize}[nosep]
  \item 与事件 (121, 0) 同模式但 truth $(p_x, p_y) = (-50, 0)$. 显示 (-100, 0) 病态点不孤立 — 在 $(-50, 0)$ 也有少数 LM 失败.
  \item Fisher $\sigma$ 在错盆地的曲率 \emph{比} (-100, 0) 病态点 \emph{小}, 因为该 truth 离病态边界更远; 即 $-50$ 的"病态" 程度比 $-100$ 轻.
\end{itemize}

\paragraph{Part II 归因.}
LM 收敛盆地失败 (Part II §9). 但发生频率(50 个 $(-50, 0)$ 事件中只 $\sim 2\%$ 失败) 远低于 $(-100, 0)$ ($\sim 12\%$).

\subsection{Profile/MCMC 跨事件诊断}
\label{sec:partIII:case_studies:profile_mcmc}

把上面 6 个事件的 Profile/MCMC 数据 (能查到的) 与 Fisher 横向对比:

\begin{table}[H]
\centering\scriptsize
\caption{Fisher / Profile / MCMC 三方法在示例事件上的区间宽度对比 (MeV/$c$). 缺失项表示该事件未抽到 Profile/MCMC 子集.}
\label{tab:partIII:case_methods_comparison}
\begin{tabular}{lrrrrrrr}
\toprule
事件 & 分量 & Fisher $\sigma$ & Fisher 半宽(68\%) & Profile 半宽(68\%) & MCMC 半宽(68\%) & MCMC ESS & MCMC Geweke \\
\midrule
\multirow{4}{*}{(655, 0)}
 & $p_x$ & 6.82 & 6.82 & 7.32 & --- & --- & --- \\
 & $p_y$ & 0.63 & 0.63 & 0.46 & --- & --- & --- \\
 & $p_z$ & 6.95 & 6.95 & 6.62 & --- & --- & --- \\
 & $|\vec p|$ & 6.45 & 6.45 & 6.45 & 4.74 & 13.8 & 6.31 \\
\midrule
\multirow{4}{*}{(178, 0)}
 & $p_x$ & 6.35 & 6.35 & 6.75 & --- & --- & --- \\
 & $p_y$ & 0.50 & 0.50 & 0.46 & --- & --- & --- \\
 & $p_z$ & 7.45 & 7.45 & 7.13 & --- & --- & --- \\
 & $|\vec p|$ & --- & --- & --- & 5.08 & 13.6 & 1.83 \\
\midrule
\multirow{4}{*}{(102, 0)}
 & $p_x$ & 6.65 & 6.65 & 5.69 & --- & --- & --- \\
 & $p_y$ & 0.49 & 0.49 & 0.40 & --- & --- & --- \\
 & $p_z$ & 7.19 & 7.19 & 6.07 & --- & --- & --- \\
 & $|\vec p|$ & --- & --- & --- & 6.51 & 13.5 & 1.54 \\
\midrule
\multirow{4}{*}{(298, 0)}
 & $p_x$ & --- & --- & 6.60 & --- & --- & --- \\
 & $p_y$ & --- & --- & 0.50 & --- & --- & --- \\
 & $p_z$ & --- & --- & 6.97 & --- & --- & --- \\
 & $|\vec p|$ & --- & --- & --- & 6.65 & 8.1 & 4.68 \\
\bottomrule
\end{tabular}
\end{table}

\paragraph{读法.}
\begin{itemize}[nosep]
  \item 三方法宽度 \emph{在同一事件内} 差距 $< 25\%$ — 与状态报告 §误差分析 中 ensemble 级的差距一致.
  \item Profile $p_y$ 比 Fisher $p_y$ 略\emph{窄} (0.40--0.50 vs.\ 0.49--0.65) — Profile 沿 $\Delta\chi^2 = 1$ 轮廓走, 而 Fisher 是局部曲率假高斯; 在 $(u,v,p)$ 参数化下两者都低估真实 $\sigma$, 但 Profile 更甚 (因为它沿 chi-square 等高线走得更紧).
  \item MCMC ESS 普遍 $13$ 左右 (链长 160 步, burn-in 80 步), Geweke $|z|$ 多 $> 2$ — 链未 mix, 但区间数字仍然有意义, 因为 walker 仍在主峰内随机游走, 16/84 分位点是经验估计的中心宽度.
\end{itemize}

\paragraph{为什么 MCMC ESS 低?}
状态报告 §误差分析 给出 ensemble 中位 ESS $\approx 12.6$ (\texttt{summary.csv}). 链长 $160 - 80 = 80$ 净样本, ESS $\approx 12$ 意味着 \emph{自相关时间} $\tau \approx 80/12 \approx 7$ 步 — 即每 7 步独立一次, 接收率 0.33 配合 walker step size $\approx \sigma_{\mathrm{Fisher}}$ 是合理的, 但链长太短. 把链长增到 $10^4$ 应该让 ESS $\sim 10^3$, 那时 MCMC 区间才真正 reliable. 这是 Part II 不在范围, 但 Part III 应记录下来作为 next step.

\subsection{所有 6 个事件的对比总览}
\label{sec:partIII:case_studies:overview}

\begin{table}[H]
\centering\scriptsize
\caption{6 个示例事件的全面对比.}
\label{tab:partIII:case_overview}
\begin{tabular}{lrrrrrrl}
\toprule
事件 & truth $(p_x, p_y)$ & 残差 $|p_x|$ & 残差 $|p_y|$ & 残差 $|p_z|$ & $\chi^2_{\mathrm{red}}$ & 类别 & 主因 \\
\midrule
(655, 0) & $(50, 0)$    & 2.1   & 0.7   & 3.5   & 0.0003 & A & 名义 \\
(482, 0) & $(-100, 20)$ & 4.4   & 1.1   & 2.4   & 2.03   & B & dE/dx + (u,v,p) bias \\
(102, 0) & $(100, -20)$ & 6.0   & 4.4   & 2.1   & 0.13   & A* & (u,v,p) $p_y$ 线性化 \\
(742, 0) & $(0, 0)$     & 148   & 0.6   & 549   & 2.04   & B & LM 错盆地 (chi-square 局部) \\
(121, 0) & $(-100, 0)$  & 514   & 2.5   & 335   & 1244   & C & 病态点 \\
(229, 0) & $(-50, 0)$   & 450   & 11.5  & 312   & 1548   & C & 同上, 减弱版 \\
\bottomrule
\end{tabular}
\end{table}
\textit{注: A* = $\chi^2_{\mathrm{red}}$ 落在 Class A 范围, 但 $p_y$ 单分量极端欠覆盖, 是 (u, v, p) 参数化偏的特殊体现.}

\paragraph{LM trace 占位.}
% TODO: 当前 PDCErrorAnalysis.cc 不持久化 LM 的迭代历史 (lambda 变化, 残差变化, 参数轨迹).
% 加入 --dump-lm-trace 选项, 把每条事件的 (iter, lambda, chi2, u, v, p) 写到
% build/test_output/.../lm_traces/event_{N}.csv 即可绘 trace 图.
% 期望对 121, 229, 742 三条 Class C 事件画 (u, p) 平面的 LM 轨迹,
% 应能直观看到它们粘在错盆地, 没有跳出.

\subsection{小结}
\label{sec:partIII:case_studies:summary}

\begin{itemize}[nosep]
  \item 6 个示例覆盖 4 种主要机制: 名义 (A), dE/dx (B), $(u, v, p)$ 线性化偏 (A*), LM 错盆地 (B/C).
  \item 状态报告 ratio 1.86 ($p_y$) 在事件 102 上具体表现为 \emph{残差 $4.4 \times \sigma$} 的极端欠覆盖.
  \item 状态报告 $|\vec p|$ 覆盖率 0.76 在事件 742, 121, 229 上具体表现为 LM 错盆地 + Fisher $\sigma$ 在错盆地里仍给出"看似自洽" 的曲率.
  \item LM trace 数据持久化是接下来的高优先级 dev item.
\end{itemize}

% =============================================================

\section{扫描研究方案与外推}
\label{sec:partIII:scans}

本章给出四组待运行的扫描研究; 大部分是 \emph{方案+外推}, 实际运行未执行. 每个扫描列出: 物理动机, 期望趋势, 待运行命令, 期望耗时, 输出 CSV 路径.

\subsection{扫描 BS-1: 磁场 $B_y$ $\pm 5\%$}
\label{sec:partIII:scans:bs1}

\paragraph{物理动机.}
SAMURAI 偶极磁场图来自实测 + 插值 (\texttt{sm\_field.dat}). 现场实测精度 $\pm 1\%$, 但插值与 RK 步进对中心场强敏感. 扫 $B_y$ 标度 $\pm 5\%$ 给出谱仪对场强标定的灵敏度.

\paragraph{期望趋势.}
RK 重建假设场图正确; 真实场强为 $B (1 + \delta_B)$ 但代码用 $B$, 导致弯曲半径 $\rho = p / (qB)$ 算错 $\delta_B$. 拟合时 LM 把这个吸收到 $\hat p$ 里, 即 $\hat p / p \approx 1 + \delta_B$. 线性外推:
\[
  \hat{|\vec p|} \;\approx\; |\vec p|^{\mathrm{truth}} \cdot (1 + \delta_B), \quad \sigma_{|\vec p|}^{\mathrm{stat}} \to \sigma_{|\vec p|}^{\mathrm{stat}}.
\]
$\sigma$ 不变, 但中心估计偏移. 当前 $|\vec p|$ 中位数 630 MeV/$c$, $\delta_B = +5\%$ 给 $\hat p \approx 661$, $\delta_B = -5\%$ 给 $\hat p \approx 599$.

\paragraph{覆盖率推论.}
\begin{itemize}[nosep]
  \item $\delta_B = +5\%$: 中心偏 $+30$ MeV/$c$, Fisher $\sigma \sim 7$ MeV/$c$, pull mean $\sim 4.3$, 覆盖率 $|\vec p|$ 应跌到 $< 5\%$.
  \item $\delta_B = -5\%$: 镜像, pull mean $\sim -4.3$.
  \item $\delta_B \pm 1\%$: pull mean $\sim 0.9$, 覆盖率从 0.76 跌到 $\sim 0.6$.
\end{itemize}

\paragraph{待运行命令.}
% TODO: B-field scan run BS-1
\begin{lstlisting}[language=bash,caption={B-field scan BS-1 命令模板},label=lst:partIII:bs1_cmd]
# 先把 sm_field.dat 复制成两份缩放副本:
python3 scripts/analysis/scale_sm_field.py --in data/input/sm_field.dat \
    --out data/input/sm_field_scaled_p5pct.dat  --scale 1.05
python3 scripts/analysis/scale_sm_field.py --in data/input/sm_field.dat \
    --out data/input/sm_field_scaled_m5pct.dat  --scale 0.95

for SCALE in p5pct m5pct; do
  ENSEMBLE_SIZE=50 PROFILE_PER_QUARTILE=8 MCMC_PER_QUARTILE=16 \
      MAGNETIC_FIELD=data/input/sm_field_scaled_${SCALE}.dat \
      OUT_ROOT=build/test_output/scan_bs1_${SCALE} \
      bash scripts/analysis/run_ensemble_coverage.sh
done
\end{lstlisting}

\paragraph{期望耗时.} 单 scale $\sim 40$ min ($\times 2 = 80$ min).
\paragraph{期望输出.} \texttt{build/test\_output/scan\_bs1\_p5pct/}, \texttt{scan\_bs1\_m5pct/} 各带完整 ensemble 输出.

\subsection{扫描 WS-1: PDC 丝坐标分辨 $\sigma_{\mathrm{wire}}$}
\label{sec:partIII:scans:ws1}

\paragraph{物理动机.}
PDC 丝坐标分辨当前默认 $\sigma_{\mathrm{wire}} \approx 0.2$ mm (对应 cluster 重心精度). 实测 PDC 数据有 $\sim 0.5$ mm; 极端情形 (cluster 中包含 noise hit) 可到 $\sim 2$ mm. 扫 $\sigma_{\mathrm{wire}} \in \{0.1, 0.3, 0.5, 1.0, 2.0\}$ mm 验证 Glückstern $\sigma_p \propto \sigma_{\mathrm{wire}}$ 关系.

\paragraph{期望趋势.}
Glückstern 公式 (Part I §7) 给出
\[
  \sigma_p \;\propto\; \sigma_{\mathrm{wire}} \cdot |\vec p| / (qBL^2).
\]
当前 $\sigma_p \approx 7$ MeV/$c$ at $\sigma_{\mathrm{wire}} \approx 0.2$ mm. 线性外推:

\begin{table}[H]
\centering\small
\caption{$\sigma_p$ vs.\ $\sigma_{\mathrm{wire}}$ 线性外推.}
\label{tab:partIII:ws1_extrapolate}
\begin{tabular}{rr}
\toprule
$\sigma_{\mathrm{wire}}$ (mm) & $\sigma_p$ (MeV/$c$, 外推) \\
\midrule
0.1 & 3.5 \\
0.3 & 10.5 \\
0.5 & 17.5 \\
1.0 & 35.0 \\
2.0 & 70.0 \\
\bottomrule
\end{tabular}
\end{table}

状态报告 §理论分辨极限 给的 0.5 mm 与 2.0 mm 端点应与上表一致 (待状态报告交叉核对).

\paragraph{待运行命令.}
% TODO: sigma_wire scan run WS-1
\begin{lstlisting}[language=bash,caption={sigma\_wire scan WS-1},label=lst:partIII:ws1_cmd]
for SW in 0.1 0.3 0.5 1.0 2.0; do
  ENSEMBLE_SIZE=50 PROFILE_PER_QUARTILE=8 MCMC_PER_QUARTILE=16 \
      PDC_SIGMA_WIRE_MM=${SW} \
      OUT_ROOT=build/test_output/scan_ws1_sw${SW//./p} \
      bash scripts/analysis/run_ensemble_coverage.sh
done
\end{lstlisting}
该脚本要求 \texttt{run\_ensemble\_coverage.sh} 解析 \texttt{PDC\_SIGMA\_WIRE\_MM} 环境变量并传给 ensemble macro; 当前没有此 hook (% TODO: 在 macro/ensemble template 中加 \\setSigmaWireMm 命令).

\paragraph{期望耗时.} 单 step $\sim 40$ min, $\times 5 = 200$ min ($\sim 3.5$ h).
\paragraph{期望输出.} \texttt{build/test\_output/scan\_ws1\_sw\{0p1,0p3,0p5,1p0,2p0\}/}.

\subsection{扫描 TS-1: 靶倾角 $\alpha_{\mathrm{tgt}}$}
\label{sec:partIII:scans:ts1}

\paragraph{物理动机.}
2026-04-19 修复后, 重建端加入 $R_y(+\alpha_{\mathrm{tgt}})$ (\texttt{PDCFrameRotation.hh}) 把 lab 系结果旋回靶系. 现行配置 $\alpha_{\mathrm{tgt}} = 3°$. 我们要验证:
\begin{enumerate}[label=(\alph*), nosep]
  \item $\alpha_{\mathrm{tgt}} = 0°$ (无旋转): 重建端 frame fix 不应改变物理结果, 即 reco 仍然在 lab = target frame. $p_x$ 残差应仍 $\sim 0$.
  \item $\alpha_{\mathrm{tgt}} = 5°$: $-p_z \sin(5°) \approx -55$ MeV/$c$ 是 truth-vs-reco 的伪偏差, 但 frame fix 把它消除, 残差仍 $\sim 0$.
\end{enumerate}
期待: 三种角度下 $p_x$ 残差半宽不变 (frame 旋转是 \emph{坐标变换}, 无物理影响), 这是 sanity check.

\paragraph{期望趋势.}
\begin{itemize}[nosep]
  \item $\alpha_{\mathrm{tgt}} = 0°, 3°, 5°$: $p_x, p_y, p_z, |\vec p|$ 三种覆盖率应 \emph{完全一致} (在 ensemble 噪声范围内).
  \item 如果 $\alpha_{\mathrm{tgt}} = 5°$ 比 $0°$ 给出不同覆盖率, 说明 PDCFrameRotation 实施有 bug (例如旋转方向错, 或者旋转应用到了部分 PDC1/PDC2 而非全 reco).
\end{itemize}

\paragraph{待运行命令.}
% TODO: tilt scan run TS-1
\begin{lstlisting}[language=bash,caption={tilt scan TS-1},label=lst:partIII:ts1_cmd]
for ALPHA in 0.0 3.0 5.0; do
  ENSEMBLE_SIZE=50 PROFILE_PER_QUARTILE=8 MCMC_PER_QUARTILE=16 \
      TARGET_TILT_DEG=${ALPHA} \
      OUT_ROOT=build/test_output/scan_ts1_alpha${ALPHA//./p} \
      bash scripts/analysis/run_ensemble_coverage.sh
done
\end{lstlisting}

\paragraph{期望耗时.} $3 \times 40 = 120$ min.
\paragraph{期望输出.} 三个 ensemble dir, 比对 \texttt{coverage\_with\_ci.csv} 的 $p_x$ 行.

\subsection{扫描 ES-1: ensemble 大小 $N$}
\label{sec:partIII:scans:es1}

\paragraph{物理动机.}
当前每 truth 点 $N = 50$, 总 $N = 744$ (15 truth bin); 每点 CP $\pm 0.13$ 略嫌粗. 验证生产应取 $N = 200$/truth ($\sim 3000$ 总).

\paragraph{期望趋势.}
覆盖率本身应不变 (ensemble 是 i.i.d. 抽样), 但 CP CI 半宽按 $1 / \sqrt{N}$ 收窄:

\begin{table}[H]
\centering\small
\caption{Ensemble 大小与 CP 半宽 (per truth bin, $\hat p = 0.68$).}
\label{tab:partIII:es1_extrapolate}
\begin{tabular}{rrr}
\toprule
$N$/truth & 总 $N$ & CP 95\% 半宽 \\
\midrule
10  & 150  & $\pm 0.30$ \\
50  & 750  & $\pm 0.13$ (当前) \\
200 & 3000 & $\pm 0.07$ \\
1000 & 15000 & $\pm 0.03$ \\
\bottomrule
\end{tabular}
\end{table}

\paragraph{待运行命令.}
% TODO: ensemble-size scaling run ES-1
\begin{lstlisting}[language=bash,caption={ensemble size scan ES-1},label=lst:partIII:es1_cmd]
for N in 10 50 200 1000; do
  ENSEMBLE_SIZE=${N} PROFILE_PER_QUARTILE=8 MCMC_PER_QUARTILE=16 \
      OUT_ROOT=build/test_output/scan_es1_n${N} \
      bash scripts/analysis/run_ensemble_coverage.sh
done
\end{lstlisting}

\paragraph{期望耗时.} $N=10$ 单核 8 min, $N=50$ 40 min (already done), $N=200$ 160 min, $N=1000$ 800 min ($\sim 13$ h). 应在多核 8 cores 上并行.

\paragraph{期望输出.}
\texttt{build/test\_output/scan\_es1\_n\{10,50,200,1000\}/}; 用同一个 \texttt{analyze\_ensemble\_coverage.py} 做 CP CI 对比.

\subsection{扫描总结}
\label{sec:partIII:scans:summary}

\begin{table}[H]
\centering\small
\caption{扫描研究优先级与状态.}
\label{tab:partIII:scan_priority}
\begin{tabular}{lllll}
\toprule
ID & 扫描 & 优先级 & 状态 & 期望产出 \\
\midrule
BS-1 & $B_y \pm 5\%$        & 中 & 待运行 & 谱仪场强灵敏度 \\
WS-1 & $\sigma_{\mathrm{wire}}$ 0.1--2.0 mm & 高 & 待运行 (需 macro hook) & Glückstern 验证 \\
TS-1 & $\alpha_{\mathrm{tgt}}$ 0--5°  & 中 & 待运行 (需 macro hook) & frame fix sanity \\
ES-1 & $N$ 10--1000           & 高 & 待运行 (无 hook 改动) & 生产 ensemble 标定 \\
\bottomrule
\end{tabular}
\end{table}

% =============================================================

\section{LM 收敛盆地研究: $(-100, 0)$ 病态点剖析}
\label{sec:partIII:basin}

状态报告 §误差分析 在 G4 涨落研究里发现, $(p_x, p_y) = (-100, 0)$ truth 点的 G4 dispersion / Fisher $\sigma$ 比例在 $p_x$ 方向达到 \textbf{3.51}, 在 $p_z$ 方向达到 \textbf{3.51} (同源, 因为是同一组 LM 失败事件主导). 这远高于其他 14 个 truth 点的 $0.5 \sim 1.5$ 范围. 本章详细诊断为什么这个特定 truth 点是 LM 收敛盆地的失败热点, 并提出两条修正路径.

\subsection{从 ensemble 数据直接读取 $(-100, 0)$ 的失败统计}
\label{sec:partIII:basin:numbers}

\begin{table}[H]
\centering\small
\caption{$(p_x, p_y) = (-100, 0)$ truth 点的 50 条事件残差统计 (从 \texttt{ensemble\_pdc\_reco\_merged.csv} 直接计算).}
\label{tab:partIII:basin_pathology_stats}
\begin{tabular}{lr}
\toprule
统计量 & 值 \\
\midrule
$N$ 总事件 & 50 \\
$N$ 残差 $|\Delta p_x| > 50$ MeV/$c$ & 10 (20\%) \\
$N$ 残差 $|\Delta p_x| > 200$ MeV/$c$ & 6 (12\%) \\
$N$ $\hat p_x < -300$ MeV/$c$ & 6 (12\%) \\
失败事件中 $\hat p_x$ 范围 & $[-614, -350]$ \\
失败事件 $\chi^2_{\mathrm{red}}$ 范围 & $[583, 1244]$ \\
失败事件 $\chi^2_{\mathrm{red}}$ 中位 & $\sim 788$ \\
$\Delta p_x$ 中位 (含失败) & $-1.2$ MeV/$c$ \\
$\Delta p_x$ 均值 (含失败) & $-51.3$ MeV/$c$ (失败拖拽) \\
$\Delta p_x$ 标准差 (含失败) & $124.9$ MeV/$c$ \\
$\Delta p_x$ p84 - p16 半宽 (鲁棒) & $\sim 7$ MeV/$c$ \\
\bottomrule
\end{tabular}
\end{table}

\paragraph{诠释.}
\begin{itemize}[nosep]
  \item 失败事件占 12\% — 远高于其他 truth 点 (0--4\%).
  \item 失败事件 $\hat p_x \in [-614, -350]$, 即 LM 把 truth $-100$ 误认成 $\sim -500$ 这个量级. 残差 $\sim -400$ MeV/$c$ 是 truth 的 4 倍.
  \item 中位残差仍 $\sim 0$ — 主峰是好的, 失败事件是 \emph{拖尾}.
  \item Fisher $\sigma$ 在主峰 $\sim 7$ MeV/$c$ 但失败事件 $\sigma$ 跳到 $\sim 60$ — Fisher 在错盆地"刻舟求剑".
\end{itemize}

\subsection{为什么是 $(-100, 0)$?}
\label{sec:partIII:basin:why}

\paragraph{网格搜索的几何.}
\texttt{PDCRkAnalysisInternal.cc} 第 392--424 行的 grid search 用
\begin{itemize}[nosep]
  \item $u$ (= $p_x / p_z$ slope at target) 网格: $u \in [u_{\mathrm{line}} - 1, u_{\mathrm{line}} + 1]$, 7 步, 步长 $\sim 0.33$;
  \item $v$ (= $p_y / p_z$): $v \in [v_{\mathrm{line}} - 0.3, v_{\mathrm{line}} + 0.3]$, 3 步, 步长 $0.3$;
  \item $|\vec p|$: $\{200, 400, 700, 1200, 2000\}$ MeV/$c$ 5 个 logspace-like 候选.
\end{itemize}
其中 $u_{\mathrm{line}}, v_{\mathrm{line}}$ 是 \emph{靶到最远 PDC 的直线方向} (\texttt{init\_dir}), 即不考虑磁场偏转的几何视线方向.

\paragraph{$(-100, 0)$ 的 $u_{\mathrm{line}}$.}
truth $p = (-100, 0, 627)$, 出射方向 $u = -100/627 \approx -0.160$. 但磁场把质子在 PDC2 处弯了 $\sim 25°$ ($\theta_{xz} \sim 0.43$ rad), PDC2 真实位置在 $u \approx -0.43 + (-0.16) = -0.59$ 附近. 因此粗网格中心 $u_{\mathrm{line}} \approx -0.59$ (由 PDC2 几何定义, 不是 truth 出射 slope).

实际计算: PDC2 在 $z \approx 4655$ mm 处, 该事件 PDC2 命中 $x \approx -2745$ mm (由表 \nameref{sec:partIII:basin:numbers} 隐含); 直线 slope $u_{\mathrm{line}} = x_{\mathrm{PDC2}} / z_{\mathrm{PDC2}} \approx -2745 / 4655 \approx -0.59$. 网格 $u \in [-1.59, 0.41]$ —— truth $u = -0.16$ 在网格中心偏右 ($+0.43$ 处, 第 4 个网格点).

\paragraph{为什么是边界?}
truth $u = -0.16$ 不是在 $u_{\mathrm{line}} \pm 1$ 网格的 \emph{物理边界}, 但它\emph{接近 chi-square 第二低盆地} 的边界. 状态报告 $§误差分析$ 中 G4 涨落的 5$\times$3 truth 网格图(图 \texttt{figures/g4\_per\_truth\_box\_*.png}) 显示在 $(-100, 0)$ 这个点 $\hat p_x$ 分布是双峰: 主峰 $\hat p_x \approx -100$ 包含 88\% 事件, 副峰 $\hat p_x \approx -500$ 包含 12\% 事件. 副峰的位置对应 LM "卡在错盆地" 的 chi-square 局部极小.

\paragraph{为什么 $(-100, 20)$ 没那么严重?}
truth $(p_x, p_y) = (-100, 20)$ 的 50 事件中 $\Delta p_x$ 半宽 $\sim 8.6$ MeV/$c$ (从 \texttt{per\_truth\_point\_g4\_vs\_fisher.csv}: g4\_halfw68\_px = 8.64), 比 $(-100, 0)$ 的 $34.8$ 要小 4 倍. 这说明 $p_y \ne 0$ 引入的 $v = p_y / p_z$ 偏离对 LM "推开错盆地" 起到了关键作用 —— 网格 $v$ 步长 0.3 在 $p_y = 0$ 时正好覆盖错盆地, 在 $p_y = \pm 20$ 时把网格推到错盆地外, LM 看不到副峰.

\subsection{Chi-square 表面拓扑示意}
\label{sec:partIII:basin:chi2_topology}

% TODO: 需要画一张 (u, p) 平面 chi-square 等高线图,
% 由 scripts/analysis/plot_chi2_landscape.py 生成 (待写).
% 输入: PDC1, PDC2, target geometry; truth (p_x, p_y, p_z) = (-100, 0, 627).
% 输出: figures/diagnostics/chi2_landscape_p100_p0.png.
% 应能直观看到两个谷(主谷 $u \sim -0.16$, 副谷 $u \sim -0.7 \sim -1.0$)
% 以及它们之间的鞍点, 标注网格采样点.

期望画面:
\begin{itemize}[nosep]
  \item 主谷 (truth 谷): $(u, p) \sim (-0.16, 627)$, $\chi^2 \sim 0$;
  \item 副谷 (病态谷): $(u, p) \sim (-0.85, 1145)$, $\chi^2 \sim 1244$ — 注意 $\chi^2$ 不为零 \emph{是因为} 该副谷是 chi-square 非零的局部极小, 不是真极小;
  \item 鞍点连接两谷, $u_{\mathrm{saddle}} \sim -0.5$;
  \item 网格点 7$\times$5 = 35 个候选 ($v$ 维度暂不画), 7 个 $u$ 候选中 4 个在主谷一侧, 3 个在副谷一侧.
\end{itemize}

\subsection{修复方案 (a): 网格不对称延展}
\label{sec:partIII:basin:fix_a}

\paragraph{思路.}
$|u_{\mathrm{line}}|$ 大时, 主谷 $u_{\mathrm{truth}}$ 与 $u_{\mathrm{line}}$ 距离 $\sim 0.4$; 副谷 $u_{\mathrm{副}}$ 与 $u_{\mathrm{line}}$ 距离也 $\sim 0.3$ (在另一侧). 当前网格 $\pm 1$ 对称延展, 副谷一定能命中. 改成 \emph{不对称} 延展, 让网格偏向 truth 谷:
\[
  u_{\min} = u_{\mathrm{line}} - 0.3, \quad u_{\max} = u_{\mathrm{line}} + 1.0.
\]
即把网格 \emph{向更靠近 truth 出射方向} 偏置 (更小的 $|u|$). 但这要求知道 truth 出射方向的先验, 不可行.

\paragraph{改进思路.}
基于物理: 弯曲方向永远把 $u_{\mathrm{line}}$ 推得比 $u_{\mathrm{truth}}$ \emph{绝对值更大} (磁场把出射 deflect 到更陡的弧). 因此网格 \emph{向 $|u|$ 更小的方向} 偏置:
\[
  u_{\min} = u_{\mathrm{line}} - 0.3 \cdot \mathrm{sign}(u_{\mathrm{line}}), \quad u_{\max} = u_{\mathrm{line}} + 1.0 \cdot \mathrm{sign}(u_{\mathrm{line}}) \cdot (-1).
\]
等价地: 网格半幅 $\pm 1$, 但 truth-side 半幅 $-1.0$, 远 truth-side 半幅 $+0.3$.

\paragraph{代码改动.}
\texttt{libs/analysis\_pdc\_reco/src/PDCRkAnalysisInternal.cc} 第 392--410 行:
\begin{lstlisting}[caption={网格不对称延展 (建议改动)}]
constexpr double kUFarHalfRange  = 1.0;   // away from target
constexpr double kUNearHalfRange = 0.3;   // toward target slope
double u_lo = u_center - (u_center >= 0 ? kUNearHalfRange : kUFarHalfRange);
double u_hi = u_center + (u_center >= 0 ? kUFarHalfRange  : kUNearHalfRange);
for (int iu = 0; iu < kGridU; ++iu) {
    const double u = u_lo + (u_hi - u_lo) * iu / (kGridU - 1);
    ...
}
\end{lstlisting}

\paragraph{预期效果.}
$(-100, 0)$ truth 点 $u_{\mathrm{line}} \approx -0.59$, 改后网格 $u \in [-0.89, +0.41]$ (truth $-0.16$ 在中心偏右). 副谷在 $u \approx -0.85$ 仍在网格内, 但 truth 谷被 4 个网格点覆盖, 副谷只 1 个. 正确盆地会以多起点压倒错盆地.

\subsection{修复方案 (b): 加入 "magnetic-kick-corrected" 起点}
\label{sec:partIII:basin:fix_b}

\paragraph{思路.}
Linear backward 起点: 用 PDC1, PDC2 做反向直线外推到 target $z$, 得到 $u_{\mathrm{back}} = (x_{\mathrm{PDC1}} - x_{\mathrm{tgt}}) / (z_{\mathrm{PDC1}} - z_{\mathrm{tgt}})$. 这忽略磁场, 但可以从这个起点开始一次 $|\vec p|$ 灵敏度估计:

设质子在 dipole 中弯过弧长 $L$, 弯曲角 $\theta = qBL / p$ (Highland-style). 用 PDC2 处的 incidence angle $\theta_{\mathrm{in}}$ 减去 $\theta$ 即得 emission angle, 再外推到 target 给出 \emph{magnetic-kick corrected} $u_{\mathrm{corr}}$.

实际计算更简单: PDC1 和 PDC2 的角度差 $\Delta\theta_{\mathrm{PDCs}} = \mathrm{atan2}(\Delta x, \Delta z)$ 给出 \emph{出 PDC 后} 的轨迹方向; 把它向 target 外推 (无场), 得到的 slope 就是 \emph{出射的} $u$, 再加上一个磁场修正项 $\sim qB\Delta z / p$, 给出 \emph{在 target 的} $u$.

\paragraph{代码改动.}
新加一个候选 $u$ 候选 (替代或补充网格中心):
\begin{lstlisting}[caption={magnetic-kick-corrected 起点 (建议改动)}]
double u_pdc_to_pdc = (track.pdc2.X() - track.pdc1.X()) /
                      (track.pdc2.Z() - track.pdc1.Z());
double dz_pdc1_target = track.pdc1.Z() - target.target_position.Z();
// magnetic kick estimate (assume qBy = -0.5 T, p = 600 MeV/c => 0.4 rad/m)
double magnetic_correction = 0.5 * dz_pdc1_target / 1500.0;  // approx
double u_back = u_pdc_to_pdc - magnetic_correction;

// add as extra candidate
candidates.push_back({{u_back, v_back, 600.0}, Evaluate(...).chi2_raw});
\end{lstlisting}

\paragraph{预期效果.}
$(-100, 0)$ 的 $u_{\mathrm{back}}$ 接近 $-0.16$ (因为 PDC1, PDC2 之间的 slope 反映出 $\hat p$ 而非 $u_{\mathrm{line}}$); 加进 candidates 后, LM 多起点必有一个落到主谷.

\subsection{第三种方案 (c): 一阶 backward Kalman}
\label{sec:partIII:basin:fix_c}

更彻底的方案: 用 Kalman 反向滤波 — 从 PDC2 倒推到 PDC1 倒推到 target, 每步用 RK4 反向积分, 得到一个相对精确的 target-frame 起点. 但这基本等价于做\emph{第二次拟合}, 工作量与替换 LM 等同, 暂不优先.

\subsection{修复后预期 ratio}
\label{sec:partIII:basin:fix_expected}

\begin{table}[H]
\centering\small
\caption{修复方案对 $(-100, 0)$ truth 点 $\Delta p_x$ 半宽的预期影响.}
\label{tab:partIII:basin_fix_table}
\begin{tabular}{lr}
\toprule
方案 & 预期 $\Delta p_x$ 半宽 (MeV/$c$) \\
\midrule
当前 & 34.8 (主峰 $\sim 5$ + 12\% 失败 $\sim 400$) \\
方案 (a) 网格不对称 & $\sim 8$ (失败率应降到 0--2\%) \\
方案 (b) magnetic-kick 起点 & $\sim 5$ (失败率 $\sim 0$, 但会引入新 LM-trap 风险) \\
方案 (a) + (b) 联合 & $\sim 5$ (主峰 + 极少失败) \\
\bottomrule
\end{tabular}
\end{table}

\subsection{副谷的物理来源: 反方向轨迹}
\label{sec:partIII:basin:reverse_trajectory}

副谷 $\hat p_x \approx -500$ 是什么物理? 一种解释是: 给定 PDC1, PDC2 命中, 存在两个不同的 (target 出射, 总动量) 配置都能解释这两个点 — 一个是 truth ($\hat p_x = -100$, $\hat p_z = 627$), 另一个是高动量反向 ($\hat p_x = -500$, $\hat p_z = 940$).

\paragraph{推导.}
质子在磁场 $B_y$ 中, 弯曲半径 $\rho = p / (q B)$. 给定 PDC1, PDC2 两点和 target, 拟合即解三圆心方程 — 圆经过三个点. 但 \emph{磁场不均匀} (SAMURAI dipole 有 fringe field), 所以"圆经过三点"不严格; 不同 $|\vec p|$ 的 RK 积分可以经过相同的 PDC2 但 source 出射方向不同.

具体: 设 truth 解给出 $u_{\mathrm{truth}} = -0.16$ at target. 副谷解给出 $u_{\mathrm{副}} \approx -0.85$ at target — 方向更陡, 但 \emph{动量更高} 让弯曲变浅, 弯过 dipole 后到达 PDC2 的位置接近. 这是磁场+动量的二维耦合带来的 chi-square 多解.

\paragraph{$p_y = 0$ 的特殊性.}
为什么 $p_y \ne 0$ 时副谷消失? 因为 $v = p_y/p_z$ 为非零时, $v$ 网格 $\{v_{\mathrm{line}} - 0.3, v_{\mathrm{line}}, v_{\mathrm{line}} + 0.3\}$ 三个候选与 truth $v$ 距离都不同; 副谷需要特定的 $(u_{\mathrm{副}}, v_{\mathrm{副}})$ 组合, $v_{\mathrm{副}} \ne 0$ 时网格不再覆盖. 在 $p_y = 0$, truth $v = 0$, $v_{\mathrm{line}} = 0$, 网格中央点正好是 truth $v$ — 但同时也在副谷的 $v$ 上 (因为副谷沿 $u$ 维度延伸, $v$ 中心仍 $\sim 0$).

\subsection{为什么 $-100, 0$ 但不是 $+100, 0$?}
\label{sec:partIII:basin:asymmetry}

表 \ref{tab:partIII:other_truth_points} 显示 $(+100, 0)$ ratio 0.57 (健康). 不对称性来自:

\begin{itemize}[nosep]
  \item SAMURAI dipole $B_y$ 极性把正电荷向 $+x$ 方向偏, 即 deuteron / proton 的 $u_{\mathrm{line}}$ 偏 $+u$. truth $p_x = -100$ 的事件出射朝 $-x$, 但 magnetic kick 把它\emph{反向}弯到 $+x$ — 即 PDC2 命中位置在 $+x$ 处, $u_{\mathrm{line}} > 0$. 但实测 (\texttt{pdc\_dispersion\_summary.csv}) 显示 $(p_x = -100, p_y = 0)$ 的 PDC2 中位 $x \approx -3500$ mm, 仍在 $-x$;
  \item 这就需要重新检查: $u_{\mathrm{line}} = -3500/4655 \approx -0.75$, 网格 $u \in [-1.75, 0.25]$. truth $u = -0.16$ 在网格中点偏右. 副谷 $u_{\mathrm{副}} \approx -0.85$ 在网格中央. 即副谷正好被网格"瞄准";
  \item truth $p_x = +100$: 出射朝 $+x$, magnetic kick 反向弯, PDC2 命中 $\sim -2900$ mm (从 \texttt{pdc\_dispersion\_summary.csv} 推断), $u_{\mathrm{line}} \approx -0.62$, 网格 $u \in [-1.62, 0.38]$, truth $u = +0.16$ 在网格右端边缘 — \emph{副谷} 在 $u \approx -0.6$ 也在网格中央, 但 truth $u$ 处于边缘也意味着 LM 容易跑出主谷, 走向副谷的几率反而小. 经验上 ratio 0.57 是健康的;
  \item 物理结论: 网格设计在 \emph{$u_{\mathrm{line}}$ 与 $u_{\mathrm{truth}}$ 同号} 时(主谷与副谷都在网格内, 副谷更靠近中心) 容易失败.
\end{itemize}

\subsection{监控诊断: LM trace 持久化}
\label{sec:partIII:basin:lm_trace}

为以后追踪 LM 失败, 应在 \texttt{PDCErrorAnalysis.cc} 加 \texttt{--dump-lm-trace} 选项, 把每条事件的:
\begin{itemize}[nosep]
  \item 网格搜索的 105 个候选 $(u, v, p, \chi^2)$;
  \item 多起点 LM 的 $\le 5$ 条迭代历史 $(\mathrm{iter}, \lambda, \chi^2, u, v, p)$;
  \item 最终选取的最优点;
\end{itemize}
持久化到 \texttt{lm\_traces/event\_\{N\}.csv}. 这给后续 visualization 提供数据基础, 也允许在数据 (815 tracks, no truth) 上识别 LM 失败 (无 truth 时只能看 $\chi^2$ + LM trace, 不能看残差).

% TODO: PDCErrorAnalysis.cc 加 --dump-lm-trace 选项 (作 dev item).

\subsection{其他可疑 truth 点}
\label{sec:partIII:basin:other_points}

从 \texttt{per\_truth\_point\_g4\_vs\_fisher.csv} 直接读出 ratio (g4 / fisher) :

\begin{table}[H]
\centering\small
\caption{各 truth 点的 g4 半宽 / Fisher $\sigma$ 比例 (从 \texttt{per\_truth\_point\_g4\_vs\_fisher.csv}, $p_x$ 列). $p_y \in \{-20, 0, 20\}$ 对应三行.}
\label{tab:partIII:other_truth_points}
\begin{tabular}{rrrr}
\toprule
truth $p_x$ & ratio $p_y = -20$ & ratio $p_y = 0$ & ratio $p_y = 20$ \\
\midrule
$-100$ & 0.48 & \textbf{3.51} & 0.77 \\
$-50$  & 0.40 & 0.38 & 0.66 \\
$0$    & 0.84 & 0.63 & 0.89 \\
$+50$  & 0.51 & 0.56 & 0.72 \\
$+100$ & 0.62 & 0.57 & 0.96 \\
\bottomrule
\end{tabular}
\end{table}

\paragraph{读法.}
\begin{itemize}[nosep]
  \item $(-100, 0)$ 是唯一 ratio $> 1$ 的 truth 点; 失败率 12\% 是孤立病态.
  \item 其他 14 个 truth 点 ratio 都在 $0.4$--$1.0$ 范围, 基本健康.
  \item ratio $< 1$ 反映 \emph{Fisher 略保守}, 这是 PDC 高斯分辨条件下的预期; 主要担忧是 $> 1$ 的方向.
\end{itemize}

\subsection{小结}
\label{sec:partIII:basin:summary}

\begin{itemize}[nosep]
  \item $(-100, 0)$ 是 LM 收敛盆地失败的孤立热点, 12\% 事件 trap 在 $\hat p_x \approx -500$ 副谷.
  \item 物理原因: $u_{\mathrm{line}}$ 距 $u_{\mathrm{truth}}$ 远, 网格对称延展正好覆盖错盆地; $p_y = 0$ 时 $v$ 网格不能推开错盆地.
  \item 修复: (a) 网格不对称延展 + (b) magnetic-kick-corrected 起点, 联合预期把 ratio 从 3.51 降到 $< 1$.
  \item LM trace 持久化是接下来追踪 LM 病态的高优先级 dev item.
\end{itemize}

% =============================================================

\subsection*{Part III 全章总结}

本部分给出了 RK 误差分析的实证诊断完整框架, 从二项 CI 的统计学基础到事件级 LM trace 的具体追踪.

\textbf{六章主要结论.}
\begin{enumerate}[label=\arabic*., nosep]
  \item \textbf{覆盖率方法学}: Clopper-Pearson 是当前 default, 比 Wilson 略保守, 比 Wald 在端点正确; $N=744$ 给 CP $\pm 0.034$ 足以诊断当前的 $p_y$ 欠覆盖, 生产 ensemble 应取 $N \ge 200$/truth.
  \item \textbf{Pull 分布}: 实测 744 事件 pull 中位 $\sim 0$ (frame fix 后无偏), $p_y$ 主峰 $\sigma_z \approx 2.48$ (Fisher 低估 $\sim 1.86\times$), $p_x, p_z$ 主峰窄但有 LM 长尾干扰; 鲁棒尺度估计 IQR/1.349 是接下来要加的指标.
  \item \textbf{$\chi^2$ 分布与失败模式}: 50 条 ($6.7\%$) $\chi^2_{\mathrm{red}} > 10$ 远超 $\chi^2_1$ 自然涨落 ($0.16\%$, $42\times$); 三类机制 (LM 错盆地, 反常 hit, MS 爆发); $\chi^2 < 5$ quality cut 应在生产分析里默认开启.
  \item \textbf{6 个示例事件}: 涵盖 4 种主要机制. 事件 102 $z_{p_y} = -8.94$ 是 (u, v, p) 偏典型, 事件 121, 229 是 (-100, 0) 病态典型, 事件 742 是 LM 错盆地中等表现.
  \item \textbf{扫描研究方案}: BS-1 (B 场), WS-1 ($\sigma_{\mathrm{wire}}$), TS-1 (靶倾角 sanity), ES-1 (ensemble 规模) 已列出命令模板与外推; ES-1 优先级最高.
  \item \textbf{$(-100, 0)$ 病态点剖析}: 12\% 事件 trap 在副谷, 修复方案 (a) 网格不对称延展 + (b) magnetic-kick-corrected 起点, 联合预期把 ratio 从 3.51 降到 $< 1$.
\end{enumerate}

\textbf{下一阶段优先 dev items.}
\begin{itemize}[nosep]
  \item LM trace 持久化 (\texttt{PDCErrorAnalysis.cc} \texttt{--dump-lm-trace} 选项);
  \item Pull 直方图绘图脚本 + IQR/MAD 鲁棒尺度;
  \item Ensemble script 加 \texttt{PDC\_SIGMA\_WIRE\_MM}, \texttt{TARGET\_TILT\_DEG}, \texttt{MAGNETIC\_FIELD} 环境变量 hook;
  \item \texttt{analyze\_ensemble\_coverage.py} 加 \texttt{--chi2-max} 选项支持 quality cut;
  \item 网格搜索不对称延展 + magnetic-kick 起点的 RK 实施 + ensemble 验证.
\end{itemize}

\textbf{与 Part I, Part II 的连接.}
\begin{itemize}[nosep]
  \item 第 \ref{sec:partIII:coverage_methodology} 章的 CP CI 推导与 Part I §2 (Cramér-Rao) 的 ensemble 验证范式互补;
  \item 第 \ref{sec:partIII:pulls} 章的 pull = $\mathcal{N}(0, 1)$ 期望与 Part I §3 (Fisher) 的协方差一致性紧密相关;
  \item 第 \ref{sec:partIII:chi2} 章的失败模式分类与 Part II §1--9 的误差源对应表完全对齐;
  \item 第 \ref{sec:partIII:case_studies} 章的 6 个示例事件直接验证了 Part II §5 (dE/dx), §8 ((u,v,p) 参数化偏), §9 (LM 收敛盆地);
  \item 第 \ref{sec:partIII:scans} 章的扫描计划是 Part II 各源的 sensitivity 验证;
  \item 第 \ref{sec:partIII:basin} 章 (-100, 0) 剖析是 Part II §9 的具体落地.
\end{itemize}

% =============================================================
% End of Part III.
% =============================================================
 加载。
%
% 不要在此文件中包含 \documentclass / \begin{document} / 序言;
% 主壳层提供 \part{第三部分: 实证诊断} 标题, ctex+xelatex 字体,
% lstlisting 中文样式, booktabs/multirow/amsmath 等宏包。
%
% 创建日期: 2026-04-26
% 作者:    Baiting Tian
% 关联文件:
%   * docs/reports/reconstruction/momentum_reco/rk/rk_reconstruction_status_20260416_zh.tex
%   * docs/superpowers/specs/2026-04-26-rk-error-analysis-deep-dive-design.md
%   * docs/reports/reconstruction/momentum_reco/rk/rk_reconstruction_bugfix_and_error_analysis_20260416.md
%
% 本文件结构 (六章):
%   1. 覆盖率方法学: Clopper-Pearson / Wilson / Wald
%   2. Pull 分布: 定义, 期望, 解读, 实测
%   3. chi^2 分布与失败模式分类
%   4. 5-10 个单事件深度分析
%   5. 扫描研究方案与外推 (多数为 % TODO)
%   6. LM 收敛盆地: (-100, 0) 病态点剖析
%
% 数据来源:
%   * track_summary.csv  (744 事件, 修复后靶系)
%   * profile_samples.csv (128 事件)
%   * bayesian_samples.csv (128 事件)
%   * g4_fluctuation/per_truth_point_g4_vs_fisher.csv
%   * g4_fluctuation/ensemble_pdc_reco_merged.csv (744 事件)
% =============================================================

\section{覆盖率方法学: Clopper-Pearson, Wilson, Wald 三选一}
\label{sec:partIII:coverage_methodology}

整个状态报告 (\StatusReport{}) 把覆盖率作为 RK 误差分析的最终判据: 把名义 68\%/95\% 的方法区间打到 744 条事件上, 数 truth 落在区间内的比例; 这个比例本身又是一个二项随机变量, 必须给一个置信区间才能下结论. 报告里默认采用 Clopper-Pearson 95\% 置信区间(以下简称 CP CI), 本章给出原因, 并把它和教科书里另两种竞争方案 (Wilson 与 Wald) 做对比.

\subsection{二项覆盖率的统计模型}
\label{sec:partIII:coverage_methodology:model}

设我们独立观察 $N$ 个事件, 每个事件的方法区间 $[\hat{\theta}_i - \delta_i^-, \hat{\theta}_i + \delta_i^+]$ 是否覆盖 truth $\theta_i^{\mathrm{truth}}$ 由指示变量
\[
  X_i \;=\; \mathbb{1}\bigl[\hat\theta_i - \delta_i^- \le \theta_i^{\mathrm{truth}} \le \hat\theta_i + \delta_i^+\bigr]
\]
描述. 在 "方法自洽" 的零假设下, $\Pr(X_i = 1) = p_0$ 与名义覆盖水平 (例如 $p_0 = 0.68$) 一致, 且对 $i$ 而言 $X_i$ 独立同分布. 总和 $K = \sum_i X_i$ 服从二项分布
\[
  K \sim \mathrm{Bin}(N, p_0).
\]
观测覆盖率是点估计 $\hat p = K / N$. CI 的目标是: 给定观测 $K = k$, 构造一个区间 $[p_L, p_U]$ 使得在所有真实的 $p$ 上, 该区间覆盖真实 $p$ 的概率不低于 $1 - \alpha$ (例如 $1 - \alpha = 0.95$).

\subsection{Wald 区间(正态近似)与失败模式}
\label{sec:partIII:coverage_methodology:wald}

最常见的写法是 Wald 正态近似:
\[
  p_{L/U}^{\mathrm{Wald}} \;=\; \hat p \;\pm\; z_{\alpha/2}\, \sqrt{\frac{\hat p (1 - \hat p)}{N}}.
\]
教科书把它印在第一页是因为它形式简单, 但它在两端附近(尤其当 $\hat p$ 靠近 $0$ 或 $1$, 或者 $N$ 不够大)有以下三个失败模式:

\begin{enumerate}[label=(\alph*), nosep]
  \item 当 $\hat p = 0$ 或 $1$ 时, 标准误差 $\sqrt{\hat p(1-\hat p)/N} = 0$, 区间退化为单点 $\{0\}$ 或 $\{1\}$. 这显然不能用作 CI ——任何有限 $p$ 都可能产生 $K = 0$ 或 $K = N$.
  \item 当 $N$ 很小 (本章场景: 状态报告中 MCMC 子集 $N = 128$, 早期试运行 $N = 32$), 区间端点可能跑出 $[0, 1]$ 的物理范围.
  \item 实际覆盖率(在所有 $p$ 上对 Wald CI 求频率)远低于名义 $1 - \alpha$. Brown, Cai, DasGupta (2001, \emph{Statistical Science}) 给出了著名的 "wiggle plot": Wald 在 $p$ 略偏离 0.5 时的覆盖率会陷入 0.85 量级的低谷, 远低于名义 0.95.
\end{enumerate}

第三点在状态报告 §误差分析 里直接致命: 修复前 $p_x$ 覆盖率 $\hat p \approx 0.05$ (744 中 38 条命中); Wald 给的区间是
\[
  [0.05 - 1.96\sqrt{0.05\cdot 0.95/744},\ 0.05 + \ldots] \approx [0.034, 0.066],
\]
看起来很窄; 但当 $\hat p$ 离 0 这么近, Wald 显著地低估了不对称性 (真实分布右偏更厉害). 我们要的是一种在 $\hat p \to 0$ 或 $\hat p \to 1$ 极端区域仍然行为正确的 CI.

\subsection{Clopper-Pearson 精确二项 CI}
\label{sec:partIII:coverage_methodology:cp}

Clopper 与 Pearson (1934, \emph{Biometrika}) 提出: 让 $[p_L, p_U]$ 同时满足
\begin{align}
  \Pr\bigl(K \ge k \mid p = p_L\bigr) &\;=\; \alpha/2, \\
  \Pr\bigl(K \le k \mid p = p_U\bigr) &\;=\; \alpha/2.
\end{align}
即 $p_L$ 是 "真实 $p$ 至少为多少时, $K \ge k$ 这件事的概率才能压到 $\alpha/2$" 的最大值; $p_U$ 是 "真实 $p$ 至多为多少时, $K \le k$ 的概率才能压到 $\alpha/2$" 的最小值. 这是一种 "反演 (invert) 二项 CDF" 的 CI 构造, 直觉上等价于把假设检验的接受域翻成参数的置信域.

\paragraph{封闭形式与 Beta-incomplete.}
利用二项 CDF 与不完全 Beta 函数的对偶
\[
  \sum_{j=0}^{k}\binom{N}{j} p^j (1-p)^{N-j} \;=\; I_{1-p}(N-k,\, k+1),
\]
可以把 CP 区间端点写成 Beta 分布分位点:
\begin{align}
  p_L &\;=\; \mathrm{Beta}^{-1}\!\bigl(\alpha/2;\ k,\ N - k + 1\bigr), \\
  p_U &\;=\; \mathrm{Beta}^{-1}\!\bigl(1 - \alpha/2;\ k + 1,\ N - k\bigr).
\end{align}
端点情形: $k = 0 \Rightarrow p_L = 0$, $k = N \Rightarrow p_U = 1$, 这正是物理上希望的 "无下/上界" 行为.

\paragraph{CP 是保守的.}
CP 的 \emph{实际} 覆盖率始终不低于 $1 - \alpha$ (这就是 "精确" 的含义), 但不一定恰等于 $1 - \alpha$. 由于二项分布的离散性, 在大多数 $p$ 值上, 实际覆盖率会略高于 $1 - \alpha$. 这意味着 CP 的区间宽度比名义宽度需要的"刚好" 95\% 略宽一些; 这个保守性恰恰是状态报告需要的: 对每条覆盖率行, "如果 CP CI 不包含 0.68, 那不是统计噪声", 这是一个强论断.

\subsection{Wilson 区间}
\label{sec:partIII:coverage_methodology:wilson}

Wilson (1927) 提出的折中方案是把 Wald 的 $\hat p$ 替换为后验 Beta 调整:
\[
  p_{L/U}^{\mathrm{Wilson}} \;=\;
  \frac{\hat p + \tfrac{z^2}{2N}}{1 + \tfrac{z^2}{N}}
  \;\pm\;
  \frac{z}{1 + \tfrac{z^2}{N}}\,
  \sqrt{\frac{\hat p (1 - \hat p)}{N} + \frac{z^2}{4 N^2}}.
\]
其中 $z = z_{\alpha/2}$ (例如 $1.96$ 对应 95\% CI). 与 Wald 相比, Wilson:
\begin{itemize}[nosep]
  \item 在 $\hat p \in [0.1, 0.9]$ 区间精度极好, 实际覆盖率非常接近名义值;
  \item 区间端点保留在 $[0, 1]$;
  \item 在 $\hat p \to 0/1$ 退化得比 Wald 慢 (但仍比 CP 快).
\end{itemize}

\subsection{$N=128$, $K=80$ 的数值算例}
\label{sec:partIII:coverage_methodology:example}

把 $\hat p \approx 0.625$ 的情形 (例如 MCMC $|\vec p|$ 95\% 覆盖率, $N=128$, $K=115$) 取近似数 $K = 80$ 出来作演算示例. 95\% 双侧:

\begin{table}[H]
\centering\small
\caption{$N=128$, $K=80$ 时三种 CI 的对比 ($\hat p = 0.625$).}
\label{tab:partIII:cp_wilson_wald_demo}
\begin{tabular}{lcccc}
\toprule
方法 & 公式入口 & $p_L$ & $p_U$ & 半宽 \\
\midrule
Wald     & $\hat p \pm 1.96\sqrt{\hat p (1-\hat p)/N}$ & 0.5412 & 0.7088 & 0.0838 \\
Wilson   & 公式 \eqref{eq:partIII:wilson} & 0.5380 & 0.7053 & 0.0837 \\
Clopper-Pearson & Beta 分位点 & 0.5365 & 0.7068 & 0.0852 \\
\bottomrule
\end{tabular}
\end{table}

\begin{equation}
\label{eq:partIII:wilson}
  p_{L/U}^{\mathrm{Wilson}} = \frac{\hat p + z^2/(2N)}{1 + z^2/N}
  \pm \frac{z\sqrt{\hat p(1-\hat p)/N + z^2/(4N^2)}}{1 + z^2/N}.
\end{equation}

% TODO: 用 scripts/analysis/analyze_ensemble_coverage.py 内部 _clopper_pearson 函数
% (它即用 scipy.stats.beta.ppf) 与 statsmodels.stats.proportion 的 wilson_score
% 对全部 4*4 覆盖率行批量出 CSV; 当前数值是手算/样例,
% 需要的运行: ENSEMBLE_SIZE=50 ./scripts/analysis/run_ensemble_coverage.sh
% 已经跑过, 这里仅缺一个三方法对比表生成器.

在 $\hat p \approx 0.625$ 这种 "中部" 区间, 三种方法差距 $< 1\%$; 真正的差距出现在尾部. 对极端 $\hat p \approx 0.05$, $N=744$, $K=38$:

\begin{table}[H]
\centering\small
\caption{$N=744$, $K=38$ 时三种 CI 的对比 ($\hat p = 0.0511$, 修复前 $p_x$ 覆盖率).}
\label{tab:partIII:cp_wilson_wald_extreme}
\begin{tabular}{lccc}
\toprule
方法 & $p_L$ & $p_U$ & 备注 \\
\midrule
Wald     & 0.0353 & 0.0669 & 端点近似对称, 偏窄 \\
Wilson   & 0.0376 & 0.0694 & 内推到 $[0, 1]$ \\
Clopper-Pearson & 0.0365 & 0.0698 & 略宽, 给出真正下界 \\
\bottomrule
\end{tabular}
\end{table}

CP 与 Wilson 的差距不大 (各 $\sim 0.001$), 但都明显比 Wald 长. 既然计算成本相同 ($\sim$ 一次 Beta 分位点计算), 选 CP 是合理的工程选择: 永远不会因为 $\hat p$ 走到极端而退化, 永远是保守的.

\subsection{CP 区间反演的几何直觉}
\label{sec:partIII:coverage_methodology:geometry}

CP 的反演构造可以画一张二维图来理解. 在 $(p, K)$ 平面上(横轴是参数 $p$, 纵轴是观测计数 $K$), 对每个 $p$ 画 $K$ 的双侧 $\alpha/2$ 接受区间 $[K_L(p), K_U(p)]$:
\begin{align}
  K_L(p) &= \min\{k : \Pr(K \ge k \mid p) \le 1 - \alpha/2\}, \\
  K_U(p) &= \max\{k : \Pr(K \le k \mid p) \le 1 - \alpha/2\}.
\end{align}
这两条上下边界形成一个"接受带"; 给定观测 $k$, 把横线 $K = k$ 与该带相交, 横向投影即得到 CP CI $[p_L, p_U]$. 离散的二项分布让边界呈"阶梯状", 这是 CP 略保守的根源 — 阶梯在某些 $p$ 上比连续构造高出 $\sim 1$ 单位.

如下伪代码反映了这一构造的数值实现 (与 \texttt{analyze\_ensemble\_coverage.py} 的 \texttt{\_clopper\_pearson} 函数一致):

\begin{lstlisting}[language=Python,caption={CP CI 的最小数值实现},label=lst:partIII:cp_python]
from scipy import stats

def clopper_pearson(k, n, alpha=0.05):
    """Clopper-Pearson exact binomial CI, two-sided 1-alpha coverage."""
    if k == 0:
        p_lo = 0.0
    else:
        p_lo = stats.beta.ppf(alpha / 2, k, n - k + 1)
    if k == n:
        p_hi = 1.0
    else:
        p_hi = stats.beta.ppf(1 - alpha / 2, k + 1, n - k)
    return p_lo, p_hi
\end{lstlisting}

\paragraph{阶梯效应数值化.}
对 $N=128$, 取 $K=87$ 与 $K=88$ (相邻整数), CP CI 的下端跳变:
\begin{itemize}[nosep]
  \item $K=87$: $p_L = 0.602$;
  \item $K=88$: $p_L = 0.610$.
\end{itemize}
即 $K$ 增 1 时 $p_L$ 跳 $\sim 0.008$. 状态报告里观测覆盖率精度因此有 \emph{量子化下限} $\sim 1/N$, 这是离散二项分布的根本约束, 不是计算误差.

\subsection{多覆盖率水平: 95\% vs.\ 99\%}
\label{sec:partIII:coverage_methodology:multilevel}

状态报告默认 95\% CP CI; 偶尔需要 99\% 来给出"几乎确定" 的结论. 把上面 $\hat p = 0.625$, $N=128$ 的算例扩展:

\begin{table}[H]
\centering\small
\caption{$N=128$, $K=80$ 时不同 $\alpha$ 的 CP CI.}
\label{tab:partIII:cp_multilevel}
\begin{tabular}{lcccc}
\toprule
$1 - \alpha$ & $z_{\alpha/2}$ (Wald 参考) & $p_L$ (CP) & $p_U$ (CP) & 半宽 \\
\midrule
68\% & 1.00 & 0.582 & 0.667 & 0.043 \\
90\% & 1.64 & 0.553 & 0.692 & 0.069 \\
95\% & 1.96 & 0.536 & 0.707 & 0.085 \\
99\% & 2.58 & 0.508 & 0.730 & 0.111 \\
\bottomrule
\end{tabular}
\end{table}

\paragraph{推论.} 99\% 比 95\% 半宽宽 $\sim 1.3\times$. 状态报告的覆盖率结论 ($p_y$ 欠覆盖 $0.42$ vs.\ $0.68$) 用 99\% CP CI 仍然成立 (差距 $0.26 \gg 0.111$). 但 MCMC 的 $|\vec p|$ 0.633 vs.\ 0.68 的差距 0.05 在 95\% CP CI 半宽 0.085 之内, 在 99\% 半宽 0.111 之内 — 这个差异 \emph{不显著}, 这是状态报告 §误差分析 中已经强调的事实.

\subsection{ensemble 数据的独立性假设}
\label{sec:partIII:coverage_methodology:independence}

CP CI 假设 $X_i$ i.i.d.. 现实情况:
\begin{itemize}[nosep]
  \item ensemble 在 15 个 truth bin 上分组 ($N=50$/bin, 总 $N=744$). 同一 truth bin 内的 50 条事件的 \emph{Geant4 输入} 完全相同, 只是随机种子不同 — 这种意义上独立.
  \item 但是, 同一 truth bin 内 LM 失败事件占 12\% (在 $-100, 0$ bin) 或 $0$\% (在其他 bin) — \emph{失败的发生与 truth bin 强相关}, 不是 i.i.d.. 这意味着把全 744 看作单一二项总样本会低估 CP CI.
  \item 更严格的处理: 对每个 truth bin 单独算 CP CI, 然后做 stratified weighted average. 当前 \texttt{analyze\_ensemble\_coverage.py} 是 pooled 处理.
\end{itemize}

\paragraph{stratified analysis 占位.}
% TODO: analyze_ensemble_coverage.py 加 --stratify-by-truth 选项,
% 输出 build/test_output/.../coverage_per_truth_bin.csv,
% 每行 (truth_px, truth_py, N, k, p_hat, cp_lo, cp_hi).

\subsection{为什么不用 bootstrap?}
\label{sec:partIII:coverage_methodology:no_bootstrap}

Bootstrap 在覆盖率比例上是完全可行的: 把 $N$ 个 $X_i$ 重抽样 $B$ 次, 每次算 $\hat p^*_b$, 取经验分位点. 但相对 CP 它有三个劣势:

\begin{itemize}[nosep]
  \item \textbf{计算成本}: $B = 10^4$ 次重抽样 vs.\ 单次 Beta 分位点; 在 60 个覆盖率行 (4 分量 $\times$ 4 方法 $\times$ 多覆盖水平 $\times$ 多 truth bin) 累积起来非微秒.
  \item \textbf{近似性}: bootstrap 的覆盖率本身仍然是渐近 $1 - \alpha$, 在小 $N$ 不"精确".
  \item \textbf{无法重复}: 每次 bootstrap 有随机种子, 报告读者复现时数字会跳; CP 是闭式的.
\end{itemize}

闭式 CP 足够保守, 故选定为状态报告标准.

\subsection{ensemble 大小 $\to$ CP 半宽的换算表}
\label{sec:partIII:coverage_methodology:N_scaling}

当 $\hat p = 0.68$ 时, 二项标准差为 $\sqrt{p(1-p)/N} = \sqrt{0.2176/N}$, Wald 半宽是 $1.96$ 倍此值; CP 半宽与 Wald 的中部值相差 $< 5\%$. 估算各 $N$ 下的 CP 95\% 半宽:

\begin{table}[H]
\centering\small
\caption{$\hat p = 0.68$ 时 $N$ 与 CP 95\% 半宽的近似关系. 列 "Wald 近似" 给 $1.96\sqrt{p(1-p)/N}$, 列 "CP 实际" 由 \texttt{scipy.stats.beta.ppf} 直接计算.}
\label{tab:partIII:N_cp_halfw}
\begin{tabular}{rrr}
\toprule
$N$ & Wald 近似 & CP 实际 \\
\midrule
32   & 0.162 & 0.180 \\
64   & 0.114 & 0.121 \\
128  & 0.081 & 0.084 \\
256  & 0.057 & 0.058 \\
744  & 0.034 & 0.034 \\
5000 & 0.013 & 0.013 \\
\bottomrule
\end{tabular}
\end{table}

\textbf{解读}: 状态报告里 Fisher/Laplace 用 $N=744$, CP 半宽 $\sim 0.034$, 即名义 0.68 与观测 $\sim 0.42$ ($p_y$) 的差距 $0.26$ 远超过 $\pm 0.034$, 是 "真信号"; Profile/MCMC 用 $N=128$, CP 半宽 $\sim 0.084$, 仍能区分 $0.42$ vs $0.68$. 但是, 一旦缩到 $N=32$, 半宽 $\sim 0.18$ 跨过差距, 任何小于 $\sim 18\%$ 的偏离都会埋没在统计噪声里.

\textbf{推论}: 生产用 $N=200$ (CP $\pm 0.07$) 已经够诊断 $\sim 10\%$ 级偏差; 当前 $N=50$ 每 truth 点(总 $N=744$) 只在每个 truth bin 内 CP $\pm 0.13$, 每点诊断略嫌粗. 下文 \nameref{sec:partIII:scans} 里 ES-1 计划把生产 ensemble 增到 $N=200$/truth 点 ($\sim 3000$ 总).

\subsection{小结}
\label{sec:partIII:coverage_methodology:summary}

\begin{itemize}[nosep]
  \item Wald 在 $\hat p \to 0/1$ 退化, 在 $\hat p \in (0.1, 0.9)$ 也轻微低覆盖, 状态报告不用.
  \item Wilson 大多数情形与 CP 接近, 但缺乏 "永远 $\ge 1 - \alpha$" 的精确性保证.
  \item Clopper-Pearson 是默认选项: 闭式, 保守, 在尾部行为正确.
  \item 当前 $N=744$ 给 CP $\pm 0.034$, 足以分辨当前的 $p_y$ 欠覆盖 $0.42$ vs.\ 名义 $0.68$; $N=128$ MCMC/Profile 子集给 CP $\pm 0.084$, 仍可诊断 $\sim 0.1$ 级偏差.
\end{itemize}

% =============================================================

\section{Pull 分布: 定义, 期望, 解读}
\label{sec:partIII:pulls}

覆盖率检验告诉你 "区间宽度对不对", 但它把整个分布折成一个标量. 真正展示 "误差棒形状是否正确" 的诊断量是 \emph{pull}. 本章给出 pull 的形式定义, 在三种缺陷模式下的标志特征, 并直接读取 \texttt{track\_summary.csv} 算出 744 事件的 pull 统计.

\subsection{Pull 定义}
\label{sec:partIII:pulls:definition}

设第 $i$ 条事件的某个动量分量(以 $p_x$ 为例)的真值为 $p_{x,i}^{\mathrm{truth}}$, 重建给出的中心估计为 $\hat p_{x,i}$, 误差分析(本节默认 Fisher) 给出的标准差为 $\sigma_{p_x, i}$. 定义事件级 pull:
\[
  z_{p_x, i} \;=\; \frac{\hat p_{x,i} - p_{x,i}^{\mathrm{truth}}}{\sigma_{p_x, i}}.
\]
注意: pull 并不等于 \emph{标准化残差} 一般意义上的 $\chi^2$ 项 —— 它取的是 \emph{重建估计与真值之差} 除以\emph{该事件的}估计 $\sigma$, 不是平均 $\sigma$. 这意味着每条事件都用自己的方差归一, 是一个真正的 "诊断这条事件" 的量.

\subsection{标准期望}
\label{sec:partIII:pulls:expected}

如果(且仅如果)以下三个条件都成立,
\begin{enumerate}[label=(\alph*), nosep]
  \item 重建估计是无偏的: $\mathbb{E}[\hat p_x | p_x^{\mathrm{truth}}] = p_x^{\mathrm{truth}}$;
  \item 误差分析给出的 $\sigma$ 与真实方差一致: $\mathrm{Var}[\hat p_x | p_x^{\mathrm{truth}}] = \sigma_{p_x}^2$;
  \item 残差是高斯分布的: $\hat p_x | p_x^{\mathrm{truth}} \sim \mathcal{N}(p_x^{\mathrm{truth}}, \sigma_{p_x}^2)$,
\end{enumerate}
那么 pull 服从标准正态:
\[
  z \sim \mathcal{N}(0, 1).
\]
推论:
\begin{itemize}[nosep]
  \item $\mathbb{E}[z] = 0$,
  \item $\mathrm{Var}[z] = 1$ 即 $\sigma_z = 1$,
  \item $\Pr(|z| \le 1) = 0.6827$, $\Pr(|z| \le 1.96) = 0.95$.
\end{itemize}
任何对 $\mathcal{N}(0,1)$ 的偏离都对应一种系统问题. 表 \ref{tab:partIII:pull_diagnostics_template} 给出三类典型缺陷的特征:

\begin{table}[H]
\centering\small
\caption{Pull 分布对应的故障模式.}
\label{tab:partIII:pull_diagnostics_template}
\begin{tabular}{llp{0.45\linewidth}}
\toprule
特征 & 说明 & 故障模式 \\
\midrule
$\bar z \ne 0$ & pull 中心不在零 & 重建有偏: 坐标系/dE/dx/对齐错误的几何位移 \\
$\sigma_z > 1$ & pull 太宽 & $\sigma$ 低估: Fisher 线性化丢了曲率, 或者 (u,v,p) 参数化压抑了某分量 \\
$\sigma_z < 1$ & pull 太窄 & $\sigma$ 高估: Fisher 把不存在的相关性 / 多次散射 / dE/dx 加进去了 \\
\bottomrule
\end{tabular}
\end{table}

\textbf{示例}: 状态报告里修复前的 $p_x$ pull 中位 $\sim -33 / 7.4 \approx -4.4$ 是 (a) 类典型 (frame mismatch); 当前 $p_y$ ratio $1.86$ $\Rightarrow \sigma_z \approx 1.86$ 是 (b) 类典型; 当前 $p_x, p_z$ ratio $\approx 0.8$ $\Rightarrow \sigma_z \approx 0.8$ 是 (c) 类轻症.

\subsection{744 事件 pull 实测}
\label{sec:partIII:pulls:measured}

下表是从 \texttt{track\_summary.csv} 直接计算的 pull 统计 (定义见上, $\sigma$ 取 \texttt{fisher\_*\_sigma} 列):

\begin{table}[H]
\centering\small
\caption{Pull 统计(744 事件, 修复后靶系); $\bar z$ = 均值; $\sigma_z$ = 标准差; $|z|\le 1$ 与 $|z|\le 1.96$ = 实测命中率, 名义为 $0.68$ 与 $0.95$.}
\label{tab:partIII:pull_measured}
\begin{tabular}{lrrrrrl}
\toprule
分量 & $\bar z$ & 中位 $z$ & $\sigma_z$ & $|z|\le 1$ & $|z|\le 1.96$ & 备注 \\
\midrule
$p_x$ & $-0.572$ & $-0.026$ & $4.216$ & $80.1\%$ & $91.0\%$ & 长尾来自 6 条 LM 失败事件 \\
$p_y$ & $-0.178$ & $-0.066$ & $2.477$ & $42.1\%$ & $69.0\%$ & 整体宽 $\sigma_z \approx 2.5$ \\
$p_z$ & $-0.293$ & $-0.068$ & $9.050$ & $77.2\%$ & $90.2\%$ & 长尾更严重(尾部 $|z| > 100$) \\
$|\vec p|$ & $-2.823$ & $-0.066$ & $59.989$ & $76.1\%$ & $90.1\%$ & 同 $p_z$, 长尾事件主导 \\
\bottomrule
\end{tabular}
\end{table}

\paragraph{读法.}
\begin{itemize}[nosep]
  \item \textbf{中位数 $\sim 0$}: 中心估计修复后 (2026-04-19 frame fix) 已经无偏, 没有 frame mismatch 残余. 三分量中位 pull 都 $|z_{\mathrm{med}}| < 0.07$.
  \item \textbf{均值 $\bar z$ 偏负}: 但拖到长尾 (LM 失败事件) 把均值压成 $\sim -0.5$ ($p_x$); $|\vec p|$ 更夸张 $\bar z = -2.8$ 因为 $|\vec p|$ 是\emph{严格非负} 的, LM 跑飞时 $\hat{|\vec p|} \to 1100$ 但 truth $\sim 635$, 这种事件单独贡献 $\Delta z \sim +75$ 但少数; 实际 6 条事件贡献了 95\% 的方差.
  \item \textbf{$\sigma_z (p_y) = 2.48$}: 没有 LM 长尾的"干净"分量, 这个 2.48 跟覆盖率比值 1.86 不完全吻合, 原因是 pull 用的是 \emph{每事件 $\sigma$} 而不是中位 $\sigma$, 而且分布不是高斯; ratio 1.86 是 \emph{鲁棒半宽 / Fisher $\sigma$ 中位}, 数学上不等价, 但定性上同向 (Fisher 低估).
  \item \textbf{$\sigma_z (p_x), \sigma_z (p_z)$ 极大但中位 $|z|\le 1$ 比例 $\sim 0.78$}: 这是 "干净事件占 $93\%$, 6 条 LM 失败事件占 $7\%$" 的双峰分布表现; 鲁棒尺度估计(中位 + IQR) 对这个分布更稳, 也对应表 §误差分析 的覆盖率数字.
\end{itemize}

\paragraph{Pull 的鲁棒尺度估计.}
\label{sec:partIII:pulls:robust}
对 LM 长尾敏感的修正方案: 取 \emph{IQR / 1.349} (高斯下与 $\sigma$ 一致, 对长尾鲁棒) 或者 MAD (median absolute deviation) $\times 1.4826$.

% TODO: 把下表的 IQR/MAD 列也加上,
% 需要在 analyze_ensemble_coverage.py 加 --robust-pull-scale 选项.
\begin{table}[H]
\centering\small
\caption{Pull 的鲁棒尺度估计(待运行, 当前为占位): IQR/1.349 与 $\sigma_z$ 直接对比.}
\label{tab:partIII:pull_robust_scale}
\begin{tabular}{lrrr}
\toprule
分量 & $\sigma_z$ (经典) & IQR/1.349 & 比值 \\
\midrule
$p_x$    & 4.216 & TODO & TODO \\
$p_y$    & 2.477 & TODO & TODO \\
$p_z$    & 9.050 & TODO & TODO \\
$|\vec p|$ & 59.99 & TODO & TODO \\
\bottomrule
\end{tabular}
\end{table}

\subsection{从 CSV 计算 pull 的代码片段}
\label{sec:partIII:pulls:code}

为可复现, 给出从 \texttt{track\_summary.csv} 直接算 pull 的 Python 片段:

\begin{lstlisting}[language=Python,caption={pull 统计的最小复现脚本},label=lst:partIII:pull_python]
import pandas as pd, numpy as np

CSV = ('build/test_output/ensemble_n50_targetframe/'
       'rk_fixed_target_pdc_only_error_analysis/track_summary.csv')
df = pd.read_csv(CSV)

# 残差与 pull
df['truth_p'] = np.sqrt(df.truth_px**2 + df.truth_py**2 + df.truth_pz**2)
for c in ['px', 'py', 'pz', 'p']:
    df[f'pull_{c}'] = (df[f'reco_{c}'] - df[f'truth_{c}']) / df[f'fisher_{c}_sigma']

for c in ['px', 'py', 'pz', 'p']:
    z = df[f'pull_{c}']
    in68 = (z.abs() <= 1.0).mean() * 100
    in95 = (z.abs() <= 1.96).mean() * 100
    print(f'{c}: mean={z.mean():+.3f} median={z.median():+.3f} '
          f'sigma={z.std():.3f}  |z|<=1: {in68:.1f}%  |z|<=1.96: {in95:.1f}%')
\end{lstlisting}

\subsection{Pull 直方图(占位)}
\label{sec:partIII:pulls:histograms}

% TODO: 需要的 PNG, 由 scripts/analysis/plot_pull_histograms.py 生成:
%   figures/diagnostics/pull_distribution_px.png
%   figures/diagnostics/pull_distribution_py.png
%   figures/diagnostics/pull_distribution_pz.png
%   figures/diagnostics/pull_distribution_p.png
% 该脚本未提交; 拟接受 --input track_summary.csv 与 --truth-cut 0.5 (排除 LM 失败事件)
% 双轴显示: 左 全部 744 事件 (含 LM 失败长尾), 右 truth-cut 后干净分布,
% 叠加 N(0,1) PDF.

四张直方图(占位)将分别显示 $p_x, p_y, p_z, |\vec p|$ 的 pull 分布, 与 $\mathcal{N}(0, 1)$ 参考曲线叠加. 期待画面:
\begin{itemize}[nosep]
  \item $p_x$ 分布: 主峰几乎与 $\mathcal{N}(0, 1)$ 重合, 但右上 (LM 失败) 的 6 条事件远远拖出 $|z| > 30$ 的尾巴;
  \item $p_y$ 分布: 主峰宽于 $\mathcal{N}(0, 1)$, $\sigma_z \sim 2.5$;
  \item $p_z$ 分布: 类似 $p_x$ 但长尾更长 (LM 失败时 $p_z$ 跑飞到 $\sim 940$ MeV/$c$);
  \item $|\vec p|$ 分布: 同 $p_z$, 主峰 $\bar z \sim 0$ 但长尾左偏(失败事件 $\hat{|p|} \gg p^{\mathrm{truth}}$).
\end{itemize}

\subsection{每事件 $\sigma$ 与 ensemble 平均 $\sigma$ 的差异}
\label{sec:partIII:pulls:per_event_vs_ensemble}

Pull 定义里 "$\sigma_i$" 是 \emph{每事件的} Fisher $\sigma$, 不是 ensemble 平均. 这有微妙的影响:

\begin{itemize}[nosep]
  \item 如果 Fisher $\sigma_i$ 在事件之间剧烈变化 (例如 $|\vec p|$ Fisher 在 $|p_x|=100$ 时 $\sim 7$, 在病态点 $\sim 60$), 而残差幅度也跟着变大, pull 仍可能保持 $\mathcal{N}(0,1)$;
  \item 反之, 如果 $\sigma_i$ 是常数但残差有 truth-bin 依赖 (例如 dE/dx 在大 $|\vec p|$ 处更大), pull 会有 truth-bin 依赖.
\end{itemize}

换言之, pull 是\emph{标准化残差}的事件级版本. 它对 "Fisher $\sigma$ 是否随事件正确缩放" 是敏感的诊断量.

\paragraph{Per-truth-bin pull 检查.}
% TODO: 把 744 事件按 15 个 truth bin 分组, 每 bin 算 pull mean, std,
% 输出 pull_per_truth_bin.csv. 期望: 14 个非病态 bin 的 pull std 接近 1
% (frame fix 后), $(-100, 0)$ bin 因 12% LM 失败拖出 std $\sim 5$.
% 命令: python3 scripts/analysis/analyze_ensemble_coverage.py \
%         --stratify-by-truth --pull-per-bin

预期结果(从 \nameref{tab:partIII:other_truth_points} ratio 列推算):
\begin{itemize}[nosep]
  \item 14 个非病态 bin: pull std $\sim 0.4$--$1.0$ (大多 Fisher 偏保守, $\sigma_z < 1$);
  \item $(-100, 0)$ bin: pull std $\sim 5$ ($p_x$, $p_z$ 双方向被 LM 失败拖大).
\end{itemize}

\subsection{Pull 与覆盖率的关系}
\label{sec:partIII:pulls:vs_coverage}

二者本质上同源: 一个分量的 68\% 覆盖率 $= \Pr(|z| \le 1)$. 对高斯分布严格等价. 实际数据由于长尾, 覆盖率(对长尾"鲁棒")与 $\sigma_z$ (对长尾敏感) 给出不同信号:

\begin{table}[H]
\centering\small
\caption{覆盖率与 $\sigma_z$ 的对比 (744 事件): 覆盖率对 LM 长尾鲁棒, $\sigma_z$ 不鲁棒.}
\label{tab:partIII:pull_vs_coverage}
\begin{tabular}{lrrr}
\toprule
分量 & 覆盖率 $|z|\le 1$ & $\sigma_z$ & 名义 $\Phi(\sigma_z)$ \\
\midrule
$p_x$ & 0.801 & 4.22 & 0.232 (高斯下) \\
$p_y$ & 0.421 & 2.48 & 0.394 (高斯下) \\
$p_z$ & 0.772 & 9.05 & 0.111 (高斯下) \\
$|\vec p|$ & 0.761 & 60.0 & 0.013 (高斯下) \\
\bottomrule
\end{tabular}
\end{table}

\textbf{解读}: $p_y$ 的 覆盖率 0.42 几乎与 "高斯下 $\sigma_z = 2.48$ 给出的 $\Phi(1/2.48) - \Phi(-1/2.48) = 0.394$" 一致 (差 $\sim 7\%$); 这说明 $p_y$ 主峰是单峰高斯, 但宽度被 Fisher 低估了 $2.48\times$. 反观 $p_x, p_z$, 经典 $\sigma_z$ 由长尾撑大到 $4 \sim 9$, 高斯换算只能给 $0.23 \sim 0.11$ 覆盖率, 但实测 $0.77 \sim 0.80$ —— 差异是因为分布不是高斯, 主峰窄, 长尾稀疏. 鲁棒指标 (中位 + IQR/1.349) 才是 $p_x, p_z$ 的可靠诊断.

\subsection{小结}
\label{sec:partIII:pulls:summary}

\begin{itemize}[nosep]
  \item Pull 是诊断 "误差棒形状是否正确" 的标准量, 期望 $\mathcal{N}(0, 1)$.
  \item 当前 744 事件: 三分量中位 pull $\sim 0$ (frame fix 后无偏), $p_y$ 主峰 $\sigma_z \approx 2.48$ ($\Rightarrow$ Fisher 低估 $\sim 1.86\times$, 与覆盖率 ratio 一致); $p_x, p_z$ 主峰窄但长尾干扰 $\sigma_z$.
  \item 接下来需要: pull 直方图绘图脚本; IQR/MAD 鲁棒尺度; truth-cut 后(去 LM 失败事件) 重新拟合主峰高斯, 应给出与覆盖率一致的 $\sigma_z$.
\end{itemize}

% =============================================================

\section{$\chi^2$ 分布与失败模式分类}
\label{sec:partIII:chi2}

\subsection{自由度 1 的 $\chi^2$ 分布是重尾的}
\label{sec:partIII:chi2:dof1}

RK 重建在 PDC1+PDC2 共两点时, 残差维度 = 4 (两个 PDC 各 $u, v$), 参数维度 = 3 ($u, v, p$ 的 $u$ 是斜率不是丝面坐标), 自由度 $N_{\mathrm{dof}} = 4 - 3 = 1$. 这是状态报告 §\(\chi^2\)~目标函数 写明的事实. 自由度 1 的 $\chi^2$ 分布密度为
\[
  f_{\chi^2_1}(x) \;=\; \frac{1}{\sqrt{2\pi x}} e^{-x/2}, \quad x > 0.
\]
在原点发散($x \to 0$ 时密度 $\to \infty$, 但积分仍收敛), 在尾部以 $e^{-x/2}/\sqrt{x}$ 衰减 —— 比 $\chi^2_2 \sim e^{-x/2}$ 慢. 数值上:

\begin{table}[H]
\centering\small
\caption{$\chi^2_1$ 分布的尾部概率 ($\chi^2_{\mathrm{red}} = \chi^2 / N_{\mathrm{dof}}$, 此处 $N_{\mathrm{dof}} = 1$).}
\label{tab:partIII:chi2_tail_theory}
\begin{tabular}{rrr}
\toprule
$\chi^2_{\mathrm{red}}$ 阈值 & $\Pr(\chi^2_1 > t)$ & $-\log_{10}\Pr$ \\
\midrule
1   & 0.3173 & 0.50 \\
2   & 0.1573 & 0.80 \\
3   & 0.0833 & 1.08 \\
4   & 0.0455 & 1.34 \\
6   & 0.0143 & 1.85 \\
10  & 0.00157 & 2.80 \\
20  & 7.7$\times 10^{-6}$ & 5.11 \\
50  & 1.5$\times 10^{-12}$ & 11.81 \\
\bottomrule
\end{tabular}
\end{table}

\textbf{推论}: 即使在零假设下(模型完美吻合), 1 自由度 $\chi^2_1 > 4$ 的事件仍占 $\sim 4.6\%$, $\chi^2_1 > 10$ 占 $\sim 0.16\%$. 在 744 事件中, 我们 \emph{期望}
\[
  N_{\chi^2 > 4} \approx 33.8, \qquad N_{\chi^2 > 10} \approx 1.17.
\]

\subsection{744 事件实测对比}
\label{sec:partIII:chi2:measured}

直接读 \texttt{track\_summary.csv} 的 \texttt{chi2\_reduced} 列:

\begin{table}[H]
\centering\small
\caption{744 事件的 $\chi^2_{\mathrm{red}}$ 分布观测值 vs.\ 理论 ($\chi^2_1$ 假设). $N_{\mathrm{exp}}$ = 744 $\times$ $\Pr(\chi^2_1 > t)$.}
\label{tab:partIII:chi2_observed}
\begin{tabular}{rrrr}
\toprule
阈值 $t$ & 观测 $N$ & 期望 $N_{\mathrm{exp}}$ & 观测/期望 \\
\midrule
2  & 91 & 117.0 & 0.78 \\
4  & 60 & 33.8  & 1.78 \\
6  & 54 & 10.6  & 5.09 \\
10 & 50 & 1.17  & 42.7 \\
20 & 42 & 0.006 & 7000 \\
\bottomrule
\end{tabular}
\end{table}

\paragraph{读法.}
\begin{itemize}[nosep]
  \item $\chi^2_{\mathrm{red}} > 2$ 出现 91 条 vs.\ 期望 117 —— \emph{低于}期望. 主峰核心拟合得比 1 自由度 $\chi^2_1$ 还紧, 提示 $N_{\mathrm{dof}}$ 可能在 \emph{有效} 上略大于 1 (例如 PDC 测量误差被夸大, 拟合"省下"自由度), 或 LM 把 $\chi^2$ 推得过于贴近 0.
  \item $\chi^2_{\mathrm{red}} > 4$ 之上, 观测开始压倒期望: $> 4$ 已经是 1.78 倍, $> 10$ 是 \textbf{42.7 倍}, $> 20$ 是 7000 倍.
  \item 50 条 $\chi^2_{\mathrm{red}} > 10$ 的事件, 从 7\% 占比看绝对不是 $\chi^2_1$ 自然涨落; 它们是另一个机制产生的 (LM 局部极小, 多次散射爆发, dE/dx 离群, 等等).
\end{itemize}

\subsection{失败模式分类}
\label{sec:partIII:chi2:taxonomy}

把 744 事件按 $\chi^2_{\mathrm{red}}$ 分成三类, 给出每类的代表事件 (从 \texttt{track\_summary.csv} 直接读取):

\begin{table}[H]
\centering\scriptsize
\caption{失败模式分类(按 $\chi^2_{\mathrm{red}}$ 分箱). 数据来自 \texttt{track\_summary.csv} 直接读取; \texttt{event\_index, track\_index} 唯一定位.}
\label{tab:partIII:taxonomy}
\begin{tabular}{lrlllll}
\toprule
类别 & $\chi^2_{\mathrm{red}}$ 范围 & 占比 & 代表事件 & 残差 ($p_x, p_y, p_z$) MeV/$c$ & Fisher $\sigma$ ($p_x, p_y, p_z$) & 物理诠释 \\
\midrule
A 正常 & $[0, 2)$ & $87.8\%$ & (626, 0) & $(-2.49, +0.66, +4.38)$ & $(6.82, 0.64, 6.91)$ & 名义重建, 无明显 dE/dx \\
A 正常 & $[0, 2)$ & --- & (428, 0) & $(+2.36, -0.97, -3.51)$ & $(7.00, 0.75, 6.22)$ & 同上, 不同 truth bin \\
B 中度 & $[2, 10]$ & $5.5\%$ & (482, 0) & $(-4.37, +1.14, +2.38)$ & $(11.16, 1.03, 5.41)$ & dE/dx 致 $p_z$ pull, $p_x$ ratio 大 \\
B 中度 & $[2, 10]$ & --- & (742, 0) & $(-148, +0.56, -549)$ & $(10.58, 1.45, 8.54)$ & 极少见: $p_z$ 跑到 77, LM "刚刚到" \\
C 严重 & $> 10$ & $6.7\%$ & (229, 0) & $(-449, +11.5, +312)$ & $(30.4, 1.98, 19.5)$ & LM trapped at $p_x \to -500$ \\
C 严重 & $> 10$ & --- & (121, 0) & $(-514, +2.49, +335)$ & $(58.5, 2.11, 38.7)$ & $(-100, 0)$ 病态点 \\
\bottomrule
\end{tabular}
\end{table}

\paragraph{Class A — 正常事件.}
$\sim 88\%$ 的事件落在这里. $\chi^2_{\mathrm{red}} \in [0, 2)$, 残差小到 Fisher $\sigma$ 量级, pull 接近 $\mathcal{N}(0, 1)$. 这些事件给出的覆盖率本应为 0.68; 当前 $p_y$ 在 Class A 内部仍欠覆盖, 因为 Class A 内 \emph{每事件} 的 $p_y$ Fisher 低估 $\sim 1.86\times$.

\paragraph{Class B — 中度失配 ($\chi^2_{\mathrm{red}} \in [2, 10]$).}
$\sim 5.5\%$. 残差不再 Fisher $\sigma$ 量级但仍有限. 主要机制:
\begin{enumerate}[label=(\arabic*), nosep]
  \item dE/dx 缺失: $p_z$ 在 1 m 空气段损失 $\sim 0.5$ MeV/$c$, 在 PDC 气体里再损失 $\sim 0.2$ MeV/$c$, 合计 $\sim 0.7$ MeV/$c$; LM 把这个能损"吸收"到 $p_x, p_z$ 的拟合里, $\chi^2$ 增大;
  \item 多次散射在 1 m 空气段超出 Fisher 假设(高斯独立 $\sigma$): 一次大角散射 $\theta \sim 5\sigma_{\mathrm{ms}}$ 让两个 PDC 的 $u$ 残差不再独立, $\chi^2$ 翻倍;
  \item $(u, v, p)$ 参数化的非线性: LM 收敛到 $(u, v, p)$ 空间的最小, 但其在 $(p_x, p_y, p_z)$ 空间未必无偏(参看 Part II §8 详述).
\end{enumerate}

\paragraph{Class C — 严重失败 ($\chi^2_{\mathrm{red}} > 10$).}
$\sim 6.7\%$, 50 条事件. 三种机制:
\begin{enumerate}[label=(\arabic*), nosep]
  \item \textbf{LM 困在错误盆地}: 网格搜索的初值 $(u, v, p)$ 落在 $\chi^2$ 二级最小附近, LM 收敛到 $\hat p_x \sim -500$, $\hat p_z \sim 940$ —— 这是状态报告 §误差分析 里描述的 "(-100, 0) 病态点". 第 \ref{sec:partIII:basin} 章详细剖析.
  \item \textbf{反常 PDC 命中}: PDC 簇心被错误识别(如多径迹串扰), 单条 hit 的 $u$ 偏离 truth $\sim 100$ mm, 残差扭曲, $\chi^2 \to 10^3$. 当前 PDC 模拟未注入此类异常, 但实际数据中确实存在.
  \item \textbf{多次散射爆发}: 1 m 空气段一次硬散射(电磁簇射或大角弹性), 让事件偏离高斯 MS 假设. 当前 Geant4 模拟开了 \texttt{G4UrbanMscModel} 含尾部, 偶然发生概率 $\sim 10^{-3}$.
\end{enumerate}

\subsection{每类事件在覆盖率统计中的贡献}
\label{sec:partIII:chi2:contribution}

\begin{table}[H]
\centering\small
\caption{各类事件在 744 事件覆盖率统计中的贡献(估算).}
\label{tab:partIII:chi2_class_contribution}
\begin{tabular}{lrr}
\toprule
类别 & 事件数 & 对总覆盖率(of $|\vec p|$) 拉低量 \\
\midrule
A 正常 & 654 & $\sim 0$ (本身贴近 0.68) \\
B 中度 & 41  & $\sim 0.005$ (覆盖 $\sim 0.5$, 比 0.68 低 0.18, $\times$ 41/744) \\
C 严重 & 49  & $\sim 0.066$ (覆盖 $\sim 0$, 比 0.68 低 0.68, $\times$ 49/744) \\
\midrule
合计    & 744 & $\sim 0.071$, 即把名义 0.68 拉到 $\sim 0.61$ \\
\bottomrule
\end{tabular}
\end{table}

但是状态报告 $|\vec p|$ 实测覆盖率是 0.76 不是 0.61 —— Fisher $\sigma$ 在 Class C 里 \emph{自适应增大} (例如事件 121 $\sigma_{p_x} = 58.5$, 事件 229 $\sigma_{p_x} = 30.4$), 部分 Class C 事件因 Fisher $\sigma$ 巨大反而被算作 "覆盖". 这是 Fisher 在大 $\chi^2$ 时数值不稳定的副作用 ——它假设 $\chi^2$ 在最优点是高斯 quadratic, 但当 LM 困在错误盆地, "最优点" 在错误位置, Hessian 给出的 $\sigma$ 也在错误尺度上.

\subsection{Wilks's theorem 与 $\chi^2$ 解读边界}
\label{sec:partIII:chi2:wilks}

Wilks (1938) 定理给出: 在零假设下, $-2 \ln (\mathcal{L}_{\mathrm{null}} / \mathcal{L}_{\mathrm{alt}})$ 渐近 $\chi^2_{\Delta k}$ 分布, 其中 $\Delta k$ 是参数维度差. 在我们的问题里:
\begin{itemize}[nosep]
  \item 备择假设 $H_1$: 自由 $(u, v, p)$ 三参数最小化 $\chi^2$, 给 $\chi^2_{\min}$;
  \item 零假设 $H_0$: 模型完美 (truth 已知, 无须拟合), 残差 $\sim \mathcal{N}(0, \sigma_{\mathrm{PDC}})$;
  \item Wilks: $\chi^2_{\min} \sim \chi^2_{N_{\mathrm{res}} - N_{\mathrm{par}}} = \chi^2_1$.
\end{itemize}

\paragraph{Wilks 失效的条件.}
Wilks 假设 (a) 模型在最优点是 regular (Hessian 正定), (b) 数据量大 (asymptotic). 在 $N_{\mathrm{dof}} = 1$ 时, "数据量大" 是个尴尬的条件 — 残差只有 4 个, 参数 3 个, 远未 asymptotic. 但是 PDC 残差是 \emph{独立 4 维高斯}, 在这种全高斯情形 Wilks 即便 $N_{\mathrm{dof}} = 1$ 也精确成立. 这是为什么状态报告中 Class A 的 $\chi^2$ 经验分布与 $\chi^2_1$ 形状一致.

\paragraph{Wilks 在 LM 失败时不成立.}
Class C 事件不满足 Wilks 假设, 因为 (a) LM 不在 $\chi^2$ 全局最小, Hessian 在错盆地内的曲率与全局最小处不同, regular 假设破裂. 这就是为什么 Class C 的 $\chi^2$ 数量级是 $10^3$ 而不是 $\chi^2_1$ 的 $\sim 1$.

\paragraph{p-value 解读.}
对 Class A 主峰内的事件, 把 $\chi^2_{\mathrm{red}}$ 转成 p-value:
\[
  p_{\mathrm{val}} = \Pr(\chi^2_1 > \chi^2_{\mathrm{obs}}).
\]
例如 $\chi^2_{\mathrm{obs}} = 0.5 \Rightarrow p_{\mathrm{val}} = 0.48$ (好); $\chi^2_{\mathrm{obs}} = 4 \Rightarrow p_{\mathrm{val}} = 0.046$ (5\% 水平边界). 在 744 事件中, p-value 小于 $0.05$ 的事件应占 $\sim 5\%$, 即 $\sim 37$ 条; 实测 60 条 $\chi^2 > 4$, 比期望多 $\sim 1.8\times$, 与表 \ref{tab:partIII:chi2_observed} 一致.

\subsection{$\chi^2$ 阈值作为事件级 quality cut}
\label{sec:partIII:chi2:cut}

实用建议: 在最终物理分析里, 加 $\chi^2_{\mathrm{red}} < 5$ 的 quality cut 可以剔除大部分 Class C, 保留 $> 95\%$ 的 Class A+B. 这种 cut 在覆盖率上的效果(模拟):

\begin{table}[H]
\centering\small
\caption{$\chi^2_{\mathrm{red}}$ cut 对覆盖率的影响(估算, 待运行).}
\label{tab:partIII:chi2_cut_estimate}
\begin{tabular}{lrrrr}
\toprule
Cut & 保留事件 & $|\vec p|$ Fisher 68\% & $p_y$ Fisher 68\% & 备注 \\
\midrule
无         & 744 & 0.761 & 0.421 & 当前数字 \\
$< 10$    & $\sim 694$ & TODO & TODO & 应消除 LM 失败 \\
$< 5$     & $\sim 690$ & TODO & TODO & 同上 + 中度 dE/dx 离群 \\
$< 2$     & $\sim 653$ & TODO & TODO & 极保守 \\
\bottomrule
\end{tabular}
\end{table}

% TODO: 上表数字需要在 analyze_ensemble_coverage.py 加 --chi2-max 选项后重跑.
% 命令: python3 scripts/analysis/analyze_ensemble_coverage.py \
%   --input-dir build/test_output/ensemble_n50_targetframe/.../ \
%   --output-dir build/test_output/ensemble_n50_targetframe/coverage_chi2cut5 \
%   --chi2-max 5
% 期望运行时间: 30 秒 (纯 Python 后处理).

\subsection{小结}
\label{sec:partIII:chi2:summary}

\begin{itemize}[nosep]
  \item $\chi^2_1$ 分布尾部重: $> 4$ 的占 $\sim 4.6\%$, $> 10$ 的占 $\sim 0.16\%$ (无机制下).
  \item 实测 744 事件中, $> 10$ 占 $6.7\%$, 比理论高 \emph{42 倍}, 是真正 "另一个机制" 在主导.
  \item 三类机制: LM 错盆地 (主要), 反常 hit, MS 爆发. 对应代表事件已列在表 \ref{tab:partIII:taxonomy}.
  \item $\chi^2$ cut 是简单粗暴但有效的 quality cut, 应在生产分析里默认开启.
\end{itemize}

% =============================================================

\section{单事件深度分析: 6 个示例}
\label{sec:partIII:case_studies}

本章逐条分析 6 个代表事件, 每条事件:
(a) 引用 \texttt{track\_summary.csv} 的原始行;
(b) 标注哪些 Part II §1--9 误差源主导其残差;
(c) 当 \texttt{profile\_samples.csv}/\texttt{bayesian\_samples.csv} 包含该事件时, 引用对应行;
(d) 讨论 LM trace 的预期形态 (实际 trace 数据当前未持久化, 需要 \texttt{PDCErrorAnalysis.cc} 加调试输出).

事件选取覆盖谱为: 标准 Class A (1 条), 中度 dE/dx (1 条), $p_y$ 显著欠覆盖 (1 条), Class C 极端 (1 条 from $(-100, 0)$ 病态点), Class C 中等 (1 条), 边界事件 (1 条).

\subsection{事件 (655, 0) — Class A 标准}
\label{sec:partIII:case_studies:event_655}

\begin{table}[H]
\centering\small
\caption{事件 (655, 0) 数据卡片.}
\begin{tabular}{ll}
\toprule
truth $(p_x, p_y, p_z)$ & $(50.0, 0.0, 627.0)$ MeV/$c$ \\
reco $(p_x, p_y, p_z)$ & $(47.91, -0.67, 630.46)$ MeV/$c$ \\
残差 $(p_x, p_y, p_z)$ & $(-2.09, -0.67, +3.46)$ MeV/$c$ \\
$\chi^2_{\mathrm{red}}$ & $0.000308$ \\
Fisher $\sigma$ $(p_x, p_y, p_z)$ & $(6.82, 0.63, 6.95)$ MeV/$c$ \\
LM 迭代数 & 0 (网格直接命中) \\
Profile 区间 ($p_x$) & $[40.59, 55.23]$, 半宽 $7.32$ \\
Profile 区间 ($p_y$) & $[-1.13, -0.21]$, 半宽 $0.46$ \\
Profile 区间 ($p_z$) & $[623.83, 637.08]$, 半宽 $6.62$ \\
MCMC 区间 ($|\vec p|$) & $[628.19, 637.67]$, 半宽 $4.74$ \\
MCMC acc / ESS / Geweke & $0.355 / 13.8 / 6.31$ \\
\bottomrule
\end{tabular}
\end{table}

\paragraph{诠释.}
\begin{itemize}[nosep]
  \item 残差全部 $< 1\sigma$, 拟合理想; $\chi^2_{\mathrm{red}} < 10^{-3}$ 提示 LM 把残差几乎压到机器精度.
  \item Profile 区间宽度($7.32, 0.46, 6.62$) 比 Fisher 略窄, 与状态报告 §误差分析方法对比表 中 Profile $\sim 0.94 \times$ Fisher 的趋势一致.
  \item MCMC ESS 仅 13.8 (链长 160 步, 80 burn-in), Geweke $z = 6.31$ 远超 $|z| > 2$ 的不收敛阈值 —— 链没有 mix; MCMC 区间数字"碰巧" 与 Profile 接近不代表后验已被恰当采样, 仅说明高斯峰在 $\hat p \sim 633$ 附近且 walker 在峰内随机游走.
  \item LM 迭代数 = 0 意味着粗网格搜索直接给出 $\chi^2 < 10^{-3}$ 的候选, 后续 LM 不需要进一步移动. 这是 Class A 的常见模式.
\end{itemize}

\paragraph{Part II 误差源归因.}
\begin{itemize}[nosep]
  \item PDC 位置分辨 (Part II §1): $\sim 200$ μm 在两个 PDC 上, 通过 Glückstern 给出 $\sigma_p \sim 6$--$7$ MeV/$c$ —— 与 Fisher 一致;
  \item MS in air (Part II §2): $\sim 0.3$ mrad over 1 m, 在 $|\vec p| = 627$ 上贡献 $\sim 1$ MeV/$c$ —— 嵌在 6.8 里;
  \item dE/dx mean (Part II §5): $\sim 0.7$ MeV/$c$ for 1 m air + 2 PDC, 残差 $-2$ MeV/$c$ 与之同号但稍大, 可能配额 $\sim 1$ MeV/$c$;
  \item 其他源: 在此事件中可忽略.
\end{itemize}

\subsection{事件 (482, 0) — Class B 中度失配}
\label{sec:partIII:case_studies:event_482}

\begin{table}[H]
\centering\small
\caption{事件 (482, 0) 数据卡片.}
\begin{tabular}{ll}
\toprule
truth $(p_x, p_y, p_z)$ & $(-100.0, +20.0, 627.0)$ MeV/$c$ \\
reco $(p_x, p_y, p_z)$ & $(-104.37, +21.14, 629.38)$ MeV/$c$ \\
残差 $(p_x, p_y, p_z)$ & $(-4.37, +1.14, +2.38)$ MeV/$c$ \\
$\chi^2_{\mathrm{red}}$ & $2.035$ \\
Fisher $\sigma$ $(p_x, p_y, p_z)$ & $(11.16, 1.03, 5.41)$ MeV/$c$ \\
LM 迭代数 & 0 \\
\bottomrule
\end{tabular}
\end{table}

\paragraph{诠释.}
\begin{itemize}[nosep]
  \item 残差 $/\sigma$ 比为 $(0.39, 1.11, 0.44)$ — $p_y$ 已经在 $1\sigma$ 边缘. 这件 (truth $p_y = 20$) 加上 (truth $p_x = -100$) 提供较高 PDC2 弯曲 $\theta_{xz} \sim -16°$, 把 dE/dx 在弯曲弧上分配的能量损失更多, $p_z$ 残差正向偏 (重建 $p_z$ 大于 truth, 因为 dE/dx 缺失意味着 LM 假设无能损, 把 momentum 分给了 $p_z$).
  \item $\chi^2_{\mathrm{red}} = 2.03$ 处于 $\chi^2_1$ 上 $\sim 16\%$ 的概率密度尾, 是 Class B 的边缘; 但 6.7\% 的 $\chi^2_{\mathrm{red}} > 2$ 占比远大于 16\%, 说明 dE/dx 系统性地压在所有 $|p_x| = 100$ 事件上.
  \item Fisher $\sigma_{p_x} = 11.16$ 略高于平均(其他 $|p_x| = 100$ 事件 $\sigma_{p_x} \sim 7$--$9$), 可能因为 LM 困在略偏离最优的位置, Hessian 在那里的曲率更平.
\end{itemize}

\paragraph{Part II 归因.}
\begin{itemize}[nosep]
  \item dE/dx 缺失 (Part II §5) — 主导, 贡献 $\sim 0.7$ MeV/$c$ 到 $p_z$ 残差;
  \item $(u, v, p)$ 参数化偏 (Part II §8) — 在大 $|p_x|$ 时尤其显著, 贡献 $\sim 1$--$2$ MeV/$c$;
  \item PDC 分辨 + MS — 嵌在 Fisher $\sigma$ 中.
\end{itemize}

\subsection{事件 (102, 0) — $p_y$ 强欠覆盖}
\label{sec:partIII:case_studies:event_102}

\begin{table}[H]
\centering\small
\caption{事件 (102, 0) 数据卡片.}
\begin{tabular}{ll}
\toprule
truth $(p_x, p_y, p_z)$ & $(100.0, -20.0, 627.0)$ MeV/$c$ \\
reco $(p_x, p_y, p_z)$ & $(105.99, -24.36, 624.85)$ MeV/$c$ \\
残差 $(p_x, p_y, p_z)$ & $(+5.99, -4.36, -2.15)$ MeV/$c$ \\
$\chi^2_{\mathrm{red}}$ & $0.133$ \\
Fisher $\sigma$ $(p_x, p_y, p_z)$ & $(6.65, 0.49, 7.19)$ MeV/$c$ \\
$z_{p_y}$ & $-8.94$ \\
Profile $p_y$ 区间 & $[-24.76, -23.96]$, 半宽 $0.40$ \\
MCMC $|\vec p|$ & $[629.03, 642.06]$, ESS $13.5$, Geweke $1.54$ \\
\bottomrule
\end{tabular}
\end{table}

\paragraph{诠释.}
\begin{itemize}[nosep]
  \item $p_y$ 残差 $-4.36$ 是 Fisher $\sigma_{p_y} = 0.49$ 的 \emph{8.9 倍} —— 强欠覆盖典型. 但是 $\chi^2_{\mathrm{red}} = 0.133$ 极低, 说明 LM 把整个残差消耗在 \emph{两个 PDC 同向} 的 $v$ 残差上, $\chi^2$ 没有 "怒"
.
  \item Profile 区间 $[-24.76, -23.96]$ 半宽 $0.40$ 仍然不包含 truth $-20$ —— Profile 与 Fisher 同方向欠覆盖, 这是 Part II §8 描述的 $(u, v, p)$ 线性化偏 $p_y$ 的直接表现.
  \item MCMC ESS 13.5, Geweke $1.54$ 偶然在阈值内 (但仍低 ESS), MCMC 也在 \emph{偏移的}峰附近游走.
\end{itemize}

\paragraph{Part II 归因.}
$p_y$ 残差 $-4.36$ 全部归到 \emph{(u, v, p) 参数化的 Jacobian 压扁} (Part II §8). 这是 "Fisher 信息矩阵在 $(u, v, p)$ 空间到 $(p_x, p_y, p_z)$ 空间转换时, $p_y$ 列 Jacobian $\partial p_y / \partial v$ 没有恢复全部曲率信息". 由于 truth $p_y = -20$ 超出 $\pm \sigma_{p_y}$ 9 倍, LM 找的最优 $v$ 与 truth $v$ 差距远大于线性化下的预期; Fisher $\sigma$ 是从最优点的 Hessian 导, 在最优点附近的曲率不能反映这种远距离偏移.

\subsection{事件 (482, 0) 边界化版本: 事件 (742, 0) — Class B 边界}
\label{sec:partIII:case_studies:event_742}

\begin{table}[H]
\centering\small
\caption{事件 (742, 0) 数据卡片. 这个事件介于 Class B 与 Class C, $\chi^2_{\mathrm{red}} \approx 2$ 但残差极大 — LM 找到了一个"chi-square 看起来满意但物理上荒谬"的解.}
\begin{tabular}{ll}
\toprule
truth $(p_x, p_y, p_z)$ & $(0.0, 0.0, 627.0)$ MeV/$c$ \\
reco $(p_x, p_y, p_z)$ & $(-147.95, +0.56, +77.50)$ MeV/$c$ \\
残差 $(p_x, p_y, p_z)$ & $(-147.95, +0.56, -549.50)$ MeV/$c$ \\
$\chi^2_{\mathrm{red}}$ & $2.040$ \\
Fisher $\sigma$ $(p_x, p_y, p_z)$ & $(10.58, 1.45, 8.54)$ MeV/$c$ \\
LM 迭代数 & 5 \\
\bottomrule
\end{tabular}
\end{table}

\paragraph{诠释.}
\begin{itemize}[nosep]
  \item 残差 $p_z = -549.5$, $|\vec p| = 167$ vs.\ truth 627 —— LM 落在 $|\vec p|$ 极低的局部极小. 但 $\chi^2_{\mathrm{red}} = 2.04$ 看起来还行, 这是\emph{反演问题不唯一}的典型征兆.
  \item Fisher $\sigma$ 在错的位置算: $\sigma_{p_z} = 8.54$ 反映的是 \emph{该 chi-square 谷里} 的二次曲率, 与正确谷($\hat p_z \sim 627$) 的曲率无关.
  \item LM 迭代了 5 次后困在这, 因为粗网格初值落在错盆地. 网格 $u_{\mathrm{center}} = 0 / 627 \to 0$, $v = 0$, $p \in \{200, 400, 700, 1200, 2000\}$ 的某一个 $\Rightarrow$ 这种 truth $(0, 0, 627)$ 居然能困入低 $|\vec p|$ 谷, 提示磁场积分 + RK 可能在某个 $u$ 上有副 $\chi^2$ 极小.
\end{itemize}

\paragraph{Part II 归因.}
\begin{itemize}[nosep]
  \item LM 收敛盆地失败 (Part II §9) — 主导;
  \item 网格搜索初始化不足: 7$\times$3$\times$5 = 105 个候选, 但 $u$ 半幅 1.0 给的网格步长 $0.33$ 在某些 truth 点会跨过 chi-square 鞍点.
\end{itemize}

\subsection{事件 (121, 0) — Class C 极端 ($-100, 0$ 病态)}
\label{sec:partIII:case_studies:event_121}

\begin{table}[H]
\centering\small
\caption{事件 (121, 0) 数据卡片. 这是 \nameref{sec:partIII:basin} 章详细分析的"病态点"代表.}
\begin{tabular}{ll}
\toprule
truth $(p_x, p_y, p_z)$ & $(-100.0, 0.0, 627.0)$ MeV/$c$ \\
reco $(p_x, p_y, p_z)$ & $(-613.70, +2.49, +961.92)$ MeV/$c$ \\
残差 $(p_x, p_y, p_z)$ & $(-513.70, +2.49, +334.92)$ MeV/$c$ \\
$\chi^2_{\mathrm{red}}$ & $1243.89$ \\
Fisher $\sigma$ $(p_x, p_y, p_z)$ & $(58.49, 2.11, 38.66)$ MeV/$c$ \\
LM 迭代数 & 0 (粗网格已经把它推到这) \\
\bottomrule
\end{tabular}
\end{table}

\paragraph{诠释.}
\begin{itemize}[nosep]
  \item 残差 $p_x = -514$ 是 truth 的 $5\times$, $|\vec p| \sim 1145$ vs.\ truth 635. 这是粗网格的某个 $(u, v, p)$ 候选直接给出的 $\chi^2 = 1244$, 而 LM 没有从这里再走 (迭代数 $= 0$ 意味着初值的 $\chi^2$ 在多起点中是最小的).
  \item Fisher $\sigma_{p_x} = 58.5$ 巨大, 因为 Hessian 在这个错盆地里曲率小; 但残差 $-514$ 仍是 $\sigma$ 的 $8.8\times$, pull $z = -8.78$.
  \item 该事件不在 \texttt{profile\_samples.csv} 中(profile 抽样按 $\chi^2$ 分位抽 32 条 $\times$ 4 quartile = 128 条, 该事件 $\chi^2$ 离群 + 不在 quartile 抽样列表里).
\end{itemize}

\paragraph{Part II 归因.}
\begin{itemize}[nosep]
  \item LM 收敛盆地失败 (Part II §9) — 几乎全部;
  \item 第六章详细剖析为什么 truth $(p_x, p_y) = (-100, 0)$ 是网格 $u$ 范围的边界点.
\end{itemize}

\subsection{事件 (229, 0) — Class C 中等}
\label{sec:partIII:case_studies:event_229}

\begin{table}[H]
\centering\small
\caption{事件 (229, 0) 数据卡片. truth $(p_x, p_y) = (-50, 0)$, 跑成了 LM 失败.}
\begin{tabular}{ll}
\toprule
truth $(p_x, p_y, p_z)$ & $(-50.0, 0.0, 627.0)$ MeV/$c$ \\
reco $(p_x, p_y, p_z)$ & $(-499.56, +11.54, +938.68)$ MeV/$c$ \\
残差 $(p_x, p_y, p_z)$ & $(-449.56, +11.54, +311.68)$ MeV/$c$ \\
$\chi^2_{\mathrm{red}}$ & $1548.20$ \\
Fisher $\sigma$ $(p_x, p_y, p_z)$ & $(30.43, 1.98, 19.48)$ MeV/$c$ \\
LM 迭代数 & 0 \\
\bottomrule
\end{tabular}
\end{table}

\paragraph{诠释.}
\begin{itemize}[nosep]
  \item 与事件 (121, 0) 同模式但 truth $(p_x, p_y) = (-50, 0)$. 显示 (-100, 0) 病态点不孤立 — 在 $(-50, 0)$ 也有少数 LM 失败.
  \item Fisher $\sigma$ 在错盆地的曲率 \emph{比} (-100, 0) 病态点 \emph{小}, 因为该 truth 离病态边界更远; 即 $-50$ 的"病态" 程度比 $-100$ 轻.
\end{itemize}

\paragraph{Part II 归因.}
LM 收敛盆地失败 (Part II §9). 但发生频率(50 个 $(-50, 0)$ 事件中只 $\sim 2\%$ 失败) 远低于 $(-100, 0)$ ($\sim 12\%$).

\subsection{Profile/MCMC 跨事件诊断}
\label{sec:partIII:case_studies:profile_mcmc}

把上面 6 个事件的 Profile/MCMC 数据 (能查到的) 与 Fisher 横向对比:

\begin{table}[H]
\centering\scriptsize
\caption{Fisher / Profile / MCMC 三方法在示例事件上的区间宽度对比 (MeV/$c$). 缺失项表示该事件未抽到 Profile/MCMC 子集.}
\label{tab:partIII:case_methods_comparison}
\begin{tabular}{lrrrrrrr}
\toprule
事件 & 分量 & Fisher $\sigma$ & Fisher 半宽(68\%) & Profile 半宽(68\%) & MCMC 半宽(68\%) & MCMC ESS & MCMC Geweke \\
\midrule
\multirow{4}{*}{(655, 0)}
 & $p_x$ & 6.82 & 6.82 & 7.32 & --- & --- & --- \\
 & $p_y$ & 0.63 & 0.63 & 0.46 & --- & --- & --- \\
 & $p_z$ & 6.95 & 6.95 & 6.62 & --- & --- & --- \\
 & $|\vec p|$ & 6.45 & 6.45 & 6.45 & 4.74 & 13.8 & 6.31 \\
\midrule
\multirow{4}{*}{(178, 0)}
 & $p_x$ & 6.35 & 6.35 & 6.75 & --- & --- & --- \\
 & $p_y$ & 0.50 & 0.50 & 0.46 & --- & --- & --- \\
 & $p_z$ & 7.45 & 7.45 & 7.13 & --- & --- & --- \\
 & $|\vec p|$ & --- & --- & --- & 5.08 & 13.6 & 1.83 \\
\midrule
\multirow{4}{*}{(102, 0)}
 & $p_x$ & 6.65 & 6.65 & 5.69 & --- & --- & --- \\
 & $p_y$ & 0.49 & 0.49 & 0.40 & --- & --- & --- \\
 & $p_z$ & 7.19 & 7.19 & 6.07 & --- & --- & --- \\
 & $|\vec p|$ & --- & --- & --- & 6.51 & 13.5 & 1.54 \\
\midrule
\multirow{4}{*}{(298, 0)}
 & $p_x$ & --- & --- & 6.60 & --- & --- & --- \\
 & $p_y$ & --- & --- & 0.50 & --- & --- & --- \\
 & $p_z$ & --- & --- & 6.97 & --- & --- & --- \\
 & $|\vec p|$ & --- & --- & --- & 6.65 & 8.1 & 4.68 \\
\bottomrule
\end{tabular}
\end{table}

\paragraph{读法.}
\begin{itemize}[nosep]
  \item 三方法宽度 \emph{在同一事件内} 差距 $< 25\%$ — 与状态报告 §误差分析 中 ensemble 级的差距一致.
  \item Profile $p_y$ 比 Fisher $p_y$ 略\emph{窄} (0.40--0.50 vs.\ 0.49--0.65) — Profile 沿 $\Delta\chi^2 = 1$ 轮廓走, 而 Fisher 是局部曲率假高斯; 在 $(u,v,p)$ 参数化下两者都低估真实 $\sigma$, 但 Profile 更甚 (因为它沿 chi-square 等高线走得更紧).
  \item MCMC ESS 普遍 $13$ 左右 (链长 160 步, burn-in 80 步), Geweke $|z|$ 多 $> 2$ — 链未 mix, 但区间数字仍然有意义, 因为 walker 仍在主峰内随机游走, 16/84 分位点是经验估计的中心宽度.
\end{itemize}

\paragraph{为什么 MCMC ESS 低?}
状态报告 §误差分析 给出 ensemble 中位 ESS $\approx 12.6$ (\texttt{summary.csv}). 链长 $160 - 80 = 80$ 净样本, ESS $\approx 12$ 意味着 \emph{自相关时间} $\tau \approx 80/12 \approx 7$ 步 — 即每 7 步独立一次, 接收率 0.33 配合 walker step size $\approx \sigma_{\mathrm{Fisher}}$ 是合理的, 但链长太短. 把链长增到 $10^4$ 应该让 ESS $\sim 10^3$, 那时 MCMC 区间才真正 reliable. 这是 Part II 不在范围, 但 Part III 应记录下来作为 next step.

\subsection{所有 6 个事件的对比总览}
\label{sec:partIII:case_studies:overview}

\begin{table}[H]
\centering\scriptsize
\caption{6 个示例事件的全面对比.}
\label{tab:partIII:case_overview}
\begin{tabular}{lrrrrrrl}
\toprule
事件 & truth $(p_x, p_y)$ & 残差 $|p_x|$ & 残差 $|p_y|$ & 残差 $|p_z|$ & $\chi^2_{\mathrm{red}}$ & 类别 & 主因 \\
\midrule
(655, 0) & $(50, 0)$    & 2.1   & 0.7   & 3.5   & 0.0003 & A & 名义 \\
(482, 0) & $(-100, 20)$ & 4.4   & 1.1   & 2.4   & 2.03   & B & dE/dx + (u,v,p) bias \\
(102, 0) & $(100, -20)$ & 6.0   & 4.4   & 2.1   & 0.13   & A* & (u,v,p) $p_y$ 线性化 \\
(742, 0) & $(0, 0)$     & 148   & 0.6   & 549   & 2.04   & B & LM 错盆地 (chi-square 局部) \\
(121, 0) & $(-100, 0)$  & 514   & 2.5   & 335   & 1244   & C & 病态点 \\
(229, 0) & $(-50, 0)$   & 450   & 11.5  & 312   & 1548   & C & 同上, 减弱版 \\
\bottomrule
\end{tabular}
\end{table}
\textit{注: A* = $\chi^2_{\mathrm{red}}$ 落在 Class A 范围, 但 $p_y$ 单分量极端欠覆盖, 是 (u, v, p) 参数化偏的特殊体现.}

\paragraph{LM trace 占位.}
% TODO: 当前 PDCErrorAnalysis.cc 不持久化 LM 的迭代历史 (lambda 变化, 残差变化, 参数轨迹).
% 加入 --dump-lm-trace 选项, 把每条事件的 (iter, lambda, chi2, u, v, p) 写到
% build/test_output/.../lm_traces/event_{N}.csv 即可绘 trace 图.
% 期望对 121, 229, 742 三条 Class C 事件画 (u, p) 平面的 LM 轨迹,
% 应能直观看到它们粘在错盆地, 没有跳出.

\subsection{小结}
\label{sec:partIII:case_studies:summary}

\begin{itemize}[nosep]
  \item 6 个示例覆盖 4 种主要机制: 名义 (A), dE/dx (B), $(u, v, p)$ 线性化偏 (A*), LM 错盆地 (B/C).
  \item 状态报告 ratio 1.86 ($p_y$) 在事件 102 上具体表现为 \emph{残差 $4.4 \times \sigma$} 的极端欠覆盖.
  \item 状态报告 $|\vec p|$ 覆盖率 0.76 在事件 742, 121, 229 上具体表现为 LM 错盆地 + Fisher $\sigma$ 在错盆地里仍给出"看似自洽" 的曲率.
  \item LM trace 数据持久化是接下来的高优先级 dev item.
\end{itemize}

% =============================================================

\section{扫描研究方案与外推}
\label{sec:partIII:scans}

本章给出四组待运行的扫描研究; 大部分是 \emph{方案+外推}, 实际运行未执行. 每个扫描列出: 物理动机, 期望趋势, 待运行命令, 期望耗时, 输出 CSV 路径.

\subsection{扫描 BS-1: 磁场 $B_y$ $\pm 5\%$}
\label{sec:partIII:scans:bs1}

\paragraph{物理动机.}
SAMURAI 偶极磁场图来自实测 + 插值 (\texttt{sm\_field.dat}). 现场实测精度 $\pm 1\%$, 但插值与 RK 步进对中心场强敏感. 扫 $B_y$ 标度 $\pm 5\%$ 给出谱仪对场强标定的灵敏度.

\paragraph{期望趋势.}
RK 重建假设场图正确; 真实场强为 $B (1 + \delta_B)$ 但代码用 $B$, 导致弯曲半径 $\rho = p / (qB)$ 算错 $\delta_B$. 拟合时 LM 把这个吸收到 $\hat p$ 里, 即 $\hat p / p \approx 1 + \delta_B$. 线性外推:
\[
  \hat{|\vec p|} \;\approx\; |\vec p|^{\mathrm{truth}} \cdot (1 + \delta_B), \quad \sigma_{|\vec p|}^{\mathrm{stat}} \to \sigma_{|\vec p|}^{\mathrm{stat}}.
\]
$\sigma$ 不变, 但中心估计偏移. 当前 $|\vec p|$ 中位数 630 MeV/$c$, $\delta_B = +5\%$ 给 $\hat p \approx 661$, $\delta_B = -5\%$ 给 $\hat p \approx 599$.

\paragraph{覆盖率推论.}
\begin{itemize}[nosep]
  \item $\delta_B = +5\%$: 中心偏 $+30$ MeV/$c$, Fisher $\sigma \sim 7$ MeV/$c$, pull mean $\sim 4.3$, 覆盖率 $|\vec p|$ 应跌到 $< 5\%$.
  \item $\delta_B = -5\%$: 镜像, pull mean $\sim -4.3$.
  \item $\delta_B \pm 1\%$: pull mean $\sim 0.9$, 覆盖率从 0.76 跌到 $\sim 0.6$.
\end{itemize}

\paragraph{待运行命令.}
% TODO: B-field scan run BS-1
\begin{lstlisting}[language=bash,caption={B-field scan BS-1 命令模板},label=lst:partIII:bs1_cmd]
# 先把 sm_field.dat 复制成两份缩放副本:
python3 scripts/analysis/scale_sm_field.py --in data/input/sm_field.dat \
    --out data/input/sm_field_scaled_p5pct.dat  --scale 1.05
python3 scripts/analysis/scale_sm_field.py --in data/input/sm_field.dat \
    --out data/input/sm_field_scaled_m5pct.dat  --scale 0.95

for SCALE in p5pct m5pct; do
  ENSEMBLE_SIZE=50 PROFILE_PER_QUARTILE=8 MCMC_PER_QUARTILE=16 \
      MAGNETIC_FIELD=data/input/sm_field_scaled_${SCALE}.dat \
      OUT_ROOT=build/test_output/scan_bs1_${SCALE} \
      bash scripts/analysis/run_ensemble_coverage.sh
done
\end{lstlisting}

\paragraph{期望耗时.} 单 scale $\sim 40$ min ($\times 2 = 80$ min).
\paragraph{期望输出.} \texttt{build/test\_output/scan\_bs1\_p5pct/}, \texttt{scan\_bs1\_m5pct/} 各带完整 ensemble 输出.

\subsection{扫描 WS-1: PDC 丝坐标分辨 $\sigma_{\mathrm{wire}}$}
\label{sec:partIII:scans:ws1}

\paragraph{物理动机.}
PDC 丝坐标分辨当前默认 $\sigma_{\mathrm{wire}} \approx 0.2$ mm (对应 cluster 重心精度). 实测 PDC 数据有 $\sim 0.5$ mm; 极端情形 (cluster 中包含 noise hit) 可到 $\sim 2$ mm. 扫 $\sigma_{\mathrm{wire}} \in \{0.1, 0.3, 0.5, 1.0, 2.0\}$ mm 验证 Glückstern $\sigma_p \propto \sigma_{\mathrm{wire}}$ 关系.

\paragraph{期望趋势.}
Glückstern 公式 (Part I §7) 给出
\[
  \sigma_p \;\propto\; \sigma_{\mathrm{wire}} \cdot |\vec p| / (qBL^2).
\]
当前 $\sigma_p \approx 7$ MeV/$c$ at $\sigma_{\mathrm{wire}} \approx 0.2$ mm. 线性外推:

\begin{table}[H]
\centering\small
\caption{$\sigma_p$ vs.\ $\sigma_{\mathrm{wire}}$ 线性外推.}
\label{tab:partIII:ws1_extrapolate}
\begin{tabular}{rr}
\toprule
$\sigma_{\mathrm{wire}}$ (mm) & $\sigma_p$ (MeV/$c$, 外推) \\
\midrule
0.1 & 3.5 \\
0.3 & 10.5 \\
0.5 & 17.5 \\
1.0 & 35.0 \\
2.0 & 70.0 \\
\bottomrule
\end{tabular}
\end{table}

状态报告 §理论分辨极限 给的 0.5 mm 与 2.0 mm 端点应与上表一致 (待状态报告交叉核对).

\paragraph{待运行命令.}
% TODO: sigma_wire scan run WS-1
\begin{lstlisting}[language=bash,caption={sigma\_wire scan WS-1},label=lst:partIII:ws1_cmd]
for SW in 0.1 0.3 0.5 1.0 2.0; do
  ENSEMBLE_SIZE=50 PROFILE_PER_QUARTILE=8 MCMC_PER_QUARTILE=16 \
      PDC_SIGMA_WIRE_MM=${SW} \
      OUT_ROOT=build/test_output/scan_ws1_sw${SW//./p} \
      bash scripts/analysis/run_ensemble_coverage.sh
done
\end{lstlisting}
该脚本要求 \texttt{run\_ensemble\_coverage.sh} 解析 \texttt{PDC\_SIGMA\_WIRE\_MM} 环境变量并传给 ensemble macro; 当前没有此 hook (% TODO: 在 macro/ensemble template 中加 \\setSigmaWireMm 命令).

\paragraph{期望耗时.} 单 step $\sim 40$ min, $\times 5 = 200$ min ($\sim 3.5$ h).
\paragraph{期望输出.} \texttt{build/test\_output/scan\_ws1\_sw\{0p1,0p3,0p5,1p0,2p0\}/}.

\subsection{扫描 TS-1: 靶倾角 $\alpha_{\mathrm{tgt}}$}
\label{sec:partIII:scans:ts1}

\paragraph{物理动机.}
2026-04-19 修复后, 重建端加入 $R_y(+\alpha_{\mathrm{tgt}})$ (\texttt{PDCFrameRotation.hh}) 把 lab 系结果旋回靶系. 现行配置 $\alpha_{\mathrm{tgt}} = 3°$. 我们要验证:
\begin{enumerate}[label=(\alph*), nosep]
  \item $\alpha_{\mathrm{tgt}} = 0°$ (无旋转): 重建端 frame fix 不应改变物理结果, 即 reco 仍然在 lab = target frame. $p_x$ 残差应仍 $\sim 0$.
  \item $\alpha_{\mathrm{tgt}} = 5°$: $-p_z \sin(5°) \approx -55$ MeV/$c$ 是 truth-vs-reco 的伪偏差, 但 frame fix 把它消除, 残差仍 $\sim 0$.
\end{enumerate}
期待: 三种角度下 $p_x$ 残差半宽不变 (frame 旋转是 \emph{坐标变换}, 无物理影响), 这是 sanity check.

\paragraph{期望趋势.}
\begin{itemize}[nosep]
  \item $\alpha_{\mathrm{tgt}} = 0°, 3°, 5°$: $p_x, p_y, p_z, |\vec p|$ 三种覆盖率应 \emph{完全一致} (在 ensemble 噪声范围内).
  \item 如果 $\alpha_{\mathrm{tgt}} = 5°$ 比 $0°$ 给出不同覆盖率, 说明 PDCFrameRotation 实施有 bug (例如旋转方向错, 或者旋转应用到了部分 PDC1/PDC2 而非全 reco).
\end{itemize}

\paragraph{待运行命令.}
% TODO: tilt scan run TS-1
\begin{lstlisting}[language=bash,caption={tilt scan TS-1},label=lst:partIII:ts1_cmd]
for ALPHA in 0.0 3.0 5.0; do
  ENSEMBLE_SIZE=50 PROFILE_PER_QUARTILE=8 MCMC_PER_QUARTILE=16 \
      TARGET_TILT_DEG=${ALPHA} \
      OUT_ROOT=build/test_output/scan_ts1_alpha${ALPHA//./p} \
      bash scripts/analysis/run_ensemble_coverage.sh
done
\end{lstlisting}

\paragraph{期望耗时.} $3 \times 40 = 120$ min.
\paragraph{期望输出.} 三个 ensemble dir, 比对 \texttt{coverage\_with\_ci.csv} 的 $p_x$ 行.

\subsection{扫描 ES-1: ensemble 大小 $N$}
\label{sec:partIII:scans:es1}

\paragraph{物理动机.}
当前每 truth 点 $N = 50$, 总 $N = 744$ (15 truth bin); 每点 CP $\pm 0.13$ 略嫌粗. 验证生产应取 $N = 200$/truth ($\sim 3000$ 总).

\paragraph{期望趋势.}
覆盖率本身应不变 (ensemble 是 i.i.d. 抽样), 但 CP CI 半宽按 $1 / \sqrt{N}$ 收窄:

\begin{table}[H]
\centering\small
\caption{Ensemble 大小与 CP 半宽 (per truth bin, $\hat p = 0.68$).}
\label{tab:partIII:es1_extrapolate}
\begin{tabular}{rrr}
\toprule
$N$/truth & 总 $N$ & CP 95\% 半宽 \\
\midrule
10  & 150  & $\pm 0.30$ \\
50  & 750  & $\pm 0.13$ (当前) \\
200 & 3000 & $\pm 0.07$ \\
1000 & 15000 & $\pm 0.03$ \\
\bottomrule
\end{tabular}
\end{table}

\paragraph{待运行命令.}
% TODO: ensemble-size scaling run ES-1
\begin{lstlisting}[language=bash,caption={ensemble size scan ES-1},label=lst:partIII:es1_cmd]
for N in 10 50 200 1000; do
  ENSEMBLE_SIZE=${N} PROFILE_PER_QUARTILE=8 MCMC_PER_QUARTILE=16 \
      OUT_ROOT=build/test_output/scan_es1_n${N} \
      bash scripts/analysis/run_ensemble_coverage.sh
done
\end{lstlisting}

\paragraph{期望耗时.} $N=10$ 单核 8 min, $N=50$ 40 min (already done), $N=200$ 160 min, $N=1000$ 800 min ($\sim 13$ h). 应在多核 8 cores 上并行.

\paragraph{期望输出.}
\texttt{build/test\_output/scan\_es1\_n\{10,50,200,1000\}/}; 用同一个 \texttt{analyze\_ensemble\_coverage.py} 做 CP CI 对比.

\subsection{扫描总结}
\label{sec:partIII:scans:summary}

\begin{table}[H]
\centering\small
\caption{扫描研究优先级与状态.}
\label{tab:partIII:scan_priority}
\begin{tabular}{lllll}
\toprule
ID & 扫描 & 优先级 & 状态 & 期望产出 \\
\midrule
BS-1 & $B_y \pm 5\%$        & 中 & 待运行 & 谱仪场强灵敏度 \\
WS-1 & $\sigma_{\mathrm{wire}}$ 0.1--2.0 mm & 高 & 待运行 (需 macro hook) & Glückstern 验证 \\
TS-1 & $\alpha_{\mathrm{tgt}}$ 0--5°  & 中 & 待运行 (需 macro hook) & frame fix sanity \\
ES-1 & $N$ 10--1000           & 高 & 待运行 (无 hook 改动) & 生产 ensemble 标定 \\
\bottomrule
\end{tabular}
\end{table}

% =============================================================

\section{LM 收敛盆地研究: $(-100, 0)$ 病态点剖析}
\label{sec:partIII:basin}

状态报告 §误差分析 在 G4 涨落研究里发现, $(p_x, p_y) = (-100, 0)$ truth 点的 G4 dispersion / Fisher $\sigma$ 比例在 $p_x$ 方向达到 \textbf{3.51}, 在 $p_z$ 方向达到 \textbf{3.51} (同源, 因为是同一组 LM 失败事件主导). 这远高于其他 14 个 truth 点的 $0.5 \sim 1.5$ 范围. 本章详细诊断为什么这个特定 truth 点是 LM 收敛盆地的失败热点, 并提出两条修正路径.

\subsection{从 ensemble 数据直接读取 $(-100, 0)$ 的失败统计}
\label{sec:partIII:basin:numbers}

\begin{table}[H]
\centering\small
\caption{$(p_x, p_y) = (-100, 0)$ truth 点的 50 条事件残差统计 (从 \texttt{ensemble\_pdc\_reco\_merged.csv} 直接计算).}
\label{tab:partIII:basin_pathology_stats}
\begin{tabular}{lr}
\toprule
统计量 & 值 \\
\midrule
$N$ 总事件 & 50 \\
$N$ 残差 $|\Delta p_x| > 50$ MeV/$c$ & 10 (20\%) \\
$N$ 残差 $|\Delta p_x| > 200$ MeV/$c$ & 6 (12\%) \\
$N$ $\hat p_x < -300$ MeV/$c$ & 6 (12\%) \\
失败事件中 $\hat p_x$ 范围 & $[-614, -350]$ \\
失败事件 $\chi^2_{\mathrm{red}}$ 范围 & $[583, 1244]$ \\
失败事件 $\chi^2_{\mathrm{red}}$ 中位 & $\sim 788$ \\
$\Delta p_x$ 中位 (含失败) & $-1.2$ MeV/$c$ \\
$\Delta p_x$ 均值 (含失败) & $-51.3$ MeV/$c$ (失败拖拽) \\
$\Delta p_x$ 标准差 (含失败) & $124.9$ MeV/$c$ \\
$\Delta p_x$ p84 - p16 半宽 (鲁棒) & $\sim 7$ MeV/$c$ \\
\bottomrule
\end{tabular}
\end{table}

\paragraph{诠释.}
\begin{itemize}[nosep]
  \item 失败事件占 12\% — 远高于其他 truth 点 (0--4\%).
  \item 失败事件 $\hat p_x \in [-614, -350]$, 即 LM 把 truth $-100$ 误认成 $\sim -500$ 这个量级. 残差 $\sim -400$ MeV/$c$ 是 truth 的 4 倍.
  \item 中位残差仍 $\sim 0$ — 主峰是好的, 失败事件是 \emph{拖尾}.
  \item Fisher $\sigma$ 在主峰 $\sim 7$ MeV/$c$ 但失败事件 $\sigma$ 跳到 $\sim 60$ — Fisher 在错盆地"刻舟求剑".
\end{itemize}

\subsection{为什么是 $(-100, 0)$?}
\label{sec:partIII:basin:why}

\paragraph{网格搜索的几何.}
\texttt{PDCRkAnalysisInternal.cc} 第 392--424 行的 grid search 用
\begin{itemize}[nosep]
  \item $u$ (= $p_x / p_z$ slope at target) 网格: $u \in [u_{\mathrm{line}} - 1, u_{\mathrm{line}} + 1]$, 7 步, 步长 $\sim 0.33$;
  \item $v$ (= $p_y / p_z$): $v \in [v_{\mathrm{line}} - 0.3, v_{\mathrm{line}} + 0.3]$, 3 步, 步长 $0.3$;
  \item $|\vec p|$: $\{200, 400, 700, 1200, 2000\}$ MeV/$c$ 5 个 logspace-like 候选.
\end{itemize}
其中 $u_{\mathrm{line}}, v_{\mathrm{line}}$ 是 \emph{靶到最远 PDC 的直线方向} (\texttt{init\_dir}), 即不考虑磁场偏转的几何视线方向.

\paragraph{$(-100, 0)$ 的 $u_{\mathrm{line}}$.}
truth $p = (-100, 0, 627)$, 出射方向 $u = -100/627 \approx -0.160$. 但磁场把质子在 PDC2 处弯了 $\sim 25°$ ($\theta_{xz} \sim 0.43$ rad), PDC2 真实位置在 $u \approx -0.43 + (-0.16) = -0.59$ 附近. 因此粗网格中心 $u_{\mathrm{line}} \approx -0.59$ (由 PDC2 几何定义, 不是 truth 出射 slope).

实际计算: PDC2 在 $z \approx 4655$ mm 处, 该事件 PDC2 命中 $x \approx -2745$ mm (由表 \nameref{sec:partIII:basin:numbers} 隐含); 直线 slope $u_{\mathrm{line}} = x_{\mathrm{PDC2}} / z_{\mathrm{PDC2}} \approx -2745 / 4655 \approx -0.59$. 网格 $u \in [-1.59, 0.41]$ —— truth $u = -0.16$ 在网格中心偏右 ($+0.43$ 处, 第 4 个网格点).

\paragraph{为什么是边界?}
truth $u = -0.16$ 不是在 $u_{\mathrm{line}} \pm 1$ 网格的 \emph{物理边界}, 但它\emph{接近 chi-square 第二低盆地} 的边界. 状态报告 $§误差分析$ 中 G4 涨落的 5$\times$3 truth 网格图(图 \texttt{figures/g4\_per\_truth\_box\_*.png}) 显示在 $(-100, 0)$ 这个点 $\hat p_x$ 分布是双峰: 主峰 $\hat p_x \approx -100$ 包含 88\% 事件, 副峰 $\hat p_x \approx -500$ 包含 12\% 事件. 副峰的位置对应 LM "卡在错盆地" 的 chi-square 局部极小.

\paragraph{为什么 $(-100, 20)$ 没那么严重?}
truth $(p_x, p_y) = (-100, 20)$ 的 50 事件中 $\Delta p_x$ 半宽 $\sim 8.6$ MeV/$c$ (从 \texttt{per\_truth\_point\_g4\_vs\_fisher.csv}: g4\_halfw68\_px = 8.64), 比 $(-100, 0)$ 的 $34.8$ 要小 4 倍. 这说明 $p_y \ne 0$ 引入的 $v = p_y / p_z$ 偏离对 LM "推开错盆地" 起到了关键作用 —— 网格 $v$ 步长 0.3 在 $p_y = 0$ 时正好覆盖错盆地, 在 $p_y = \pm 20$ 时把网格推到错盆地外, LM 看不到副峰.

\subsection{Chi-square 表面拓扑示意}
\label{sec:partIII:basin:chi2_topology}

% TODO: 需要画一张 (u, p) 平面 chi-square 等高线图,
% 由 scripts/analysis/plot_chi2_landscape.py 生成 (待写).
% 输入: PDC1, PDC2, target geometry; truth (p_x, p_y, p_z) = (-100, 0, 627).
% 输出: figures/diagnostics/chi2_landscape_p100_p0.png.
% 应能直观看到两个谷(主谷 $u \sim -0.16$, 副谷 $u \sim -0.7 \sim -1.0$)
% 以及它们之间的鞍点, 标注网格采样点.

期望画面:
\begin{itemize}[nosep]
  \item 主谷 (truth 谷): $(u, p) \sim (-0.16, 627)$, $\chi^2 \sim 0$;
  \item 副谷 (病态谷): $(u, p) \sim (-0.85, 1145)$, $\chi^2 \sim 1244$ — 注意 $\chi^2$ 不为零 \emph{是因为} 该副谷是 chi-square 非零的局部极小, 不是真极小;
  \item 鞍点连接两谷, $u_{\mathrm{saddle}} \sim -0.5$;
  \item 网格点 7$\times$5 = 35 个候选 ($v$ 维度暂不画), 7 个 $u$ 候选中 4 个在主谷一侧, 3 个在副谷一侧.
\end{itemize}

\subsection{修复方案 (a): 网格不对称延展}
\label{sec:partIII:basin:fix_a}

\paragraph{思路.}
$|u_{\mathrm{line}}|$ 大时, 主谷 $u_{\mathrm{truth}}$ 与 $u_{\mathrm{line}}$ 距离 $\sim 0.4$; 副谷 $u_{\mathrm{副}}$ 与 $u_{\mathrm{line}}$ 距离也 $\sim 0.3$ (在另一侧). 当前网格 $\pm 1$ 对称延展, 副谷一定能命中. 改成 \emph{不对称} 延展, 让网格偏向 truth 谷:
\[
  u_{\min} = u_{\mathrm{line}} - 0.3, \quad u_{\max} = u_{\mathrm{line}} + 1.0.
\]
即把网格 \emph{向更靠近 truth 出射方向} 偏置 (更小的 $|u|$). 但这要求知道 truth 出射方向的先验, 不可行.

\paragraph{改进思路.}
基于物理: 弯曲方向永远把 $u_{\mathrm{line}}$ 推得比 $u_{\mathrm{truth}}$ \emph{绝对值更大} (磁场把出射 deflect 到更陡的弧). 因此网格 \emph{向 $|u|$ 更小的方向} 偏置:
\[
  u_{\min} = u_{\mathrm{line}} - 0.3 \cdot \mathrm{sign}(u_{\mathrm{line}}), \quad u_{\max} = u_{\mathrm{line}} + 1.0 \cdot \mathrm{sign}(u_{\mathrm{line}}) \cdot (-1).
\]
等价地: 网格半幅 $\pm 1$, 但 truth-side 半幅 $-1.0$, 远 truth-side 半幅 $+0.3$.

\paragraph{代码改动.}
\texttt{libs/analysis\_pdc\_reco/src/PDCRkAnalysisInternal.cc} 第 392--410 行:
\begin{lstlisting}[caption={网格不对称延展 (建议改动)}]
constexpr double kUFarHalfRange  = 1.0;   // away from target
constexpr double kUNearHalfRange = 0.3;   // toward target slope
double u_lo = u_center - (u_center >= 0 ? kUNearHalfRange : kUFarHalfRange);
double u_hi = u_center + (u_center >= 0 ? kUFarHalfRange  : kUNearHalfRange);
for (int iu = 0; iu < kGridU; ++iu) {
    const double u = u_lo + (u_hi - u_lo) * iu / (kGridU - 1);
    ...
}
\end{lstlisting}

\paragraph{预期效果.}
$(-100, 0)$ truth 点 $u_{\mathrm{line}} \approx -0.59$, 改后网格 $u \in [-0.89, +0.41]$ (truth $-0.16$ 在中心偏右). 副谷在 $u \approx -0.85$ 仍在网格内, 但 truth 谷被 4 个网格点覆盖, 副谷只 1 个. 正确盆地会以多起点压倒错盆地.

\subsection{修复方案 (b): 加入 "magnetic-kick-corrected" 起点}
\label{sec:partIII:basin:fix_b}

\paragraph{思路.}
Linear backward 起点: 用 PDC1, PDC2 做反向直线外推到 target $z$, 得到 $u_{\mathrm{back}} = (x_{\mathrm{PDC1}} - x_{\mathrm{tgt}}) / (z_{\mathrm{PDC1}} - z_{\mathrm{tgt}})$. 这忽略磁场, 但可以从这个起点开始一次 $|\vec p|$ 灵敏度估计:

设质子在 dipole 中弯过弧长 $L$, 弯曲角 $\theta = qBL / p$ (Highland-style). 用 PDC2 处的 incidence angle $\theta_{\mathrm{in}}$ 减去 $\theta$ 即得 emission angle, 再外推到 target 给出 \emph{magnetic-kick corrected} $u_{\mathrm{corr}}$.

实际计算更简单: PDC1 和 PDC2 的角度差 $\Delta\theta_{\mathrm{PDCs}} = \mathrm{atan2}(\Delta x, \Delta z)$ 给出 \emph{出 PDC 后} 的轨迹方向; 把它向 target 外推 (无场), 得到的 slope 就是 \emph{出射的} $u$, 再加上一个磁场修正项 $\sim qB\Delta z / p$, 给出 \emph{在 target 的} $u$.

\paragraph{代码改动.}
新加一个候选 $u$ 候选 (替代或补充网格中心):
\begin{lstlisting}[caption={magnetic-kick-corrected 起点 (建议改动)}]
double u_pdc_to_pdc = (track.pdc2.X() - track.pdc1.X()) /
                      (track.pdc2.Z() - track.pdc1.Z());
double dz_pdc1_target = track.pdc1.Z() - target.target_position.Z();
// magnetic kick estimate (assume qBy = -0.5 T, p = 600 MeV/c => 0.4 rad/m)
double magnetic_correction = 0.5 * dz_pdc1_target / 1500.0;  // approx
double u_back = u_pdc_to_pdc - magnetic_correction;

// add as extra candidate
candidates.push_back({{u_back, v_back, 600.0}, Evaluate(...).chi2_raw});
\end{lstlisting}

\paragraph{预期效果.}
$(-100, 0)$ 的 $u_{\mathrm{back}}$ 接近 $-0.16$ (因为 PDC1, PDC2 之间的 slope 反映出 $\hat p$ 而非 $u_{\mathrm{line}}$); 加进 candidates 后, LM 多起点必有一个落到主谷.

\subsection{第三种方案 (c): 一阶 backward Kalman}
\label{sec:partIII:basin:fix_c}

更彻底的方案: 用 Kalman 反向滤波 — 从 PDC2 倒推到 PDC1 倒推到 target, 每步用 RK4 反向积分, 得到一个相对精确的 target-frame 起点. 但这基本等价于做\emph{第二次拟合}, 工作量与替换 LM 等同, 暂不优先.

\subsection{修复后预期 ratio}
\label{sec:partIII:basin:fix_expected}

\begin{table}[H]
\centering\small
\caption{修复方案对 $(-100, 0)$ truth 点 $\Delta p_x$ 半宽的预期影响.}
\label{tab:partIII:basin_fix_table}
\begin{tabular}{lr}
\toprule
方案 & 预期 $\Delta p_x$ 半宽 (MeV/$c$) \\
\midrule
当前 & 34.8 (主峰 $\sim 5$ + 12\% 失败 $\sim 400$) \\
方案 (a) 网格不对称 & $\sim 8$ (失败率应降到 0--2\%) \\
方案 (b) magnetic-kick 起点 & $\sim 5$ (失败率 $\sim 0$, 但会引入新 LM-trap 风险) \\
方案 (a) + (b) 联合 & $\sim 5$ (主峰 + 极少失败) \\
\bottomrule
\end{tabular}
\end{table}

\subsection{副谷的物理来源: 反方向轨迹}
\label{sec:partIII:basin:reverse_trajectory}

副谷 $\hat p_x \approx -500$ 是什么物理? 一种解释是: 给定 PDC1, PDC2 命中, 存在两个不同的 (target 出射, 总动量) 配置都能解释这两个点 — 一个是 truth ($\hat p_x = -100$, $\hat p_z = 627$), 另一个是高动量反向 ($\hat p_x = -500$, $\hat p_z = 940$).

\paragraph{推导.}
质子在磁场 $B_y$ 中, 弯曲半径 $\rho = p / (q B)$. 给定 PDC1, PDC2 两点和 target, 拟合即解三圆心方程 — 圆经过三个点. 但 \emph{磁场不均匀} (SAMURAI dipole 有 fringe field), 所以"圆经过三点"不严格; 不同 $|\vec p|$ 的 RK 积分可以经过相同的 PDC2 但 source 出射方向不同.

具体: 设 truth 解给出 $u_{\mathrm{truth}} = -0.16$ at target. 副谷解给出 $u_{\mathrm{副}} \approx -0.85$ at target — 方向更陡, 但 \emph{动量更高} 让弯曲变浅, 弯过 dipole 后到达 PDC2 的位置接近. 这是磁场+动量的二维耦合带来的 chi-square 多解.

\paragraph{$p_y = 0$ 的特殊性.}
为什么 $p_y \ne 0$ 时副谷消失? 因为 $v = p_y/p_z$ 为非零时, $v$ 网格 $\{v_{\mathrm{line}} - 0.3, v_{\mathrm{line}}, v_{\mathrm{line}} + 0.3\}$ 三个候选与 truth $v$ 距离都不同; 副谷需要特定的 $(u_{\mathrm{副}}, v_{\mathrm{副}})$ 组合, $v_{\mathrm{副}} \ne 0$ 时网格不再覆盖. 在 $p_y = 0$, truth $v = 0$, $v_{\mathrm{line}} = 0$, 网格中央点正好是 truth $v$ — 但同时也在副谷的 $v$ 上 (因为副谷沿 $u$ 维度延伸, $v$ 中心仍 $\sim 0$).

\subsection{为什么 $-100, 0$ 但不是 $+100, 0$?}
\label{sec:partIII:basin:asymmetry}

表 \ref{tab:partIII:other_truth_points} 显示 $(+100, 0)$ ratio 0.57 (健康). 不对称性来自:

\begin{itemize}[nosep]
  \item SAMURAI dipole $B_y$ 极性把正电荷向 $+x$ 方向偏, 即 deuteron / proton 的 $u_{\mathrm{line}}$ 偏 $+u$. truth $p_x = -100$ 的事件出射朝 $-x$, 但 magnetic kick 把它\emph{反向}弯到 $+x$ — 即 PDC2 命中位置在 $+x$ 处, $u_{\mathrm{line}} > 0$. 但实测 (\texttt{pdc\_dispersion\_summary.csv}) 显示 $(p_x = -100, p_y = 0)$ 的 PDC2 中位 $x \approx -3500$ mm, 仍在 $-x$;
  \item 这就需要重新检查: $u_{\mathrm{line}} = -3500/4655 \approx -0.75$, 网格 $u \in [-1.75, 0.25]$. truth $u = -0.16$ 在网格中点偏右. 副谷 $u_{\mathrm{副}} \approx -0.85$ 在网格中央. 即副谷正好被网格"瞄准";
  \item truth $p_x = +100$: 出射朝 $+x$, magnetic kick 反向弯, PDC2 命中 $\sim -2900$ mm (从 \texttt{pdc\_dispersion\_summary.csv} 推断), $u_{\mathrm{line}} \approx -0.62$, 网格 $u \in [-1.62, 0.38]$, truth $u = +0.16$ 在网格右端边缘 — \emph{副谷} 在 $u \approx -0.6$ 也在网格中央, 但 truth $u$ 处于边缘也意味着 LM 容易跑出主谷, 走向副谷的几率反而小. 经验上 ratio 0.57 是健康的;
  \item 物理结论: 网格设计在 \emph{$u_{\mathrm{line}}$ 与 $u_{\mathrm{truth}}$ 同号} 时(主谷与副谷都在网格内, 副谷更靠近中心) 容易失败.
\end{itemize}

\subsection{监控诊断: LM trace 持久化}
\label{sec:partIII:basin:lm_trace}

为以后追踪 LM 失败, 应在 \texttt{PDCErrorAnalysis.cc} 加 \texttt{--dump-lm-trace} 选项, 把每条事件的:
\begin{itemize}[nosep]
  \item 网格搜索的 105 个候选 $(u, v, p, \chi^2)$;
  \item 多起点 LM 的 $\le 5$ 条迭代历史 $(\mathrm{iter}, \lambda, \chi^2, u, v, p)$;
  \item 最终选取的最优点;
\end{itemize}
持久化到 \texttt{lm\_traces/event\_\{N\}.csv}. 这给后续 visualization 提供数据基础, 也允许在数据 (815 tracks, no truth) 上识别 LM 失败 (无 truth 时只能看 $\chi^2$ + LM trace, 不能看残差).

% TODO: PDCErrorAnalysis.cc 加 --dump-lm-trace 选项 (作 dev item).

\subsection{其他可疑 truth 点}
\label{sec:partIII:basin:other_points}

从 \texttt{per\_truth\_point\_g4\_vs\_fisher.csv} 直接读出 ratio (g4 / fisher) :

\begin{table}[H]
\centering\small
\caption{各 truth 点的 g4 半宽 / Fisher $\sigma$ 比例 (从 \texttt{per\_truth\_point\_g4\_vs\_fisher.csv}, $p_x$ 列). $p_y \in \{-20, 0, 20\}$ 对应三行.}
\label{tab:partIII:other_truth_points}
\begin{tabular}{rrrr}
\toprule
truth $p_x$ & ratio $p_y = -20$ & ratio $p_y = 0$ & ratio $p_y = 20$ \\
\midrule
$-100$ & 0.48 & \textbf{3.51} & 0.77 \\
$-50$  & 0.40 & 0.38 & 0.66 \\
$0$    & 0.84 & 0.63 & 0.89 \\
$+50$  & 0.51 & 0.56 & 0.72 \\
$+100$ & 0.62 & 0.57 & 0.96 \\
\bottomrule
\end{tabular}
\end{table}

\paragraph{读法.}
\begin{itemize}[nosep]
  \item $(-100, 0)$ 是唯一 ratio $> 1$ 的 truth 点; 失败率 12\% 是孤立病态.
  \item 其他 14 个 truth 点 ratio 都在 $0.4$--$1.0$ 范围, 基本健康.
  \item ratio $< 1$ 反映 \emph{Fisher 略保守}, 这是 PDC 高斯分辨条件下的预期; 主要担忧是 $> 1$ 的方向.
\end{itemize}

\subsection{小结}
\label{sec:partIII:basin:summary}

\begin{itemize}[nosep]
  \item $(-100, 0)$ 是 LM 收敛盆地失败的孤立热点, 12\% 事件 trap 在 $\hat p_x \approx -500$ 副谷.
  \item 物理原因: $u_{\mathrm{line}}$ 距 $u_{\mathrm{truth}}$ 远, 网格对称延展正好覆盖错盆地; $p_y = 0$ 时 $v$ 网格不能推开错盆地.
  \item 修复: (a) 网格不对称延展 + (b) magnetic-kick-corrected 起点, 联合预期把 ratio 从 3.51 降到 $< 1$.
  \item LM trace 持久化是接下来追踪 LM 病态的高优先级 dev item.
\end{itemize}

% =============================================================

\subsection*{Part III 全章总结}

本部分给出了 RK 误差分析的实证诊断完整框架, 从二项 CI 的统计学基础到事件级 LM trace 的具体追踪.

\textbf{六章主要结论.}
\begin{enumerate}[label=\arabic*., nosep]
  \item \textbf{覆盖率方法学}: Clopper-Pearson 是当前 default, 比 Wilson 略保守, 比 Wald 在端点正确; $N=744$ 给 CP $\pm 0.034$ 足以诊断当前的 $p_y$ 欠覆盖, 生产 ensemble 应取 $N \ge 200$/truth.
  \item \textbf{Pull 分布}: 实测 744 事件 pull 中位 $\sim 0$ (frame fix 后无偏), $p_y$ 主峰 $\sigma_z \approx 2.48$ (Fisher 低估 $\sim 1.86\times$), $p_x, p_z$ 主峰窄但有 LM 长尾干扰; 鲁棒尺度估计 IQR/1.349 是接下来要加的指标.
  \item \textbf{$\chi^2$ 分布与失败模式}: 50 条 ($6.7\%$) $\chi^2_{\mathrm{red}} > 10$ 远超 $\chi^2_1$ 自然涨落 ($0.16\%$, $42\times$); 三类机制 (LM 错盆地, 反常 hit, MS 爆发); $\chi^2 < 5$ quality cut 应在生产分析里默认开启.
  \item \textbf{6 个示例事件}: 涵盖 4 种主要机制. 事件 102 $z_{p_y} = -8.94$ 是 (u, v, p) 偏典型, 事件 121, 229 是 (-100, 0) 病态典型, 事件 742 是 LM 错盆地中等表现.
  \item \textbf{扫描研究方案}: BS-1 (B 场), WS-1 ($\sigma_{\mathrm{wire}}$), TS-1 (靶倾角 sanity), ES-1 (ensemble 规模) 已列出命令模板与外推; ES-1 优先级最高.
  \item \textbf{$(-100, 0)$ 病态点剖析}: 12\% 事件 trap 在副谷, 修复方案 (a) 网格不对称延展 + (b) magnetic-kick-corrected 起点, 联合预期把 ratio 从 3.51 降到 $< 1$.
\end{enumerate}

\textbf{下一阶段优先 dev items.}
\begin{itemize}[nosep]
  \item LM trace 持久化 (\texttt{PDCErrorAnalysis.cc} \texttt{--dump-lm-trace} 选项);
  \item Pull 直方图绘图脚本 + IQR/MAD 鲁棒尺度;
  \item Ensemble script 加 \texttt{PDC\_SIGMA\_WIRE\_MM}, \texttt{TARGET\_TILT\_DEG}, \texttt{MAGNETIC\_FIELD} 环境变量 hook;
  \item \texttt{analyze\_ensemble\_coverage.py} 加 \texttt{--chi2-max} 选项支持 quality cut;
  \item 网格搜索不对称延展 + magnetic-kick 起点的 RK 实施 + ensemble 验证.
\end{itemize}

\textbf{与 Part I, Part II 的连接.}
\begin{itemize}[nosep]
  \item 第 \ref{sec:partIII:coverage_methodology} 章的 CP CI 推导与 Part I §2 (Cramér-Rao) 的 ensemble 验证范式互补;
  \item 第 \ref{sec:partIII:pulls} 章的 pull = $\mathcal{N}(0, 1)$ 期望与 Part I §3 (Fisher) 的协方差一致性紧密相关;
  \item 第 \ref{sec:partIII:chi2} 章的失败模式分类与 Part II §1--9 的误差源对应表完全对齐;
  \item 第 \ref{sec:partIII:case_studies} 章的 6 个示例事件直接验证了 Part II §5 (dE/dx), §8 ((u,v,p) 参数化偏), §9 (LM 收敛盆地);
  \item 第 \ref{sec:partIII:scans} 章的扫描计划是 Part II 各源的 sensitivity 验证;
  \item 第 \ref{sec:partIII:basin} 章 (-100, 0) 剖析是 Part II §9 的具体落地.
\end{itemize}

% =============================================================
% End of Part III.
% =============================================================
 加载。
%
% 不要在此文件中包含 \documentclass / \begin{document} / 序言;
% 主壳层提供 \part{第三部分: 实证诊断} 标题, ctex+xelatex 字体,
% lstlisting 中文样式, booktabs/multirow/amsmath 等宏包。
%
% 创建日期: 2026-04-26
% 作者:    Baiting Tian
% 关联文件:
%   * docs/reports/reconstruction/momentum_reco/rk/rk_reconstruction_status_20260416_zh.tex
%   * docs/superpowers/specs/2026-04-26-rk-error-analysis-deep-dive-design.md
%   * docs/reports/reconstruction/momentum_reco/rk/rk_reconstruction_bugfix_and_error_analysis_20260416.md
%
% 本文件结构 (六章):
%   1. 覆盖率方法学: Clopper-Pearson / Wilson / Wald
%   2. Pull 分布: 定义, 期望, 解读, 实测
%   3. chi^2 分布与失败模式分类
%   4. 5-10 个单事件深度分析
%   5. 扫描研究方案与外推 (多数为 % TODO)
%   6. LM 收敛盆地: (-100, 0) 病态点剖析
%
% 数据来源:
%   * track_summary.csv  (744 事件, 修复后靶系)
%   * profile_samples.csv (128 事件)
%   * bayesian_samples.csv (128 事件)
%   * g4_fluctuation/per_truth_point_g4_vs_fisher.csv
%   * g4_fluctuation/ensemble_pdc_reco_merged.csv (744 事件)
% =============================================================

\section{覆盖率方法学: Clopper-Pearson, Wilson, Wald 三选一}
\label{sec:partIII:coverage_methodology}

整个状态报告 (\StatusReport{}) 把覆盖率作为 RK 误差分析的最终判据: 把名义 68\%/95\% 的方法区间打到 744 条事件上, 数 truth 落在区间内的比例; 这个比例本身又是一个二项随机变量, 必须给一个置信区间才能下结论. 报告里默认采用 Clopper-Pearson 95\% 置信区间(以下简称 CP CI), 本章给出原因, 并把它和教科书里另两种竞争方案 (Wilson 与 Wald) 做对比.

\subsection{二项覆盖率的统计模型}
\label{sec:partIII:coverage_methodology:model}

设我们独立观察 $N$ 个事件, 每个事件的方法区间 $[\hat{\theta}_i - \delta_i^-, \hat{\theta}_i + \delta_i^+]$ 是否覆盖 truth $\theta_i^{\mathrm{truth}}$ 由指示变量
\[
  X_i \;=\; \mathbb{1}\bigl[\hat\theta_i - \delta_i^- \le \theta_i^{\mathrm{truth}} \le \hat\theta_i + \delta_i^+\bigr]
\]
描述. 在 "方法自洽" 的零假设下, $\Pr(X_i = 1) = p_0$ 与名义覆盖水平 (例如 $p_0 = 0.68$) 一致, 且对 $i$ 而言 $X_i$ 独立同分布. 总和 $K = \sum_i X_i$ 服从二项分布
\[
  K \sim \mathrm{Bin}(N, p_0).
\]
观测覆盖率是点估计 $\hat p = K / N$. CI 的目标是: 给定观测 $K = k$, 构造一个区间 $[p_L, p_U]$ 使得在所有真实的 $p$ 上, 该区间覆盖真实 $p$ 的概率不低于 $1 - \alpha$ (例如 $1 - \alpha = 0.95$).

\subsection{Wald 区间(正态近似)与失败模式}
\label{sec:partIII:coverage_methodology:wald}

最常见的写法是 Wald 正态近似:
\[
  p_{L/U}^{\mathrm{Wald}} \;=\; \hat p \;\pm\; z_{\alpha/2}\, \sqrt{\frac{\hat p (1 - \hat p)}{N}}.
\]
教科书把它印在第一页是因为它形式简单, 但它在两端附近(尤其当 $\hat p$ 靠近 $0$ 或 $1$, 或者 $N$ 不够大)有以下三个失败模式:

\begin{enumerate}[label=(\alph*), nosep]
  \item 当 $\hat p = 0$ 或 $1$ 时, 标准误差 $\sqrt{\hat p(1-\hat p)/N} = 0$, 区间退化为单点 $\{0\}$ 或 $\{1\}$. 这显然不能用作 CI ——任何有限 $p$ 都可能产生 $K = 0$ 或 $K = N$.
  \item 当 $N$ 很小 (本章场景: 状态报告中 MCMC 子集 $N = 128$, 早期试运行 $N = 32$), 区间端点可能跑出 $[0, 1]$ 的物理范围.
  \item 实际覆盖率(在所有 $p$ 上对 Wald CI 求频率)远低于名义 $1 - \alpha$. Brown, Cai, DasGupta (2001, \emph{Statistical Science}) 给出了著名的 "wiggle plot": Wald 在 $p$ 略偏离 0.5 时的覆盖率会陷入 0.85 量级的低谷, 远低于名义 0.95.
\end{enumerate}

第三点在状态报告 §误差分析 里直接致命: 修复前 $p_x$ 覆盖率 $\hat p \approx 0.05$ (744 中 38 条命中); Wald 给的区间是
\[
  [0.05 - 1.96\sqrt{0.05\cdot 0.95/744},\ 0.05 + \ldots] \approx [0.034, 0.066],
\]
看起来很窄; 但当 $\hat p$ 离 0 这么近, Wald 显著地低估了不对称性 (真实分布右偏更厉害). 我们要的是一种在 $\hat p \to 0$ 或 $\hat p \to 1$ 极端区域仍然行为正确的 CI.

\subsection{Clopper-Pearson 精确二项 CI}
\label{sec:partIII:coverage_methodology:cp}

Clopper 与 Pearson (1934, \emph{Biometrika}) 提出: 让 $[p_L, p_U]$ 同时满足
\begin{align}
  \Pr\bigl(K \ge k \mid p = p_L\bigr) &\;=\; \alpha/2, \\
  \Pr\bigl(K \le k \mid p = p_U\bigr) &\;=\; \alpha/2.
\end{align}
即 $p_L$ 是 "真实 $p$ 至少为多少时, $K \ge k$ 这件事的概率才能压到 $\alpha/2$" 的最大值; $p_U$ 是 "真实 $p$ 至多为多少时, $K \le k$ 的概率才能压到 $\alpha/2$" 的最小值. 这是一种 "反演 (invert) 二项 CDF" 的 CI 构造, 直觉上等价于把假设检验的接受域翻成参数的置信域.

\paragraph{封闭形式与 Beta-incomplete.}
利用二项 CDF 与不完全 Beta 函数的对偶
\[
  \sum_{j=0}^{k}\binom{N}{j} p^j (1-p)^{N-j} \;=\; I_{1-p}(N-k,\, k+1),
\]
可以把 CP 区间端点写成 Beta 分布分位点:
\begin{align}
  p_L &\;=\; \mathrm{Beta}^{-1}\!\bigl(\alpha/2;\ k,\ N - k + 1\bigr), \\
  p_U &\;=\; \mathrm{Beta}^{-1}\!\bigl(1 - \alpha/2;\ k + 1,\ N - k\bigr).
\end{align}
端点情形: $k = 0 \Rightarrow p_L = 0$, $k = N \Rightarrow p_U = 1$, 这正是物理上希望的 "无下/上界" 行为.

\paragraph{CP 是保守的.}
CP 的 \emph{实际} 覆盖率始终不低于 $1 - \alpha$ (这就是 "精确" 的含义), 但不一定恰等于 $1 - \alpha$. 由于二项分布的离散性, 在大多数 $p$ 值上, 实际覆盖率会略高于 $1 - \alpha$. 这意味着 CP 的区间宽度比名义宽度需要的"刚好" 95\% 略宽一些; 这个保守性恰恰是状态报告需要的: 对每条覆盖率行, "如果 CP CI 不包含 0.68, 那不是统计噪声", 这是一个强论断.

\subsection{Wilson 区间}
\label{sec:partIII:coverage_methodology:wilson}

Wilson (1927) 提出的折中方案是把 Wald 的 $\hat p$ 替换为后验 Beta 调整:
\[
  p_{L/U}^{\mathrm{Wilson}} \;=\;
  \frac{\hat p + \tfrac{z^2}{2N}}{1 + \tfrac{z^2}{N}}
  \;\pm\;
  \frac{z}{1 + \tfrac{z^2}{N}}\,
  \sqrt{\frac{\hat p (1 - \hat p)}{N} + \frac{z^2}{4 N^2}}.
\]
其中 $z = z_{\alpha/2}$ (例如 $1.96$ 对应 95\% CI). 与 Wald 相比, Wilson:
\begin{itemize}[nosep]
  \item 在 $\hat p \in [0.1, 0.9]$ 区间精度极好, 实际覆盖率非常接近名义值;
  \item 区间端点保留在 $[0, 1]$;
  \item 在 $\hat p \to 0/1$ 退化得比 Wald 慢 (但仍比 CP 快).
\end{itemize}

\subsection{$N=128$, $K=80$ 的数值算例}
\label{sec:partIII:coverage_methodology:example}

把 $\hat p \approx 0.625$ 的情形 (例如 MCMC $|\vec p|$ 95\% 覆盖率, $N=128$, $K=115$) 取近似数 $K = 80$ 出来作演算示例. 95\% 双侧:

\begin{table}[H]
\centering\small
\caption{$N=128$, $K=80$ 时三种 CI 的对比 ($\hat p = 0.625$).}
\label{tab:partIII:cp_wilson_wald_demo}
\begin{tabular}{lcccc}
\toprule
方法 & 公式入口 & $p_L$ & $p_U$ & 半宽 \\
\midrule
Wald     & $\hat p \pm 1.96\sqrt{\hat p (1-\hat p)/N}$ & 0.5412 & 0.7088 & 0.0838 \\
Wilson   & 公式 \eqref{eq:partIII:wilson} & 0.5380 & 0.7053 & 0.0837 \\
Clopper-Pearson & Beta 分位点 & 0.5365 & 0.7068 & 0.0852 \\
\bottomrule
\end{tabular}
\end{table}

\begin{equation}
\label{eq:partIII:wilson}
  p_{L/U}^{\mathrm{Wilson}} = \frac{\hat p + z^2/(2N)}{1 + z^2/N}
  \pm \frac{z\sqrt{\hat p(1-\hat p)/N + z^2/(4N^2)}}{1 + z^2/N}.
\end{equation}

% TODO: 用 scripts/analysis/analyze_ensemble_coverage.py 内部 _clopper_pearson 函数
% (它即用 scipy.stats.beta.ppf) 与 statsmodels.stats.proportion 的 wilson_score
% 对全部 4*4 覆盖率行批量出 CSV; 当前数值是手算/样例,
% 需要的运行: ENSEMBLE_SIZE=50 ./scripts/analysis/run_ensemble_coverage.sh
% 已经跑过, 这里仅缺一个三方法对比表生成器.

在 $\hat p \approx 0.625$ 这种 "中部" 区间, 三种方法差距 $< 1\%$; 真正的差距出现在尾部. 对极端 $\hat p \approx 0.05$, $N=744$, $K=38$:

\begin{table}[H]
\centering\small
\caption{$N=744$, $K=38$ 时三种 CI 的对比 ($\hat p = 0.0511$, 修复前 $p_x$ 覆盖率).}
\label{tab:partIII:cp_wilson_wald_extreme}
\begin{tabular}{lccc}
\toprule
方法 & $p_L$ & $p_U$ & 备注 \\
\midrule
Wald     & 0.0353 & 0.0669 & 端点近似对称, 偏窄 \\
Wilson   & 0.0376 & 0.0694 & 内推到 $[0, 1]$ \\
Clopper-Pearson & 0.0365 & 0.0698 & 略宽, 给出真正下界 \\
\bottomrule
\end{tabular}
\end{table}

CP 与 Wilson 的差距不大 (各 $\sim 0.001$), 但都明显比 Wald 长. 既然计算成本相同 ($\sim$ 一次 Beta 分位点计算), 选 CP 是合理的工程选择: 永远不会因为 $\hat p$ 走到极端而退化, 永远是保守的.

\subsection{CP 区间反演的几何直觉}
\label{sec:partIII:coverage_methodology:geometry}

CP 的反演构造可以画一张二维图来理解. 在 $(p, K)$ 平面上(横轴是参数 $p$, 纵轴是观测计数 $K$), 对每个 $p$ 画 $K$ 的双侧 $\alpha/2$ 接受区间 $[K_L(p), K_U(p)]$:
\begin{align}
  K_L(p) &= \min\{k : \Pr(K \ge k \mid p) \le 1 - \alpha/2\}, \\
  K_U(p) &= \max\{k : \Pr(K \le k \mid p) \le 1 - \alpha/2\}.
\end{align}
这两条上下边界形成一个"接受带"; 给定观测 $k$, 把横线 $K = k$ 与该带相交, 横向投影即得到 CP CI $[p_L, p_U]$. 离散的二项分布让边界呈"阶梯状", 这是 CP 略保守的根源 — 阶梯在某些 $p$ 上比连续构造高出 $\sim 1$ 单位.

如下伪代码反映了这一构造的数值实现 (与 \texttt{analyze\_ensemble\_coverage.py} 的 \texttt{\_clopper\_pearson} 函数一致):

\begin{lstlisting}[language=Python,caption={CP CI 的最小数值实现},label=lst:partIII:cp_python]
from scipy import stats

def clopper_pearson(k, n, alpha=0.05):
    """Clopper-Pearson exact binomial CI, two-sided 1-alpha coverage."""
    if k == 0:
        p_lo = 0.0
    else:
        p_lo = stats.beta.ppf(alpha / 2, k, n - k + 1)
    if k == n:
        p_hi = 1.0
    else:
        p_hi = stats.beta.ppf(1 - alpha / 2, k + 1, n - k)
    return p_lo, p_hi
\end{lstlisting}

\paragraph{阶梯效应数值化.}
对 $N=128$, 取 $K=87$ 与 $K=88$ (相邻整数), CP CI 的下端跳变:
\begin{itemize}[nosep]
  \item $K=87$: $p_L = 0.602$;
  \item $K=88$: $p_L = 0.610$.
\end{itemize}
即 $K$ 增 1 时 $p_L$ 跳 $\sim 0.008$. 状态报告里观测覆盖率精度因此有 \emph{量子化下限} $\sim 1/N$, 这是离散二项分布的根本约束, 不是计算误差.

\subsection{多覆盖率水平: 95\% vs.\ 99\%}
\label{sec:partIII:coverage_methodology:multilevel}

状态报告默认 95\% CP CI; 偶尔需要 99\% 来给出"几乎确定" 的结论. 把上面 $\hat p = 0.625$, $N=128$ 的算例扩展:

\begin{table}[H]
\centering\small
\caption{$N=128$, $K=80$ 时不同 $\alpha$ 的 CP CI.}
\label{tab:partIII:cp_multilevel}
\begin{tabular}{lcccc}
\toprule
$1 - \alpha$ & $z_{\alpha/2}$ (Wald 参考) & $p_L$ (CP) & $p_U$ (CP) & 半宽 \\
\midrule
68\% & 1.00 & 0.582 & 0.667 & 0.043 \\
90\% & 1.64 & 0.553 & 0.692 & 0.069 \\
95\% & 1.96 & 0.536 & 0.707 & 0.085 \\
99\% & 2.58 & 0.508 & 0.730 & 0.111 \\
\bottomrule
\end{tabular}
\end{table}

\paragraph{推论.} 99\% 比 95\% 半宽宽 $\sim 1.3\times$. 状态报告的覆盖率结论 ($p_y$ 欠覆盖 $0.42$ vs.\ $0.68$) 用 99\% CP CI 仍然成立 (差距 $0.26 \gg 0.111$). 但 MCMC 的 $|\vec p|$ 0.633 vs.\ 0.68 的差距 0.05 在 95\% CP CI 半宽 0.085 之内, 在 99\% 半宽 0.111 之内 — 这个差异 \emph{不显著}, 这是状态报告 §误差分析 中已经强调的事实.

\subsection{ensemble 数据的独立性假设}
\label{sec:partIII:coverage_methodology:independence}

CP CI 假设 $X_i$ i.i.d.. 现实情况:
\begin{itemize}[nosep]
  \item ensemble 在 15 个 truth bin 上分组 ($N=50$/bin, 总 $N=744$). 同一 truth bin 内的 50 条事件的 \emph{Geant4 输入} 完全相同, 只是随机种子不同 — 这种意义上独立.
  \item 但是, 同一 truth bin 内 LM 失败事件占 12\% (在 $-100, 0$ bin) 或 $0$\% (在其他 bin) — \emph{失败的发生与 truth bin 强相关}, 不是 i.i.d.. 这意味着把全 744 看作单一二项总样本会低估 CP CI.
  \item 更严格的处理: 对每个 truth bin 单独算 CP CI, 然后做 stratified weighted average. 当前 \texttt{analyze\_ensemble\_coverage.py} 是 pooled 处理.
\end{itemize}

\paragraph{stratified analysis 占位.}
% TODO: analyze_ensemble_coverage.py 加 --stratify-by-truth 选项,
% 输出 build/test_output/.../coverage_per_truth_bin.csv,
% 每行 (truth_px, truth_py, N, k, p_hat, cp_lo, cp_hi).

\subsection{为什么不用 bootstrap?}
\label{sec:partIII:coverage_methodology:no_bootstrap}

Bootstrap 在覆盖率比例上是完全可行的: 把 $N$ 个 $X_i$ 重抽样 $B$ 次, 每次算 $\hat p^*_b$, 取经验分位点. 但相对 CP 它有三个劣势:

\begin{itemize}[nosep]
  \item \textbf{计算成本}: $B = 10^4$ 次重抽样 vs.\ 单次 Beta 分位点; 在 60 个覆盖率行 (4 分量 $\times$ 4 方法 $\times$ 多覆盖水平 $\times$ 多 truth bin) 累积起来非微秒.
  \item \textbf{近似性}: bootstrap 的覆盖率本身仍然是渐近 $1 - \alpha$, 在小 $N$ 不"精确".
  \item \textbf{无法重复}: 每次 bootstrap 有随机种子, 报告读者复现时数字会跳; CP 是闭式的.
\end{itemize}

闭式 CP 足够保守, 故选定为状态报告标准.

\subsection{ensemble 大小 $\to$ CP 半宽的换算表}
\label{sec:partIII:coverage_methodology:N_scaling}

当 $\hat p = 0.68$ 时, 二项标准差为 $\sqrt{p(1-p)/N} = \sqrt{0.2176/N}$, Wald 半宽是 $1.96$ 倍此值; CP 半宽与 Wald 的中部值相差 $< 5\%$. 估算各 $N$ 下的 CP 95\% 半宽:

\begin{table}[H]
\centering\small
\caption{$\hat p = 0.68$ 时 $N$ 与 CP 95\% 半宽的近似关系. 列 "Wald 近似" 给 $1.96\sqrt{p(1-p)/N}$, 列 "CP 实际" 由 \texttt{scipy.stats.beta.ppf} 直接计算.}
\label{tab:partIII:N_cp_halfw}
\begin{tabular}{rrr}
\toprule
$N$ & Wald 近似 & CP 实际 \\
\midrule
32   & 0.162 & 0.180 \\
64   & 0.114 & 0.121 \\
128  & 0.081 & 0.084 \\
256  & 0.057 & 0.058 \\
744  & 0.034 & 0.034 \\
5000 & 0.013 & 0.013 \\
\bottomrule
\end{tabular}
\end{table}

\textbf{解读}: 状态报告里 Fisher/Laplace 用 $N=744$, CP 半宽 $\sim 0.034$, 即名义 0.68 与观测 $\sim 0.42$ ($p_y$) 的差距 $0.26$ 远超过 $\pm 0.034$, 是 "真信号"; Profile/MCMC 用 $N=128$, CP 半宽 $\sim 0.084$, 仍能区分 $0.42$ vs $0.68$. 但是, 一旦缩到 $N=32$, 半宽 $\sim 0.18$ 跨过差距, 任何小于 $\sim 18\%$ 的偏离都会埋没在统计噪声里.

\textbf{推论}: 生产用 $N=200$ (CP $\pm 0.07$) 已经够诊断 $\sim 10\%$ 级偏差; 当前 $N=50$ 每 truth 点(总 $N=744$) 只在每个 truth bin 内 CP $\pm 0.13$, 每点诊断略嫌粗. 下文 \nameref{sec:partIII:scans} 里 ES-1 计划把生产 ensemble 增到 $N=200$/truth 点 ($\sim 3000$ 总).

\subsection{小结}
\label{sec:partIII:coverage_methodology:summary}

\begin{itemize}[nosep]
  \item Wald 在 $\hat p \to 0/1$ 退化, 在 $\hat p \in (0.1, 0.9)$ 也轻微低覆盖, 状态报告不用.
  \item Wilson 大多数情形与 CP 接近, 但缺乏 "永远 $\ge 1 - \alpha$" 的精确性保证.
  \item Clopper-Pearson 是默认选项: 闭式, 保守, 在尾部行为正确.
  \item 当前 $N=744$ 给 CP $\pm 0.034$, 足以分辨当前的 $p_y$ 欠覆盖 $0.42$ vs.\ 名义 $0.68$; $N=128$ MCMC/Profile 子集给 CP $\pm 0.084$, 仍可诊断 $\sim 0.1$ 级偏差.
\end{itemize}

% =============================================================

\section{Pull 分布: 定义, 期望, 解读}
\label{sec:partIII:pulls}

覆盖率检验告诉你 "区间宽度对不对", 但它把整个分布折成一个标量. 真正展示 "误差棒形状是否正确" 的诊断量是 \emph{pull}. 本章给出 pull 的形式定义, 在三种缺陷模式下的标志特征, 并直接读取 \texttt{track\_summary.csv} 算出 744 事件的 pull 统计.

\subsection{Pull 定义}
\label{sec:partIII:pulls:definition}

设第 $i$ 条事件的某个动量分量(以 $p_x$ 为例)的真值为 $p_{x,i}^{\mathrm{truth}}$, 重建给出的中心估计为 $\hat p_{x,i}$, 误差分析(本节默认 Fisher) 给出的标准差为 $\sigma_{p_x, i}$. 定义事件级 pull:
\[
  z_{p_x, i} \;=\; \frac{\hat p_{x,i} - p_{x,i}^{\mathrm{truth}}}{\sigma_{p_x, i}}.
\]
注意: pull 并不等于 \emph{标准化残差} 一般意义上的 $\chi^2$ 项 —— 它取的是 \emph{重建估计与真值之差} 除以\emph{该事件的}估计 $\sigma$, 不是平均 $\sigma$. 这意味着每条事件都用自己的方差归一, 是一个真正的 "诊断这条事件" 的量.

\subsection{标准期望}
\label{sec:partIII:pulls:expected}

如果(且仅如果)以下三个条件都成立,
\begin{enumerate}[label=(\alph*), nosep]
  \item 重建估计是无偏的: $\mathbb{E}[\hat p_x | p_x^{\mathrm{truth}}] = p_x^{\mathrm{truth}}$;
  \item 误差分析给出的 $\sigma$ 与真实方差一致: $\mathrm{Var}[\hat p_x | p_x^{\mathrm{truth}}] = \sigma_{p_x}^2$;
  \item 残差是高斯分布的: $\hat p_x | p_x^{\mathrm{truth}} \sim \mathcal{N}(p_x^{\mathrm{truth}}, \sigma_{p_x}^2)$,
\end{enumerate}
那么 pull 服从标准正态:
\[
  z \sim \mathcal{N}(0, 1).
\]
推论:
\begin{itemize}[nosep]
  \item $\mathbb{E}[z] = 0$,
  \item $\mathrm{Var}[z] = 1$ 即 $\sigma_z = 1$,
  \item $\Pr(|z| \le 1) = 0.6827$, $\Pr(|z| \le 1.96) = 0.95$.
\end{itemize}
任何对 $\mathcal{N}(0,1)$ 的偏离都对应一种系统问题. 表 \ref{tab:partIII:pull_diagnostics_template} 给出三类典型缺陷的特征:

\begin{table}[H]
\centering\small
\caption{Pull 分布对应的故障模式.}
\label{tab:partIII:pull_diagnostics_template}
\begin{tabular}{llp{0.45\linewidth}}
\toprule
特征 & 说明 & 故障模式 \\
\midrule
$\bar z \ne 0$ & pull 中心不在零 & 重建有偏: 坐标系/dE/dx/对齐错误的几何位移 \\
$\sigma_z > 1$ & pull 太宽 & $\sigma$ 低估: Fisher 线性化丢了曲率, 或者 (u,v,p) 参数化压抑了某分量 \\
$\sigma_z < 1$ & pull 太窄 & $\sigma$ 高估: Fisher 把不存在的相关性 / 多次散射 / dE/dx 加进去了 \\
\bottomrule
\end{tabular}
\end{table}

\textbf{示例}: 状态报告里修复前的 $p_x$ pull 中位 $\sim -33 / 7.4 \approx -4.4$ 是 (a) 类典型 (frame mismatch); 当前 $p_y$ ratio $1.86$ $\Rightarrow \sigma_z \approx 1.86$ 是 (b) 类典型; 当前 $p_x, p_z$ ratio $\approx 0.8$ $\Rightarrow \sigma_z \approx 0.8$ 是 (c) 类轻症.

\subsection{744 事件 pull 实测}
\label{sec:partIII:pulls:measured}

下表是从 \texttt{track\_summary.csv} 直接计算的 pull 统计 (定义见上, $\sigma$ 取 \texttt{fisher\_*\_sigma} 列):

\begin{table}[H]
\centering\small
\caption{Pull 统计(744 事件, 修复后靶系); $\bar z$ = 均值; $\sigma_z$ = 标准差; $|z|\le 1$ 与 $|z|\le 1.96$ = 实测命中率, 名义为 $0.68$ 与 $0.95$.}
\label{tab:partIII:pull_measured}
\begin{tabular}{lrrrrrl}
\toprule
分量 & $\bar z$ & 中位 $z$ & $\sigma_z$ & $|z|\le 1$ & $|z|\le 1.96$ & 备注 \\
\midrule
$p_x$ & $-0.572$ & $-0.026$ & $4.216$ & $80.1\%$ & $91.0\%$ & 长尾来自 6 条 LM 失败事件 \\
$p_y$ & $-0.178$ & $-0.066$ & $2.477$ & $42.1\%$ & $69.0\%$ & 整体宽 $\sigma_z \approx 2.5$ \\
$p_z$ & $-0.293$ & $-0.068$ & $9.050$ & $77.2\%$ & $90.2\%$ & 长尾更严重(尾部 $|z| > 100$) \\
$|\vec p|$ & $-2.823$ & $-0.066$ & $59.989$ & $76.1\%$ & $90.1\%$ & 同 $p_z$, 长尾事件主导 \\
\bottomrule
\end{tabular}
\end{table}

\paragraph{读法.}
\begin{itemize}[nosep]
  \item \textbf{中位数 $\sim 0$}: 中心估计修复后 (2026-04-19 frame fix) 已经无偏, 没有 frame mismatch 残余. 三分量中位 pull 都 $|z_{\mathrm{med}}| < 0.07$.
  \item \textbf{均值 $\bar z$ 偏负}: 但拖到长尾 (LM 失败事件) 把均值压成 $\sim -0.5$ ($p_x$); $|\vec p|$ 更夸张 $\bar z = -2.8$ 因为 $|\vec p|$ 是\emph{严格非负} 的, LM 跑飞时 $\hat{|\vec p|} \to 1100$ 但 truth $\sim 635$, 这种事件单独贡献 $\Delta z \sim +75$ 但少数; 实际 6 条事件贡献了 95\% 的方差.
  \item \textbf{$\sigma_z (p_y) = 2.48$}: 没有 LM 长尾的"干净"分量, 这个 2.48 跟覆盖率比值 1.86 不完全吻合, 原因是 pull 用的是 \emph{每事件 $\sigma$} 而不是中位 $\sigma$, 而且分布不是高斯; ratio 1.86 是 \emph{鲁棒半宽 / Fisher $\sigma$ 中位}, 数学上不等价, 但定性上同向 (Fisher 低估).
  \item \textbf{$\sigma_z (p_x), \sigma_z (p_z)$ 极大但中位 $|z|\le 1$ 比例 $\sim 0.78$}: 这是 "干净事件占 $93\%$, 6 条 LM 失败事件占 $7\%$" 的双峰分布表现; 鲁棒尺度估计(中位 + IQR) 对这个分布更稳, 也对应表 §误差分析 的覆盖率数字.
\end{itemize}

\paragraph{Pull 的鲁棒尺度估计.}
\label{sec:partIII:pulls:robust}
对 LM 长尾敏感的修正方案: 取 \emph{IQR / 1.349} (高斯下与 $\sigma$ 一致, 对长尾鲁棒) 或者 MAD (median absolute deviation) $\times 1.4826$.

% TODO: 把下表的 IQR/MAD 列也加上,
% 需要在 analyze_ensemble_coverage.py 加 --robust-pull-scale 选项.
\begin{table}[H]
\centering\small
\caption{Pull 的鲁棒尺度估计(待运行, 当前为占位): IQR/1.349 与 $\sigma_z$ 直接对比.}
\label{tab:partIII:pull_robust_scale}
\begin{tabular}{lrrr}
\toprule
分量 & $\sigma_z$ (经典) & IQR/1.349 & 比值 \\
\midrule
$p_x$    & 4.216 & TODO & TODO \\
$p_y$    & 2.477 & TODO & TODO \\
$p_z$    & 9.050 & TODO & TODO \\
$|\vec p|$ & 59.99 & TODO & TODO \\
\bottomrule
\end{tabular}
\end{table}

\subsection{从 CSV 计算 pull 的代码片段}
\label{sec:partIII:pulls:code}

为可复现, 给出从 \texttt{track\_summary.csv} 直接算 pull 的 Python 片段:

\begin{lstlisting}[language=Python,caption={pull 统计的最小复现脚本},label=lst:partIII:pull_python]
import pandas as pd, numpy as np

CSV = ('build/test_output/ensemble_n50_targetframe/'
       'rk_fixed_target_pdc_only_error_analysis/track_summary.csv')
df = pd.read_csv(CSV)

# 残差与 pull
df['truth_p'] = np.sqrt(df.truth_px**2 + df.truth_py**2 + df.truth_pz**2)
for c in ['px', 'py', 'pz', 'p']:
    df[f'pull_{c}'] = (df[f'reco_{c}'] - df[f'truth_{c}']) / df[f'fisher_{c}_sigma']

for c in ['px', 'py', 'pz', 'p']:
    z = df[f'pull_{c}']
    in68 = (z.abs() <= 1.0).mean() * 100
    in95 = (z.abs() <= 1.96).mean() * 100
    print(f'{c}: mean={z.mean():+.3f} median={z.median():+.3f} '
          f'sigma={z.std():.3f}  |z|<=1: {in68:.1f}%  |z|<=1.96: {in95:.1f}%')
\end{lstlisting}

\subsection{Pull 直方图(占位)}
\label{sec:partIII:pulls:histograms}

% TODO: 需要的 PNG, 由 scripts/analysis/plot_pull_histograms.py 生成:
%   figures/diagnostics/pull_distribution_px.png
%   figures/diagnostics/pull_distribution_py.png
%   figures/diagnostics/pull_distribution_pz.png
%   figures/diagnostics/pull_distribution_p.png
% 该脚本未提交; 拟接受 --input track_summary.csv 与 --truth-cut 0.5 (排除 LM 失败事件)
% 双轴显示: 左 全部 744 事件 (含 LM 失败长尾), 右 truth-cut 后干净分布,
% 叠加 N(0,1) PDF.

四张直方图(占位)将分别显示 $p_x, p_y, p_z, |\vec p|$ 的 pull 分布, 与 $\mathcal{N}(0, 1)$ 参考曲线叠加. 期待画面:
\begin{itemize}[nosep]
  \item $p_x$ 分布: 主峰几乎与 $\mathcal{N}(0, 1)$ 重合, 但右上 (LM 失败) 的 6 条事件远远拖出 $|z| > 30$ 的尾巴;
  \item $p_y$ 分布: 主峰宽于 $\mathcal{N}(0, 1)$, $\sigma_z \sim 2.5$;
  \item $p_z$ 分布: 类似 $p_x$ 但长尾更长 (LM 失败时 $p_z$ 跑飞到 $\sim 940$ MeV/$c$);
  \item $|\vec p|$ 分布: 同 $p_z$, 主峰 $\bar z \sim 0$ 但长尾左偏(失败事件 $\hat{|p|} \gg p^{\mathrm{truth}}$).
\end{itemize}

\subsection{每事件 $\sigma$ 与 ensemble 平均 $\sigma$ 的差异}
\label{sec:partIII:pulls:per_event_vs_ensemble}

Pull 定义里 "$\sigma_i$" 是 \emph{每事件的} Fisher $\sigma$, 不是 ensemble 平均. 这有微妙的影响:

\begin{itemize}[nosep]
  \item 如果 Fisher $\sigma_i$ 在事件之间剧烈变化 (例如 $|\vec p|$ Fisher 在 $|p_x|=100$ 时 $\sim 7$, 在病态点 $\sim 60$), 而残差幅度也跟着变大, pull 仍可能保持 $\mathcal{N}(0,1)$;
  \item 反之, 如果 $\sigma_i$ 是常数但残差有 truth-bin 依赖 (例如 dE/dx 在大 $|\vec p|$ 处更大), pull 会有 truth-bin 依赖.
\end{itemize}

换言之, pull 是\emph{标准化残差}的事件级版本. 它对 "Fisher $\sigma$ 是否随事件正确缩放" 是敏感的诊断量.

\paragraph{Per-truth-bin pull 检查.}
% TODO: 把 744 事件按 15 个 truth bin 分组, 每 bin 算 pull mean, std,
% 输出 pull_per_truth_bin.csv. 期望: 14 个非病态 bin 的 pull std 接近 1
% (frame fix 后), $(-100, 0)$ bin 因 12% LM 失败拖出 std $\sim 5$.
% 命令: python3 scripts/analysis/analyze_ensemble_coverage.py \
%         --stratify-by-truth --pull-per-bin

预期结果(从 \nameref{tab:partIII:other_truth_points} ratio 列推算):
\begin{itemize}[nosep]
  \item 14 个非病态 bin: pull std $\sim 0.4$--$1.0$ (大多 Fisher 偏保守, $\sigma_z < 1$);
  \item $(-100, 0)$ bin: pull std $\sim 5$ ($p_x$, $p_z$ 双方向被 LM 失败拖大).
\end{itemize}

\subsection{Pull 与覆盖率的关系}
\label{sec:partIII:pulls:vs_coverage}

二者本质上同源: 一个分量的 68\% 覆盖率 $= \Pr(|z| \le 1)$. 对高斯分布严格等价. 实际数据由于长尾, 覆盖率(对长尾"鲁棒")与 $\sigma_z$ (对长尾敏感) 给出不同信号:

\begin{table}[H]
\centering\small
\caption{覆盖率与 $\sigma_z$ 的对比 (744 事件): 覆盖率对 LM 长尾鲁棒, $\sigma_z$ 不鲁棒.}
\label{tab:partIII:pull_vs_coverage}
\begin{tabular}{lrrr}
\toprule
分量 & 覆盖率 $|z|\le 1$ & $\sigma_z$ & 名义 $\Phi(\sigma_z)$ \\
\midrule
$p_x$ & 0.801 & 4.22 & 0.232 (高斯下) \\
$p_y$ & 0.421 & 2.48 & 0.394 (高斯下) \\
$p_z$ & 0.772 & 9.05 & 0.111 (高斯下) \\
$|\vec p|$ & 0.761 & 60.0 & 0.013 (高斯下) \\
\bottomrule
\end{tabular}
\end{table}

\textbf{解读}: $p_y$ 的 覆盖率 0.42 几乎与 "高斯下 $\sigma_z = 2.48$ 给出的 $\Phi(1/2.48) - \Phi(-1/2.48) = 0.394$" 一致 (差 $\sim 7\%$); 这说明 $p_y$ 主峰是单峰高斯, 但宽度被 Fisher 低估了 $2.48\times$. 反观 $p_x, p_z$, 经典 $\sigma_z$ 由长尾撑大到 $4 \sim 9$, 高斯换算只能给 $0.23 \sim 0.11$ 覆盖率, 但实测 $0.77 \sim 0.80$ —— 差异是因为分布不是高斯, 主峰窄, 长尾稀疏. 鲁棒指标 (中位 + IQR/1.349) 才是 $p_x, p_z$ 的可靠诊断.

\subsection{小结}
\label{sec:partIII:pulls:summary}

\begin{itemize}[nosep]
  \item Pull 是诊断 "误差棒形状是否正确" 的标准量, 期望 $\mathcal{N}(0, 1)$.
  \item 当前 744 事件: 三分量中位 pull $\sim 0$ (frame fix 后无偏), $p_y$ 主峰 $\sigma_z \approx 2.48$ ($\Rightarrow$ Fisher 低估 $\sim 1.86\times$, 与覆盖率 ratio 一致); $p_x, p_z$ 主峰窄但长尾干扰 $\sigma_z$.
  \item 接下来需要: pull 直方图绘图脚本; IQR/MAD 鲁棒尺度; truth-cut 后(去 LM 失败事件) 重新拟合主峰高斯, 应给出与覆盖率一致的 $\sigma_z$.
\end{itemize}

% =============================================================

\section{$\chi^2$ 分布与失败模式分类}
\label{sec:partIII:chi2}

\subsection{自由度 1 的 $\chi^2$ 分布是重尾的}
\label{sec:partIII:chi2:dof1}

RK 重建在 PDC1+PDC2 共两点时, 残差维度 = 4 (两个 PDC 各 $u, v$), 参数维度 = 3 ($u, v, p$ 的 $u$ 是斜率不是丝面坐标), 自由度 $N_{\mathrm{dof}} = 4 - 3 = 1$. 这是状态报告 §\(\chi^2\)~目标函数 写明的事实. 自由度 1 的 $\chi^2$ 分布密度为
\[
  f_{\chi^2_1}(x) \;=\; \frac{1}{\sqrt{2\pi x}} e^{-x/2}, \quad x > 0.
\]
在原点发散($x \to 0$ 时密度 $\to \infty$, 但积分仍收敛), 在尾部以 $e^{-x/2}/\sqrt{x}$ 衰减 —— 比 $\chi^2_2 \sim e^{-x/2}$ 慢. 数值上:

\begin{table}[H]
\centering\small
\caption{$\chi^2_1$ 分布的尾部概率 ($\chi^2_{\mathrm{red}} = \chi^2 / N_{\mathrm{dof}}$, 此处 $N_{\mathrm{dof}} = 1$).}
\label{tab:partIII:chi2_tail_theory}
\begin{tabular}{rrr}
\toprule
$\chi^2_{\mathrm{red}}$ 阈值 & $\Pr(\chi^2_1 > t)$ & $-\log_{10}\Pr$ \\
\midrule
1   & 0.3173 & 0.50 \\
2   & 0.1573 & 0.80 \\
3   & 0.0833 & 1.08 \\
4   & 0.0455 & 1.34 \\
6   & 0.0143 & 1.85 \\
10  & 0.00157 & 2.80 \\
20  & 7.7$\times 10^{-6}$ & 5.11 \\
50  & 1.5$\times 10^{-12}$ & 11.81 \\
\bottomrule
\end{tabular}
\end{table}

\textbf{推论}: 即使在零假设下(模型完美吻合), 1 自由度 $\chi^2_1 > 4$ 的事件仍占 $\sim 4.6\%$, $\chi^2_1 > 10$ 占 $\sim 0.16\%$. 在 744 事件中, 我们 \emph{期望}
\[
  N_{\chi^2 > 4} \approx 33.8, \qquad N_{\chi^2 > 10} \approx 1.17.
\]

\subsection{744 事件实测对比}
\label{sec:partIII:chi2:measured}

直接读 \texttt{track\_summary.csv} 的 \texttt{chi2\_reduced} 列:

\begin{table}[H]
\centering\small
\caption{744 事件的 $\chi^2_{\mathrm{red}}$ 分布观测值 vs.\ 理论 ($\chi^2_1$ 假设). $N_{\mathrm{exp}}$ = 744 $\times$ $\Pr(\chi^2_1 > t)$.}
\label{tab:partIII:chi2_observed}
\begin{tabular}{rrrr}
\toprule
阈值 $t$ & 观测 $N$ & 期望 $N_{\mathrm{exp}}$ & 观测/期望 \\
\midrule
2  & 91 & 117.0 & 0.78 \\
4  & 60 & 33.8  & 1.78 \\
6  & 54 & 10.6  & 5.09 \\
10 & 50 & 1.17  & 42.7 \\
20 & 42 & 0.006 & 7000 \\
\bottomrule
\end{tabular}
\end{table}

\paragraph{读法.}
\begin{itemize}[nosep]
  \item $\chi^2_{\mathrm{red}} > 2$ 出现 91 条 vs.\ 期望 117 —— \emph{低于}期望. 主峰核心拟合得比 1 自由度 $\chi^2_1$ 还紧, 提示 $N_{\mathrm{dof}}$ 可能在 \emph{有效} 上略大于 1 (例如 PDC 测量误差被夸大, 拟合"省下"自由度), 或 LM 把 $\chi^2$ 推得过于贴近 0.
  \item $\chi^2_{\mathrm{red}} > 4$ 之上, 观测开始压倒期望: $> 4$ 已经是 1.78 倍, $> 10$ 是 \textbf{42.7 倍}, $> 20$ 是 7000 倍.
  \item 50 条 $\chi^2_{\mathrm{red}} > 10$ 的事件, 从 7\% 占比看绝对不是 $\chi^2_1$ 自然涨落; 它们是另一个机制产生的 (LM 局部极小, 多次散射爆发, dE/dx 离群, 等等).
\end{itemize}

\subsection{失败模式分类}
\label{sec:partIII:chi2:taxonomy}

把 744 事件按 $\chi^2_{\mathrm{red}}$ 分成三类, 给出每类的代表事件 (从 \texttt{track\_summary.csv} 直接读取):

\begin{table}[H]
\centering\scriptsize
\caption{失败模式分类(按 $\chi^2_{\mathrm{red}}$ 分箱). 数据来自 \texttt{track\_summary.csv} 直接读取; \texttt{event\_index, track\_index} 唯一定位.}
\label{tab:partIII:taxonomy}
\begin{tabular}{lrlllll}
\toprule
类别 & $\chi^2_{\mathrm{red}}$ 范围 & 占比 & 代表事件 & 残差 ($p_x, p_y, p_z$) MeV/$c$ & Fisher $\sigma$ ($p_x, p_y, p_z$) & 物理诠释 \\
\midrule
A 正常 & $[0, 2)$ & $87.8\%$ & (626, 0) & $(-2.49, +0.66, +4.38)$ & $(6.82, 0.64, 6.91)$ & 名义重建, 无明显 dE/dx \\
A 正常 & $[0, 2)$ & --- & (428, 0) & $(+2.36, -0.97, -3.51)$ & $(7.00, 0.75, 6.22)$ & 同上, 不同 truth bin \\
B 中度 & $[2, 10]$ & $5.5\%$ & (482, 0) & $(-4.37, +1.14, +2.38)$ & $(11.16, 1.03, 5.41)$ & dE/dx 致 $p_z$ pull, $p_x$ ratio 大 \\
B 中度 & $[2, 10]$ & --- & (742, 0) & $(-148, +0.56, -549)$ & $(10.58, 1.45, 8.54)$ & 极少见: $p_z$ 跑到 77, LM "刚刚到" \\
C 严重 & $> 10$ & $6.7\%$ & (229, 0) & $(-449, +11.5, +312)$ & $(30.4, 1.98, 19.5)$ & LM trapped at $p_x \to -500$ \\
C 严重 & $> 10$ & --- & (121, 0) & $(-514, +2.49, +335)$ & $(58.5, 2.11, 38.7)$ & $(-100, 0)$ 病态点 \\
\bottomrule
\end{tabular}
\end{table}

\paragraph{Class A — 正常事件.}
$\sim 88\%$ 的事件落在这里. $\chi^2_{\mathrm{red}} \in [0, 2)$, 残差小到 Fisher $\sigma$ 量级, pull 接近 $\mathcal{N}(0, 1)$. 这些事件给出的覆盖率本应为 0.68; 当前 $p_y$ 在 Class A 内部仍欠覆盖, 因为 Class A 内 \emph{每事件} 的 $p_y$ Fisher 低估 $\sim 1.86\times$.

\paragraph{Class B — 中度失配 ($\chi^2_{\mathrm{red}} \in [2, 10]$).}
$\sim 5.5\%$. 残差不再 Fisher $\sigma$ 量级但仍有限. 主要机制:
\begin{enumerate}[label=(\arabic*), nosep]
  \item dE/dx 缺失: $p_z$ 在 1 m 空气段损失 $\sim 0.5$ MeV/$c$, 在 PDC 气体里再损失 $\sim 0.2$ MeV/$c$, 合计 $\sim 0.7$ MeV/$c$; LM 把这个能损"吸收"到 $p_x, p_z$ 的拟合里, $\chi^2$ 增大;
  \item 多次散射在 1 m 空气段超出 Fisher 假设(高斯独立 $\sigma$): 一次大角散射 $\theta \sim 5\sigma_{\mathrm{ms}}$ 让两个 PDC 的 $u$ 残差不再独立, $\chi^2$ 翻倍;
  \item $(u, v, p)$ 参数化的非线性: LM 收敛到 $(u, v, p)$ 空间的最小, 但其在 $(p_x, p_y, p_z)$ 空间未必无偏(参看 Part II §8 详述).
\end{enumerate}

\paragraph{Class C — 严重失败 ($\chi^2_{\mathrm{red}} > 10$).}
$\sim 6.7\%$, 50 条事件. 三种机制:
\begin{enumerate}[label=(\arabic*), nosep]
  \item \textbf{LM 困在错误盆地}: 网格搜索的初值 $(u, v, p)$ 落在 $\chi^2$ 二级最小附近, LM 收敛到 $\hat p_x \sim -500$, $\hat p_z \sim 940$ —— 这是状态报告 §误差分析 里描述的 "(-100, 0) 病态点". 第 \ref{sec:partIII:basin} 章详细剖析.
  \item \textbf{反常 PDC 命中}: PDC 簇心被错误识别(如多径迹串扰), 单条 hit 的 $u$ 偏离 truth $\sim 100$ mm, 残差扭曲, $\chi^2 \to 10^3$. 当前 PDC 模拟未注入此类异常, 但实际数据中确实存在.
  \item \textbf{多次散射爆发}: 1 m 空气段一次硬散射(电磁簇射或大角弹性), 让事件偏离高斯 MS 假设. 当前 Geant4 模拟开了 \texttt{G4UrbanMscModel} 含尾部, 偶然发生概率 $\sim 10^{-3}$.
\end{enumerate}

\subsection{每类事件在覆盖率统计中的贡献}
\label{sec:partIII:chi2:contribution}

\begin{table}[H]
\centering\small
\caption{各类事件在 744 事件覆盖率统计中的贡献(估算).}
\label{tab:partIII:chi2_class_contribution}
\begin{tabular}{lrr}
\toprule
类别 & 事件数 & 对总覆盖率(of $|\vec p|$) 拉低量 \\
\midrule
A 正常 & 654 & $\sim 0$ (本身贴近 0.68) \\
B 中度 & 41  & $\sim 0.005$ (覆盖 $\sim 0.5$, 比 0.68 低 0.18, $\times$ 41/744) \\
C 严重 & 49  & $\sim 0.066$ (覆盖 $\sim 0$, 比 0.68 低 0.68, $\times$ 49/744) \\
\midrule
合计    & 744 & $\sim 0.071$, 即把名义 0.68 拉到 $\sim 0.61$ \\
\bottomrule
\end{tabular}
\end{table}

但是状态报告 $|\vec p|$ 实测覆盖率是 0.76 不是 0.61 —— Fisher $\sigma$ 在 Class C 里 \emph{自适应增大} (例如事件 121 $\sigma_{p_x} = 58.5$, 事件 229 $\sigma_{p_x} = 30.4$), 部分 Class C 事件因 Fisher $\sigma$ 巨大反而被算作 "覆盖". 这是 Fisher 在大 $\chi^2$ 时数值不稳定的副作用 ——它假设 $\chi^2$ 在最优点是高斯 quadratic, 但当 LM 困在错误盆地, "最优点" 在错误位置, Hessian 给出的 $\sigma$ 也在错误尺度上.

\subsection{Wilks's theorem 与 $\chi^2$ 解读边界}
\label{sec:partIII:chi2:wilks}

Wilks (1938) 定理给出: 在零假设下, $-2 \ln (\mathcal{L}_{\mathrm{null}} / \mathcal{L}_{\mathrm{alt}})$ 渐近 $\chi^2_{\Delta k}$ 分布, 其中 $\Delta k$ 是参数维度差. 在我们的问题里:
\begin{itemize}[nosep]
  \item 备择假设 $H_1$: 自由 $(u, v, p)$ 三参数最小化 $\chi^2$, 给 $\chi^2_{\min}$;
  \item 零假设 $H_0$: 模型完美 (truth 已知, 无须拟合), 残差 $\sim \mathcal{N}(0, \sigma_{\mathrm{PDC}})$;
  \item Wilks: $\chi^2_{\min} \sim \chi^2_{N_{\mathrm{res}} - N_{\mathrm{par}}} = \chi^2_1$.
\end{itemize}

\paragraph{Wilks 失效的条件.}
Wilks 假设 (a) 模型在最优点是 regular (Hessian 正定), (b) 数据量大 (asymptotic). 在 $N_{\mathrm{dof}} = 1$ 时, "数据量大" 是个尴尬的条件 — 残差只有 4 个, 参数 3 个, 远未 asymptotic. 但是 PDC 残差是 \emph{独立 4 维高斯}, 在这种全高斯情形 Wilks 即便 $N_{\mathrm{dof}} = 1$ 也精确成立. 这是为什么状态报告中 Class A 的 $\chi^2$ 经验分布与 $\chi^2_1$ 形状一致.

\paragraph{Wilks 在 LM 失败时不成立.}
Class C 事件不满足 Wilks 假设, 因为 (a) LM 不在 $\chi^2$ 全局最小, Hessian 在错盆地内的曲率与全局最小处不同, regular 假设破裂. 这就是为什么 Class C 的 $\chi^2$ 数量级是 $10^3$ 而不是 $\chi^2_1$ 的 $\sim 1$.

\paragraph{p-value 解读.}
对 Class A 主峰内的事件, 把 $\chi^2_{\mathrm{red}}$ 转成 p-value:
\[
  p_{\mathrm{val}} = \Pr(\chi^2_1 > \chi^2_{\mathrm{obs}}).
\]
例如 $\chi^2_{\mathrm{obs}} = 0.5 \Rightarrow p_{\mathrm{val}} = 0.48$ (好); $\chi^2_{\mathrm{obs}} = 4 \Rightarrow p_{\mathrm{val}} = 0.046$ (5\% 水平边界). 在 744 事件中, p-value 小于 $0.05$ 的事件应占 $\sim 5\%$, 即 $\sim 37$ 条; 实测 60 条 $\chi^2 > 4$, 比期望多 $\sim 1.8\times$, 与表 \ref{tab:partIII:chi2_observed} 一致.

\subsection{$\chi^2$ 阈值作为事件级 quality cut}
\label{sec:partIII:chi2:cut}

实用建议: 在最终物理分析里, 加 $\chi^2_{\mathrm{red}} < 5$ 的 quality cut 可以剔除大部分 Class C, 保留 $> 95\%$ 的 Class A+B. 这种 cut 在覆盖率上的效果(模拟):

\begin{table}[H]
\centering\small
\caption{$\chi^2_{\mathrm{red}}$ cut 对覆盖率的影响(估算, 待运行).}
\label{tab:partIII:chi2_cut_estimate}
\begin{tabular}{lrrrr}
\toprule
Cut & 保留事件 & $|\vec p|$ Fisher 68\% & $p_y$ Fisher 68\% & 备注 \\
\midrule
无         & 744 & 0.761 & 0.421 & 当前数字 \\
$< 10$    & $\sim 694$ & TODO & TODO & 应消除 LM 失败 \\
$< 5$     & $\sim 690$ & TODO & TODO & 同上 + 中度 dE/dx 离群 \\
$< 2$     & $\sim 653$ & TODO & TODO & 极保守 \\
\bottomrule
\end{tabular}
\end{table}

% TODO: 上表数字需要在 analyze_ensemble_coverage.py 加 --chi2-max 选项后重跑.
% 命令: python3 scripts/analysis/analyze_ensemble_coverage.py \
%   --input-dir build/test_output/ensemble_n50_targetframe/.../ \
%   --output-dir build/test_output/ensemble_n50_targetframe/coverage_chi2cut5 \
%   --chi2-max 5
% 期望运行时间: 30 秒 (纯 Python 后处理).

\subsection{小结}
\label{sec:partIII:chi2:summary}

\begin{itemize}[nosep]
  \item $\chi^2_1$ 分布尾部重: $> 4$ 的占 $\sim 4.6\%$, $> 10$ 的占 $\sim 0.16\%$ (无机制下).
  \item 实测 744 事件中, $> 10$ 占 $6.7\%$, 比理论高 \emph{42 倍}, 是真正 "另一个机制" 在主导.
  \item 三类机制: LM 错盆地 (主要), 反常 hit, MS 爆发. 对应代表事件已列在表 \ref{tab:partIII:taxonomy}.
  \item $\chi^2$ cut 是简单粗暴但有效的 quality cut, 应在生产分析里默认开启.
\end{itemize}

% =============================================================

\section{单事件深度分析: 6 个示例}
\label{sec:partIII:case_studies}

本章逐条分析 6 个代表事件, 每条事件:
(a) 引用 \texttt{track\_summary.csv} 的原始行;
(b) 标注哪些 Part II §1--9 误差源主导其残差;
(c) 当 \texttt{profile\_samples.csv}/\texttt{bayesian\_samples.csv} 包含该事件时, 引用对应行;
(d) 讨论 LM trace 的预期形态 (实际 trace 数据当前未持久化, 需要 \texttt{PDCErrorAnalysis.cc} 加调试输出).

事件选取覆盖谱为: 标准 Class A (1 条), 中度 dE/dx (1 条), $p_y$ 显著欠覆盖 (1 条), Class C 极端 (1 条 from $(-100, 0)$ 病态点), Class C 中等 (1 条), 边界事件 (1 条).

\subsection{事件 (655, 0) — Class A 标准}
\label{sec:partIII:case_studies:event_655}

\begin{table}[H]
\centering\small
\caption{事件 (655, 0) 数据卡片.}
\begin{tabular}{ll}
\toprule
truth $(p_x, p_y, p_z)$ & $(50.0, 0.0, 627.0)$ MeV/$c$ \\
reco $(p_x, p_y, p_z)$ & $(47.91, -0.67, 630.46)$ MeV/$c$ \\
残差 $(p_x, p_y, p_z)$ & $(-2.09, -0.67, +3.46)$ MeV/$c$ \\
$\chi^2_{\mathrm{red}}$ & $0.000308$ \\
Fisher $\sigma$ $(p_x, p_y, p_z)$ & $(6.82, 0.63, 6.95)$ MeV/$c$ \\
LM 迭代数 & 0 (网格直接命中) \\
Profile 区间 ($p_x$) & $[40.59, 55.23]$, 半宽 $7.32$ \\
Profile 区间 ($p_y$) & $[-1.13, -0.21]$, 半宽 $0.46$ \\
Profile 区间 ($p_z$) & $[623.83, 637.08]$, 半宽 $6.62$ \\
MCMC 区间 ($|\vec p|$) & $[628.19, 637.67]$, 半宽 $4.74$ \\
MCMC acc / ESS / Geweke & $0.355 / 13.8 / 6.31$ \\
\bottomrule
\end{tabular}
\end{table}

\paragraph{诠释.}
\begin{itemize}[nosep]
  \item 残差全部 $< 1\sigma$, 拟合理想; $\chi^2_{\mathrm{red}} < 10^{-3}$ 提示 LM 把残差几乎压到机器精度.
  \item Profile 区间宽度($7.32, 0.46, 6.62$) 比 Fisher 略窄, 与状态报告 §误差分析方法对比表 中 Profile $\sim 0.94 \times$ Fisher 的趋势一致.
  \item MCMC ESS 仅 13.8 (链长 160 步, 80 burn-in), Geweke $z = 6.31$ 远超 $|z| > 2$ 的不收敛阈值 —— 链没有 mix; MCMC 区间数字"碰巧" 与 Profile 接近不代表后验已被恰当采样, 仅说明高斯峰在 $\hat p \sim 633$ 附近且 walker 在峰内随机游走.
  \item LM 迭代数 = 0 意味着粗网格搜索直接给出 $\chi^2 < 10^{-3}$ 的候选, 后续 LM 不需要进一步移动. 这是 Class A 的常见模式.
\end{itemize}

\paragraph{Part II 误差源归因.}
\begin{itemize}[nosep]
  \item PDC 位置分辨 (Part II §1): $\sim 200$ μm 在两个 PDC 上, 通过 Glückstern 给出 $\sigma_p \sim 6$--$7$ MeV/$c$ —— 与 Fisher 一致;
  \item MS in air (Part II §2): $\sim 0.3$ mrad over 1 m, 在 $|\vec p| = 627$ 上贡献 $\sim 1$ MeV/$c$ —— 嵌在 6.8 里;
  \item dE/dx mean (Part II §5): $\sim 0.7$ MeV/$c$ for 1 m air + 2 PDC, 残差 $-2$ MeV/$c$ 与之同号但稍大, 可能配额 $\sim 1$ MeV/$c$;
  \item 其他源: 在此事件中可忽略.
\end{itemize}

\subsection{事件 (482, 0) — Class B 中度失配}
\label{sec:partIII:case_studies:event_482}

\begin{table}[H]
\centering\small
\caption{事件 (482, 0) 数据卡片.}
\begin{tabular}{ll}
\toprule
truth $(p_x, p_y, p_z)$ & $(-100.0, +20.0, 627.0)$ MeV/$c$ \\
reco $(p_x, p_y, p_z)$ & $(-104.37, +21.14, 629.38)$ MeV/$c$ \\
残差 $(p_x, p_y, p_z)$ & $(-4.37, +1.14, +2.38)$ MeV/$c$ \\
$\chi^2_{\mathrm{red}}$ & $2.035$ \\
Fisher $\sigma$ $(p_x, p_y, p_z)$ & $(11.16, 1.03, 5.41)$ MeV/$c$ \\
LM 迭代数 & 0 \\
\bottomrule
\end{tabular}
\end{table}

\paragraph{诠释.}
\begin{itemize}[nosep]
  \item 残差 $/\sigma$ 比为 $(0.39, 1.11, 0.44)$ — $p_y$ 已经在 $1\sigma$ 边缘. 这件 (truth $p_y = 20$) 加上 (truth $p_x = -100$) 提供较高 PDC2 弯曲 $\theta_{xz} \sim -16°$, 把 dE/dx 在弯曲弧上分配的能量损失更多, $p_z$ 残差正向偏 (重建 $p_z$ 大于 truth, 因为 dE/dx 缺失意味着 LM 假设无能损, 把 momentum 分给了 $p_z$).
  \item $\chi^2_{\mathrm{red}} = 2.03$ 处于 $\chi^2_1$ 上 $\sim 16\%$ 的概率密度尾, 是 Class B 的边缘; 但 6.7\% 的 $\chi^2_{\mathrm{red}} > 2$ 占比远大于 16\%, 说明 dE/dx 系统性地压在所有 $|p_x| = 100$ 事件上.
  \item Fisher $\sigma_{p_x} = 11.16$ 略高于平均(其他 $|p_x| = 100$ 事件 $\sigma_{p_x} \sim 7$--$9$), 可能因为 LM 困在略偏离最优的位置, Hessian 在那里的曲率更平.
\end{itemize}

\paragraph{Part II 归因.}
\begin{itemize}[nosep]
  \item dE/dx 缺失 (Part II §5) — 主导, 贡献 $\sim 0.7$ MeV/$c$ 到 $p_z$ 残差;
  \item $(u, v, p)$ 参数化偏 (Part II §8) — 在大 $|p_x|$ 时尤其显著, 贡献 $\sim 1$--$2$ MeV/$c$;
  \item PDC 分辨 + MS — 嵌在 Fisher $\sigma$ 中.
\end{itemize}

\subsection{事件 (102, 0) — $p_y$ 强欠覆盖}
\label{sec:partIII:case_studies:event_102}

\begin{table}[H]
\centering\small
\caption{事件 (102, 0) 数据卡片.}
\begin{tabular}{ll}
\toprule
truth $(p_x, p_y, p_z)$ & $(100.0, -20.0, 627.0)$ MeV/$c$ \\
reco $(p_x, p_y, p_z)$ & $(105.99, -24.36, 624.85)$ MeV/$c$ \\
残差 $(p_x, p_y, p_z)$ & $(+5.99, -4.36, -2.15)$ MeV/$c$ \\
$\chi^2_{\mathrm{red}}$ & $0.133$ \\
Fisher $\sigma$ $(p_x, p_y, p_z)$ & $(6.65, 0.49, 7.19)$ MeV/$c$ \\
$z_{p_y}$ & $-8.94$ \\
Profile $p_y$ 区间 & $[-24.76, -23.96]$, 半宽 $0.40$ \\
MCMC $|\vec p|$ & $[629.03, 642.06]$, ESS $13.5$, Geweke $1.54$ \\
\bottomrule
\end{tabular}
\end{table}

\paragraph{诠释.}
\begin{itemize}[nosep]
  \item $p_y$ 残差 $-4.36$ 是 Fisher $\sigma_{p_y} = 0.49$ 的 \emph{8.9 倍} —— 强欠覆盖典型. 但是 $\chi^2_{\mathrm{red}} = 0.133$ 极低, 说明 LM 把整个残差消耗在 \emph{两个 PDC 同向} 的 $v$ 残差上, $\chi^2$ 没有 "怒"
.
  \item Profile 区间 $[-24.76, -23.96]$ 半宽 $0.40$ 仍然不包含 truth $-20$ —— Profile 与 Fisher 同方向欠覆盖, 这是 Part II §8 描述的 $(u, v, p)$ 线性化偏 $p_y$ 的直接表现.
  \item MCMC ESS 13.5, Geweke $1.54$ 偶然在阈值内 (但仍低 ESS), MCMC 也在 \emph{偏移的}峰附近游走.
\end{itemize}

\paragraph{Part II 归因.}
$p_y$ 残差 $-4.36$ 全部归到 \emph{(u, v, p) 参数化的 Jacobian 压扁} (Part II §8). 这是 "Fisher 信息矩阵在 $(u, v, p)$ 空间到 $(p_x, p_y, p_z)$ 空间转换时, $p_y$ 列 Jacobian $\partial p_y / \partial v$ 没有恢复全部曲率信息". 由于 truth $p_y = -20$ 超出 $\pm \sigma_{p_y}$ 9 倍, LM 找的最优 $v$ 与 truth $v$ 差距远大于线性化下的预期; Fisher $\sigma$ 是从最优点的 Hessian 导, 在最优点附近的曲率不能反映这种远距离偏移.

\subsection{事件 (482, 0) 边界化版本: 事件 (742, 0) — Class B 边界}
\label{sec:partIII:case_studies:event_742}

\begin{table}[H]
\centering\small
\caption{事件 (742, 0) 数据卡片. 这个事件介于 Class B 与 Class C, $\chi^2_{\mathrm{red}} \approx 2$ 但残差极大 — LM 找到了一个"chi-square 看起来满意但物理上荒谬"的解.}
\begin{tabular}{ll}
\toprule
truth $(p_x, p_y, p_z)$ & $(0.0, 0.0, 627.0)$ MeV/$c$ \\
reco $(p_x, p_y, p_z)$ & $(-147.95, +0.56, +77.50)$ MeV/$c$ \\
残差 $(p_x, p_y, p_z)$ & $(-147.95, +0.56, -549.50)$ MeV/$c$ \\
$\chi^2_{\mathrm{red}}$ & $2.040$ \\
Fisher $\sigma$ $(p_x, p_y, p_z)$ & $(10.58, 1.45, 8.54)$ MeV/$c$ \\
LM 迭代数 & 5 \\
\bottomrule
\end{tabular}
\end{table}

\paragraph{诠释.}
\begin{itemize}[nosep]
  \item 残差 $p_z = -549.5$, $|\vec p| = 167$ vs.\ truth 627 —— LM 落在 $|\vec p|$ 极低的局部极小. 但 $\chi^2_{\mathrm{red}} = 2.04$ 看起来还行, 这是\emph{反演问题不唯一}的典型征兆.
  \item Fisher $\sigma$ 在错的位置算: $\sigma_{p_z} = 8.54$ 反映的是 \emph{该 chi-square 谷里} 的二次曲率, 与正确谷($\hat p_z \sim 627$) 的曲率无关.
  \item LM 迭代了 5 次后困在这, 因为粗网格初值落在错盆地. 网格 $u_{\mathrm{center}} = 0 / 627 \to 0$, $v = 0$, $p \in \{200, 400, 700, 1200, 2000\}$ 的某一个 $\Rightarrow$ 这种 truth $(0, 0, 627)$ 居然能困入低 $|\vec p|$ 谷, 提示磁场积分 + RK 可能在某个 $u$ 上有副 $\chi^2$ 极小.
\end{itemize}

\paragraph{Part II 归因.}
\begin{itemize}[nosep]
  \item LM 收敛盆地失败 (Part II §9) — 主导;
  \item 网格搜索初始化不足: 7$\times$3$\times$5 = 105 个候选, 但 $u$ 半幅 1.0 给的网格步长 $0.33$ 在某些 truth 点会跨过 chi-square 鞍点.
\end{itemize}

\subsection{事件 (121, 0) — Class C 极端 ($-100, 0$ 病态)}
\label{sec:partIII:case_studies:event_121}

\begin{table}[H]
\centering\small
\caption{事件 (121, 0) 数据卡片. 这是 \nameref{sec:partIII:basin} 章详细分析的"病态点"代表.}
\begin{tabular}{ll}
\toprule
truth $(p_x, p_y, p_z)$ & $(-100.0, 0.0, 627.0)$ MeV/$c$ \\
reco $(p_x, p_y, p_z)$ & $(-613.70, +2.49, +961.92)$ MeV/$c$ \\
残差 $(p_x, p_y, p_z)$ & $(-513.70, +2.49, +334.92)$ MeV/$c$ \\
$\chi^2_{\mathrm{red}}$ & $1243.89$ \\
Fisher $\sigma$ $(p_x, p_y, p_z)$ & $(58.49, 2.11, 38.66)$ MeV/$c$ \\
LM 迭代数 & 0 (粗网格已经把它推到这) \\
\bottomrule
\end{tabular}
\end{table}

\paragraph{诠释.}
\begin{itemize}[nosep]
  \item 残差 $p_x = -514$ 是 truth 的 $5\times$, $|\vec p| \sim 1145$ vs.\ truth 635. 这是粗网格的某个 $(u, v, p)$ 候选直接给出的 $\chi^2 = 1244$, 而 LM 没有从这里再走 (迭代数 $= 0$ 意味着初值的 $\chi^2$ 在多起点中是最小的).
  \item Fisher $\sigma_{p_x} = 58.5$ 巨大, 因为 Hessian 在这个错盆地里曲率小; 但残差 $-514$ 仍是 $\sigma$ 的 $8.8\times$, pull $z = -8.78$.
  \item 该事件不在 \texttt{profile\_samples.csv} 中(profile 抽样按 $\chi^2$ 分位抽 32 条 $\times$ 4 quartile = 128 条, 该事件 $\chi^2$ 离群 + 不在 quartile 抽样列表里).
\end{itemize}

\paragraph{Part II 归因.}
\begin{itemize}[nosep]
  \item LM 收敛盆地失败 (Part II §9) — 几乎全部;
  \item 第六章详细剖析为什么 truth $(p_x, p_y) = (-100, 0)$ 是网格 $u$ 范围的边界点.
\end{itemize}

\subsection{事件 (229, 0) — Class C 中等}
\label{sec:partIII:case_studies:event_229}

\begin{table}[H]
\centering\small
\caption{事件 (229, 0) 数据卡片. truth $(p_x, p_y) = (-50, 0)$, 跑成了 LM 失败.}
\begin{tabular}{ll}
\toprule
truth $(p_x, p_y, p_z)$ & $(-50.0, 0.0, 627.0)$ MeV/$c$ \\
reco $(p_x, p_y, p_z)$ & $(-499.56, +11.54, +938.68)$ MeV/$c$ \\
残差 $(p_x, p_y, p_z)$ & $(-449.56, +11.54, +311.68)$ MeV/$c$ \\
$\chi^2_{\mathrm{red}}$ & $1548.20$ \\
Fisher $\sigma$ $(p_x, p_y, p_z)$ & $(30.43, 1.98, 19.48)$ MeV/$c$ \\
LM 迭代数 & 0 \\
\bottomrule
\end{tabular}
\end{table}

\paragraph{诠释.}
\begin{itemize}[nosep]
  \item 与事件 (121, 0) 同模式但 truth $(p_x, p_y) = (-50, 0)$. 显示 (-100, 0) 病态点不孤立 — 在 $(-50, 0)$ 也有少数 LM 失败.
  \item Fisher $\sigma$ 在错盆地的曲率 \emph{比} (-100, 0) 病态点 \emph{小}, 因为该 truth 离病态边界更远; 即 $-50$ 的"病态" 程度比 $-100$ 轻.
\end{itemize}

\paragraph{Part II 归因.}
LM 收敛盆地失败 (Part II §9). 但发生频率(50 个 $(-50, 0)$ 事件中只 $\sim 2\%$ 失败) 远低于 $(-100, 0)$ ($\sim 12\%$).

\subsection{Profile/MCMC 跨事件诊断}
\label{sec:partIII:case_studies:profile_mcmc}

把上面 6 个事件的 Profile/MCMC 数据 (能查到的) 与 Fisher 横向对比:

\begin{table}[H]
\centering\scriptsize
\caption{Fisher / Profile / MCMC 三方法在示例事件上的区间宽度对比 (MeV/$c$). 缺失项表示该事件未抽到 Profile/MCMC 子集.}
\label{tab:partIII:case_methods_comparison}
\begin{tabular}{lrrrrrrr}
\toprule
事件 & 分量 & Fisher $\sigma$ & Fisher 半宽(68\%) & Profile 半宽(68\%) & MCMC 半宽(68\%) & MCMC ESS & MCMC Geweke \\
\midrule
\multirow{4}{*}{(655, 0)}
 & $p_x$ & 6.82 & 6.82 & 7.32 & --- & --- & --- \\
 & $p_y$ & 0.63 & 0.63 & 0.46 & --- & --- & --- \\
 & $p_z$ & 6.95 & 6.95 & 6.62 & --- & --- & --- \\
 & $|\vec p|$ & 6.45 & 6.45 & 6.45 & 4.74 & 13.8 & 6.31 \\
\midrule
\multirow{4}{*}{(178, 0)}
 & $p_x$ & 6.35 & 6.35 & 6.75 & --- & --- & --- \\
 & $p_y$ & 0.50 & 0.50 & 0.46 & --- & --- & --- \\
 & $p_z$ & 7.45 & 7.45 & 7.13 & --- & --- & --- \\
 & $|\vec p|$ & --- & --- & --- & 5.08 & 13.6 & 1.83 \\
\midrule
\multirow{4}{*}{(102, 0)}
 & $p_x$ & 6.65 & 6.65 & 5.69 & --- & --- & --- \\
 & $p_y$ & 0.49 & 0.49 & 0.40 & --- & --- & --- \\
 & $p_z$ & 7.19 & 7.19 & 6.07 & --- & --- & --- \\
 & $|\vec p|$ & --- & --- & --- & 6.51 & 13.5 & 1.54 \\
\midrule
\multirow{4}{*}{(298, 0)}
 & $p_x$ & --- & --- & 6.60 & --- & --- & --- \\
 & $p_y$ & --- & --- & 0.50 & --- & --- & --- \\
 & $p_z$ & --- & --- & 6.97 & --- & --- & --- \\
 & $|\vec p|$ & --- & --- & --- & 6.65 & 8.1 & 4.68 \\
\bottomrule
\end{tabular}
\end{table}

\paragraph{读法.}
\begin{itemize}[nosep]
  \item 三方法宽度 \emph{在同一事件内} 差距 $< 25\%$ — 与状态报告 §误差分析 中 ensemble 级的差距一致.
  \item Profile $p_y$ 比 Fisher $p_y$ 略\emph{窄} (0.40--0.50 vs.\ 0.49--0.65) — Profile 沿 $\Delta\chi^2 = 1$ 轮廓走, 而 Fisher 是局部曲率假高斯; 在 $(u,v,p)$ 参数化下两者都低估真实 $\sigma$, 但 Profile 更甚 (因为它沿 chi-square 等高线走得更紧).
  \item MCMC ESS 普遍 $13$ 左右 (链长 160 步, burn-in 80 步), Geweke $|z|$ 多 $> 2$ — 链未 mix, 但区间数字仍然有意义, 因为 walker 仍在主峰内随机游走, 16/84 分位点是经验估计的中心宽度.
\end{itemize}

\paragraph{为什么 MCMC ESS 低?}
状态报告 §误差分析 给出 ensemble 中位 ESS $\approx 12.6$ (\texttt{summary.csv}). 链长 $160 - 80 = 80$ 净样本, ESS $\approx 12$ 意味着 \emph{自相关时间} $\tau \approx 80/12 \approx 7$ 步 — 即每 7 步独立一次, 接收率 0.33 配合 walker step size $\approx \sigma_{\mathrm{Fisher}}$ 是合理的, 但链长太短. 把链长增到 $10^4$ 应该让 ESS $\sim 10^3$, 那时 MCMC 区间才真正 reliable. 这是 Part II 不在范围, 但 Part III 应记录下来作为 next step.

\subsection{所有 6 个事件的对比总览}
\label{sec:partIII:case_studies:overview}

\begin{table}[H]
\centering\scriptsize
\caption{6 个示例事件的全面对比.}
\label{tab:partIII:case_overview}
\begin{tabular}{lrrrrrrl}
\toprule
事件 & truth $(p_x, p_y)$ & 残差 $|p_x|$ & 残差 $|p_y|$ & 残差 $|p_z|$ & $\chi^2_{\mathrm{red}}$ & 类别 & 主因 \\
\midrule
(655, 0) & $(50, 0)$    & 2.1   & 0.7   & 3.5   & 0.0003 & A & 名义 \\
(482, 0) & $(-100, 20)$ & 4.4   & 1.1   & 2.4   & 2.03   & B & dE/dx + (u,v,p) bias \\
(102, 0) & $(100, -20)$ & 6.0   & 4.4   & 2.1   & 0.13   & A* & (u,v,p) $p_y$ 线性化 \\
(742, 0) & $(0, 0)$     & 148   & 0.6   & 549   & 2.04   & B & LM 错盆地 (chi-square 局部) \\
(121, 0) & $(-100, 0)$  & 514   & 2.5   & 335   & 1244   & C & 病态点 \\
(229, 0) & $(-50, 0)$   & 450   & 11.5  & 312   & 1548   & C & 同上, 减弱版 \\
\bottomrule
\end{tabular}
\end{table}
\textit{注: A* = $\chi^2_{\mathrm{red}}$ 落在 Class A 范围, 但 $p_y$ 单分量极端欠覆盖, 是 (u, v, p) 参数化偏的特殊体现.}

\paragraph{LM trace 占位.}
% TODO: 当前 PDCErrorAnalysis.cc 不持久化 LM 的迭代历史 (lambda 变化, 残差变化, 参数轨迹).
% 加入 --dump-lm-trace 选项, 把每条事件的 (iter, lambda, chi2, u, v, p) 写到
% build/test_output/.../lm_traces/event_{N}.csv 即可绘 trace 图.
% 期望对 121, 229, 742 三条 Class C 事件画 (u, p) 平面的 LM 轨迹,
% 应能直观看到它们粘在错盆地, 没有跳出.

\subsection{小结}
\label{sec:partIII:case_studies:summary}

\begin{itemize}[nosep]
  \item 6 个示例覆盖 4 种主要机制: 名义 (A), dE/dx (B), $(u, v, p)$ 线性化偏 (A*), LM 错盆地 (B/C).
  \item 状态报告 ratio 1.86 ($p_y$) 在事件 102 上具体表现为 \emph{残差 $4.4 \times \sigma$} 的极端欠覆盖.
  \item 状态报告 $|\vec p|$ 覆盖率 0.76 在事件 742, 121, 229 上具体表现为 LM 错盆地 + Fisher $\sigma$ 在错盆地里仍给出"看似自洽" 的曲率.
  \item LM trace 数据持久化是接下来的高优先级 dev item.
\end{itemize}

% =============================================================

\section{扫描研究方案与外推}
\label{sec:partIII:scans}

本章给出四组待运行的扫描研究; 大部分是 \emph{方案+外推}, 实际运行未执行. 每个扫描列出: 物理动机, 期望趋势, 待运行命令, 期望耗时, 输出 CSV 路径.

\subsection{扫描 BS-1: 磁场 $B_y$ $\pm 5\%$}
\label{sec:partIII:scans:bs1}

\paragraph{物理动机.}
SAMURAI 偶极磁场图来自实测 + 插值 (\texttt{sm\_field.dat}). 现场实测精度 $\pm 1\%$, 但插值与 RK 步进对中心场强敏感. 扫 $B_y$ 标度 $\pm 5\%$ 给出谱仪对场强标定的灵敏度.

\paragraph{期望趋势.}
RK 重建假设场图正确; 真实场强为 $B (1 + \delta_B)$ 但代码用 $B$, 导致弯曲半径 $\rho = p / (qB)$ 算错 $\delta_B$. 拟合时 LM 把这个吸收到 $\hat p$ 里, 即 $\hat p / p \approx 1 + \delta_B$. 线性外推:
\[
  \hat{|\vec p|} \;\approx\; |\vec p|^{\mathrm{truth}} \cdot (1 + \delta_B), \quad \sigma_{|\vec p|}^{\mathrm{stat}} \to \sigma_{|\vec p|}^{\mathrm{stat}}.
\]
$\sigma$ 不变, 但中心估计偏移. 当前 $|\vec p|$ 中位数 630 MeV/$c$, $\delta_B = +5\%$ 给 $\hat p \approx 661$, $\delta_B = -5\%$ 给 $\hat p \approx 599$.

\paragraph{覆盖率推论.}
\begin{itemize}[nosep]
  \item $\delta_B = +5\%$: 中心偏 $+30$ MeV/$c$, Fisher $\sigma \sim 7$ MeV/$c$, pull mean $\sim 4.3$, 覆盖率 $|\vec p|$ 应跌到 $< 5\%$.
  \item $\delta_B = -5\%$: 镜像, pull mean $\sim -4.3$.
  \item $\delta_B \pm 1\%$: pull mean $\sim 0.9$, 覆盖率从 0.76 跌到 $\sim 0.6$.
\end{itemize}

\paragraph{待运行命令.}
% TODO: B-field scan run BS-1
\begin{lstlisting}[language=bash,caption={B-field scan BS-1 命令模板},label=lst:partIII:bs1_cmd]
# 先把 sm_field.dat 复制成两份缩放副本:
python3 scripts/analysis/scale_sm_field.py --in data/input/sm_field.dat \
    --out data/input/sm_field_scaled_p5pct.dat  --scale 1.05
python3 scripts/analysis/scale_sm_field.py --in data/input/sm_field.dat \
    --out data/input/sm_field_scaled_m5pct.dat  --scale 0.95

for SCALE in p5pct m5pct; do
  ENSEMBLE_SIZE=50 PROFILE_PER_QUARTILE=8 MCMC_PER_QUARTILE=16 \
      MAGNETIC_FIELD=data/input/sm_field_scaled_${SCALE}.dat \
      OUT_ROOT=build/test_output/scan_bs1_${SCALE} \
      bash scripts/analysis/run_ensemble_coverage.sh
done
\end{lstlisting}

\paragraph{期望耗时.} 单 scale $\sim 40$ min ($\times 2 = 80$ min).
\paragraph{期望输出.} \texttt{build/test\_output/scan\_bs1\_p5pct/}, \texttt{scan\_bs1\_m5pct/} 各带完整 ensemble 输出.

\subsection{扫描 WS-1: PDC 丝坐标分辨 $\sigma_{\mathrm{wire}}$}
\label{sec:partIII:scans:ws1}

\paragraph{物理动机.}
PDC 丝坐标分辨当前默认 $\sigma_{\mathrm{wire}} \approx 0.2$ mm (对应 cluster 重心精度). 实测 PDC 数据有 $\sim 0.5$ mm; 极端情形 (cluster 中包含 noise hit) 可到 $\sim 2$ mm. 扫 $\sigma_{\mathrm{wire}} \in \{0.1, 0.3, 0.5, 1.0, 2.0\}$ mm 验证 Glückstern $\sigma_p \propto \sigma_{\mathrm{wire}}$ 关系.

\paragraph{期望趋势.}
Glückstern 公式 (Part I §7) 给出
\[
  \sigma_p \;\propto\; \sigma_{\mathrm{wire}} \cdot |\vec p| / (qBL^2).
\]
当前 $\sigma_p \approx 7$ MeV/$c$ at $\sigma_{\mathrm{wire}} \approx 0.2$ mm. 线性外推:

\begin{table}[H]
\centering\small
\caption{$\sigma_p$ vs.\ $\sigma_{\mathrm{wire}}$ 线性外推.}
\label{tab:partIII:ws1_extrapolate}
\begin{tabular}{rr}
\toprule
$\sigma_{\mathrm{wire}}$ (mm) & $\sigma_p$ (MeV/$c$, 外推) \\
\midrule
0.1 & 3.5 \\
0.3 & 10.5 \\
0.5 & 17.5 \\
1.0 & 35.0 \\
2.0 & 70.0 \\
\bottomrule
\end{tabular}
\end{table}

状态报告 §理论分辨极限 给的 0.5 mm 与 2.0 mm 端点应与上表一致 (待状态报告交叉核对).

\paragraph{待运行命令.}
% TODO: sigma_wire scan run WS-1
\begin{lstlisting}[language=bash,caption={sigma\_wire scan WS-1},label=lst:partIII:ws1_cmd]
for SW in 0.1 0.3 0.5 1.0 2.0; do
  ENSEMBLE_SIZE=50 PROFILE_PER_QUARTILE=8 MCMC_PER_QUARTILE=16 \
      PDC_SIGMA_WIRE_MM=${SW} \
      OUT_ROOT=build/test_output/scan_ws1_sw${SW//./p} \
      bash scripts/analysis/run_ensemble_coverage.sh
done
\end{lstlisting}
该脚本要求 \texttt{run\_ensemble\_coverage.sh} 解析 \texttt{PDC\_SIGMA\_WIRE\_MM} 环境变量并传给 ensemble macro; 当前没有此 hook (% TODO: 在 macro/ensemble template 中加 \\setSigmaWireMm 命令).

\paragraph{期望耗时.} 单 step $\sim 40$ min, $\times 5 = 200$ min ($\sim 3.5$ h).
\paragraph{期望输出.} \texttt{build/test\_output/scan\_ws1\_sw\{0p1,0p3,0p5,1p0,2p0\}/}.

\subsection{扫描 TS-1: 靶倾角 $\alpha_{\mathrm{tgt}}$}
\label{sec:partIII:scans:ts1}

\paragraph{物理动机.}
2026-04-19 修复后, 重建端加入 $R_y(+\alpha_{\mathrm{tgt}})$ (\texttt{PDCFrameRotation.hh}) 把 lab 系结果旋回靶系. 现行配置 $\alpha_{\mathrm{tgt}} = 3°$. 我们要验证:
\begin{enumerate}[label=(\alph*), nosep]
  \item $\alpha_{\mathrm{tgt}} = 0°$ (无旋转): 重建端 frame fix 不应改变物理结果, 即 reco 仍然在 lab = target frame. $p_x$ 残差应仍 $\sim 0$.
  \item $\alpha_{\mathrm{tgt}} = 5°$: $-p_z \sin(5°) \approx -55$ MeV/$c$ 是 truth-vs-reco 的伪偏差, 但 frame fix 把它消除, 残差仍 $\sim 0$.
\end{enumerate}
期待: 三种角度下 $p_x$ 残差半宽不变 (frame 旋转是 \emph{坐标变换}, 无物理影响), 这是 sanity check.

\paragraph{期望趋势.}
\begin{itemize}[nosep]
  \item $\alpha_{\mathrm{tgt}} = 0°, 3°, 5°$: $p_x, p_y, p_z, |\vec p|$ 三种覆盖率应 \emph{完全一致} (在 ensemble 噪声范围内).
  \item 如果 $\alpha_{\mathrm{tgt}} = 5°$ 比 $0°$ 给出不同覆盖率, 说明 PDCFrameRotation 实施有 bug (例如旋转方向错, 或者旋转应用到了部分 PDC1/PDC2 而非全 reco).
\end{itemize}

\paragraph{待运行命令.}
% TODO: tilt scan run TS-1
\begin{lstlisting}[language=bash,caption={tilt scan TS-1},label=lst:partIII:ts1_cmd]
for ALPHA in 0.0 3.0 5.0; do
  ENSEMBLE_SIZE=50 PROFILE_PER_QUARTILE=8 MCMC_PER_QUARTILE=16 \
      TARGET_TILT_DEG=${ALPHA} \
      OUT_ROOT=build/test_output/scan_ts1_alpha${ALPHA//./p} \
      bash scripts/analysis/run_ensemble_coverage.sh
done
\end{lstlisting}

\paragraph{期望耗时.} $3 \times 40 = 120$ min.
\paragraph{期望输出.} 三个 ensemble dir, 比对 \texttt{coverage\_with\_ci.csv} 的 $p_x$ 行.

\subsection{扫描 ES-1: ensemble 大小 $N$}
\label{sec:partIII:scans:es1}

\paragraph{物理动机.}
当前每 truth 点 $N = 50$, 总 $N = 744$ (15 truth bin); 每点 CP $\pm 0.13$ 略嫌粗. 验证生产应取 $N = 200$/truth ($\sim 3000$ 总).

\paragraph{期望趋势.}
覆盖率本身应不变 (ensemble 是 i.i.d. 抽样), 但 CP CI 半宽按 $1 / \sqrt{N}$ 收窄:

\begin{table}[H]
\centering\small
\caption{Ensemble 大小与 CP 半宽 (per truth bin, $\hat p = 0.68$).}
\label{tab:partIII:es1_extrapolate}
\begin{tabular}{rrr}
\toprule
$N$/truth & 总 $N$ & CP 95\% 半宽 \\
\midrule
10  & 150  & $\pm 0.30$ \\
50  & 750  & $\pm 0.13$ (当前) \\
200 & 3000 & $\pm 0.07$ \\
1000 & 15000 & $\pm 0.03$ \\
\bottomrule
\end{tabular}
\end{table}

\paragraph{待运行命令.}
% TODO: ensemble-size scaling run ES-1
\begin{lstlisting}[language=bash,caption={ensemble size scan ES-1},label=lst:partIII:es1_cmd]
for N in 10 50 200 1000; do
  ENSEMBLE_SIZE=${N} PROFILE_PER_QUARTILE=8 MCMC_PER_QUARTILE=16 \
      OUT_ROOT=build/test_output/scan_es1_n${N} \
      bash scripts/analysis/run_ensemble_coverage.sh
done
\end{lstlisting}

\paragraph{期望耗时.} $N=10$ 单核 8 min, $N=50$ 40 min (already done), $N=200$ 160 min, $N=1000$ 800 min ($\sim 13$ h). 应在多核 8 cores 上并行.

\paragraph{期望输出.}
\texttt{build/test\_output/scan\_es1\_n\{10,50,200,1000\}/}; 用同一个 \texttt{analyze\_ensemble\_coverage.py} 做 CP CI 对比.

\subsection{扫描总结}
\label{sec:partIII:scans:summary}

\begin{table}[H]
\centering\small
\caption{扫描研究优先级与状态.}
\label{tab:partIII:scan_priority}
\begin{tabular}{lllll}
\toprule
ID & 扫描 & 优先级 & 状态 & 期望产出 \\
\midrule
BS-1 & $B_y \pm 5\%$        & 中 & 待运行 & 谱仪场强灵敏度 \\
WS-1 & $\sigma_{\mathrm{wire}}$ 0.1--2.0 mm & 高 & 待运行 (需 macro hook) & Glückstern 验证 \\
TS-1 & $\alpha_{\mathrm{tgt}}$ 0--5°  & 中 & 待运行 (需 macro hook) & frame fix sanity \\
ES-1 & $N$ 10--1000           & 高 & 待运行 (无 hook 改动) & 生产 ensemble 标定 \\
\bottomrule
\end{tabular}
\end{table}

% =============================================================

\section{LM 收敛盆地研究: $(-100, 0)$ 病态点剖析}
\label{sec:partIII:basin}

状态报告 §误差分析 在 G4 涨落研究里发现, $(p_x, p_y) = (-100, 0)$ truth 点的 G4 dispersion / Fisher $\sigma$ 比例在 $p_x$ 方向达到 \textbf{3.51}, 在 $p_z$ 方向达到 \textbf{3.51} (同源, 因为是同一组 LM 失败事件主导). 这远高于其他 14 个 truth 点的 $0.5 \sim 1.5$ 范围. 本章详细诊断为什么这个特定 truth 点是 LM 收敛盆地的失败热点, 并提出两条修正路径.

\subsection{从 ensemble 数据直接读取 $(-100, 0)$ 的失败统计}
\label{sec:partIII:basin:numbers}

\begin{table}[H]
\centering\small
\caption{$(p_x, p_y) = (-100, 0)$ truth 点的 50 条事件残差统计 (从 \texttt{ensemble\_pdc\_reco\_merged.csv} 直接计算).}
\label{tab:partIII:basin_pathology_stats}
\begin{tabular}{lr}
\toprule
统计量 & 值 \\
\midrule
$N$ 总事件 & 50 \\
$N$ 残差 $|\Delta p_x| > 50$ MeV/$c$ & 10 (20\%) \\
$N$ 残差 $|\Delta p_x| > 200$ MeV/$c$ & 6 (12\%) \\
$N$ $\hat p_x < -300$ MeV/$c$ & 6 (12\%) \\
失败事件中 $\hat p_x$ 范围 & $[-614, -350]$ \\
失败事件 $\chi^2_{\mathrm{red}}$ 范围 & $[583, 1244]$ \\
失败事件 $\chi^2_{\mathrm{red}}$ 中位 & $\sim 788$ \\
$\Delta p_x$ 中位 (含失败) & $-1.2$ MeV/$c$ \\
$\Delta p_x$ 均值 (含失败) & $-51.3$ MeV/$c$ (失败拖拽) \\
$\Delta p_x$ 标准差 (含失败) & $124.9$ MeV/$c$ \\
$\Delta p_x$ p84 - p16 半宽 (鲁棒) & $\sim 7$ MeV/$c$ \\
\bottomrule
\end{tabular}
\end{table}

\paragraph{诠释.}
\begin{itemize}[nosep]
  \item 失败事件占 12\% — 远高于其他 truth 点 (0--4\%).
  \item 失败事件 $\hat p_x \in [-614, -350]$, 即 LM 把 truth $-100$ 误认成 $\sim -500$ 这个量级. 残差 $\sim -400$ MeV/$c$ 是 truth 的 4 倍.
  \item 中位残差仍 $\sim 0$ — 主峰是好的, 失败事件是 \emph{拖尾}.
  \item Fisher $\sigma$ 在主峰 $\sim 7$ MeV/$c$ 但失败事件 $\sigma$ 跳到 $\sim 60$ — Fisher 在错盆地"刻舟求剑".
\end{itemize}

\subsection{为什么是 $(-100, 0)$?}
\label{sec:partIII:basin:why}

\paragraph{网格搜索的几何.}
\texttt{PDCRkAnalysisInternal.cc} 第 392--424 行的 grid search 用
\begin{itemize}[nosep]
  \item $u$ (= $p_x / p_z$ slope at target) 网格: $u \in [u_{\mathrm{line}} - 1, u_{\mathrm{line}} + 1]$, 7 步, 步长 $\sim 0.33$;
  \item $v$ (= $p_y / p_z$): $v \in [v_{\mathrm{line}} - 0.3, v_{\mathrm{line}} + 0.3]$, 3 步, 步长 $0.3$;
  \item $|\vec p|$: $\{200, 400, 700, 1200, 2000\}$ MeV/$c$ 5 个 logspace-like 候选.
\end{itemize}
其中 $u_{\mathrm{line}}, v_{\mathrm{line}}$ 是 \emph{靶到最远 PDC 的直线方向} (\texttt{init\_dir}), 即不考虑磁场偏转的几何视线方向.

\paragraph{$(-100, 0)$ 的 $u_{\mathrm{line}}$.}
truth $p = (-100, 0, 627)$, 出射方向 $u = -100/627 \approx -0.160$. 但磁场把质子在 PDC2 处弯了 $\sim 25°$ ($\theta_{xz} \sim 0.43$ rad), PDC2 真实位置在 $u \approx -0.43 + (-0.16) = -0.59$ 附近. 因此粗网格中心 $u_{\mathrm{line}} \approx -0.59$ (由 PDC2 几何定义, 不是 truth 出射 slope).

实际计算: PDC2 在 $z \approx 4655$ mm 处, 该事件 PDC2 命中 $x \approx -2745$ mm (由表 \nameref{sec:partIII:basin:numbers} 隐含); 直线 slope $u_{\mathrm{line}} = x_{\mathrm{PDC2}} / z_{\mathrm{PDC2}} \approx -2745 / 4655 \approx -0.59$. 网格 $u \in [-1.59, 0.41]$ —— truth $u = -0.16$ 在网格中心偏右 ($+0.43$ 处, 第 4 个网格点).

\paragraph{为什么是边界?}
truth $u = -0.16$ 不是在 $u_{\mathrm{line}} \pm 1$ 网格的 \emph{物理边界}, 但它\emph{接近 chi-square 第二低盆地} 的边界. 状态报告 $§误差分析$ 中 G4 涨落的 5$\times$3 truth 网格图(图 \texttt{figures/g4\_per\_truth\_box\_*.png}) 显示在 $(-100, 0)$ 这个点 $\hat p_x$ 分布是双峰: 主峰 $\hat p_x \approx -100$ 包含 88\% 事件, 副峰 $\hat p_x \approx -500$ 包含 12\% 事件. 副峰的位置对应 LM "卡在错盆地" 的 chi-square 局部极小.

\paragraph{为什么 $(-100, 20)$ 没那么严重?}
truth $(p_x, p_y) = (-100, 20)$ 的 50 事件中 $\Delta p_x$ 半宽 $\sim 8.6$ MeV/$c$ (从 \texttt{per\_truth\_point\_g4\_vs\_fisher.csv}: g4\_halfw68\_px = 8.64), 比 $(-100, 0)$ 的 $34.8$ 要小 4 倍. 这说明 $p_y \ne 0$ 引入的 $v = p_y / p_z$ 偏离对 LM "推开错盆地" 起到了关键作用 —— 网格 $v$ 步长 0.3 在 $p_y = 0$ 时正好覆盖错盆地, 在 $p_y = \pm 20$ 时把网格推到错盆地外, LM 看不到副峰.

\subsection{Chi-square 表面拓扑示意}
\label{sec:partIII:basin:chi2_topology}

% TODO: 需要画一张 (u, p) 平面 chi-square 等高线图,
% 由 scripts/analysis/plot_chi2_landscape.py 生成 (待写).
% 输入: PDC1, PDC2, target geometry; truth (p_x, p_y, p_z) = (-100, 0, 627).
% 输出: figures/diagnostics/chi2_landscape_p100_p0.png.
% 应能直观看到两个谷(主谷 $u \sim -0.16$, 副谷 $u \sim -0.7 \sim -1.0$)
% 以及它们之间的鞍点, 标注网格采样点.

期望画面:
\begin{itemize}[nosep]
  \item 主谷 (truth 谷): $(u, p) \sim (-0.16, 627)$, $\chi^2 \sim 0$;
  \item 副谷 (病态谷): $(u, p) \sim (-0.85, 1145)$, $\chi^2 \sim 1244$ — 注意 $\chi^2$ 不为零 \emph{是因为} 该副谷是 chi-square 非零的局部极小, 不是真极小;
  \item 鞍点连接两谷, $u_{\mathrm{saddle}} \sim -0.5$;
  \item 网格点 7$\times$5 = 35 个候选 ($v$ 维度暂不画), 7 个 $u$ 候选中 4 个在主谷一侧, 3 个在副谷一侧.
\end{itemize}

\subsection{修复方案 (a): 网格不对称延展}
\label{sec:partIII:basin:fix_a}

\paragraph{思路.}
$|u_{\mathrm{line}}|$ 大时, 主谷 $u_{\mathrm{truth}}$ 与 $u_{\mathrm{line}}$ 距离 $\sim 0.4$; 副谷 $u_{\mathrm{副}}$ 与 $u_{\mathrm{line}}$ 距离也 $\sim 0.3$ (在另一侧). 当前网格 $\pm 1$ 对称延展, 副谷一定能命中. 改成 \emph{不对称} 延展, 让网格偏向 truth 谷:
\[
  u_{\min} = u_{\mathrm{line}} - 0.3, \quad u_{\max} = u_{\mathrm{line}} + 1.0.
\]
即把网格 \emph{向更靠近 truth 出射方向} 偏置 (更小的 $|u|$). 但这要求知道 truth 出射方向的先验, 不可行.

\paragraph{改进思路.}
基于物理: 弯曲方向永远把 $u_{\mathrm{line}}$ 推得比 $u_{\mathrm{truth}}$ \emph{绝对值更大} (磁场把出射 deflect 到更陡的弧). 因此网格 \emph{向 $|u|$ 更小的方向} 偏置:
\[
  u_{\min} = u_{\mathrm{line}} - 0.3 \cdot \mathrm{sign}(u_{\mathrm{line}}), \quad u_{\max} = u_{\mathrm{line}} + 1.0 \cdot \mathrm{sign}(u_{\mathrm{line}}) \cdot (-1).
\]
等价地: 网格半幅 $\pm 1$, 但 truth-side 半幅 $-1.0$, 远 truth-side 半幅 $+0.3$.

\paragraph{代码改动.}
\texttt{libs/analysis\_pdc\_reco/src/PDCRkAnalysisInternal.cc} 第 392--410 行:
\begin{lstlisting}[caption={网格不对称延展 (建议改动)}]
constexpr double kUFarHalfRange  = 1.0;   // away from target
constexpr double kUNearHalfRange = 0.3;   // toward target slope
double u_lo = u_center - (u_center >= 0 ? kUNearHalfRange : kUFarHalfRange);
double u_hi = u_center + (u_center >= 0 ? kUFarHalfRange  : kUNearHalfRange);
for (int iu = 0; iu < kGridU; ++iu) {
    const double u = u_lo + (u_hi - u_lo) * iu / (kGridU - 1);
    ...
}
\end{lstlisting}

\paragraph{预期效果.}
$(-100, 0)$ truth 点 $u_{\mathrm{line}} \approx -0.59$, 改后网格 $u \in [-0.89, +0.41]$ (truth $-0.16$ 在中心偏右). 副谷在 $u \approx -0.85$ 仍在网格内, 但 truth 谷被 4 个网格点覆盖, 副谷只 1 个. 正确盆地会以多起点压倒错盆地.

\subsection{修复方案 (b): 加入 "magnetic-kick-corrected" 起点}
\label{sec:partIII:basin:fix_b}

\paragraph{思路.}
Linear backward 起点: 用 PDC1, PDC2 做反向直线外推到 target $z$, 得到 $u_{\mathrm{back}} = (x_{\mathrm{PDC1}} - x_{\mathrm{tgt}}) / (z_{\mathrm{PDC1}} - z_{\mathrm{tgt}})$. 这忽略磁场, 但可以从这个起点开始一次 $|\vec p|$ 灵敏度估计:

设质子在 dipole 中弯过弧长 $L$, 弯曲角 $\theta = qBL / p$ (Highland-style). 用 PDC2 处的 incidence angle $\theta_{\mathrm{in}}$ 减去 $\theta$ 即得 emission angle, 再外推到 target 给出 \emph{magnetic-kick corrected} $u_{\mathrm{corr}}$.

实际计算更简单: PDC1 和 PDC2 的角度差 $\Delta\theta_{\mathrm{PDCs}} = \mathrm{atan2}(\Delta x, \Delta z)$ 给出 \emph{出 PDC 后} 的轨迹方向; 把它向 target 外推 (无场), 得到的 slope 就是 \emph{出射的} $u$, 再加上一个磁场修正项 $\sim qB\Delta z / p$, 给出 \emph{在 target 的} $u$.

\paragraph{代码改动.}
新加一个候选 $u$ 候选 (替代或补充网格中心):
\begin{lstlisting}[caption={magnetic-kick-corrected 起点 (建议改动)}]
double u_pdc_to_pdc = (track.pdc2.X() - track.pdc1.X()) /
                      (track.pdc2.Z() - track.pdc1.Z());
double dz_pdc1_target = track.pdc1.Z() - target.target_position.Z();
// magnetic kick estimate (assume qBy = -0.5 T, p = 600 MeV/c => 0.4 rad/m)
double magnetic_correction = 0.5 * dz_pdc1_target / 1500.0;  // approx
double u_back = u_pdc_to_pdc - magnetic_correction;

// add as extra candidate
candidates.push_back({{u_back, v_back, 600.0}, Evaluate(...).chi2_raw});
\end{lstlisting}

\paragraph{预期效果.}
$(-100, 0)$ 的 $u_{\mathrm{back}}$ 接近 $-0.16$ (因为 PDC1, PDC2 之间的 slope 反映出 $\hat p$ 而非 $u_{\mathrm{line}}$); 加进 candidates 后, LM 多起点必有一个落到主谷.

\subsection{第三种方案 (c): 一阶 backward Kalman}
\label{sec:partIII:basin:fix_c}

更彻底的方案: 用 Kalman 反向滤波 — 从 PDC2 倒推到 PDC1 倒推到 target, 每步用 RK4 反向积分, 得到一个相对精确的 target-frame 起点. 但这基本等价于做\emph{第二次拟合}, 工作量与替换 LM 等同, 暂不优先.

\subsection{修复后预期 ratio}
\label{sec:partIII:basin:fix_expected}

\begin{table}[H]
\centering\small
\caption{修复方案对 $(-100, 0)$ truth 点 $\Delta p_x$ 半宽的预期影响.}
\label{tab:partIII:basin_fix_table}
\begin{tabular}{lr}
\toprule
方案 & 预期 $\Delta p_x$ 半宽 (MeV/$c$) \\
\midrule
当前 & 34.8 (主峰 $\sim 5$ + 12\% 失败 $\sim 400$) \\
方案 (a) 网格不对称 & $\sim 8$ (失败率应降到 0--2\%) \\
方案 (b) magnetic-kick 起点 & $\sim 5$ (失败率 $\sim 0$, 但会引入新 LM-trap 风险) \\
方案 (a) + (b) 联合 & $\sim 5$ (主峰 + 极少失败) \\
\bottomrule
\end{tabular}
\end{table}

\subsection{副谷的物理来源: 反方向轨迹}
\label{sec:partIII:basin:reverse_trajectory}

副谷 $\hat p_x \approx -500$ 是什么物理? 一种解释是: 给定 PDC1, PDC2 命中, 存在两个不同的 (target 出射, 总动量) 配置都能解释这两个点 — 一个是 truth ($\hat p_x = -100$, $\hat p_z = 627$), 另一个是高动量反向 ($\hat p_x = -500$, $\hat p_z = 940$).

\paragraph{推导.}
质子在磁场 $B_y$ 中, 弯曲半径 $\rho = p / (q B)$. 给定 PDC1, PDC2 两点和 target, 拟合即解三圆心方程 — 圆经过三个点. 但 \emph{磁场不均匀} (SAMURAI dipole 有 fringe field), 所以"圆经过三点"不严格; 不同 $|\vec p|$ 的 RK 积分可以经过相同的 PDC2 但 source 出射方向不同.

具体: 设 truth 解给出 $u_{\mathrm{truth}} = -0.16$ at target. 副谷解给出 $u_{\mathrm{副}} \approx -0.85$ at target — 方向更陡, 但 \emph{动量更高} 让弯曲变浅, 弯过 dipole 后到达 PDC2 的位置接近. 这是磁场+动量的二维耦合带来的 chi-square 多解.

\paragraph{$p_y = 0$ 的特殊性.}
为什么 $p_y \ne 0$ 时副谷消失? 因为 $v = p_y/p_z$ 为非零时, $v$ 网格 $\{v_{\mathrm{line}} - 0.3, v_{\mathrm{line}}, v_{\mathrm{line}} + 0.3\}$ 三个候选与 truth $v$ 距离都不同; 副谷需要特定的 $(u_{\mathrm{副}}, v_{\mathrm{副}})$ 组合, $v_{\mathrm{副}} \ne 0$ 时网格不再覆盖. 在 $p_y = 0$, truth $v = 0$, $v_{\mathrm{line}} = 0$, 网格中央点正好是 truth $v$ — 但同时也在副谷的 $v$ 上 (因为副谷沿 $u$ 维度延伸, $v$ 中心仍 $\sim 0$).

\subsection{为什么 $-100, 0$ 但不是 $+100, 0$?}
\label{sec:partIII:basin:asymmetry}

表 \ref{tab:partIII:other_truth_points} 显示 $(+100, 0)$ ratio 0.57 (健康). 不对称性来自:

\begin{itemize}[nosep]
  \item SAMURAI dipole $B_y$ 极性把正电荷向 $+x$ 方向偏, 即 deuteron / proton 的 $u_{\mathrm{line}}$ 偏 $+u$. truth $p_x = -100$ 的事件出射朝 $-x$, 但 magnetic kick 把它\emph{反向}弯到 $+x$ — 即 PDC2 命中位置在 $+x$ 处, $u_{\mathrm{line}} > 0$. 但实测 (\texttt{pdc\_dispersion\_summary.csv}) 显示 $(p_x = -100, p_y = 0)$ 的 PDC2 中位 $x \approx -3500$ mm, 仍在 $-x$;
  \item 这就需要重新检查: $u_{\mathrm{line}} = -3500/4655 \approx -0.75$, 网格 $u \in [-1.75, 0.25]$. truth $u = -0.16$ 在网格中点偏右. 副谷 $u_{\mathrm{副}} \approx -0.85$ 在网格中央. 即副谷正好被网格"瞄准";
  \item truth $p_x = +100$: 出射朝 $+x$, magnetic kick 反向弯, PDC2 命中 $\sim -2900$ mm (从 \texttt{pdc\_dispersion\_summary.csv} 推断), $u_{\mathrm{line}} \approx -0.62$, 网格 $u \in [-1.62, 0.38]$, truth $u = +0.16$ 在网格右端边缘 — \emph{副谷} 在 $u \approx -0.6$ 也在网格中央, 但 truth $u$ 处于边缘也意味着 LM 容易跑出主谷, 走向副谷的几率反而小. 经验上 ratio 0.57 是健康的;
  \item 物理结论: 网格设计在 \emph{$u_{\mathrm{line}}$ 与 $u_{\mathrm{truth}}$ 同号} 时(主谷与副谷都在网格内, 副谷更靠近中心) 容易失败.
\end{itemize}

\subsection{监控诊断: LM trace 持久化}
\label{sec:partIII:basin:lm_trace}

为以后追踪 LM 失败, 应在 \texttt{PDCErrorAnalysis.cc} 加 \texttt{--dump-lm-trace} 选项, 把每条事件的:
\begin{itemize}[nosep]
  \item 网格搜索的 105 个候选 $(u, v, p, \chi^2)$;
  \item 多起点 LM 的 $\le 5$ 条迭代历史 $(\mathrm{iter}, \lambda, \chi^2, u, v, p)$;
  \item 最终选取的最优点;
\end{itemize}
持久化到 \texttt{lm\_traces/event\_\{N\}.csv}. 这给后续 visualization 提供数据基础, 也允许在数据 (815 tracks, no truth) 上识别 LM 失败 (无 truth 时只能看 $\chi^2$ + LM trace, 不能看残差).

% TODO: PDCErrorAnalysis.cc 加 --dump-lm-trace 选项 (作 dev item).

\subsection{其他可疑 truth 点}
\label{sec:partIII:basin:other_points}

从 \texttt{per\_truth\_point\_g4\_vs\_fisher.csv} 直接读出 ratio (g4 / fisher) :

\begin{table}[H]
\centering\small
\caption{各 truth 点的 g4 半宽 / Fisher $\sigma$ 比例 (从 \texttt{per\_truth\_point\_g4\_vs\_fisher.csv}, $p_x$ 列). $p_y \in \{-20, 0, 20\}$ 对应三行.}
\label{tab:partIII:other_truth_points}
\begin{tabular}{rrrr}
\toprule
truth $p_x$ & ratio $p_y = -20$ & ratio $p_y = 0$ & ratio $p_y = 20$ \\
\midrule
$-100$ & 0.48 & \textbf{3.51} & 0.77 \\
$-50$  & 0.40 & 0.38 & 0.66 \\
$0$    & 0.84 & 0.63 & 0.89 \\
$+50$  & 0.51 & 0.56 & 0.72 \\
$+100$ & 0.62 & 0.57 & 0.96 \\
\bottomrule
\end{tabular}
\end{table}

\paragraph{读法.}
\begin{itemize}[nosep]
  \item $(-100, 0)$ 是唯一 ratio $> 1$ 的 truth 点; 失败率 12\% 是孤立病态.
  \item 其他 14 个 truth 点 ratio 都在 $0.4$--$1.0$ 范围, 基本健康.
  \item ratio $< 1$ 反映 \emph{Fisher 略保守}, 这是 PDC 高斯分辨条件下的预期; 主要担忧是 $> 1$ 的方向.
\end{itemize}

\subsection{小结}
\label{sec:partIII:basin:summary}

\begin{itemize}[nosep]
  \item $(-100, 0)$ 是 LM 收敛盆地失败的孤立热点, 12\% 事件 trap 在 $\hat p_x \approx -500$ 副谷.
  \item 物理原因: $u_{\mathrm{line}}$ 距 $u_{\mathrm{truth}}$ 远, 网格对称延展正好覆盖错盆地; $p_y = 0$ 时 $v$ 网格不能推开错盆地.
  \item 修复: (a) 网格不对称延展 + (b) magnetic-kick-corrected 起点, 联合预期把 ratio 从 3.51 降到 $< 1$.
  \item LM trace 持久化是接下来追踪 LM 病态的高优先级 dev item.
\end{itemize}

% =============================================================

\subsection*{Part III 全章总结}

本部分给出了 RK 误差分析的实证诊断完整框架, 从二项 CI 的统计学基础到事件级 LM trace 的具体追踪.

\textbf{六章主要结论.}
\begin{enumerate}[label=\arabic*., nosep]
  \item \textbf{覆盖率方法学}: Clopper-Pearson 是当前 default, 比 Wilson 略保守, 比 Wald 在端点正确; $N=744$ 给 CP $\pm 0.034$ 足以诊断当前的 $p_y$ 欠覆盖, 生产 ensemble 应取 $N \ge 200$/truth.
  \item \textbf{Pull 分布}: 实测 744 事件 pull 中位 $\sim 0$ (frame fix 后无偏), $p_y$ 主峰 $\sigma_z \approx 2.48$ (Fisher 低估 $\sim 1.86\times$), $p_x, p_z$ 主峰窄但有 LM 长尾干扰; 鲁棒尺度估计 IQR/1.349 是接下来要加的指标.
  \item \textbf{$\chi^2$ 分布与失败模式}: 50 条 ($6.7\%$) $\chi^2_{\mathrm{red}} > 10$ 远超 $\chi^2_1$ 自然涨落 ($0.16\%$, $42\times$); 三类机制 (LM 错盆地, 反常 hit, MS 爆发); $\chi^2 < 5$ quality cut 应在生产分析里默认开启.
  \item \textbf{6 个示例事件}: 涵盖 4 种主要机制. 事件 102 $z_{p_y} = -8.94$ 是 (u, v, p) 偏典型, 事件 121, 229 是 (-100, 0) 病态典型, 事件 742 是 LM 错盆地中等表现.
  \item \textbf{扫描研究方案}: BS-1 (B 场), WS-1 ($\sigma_{\mathrm{wire}}$), TS-1 (靶倾角 sanity), ES-1 (ensemble 规模) 已列出命令模板与外推; ES-1 优先级最高.
  \item \textbf{$(-100, 0)$ 病态点剖析}: 12\% 事件 trap 在副谷, 修复方案 (a) 网格不对称延展 + (b) magnetic-kick-corrected 起点, 联合预期把 ratio 从 3.51 降到 $< 1$.
\end{enumerate}

\textbf{下一阶段优先 dev items.}
\begin{itemize}[nosep]
  \item LM trace 持久化 (\texttt{PDCErrorAnalysis.cc} \texttt{--dump-lm-trace} 选项);
  \item Pull 直方图绘图脚本 + IQR/MAD 鲁棒尺度;
  \item Ensemble script 加 \texttt{PDC\_SIGMA\_WIRE\_MM}, \texttt{TARGET\_TILT\_DEG}, \texttt{MAGNETIC\_FIELD} 环境变量 hook;
  \item \texttt{analyze\_ensemble\_coverage.py} 加 \texttt{--chi2-max} 选项支持 quality cut;
  \item 网格搜索不对称延展 + magnetic-kick 起点的 RK 实施 + ensemble 验证.
\end{itemize}

\textbf{与 Part I, Part II 的连接.}
\begin{itemize}[nosep]
  \item 第 \ref{sec:partIII:coverage_methodology} 章的 CP CI 推导与 Part I §2 (Cramér-Rao) 的 ensemble 验证范式互补;
  \item 第 \ref{sec:partIII:pulls} 章的 pull = $\mathcal{N}(0, 1)$ 期望与 Part I §3 (Fisher) 的协方差一致性紧密相关;
  \item 第 \ref{sec:partIII:chi2} 章的失败模式分类与 Part II §1--9 的误差源对应表完全对齐;
  \item 第 \ref{sec:partIII:case_studies} 章的 6 个示例事件直接验证了 Part II §5 (dE/dx), §8 ((u,v,p) 参数化偏), §9 (LM 收敛盆地);
  \item 第 \ref{sec:partIII:scans} 章的扫描计划是 Part II 各源的 sensitivity 验证;
  \item 第 \ref{sec:partIII:basin} 章 (-100, 0) 剖析是 Part II §9 的具体落地.
\end{itemize}

% =============================================================
% End of Part III.
% =============================================================


\newpage

% ============================================================
% 附录
% ============================================================
\appendix

\section{复现说明}
\label{app:reproduce}

\subsection{环境固定}
\begin{itemize}[nosep]
  \item 仓库：\texttt{smsimulator5.5}，分支 \texttt{main}（撰写时 commit
        \texttt{e363e4b} 之后）。
  \item Python：\texttt{micromamba activate anaroot-env}。
  \item Geant4：与 SAMURAI 主仓约定一致；磁场表
        \texttt{configs/simulation/.../180703-1,15T-3000.table}。
  \item 几何宏：
        \texttt{configs/simulation/DbeamTest/detailMag1to1.2T/geometry\_B115T\_pdcOptimized\_20260227\_target3deg.mac}。
\end{itemize}

\subsection{主要 CSV / 数据来源}
\label{app:csvs}

本报告引用的 CSV 在以下两个根目录下；详细 schema 见状态报告 § 数据来源与复现脚本。
\begin{itemize}[nosep]
  \item \texttt{\DeepDiveData/} —— 744 事件 ensemble，靶系，2026-04-19 修复后。
        子目录 \texttt{rk\_fixed\_target\_pdc\_only\_error\_analysis/} 下：
        \texttt{summary.csv}、\texttt{track\_summary.csv}（744 行）、
        \texttt{profile\_samples.csv}（128 行）、
        \texttt{bayesian\_samples.csv}（128 行）、
        \texttt{validation\_summary.csv}。
        子目录 \texttt{g4\_fluctuation/} 下：
        \texttt{ensemble\_pdc\_reco\_merged.csv}、
        \texttt{pdc\_dispersion\_summary.csv}、
        \texttt{reco\_momentum\_dispersion\_summary.csv}、
        \texttt{per\_truth\_point\_g4\_vs\_fisher.csv}。
  \item \texttt{\AnaRoot/} —— 815 track，缺 truth；用于状态报告
        method comparison §对比。本报告不引用其覆盖率。
\end{itemize}

\subsection{主管线命令}
\begin{lstlisting}[language=bash]
# Build (incremental)
micromamba activate anaroot-env
mkdir -p build && cd build
cmake .. -DCMAKE_BUILD_TYPE=Release -DBUILD_TESTS=ON \
         -DBUILD_APPS=ON -DBUILD_ANALYSIS=ON -DWITH_ANAROOT=ON
make -j$(nproc)

# 744-event ensemble (sim + reco + 误差分析,单核约 40 分钟)
ENSEMBLE_SIZE=50 PROFILE_PER_QUARTILE=8 MCMC_PER_QUARTILE=16 \
    OUT_ROOT=build/test_output/ensemble_n50_targetframe \
    bash scripts/analysis/run_ensemble_coverage.sh

# 覆盖率聚合
python3 scripts/analysis/analyze_ensemble_coverage.py \
    --input-dir  build/test_output/ensemble_n50_targetframe/rk_fixed_target_pdc_only_error_analysis \
    --output-dir build/test_output/ensemble_n50_targetframe/coverage_final
\end{lstlisting}

\section{符号 \(\to\) 代码标识对照（glossary）}
\label{app:glossary}

\begin{table}[H]
\centering\footnotesize
\caption{本报告主要数学符号与代码标识对照}
\label{tab:glossary}
\begin{tabular}{lll}
\toprule
符号 & 含义 & 代码标识 \\
\midrule
$\vec\alpha$ & 拟合参数矢量 & \texttt{RkParameters::data} (PDCRecoTypes.hh) \\
$u, v$ & 出射方向正切 $p_x/p_z, p_y/p_z$ & \texttt{params.u}, \texttt{params.v} \\
$p$ & 动量模 $|\vec p|$ & \texttt{params.p} \\
$\delta x, \delta y$ & 靶点位置偏移 (mm) & \texttt{params.dx}, \texttt{dy} \\
$\vec p$ & 物理动量 $(p_x, p_y, p_z)$ & \texttt{recoP4.Vect()} \\
$\hat u, \hat v$ & PDC 丝方向单位矢量 & \texttt{PDCSimAna::WireDir} \\
$\theta$ & PDC 倾斜角 (57$^\circ$) & \texttt{kPdcTiltAngle} \\
$\alpha_\mathrm{tgt}$ & 靶倾角 (3$^\circ$) & \texttt{kTargetTiltDeg} (PDCFrameRotation.hh) \\
$r_i$ & 原始残差 & \texttt{computeRawResiduals} \\
$\vec w$ & 白化残差 & \texttt{computeWhitenedResiduals} \\
$L$ & Cholesky 因子 (白化矩阵) & \texttt{whiteningCholesky} \\
$J$ & 残差 Jacobian & \texttt{computeJacobian} (PDCRkAnalysisInternal.cc) \\
$\mathcal I = J^\top J$ & Fisher 信息矩阵 & \texttt{fisherInformation} \\
$\Sigma_{\vec\alpha}$ & 参数协方差 & \texttt{paramCovariance} \\
$M$ & 动量空间 Jacobian & \texttt{momentumJacobian} \\
$\Sigma_{\vec p}$ & 动量协方差 (3$\times$3) & \texttt{momentumCovariance} (PDCErrorAnalysis.cc) \\
$\chi^2$ & 加权残差平方和 & \texttt{Evaluate}, \texttt{chi2} \\
$\kappa$ & Fisher 矩阵条件数 & \texttt{fisherConditionNumber} \\
ESS & 有效样本量 & \texttt{computeESS} \\
$z$-Geweke & MCMC 收敛诊断 & \texttt{computeGewekeZ} \\
\bottomrule
\end{tabular}
\end{table}

\section{开放问题与待执行实验清单}
\label{app:open}

\begin{table}[H]
\centering\footnotesize
\caption{本报告中所有 \texttt{\% TODO} 标记汇总; 优先级 P0 = blocker, P3 = nice-to-have}
\label{tab:todos}
\begin{tabular}{cclll}
\toprule
ID & 优先级 & 来源章节 & 内容 & 估计成本 \\
\midrule
R-1 & P0 & II §5  & RK4 中加入 Bethe-Bloch 连续 $dE/dx$
                  & 1--2 周代码 + 1 次 ensemble \\
R-2 & P0 & II §10 & MS 过程噪声并入 $\Sigma_\mathrm{meas}$ (Kalman)
                  & 1--2 周代码 + 1 次 ensemble \\
R-3 & P1 & II §8  & Cartesian $(p_x,p_y,p_z)$ LM 拟合分支
                  & 已立项；写完即可 \\
R-4 & P1 & II §9  & 非对称扩展粗格 + 反向初值
                  & 半周代码 \\
R-5 & P1 & III §3 & LM trace 持久化 (\texttt{lm\_traces/event\_N.csv})
                  & 1 天代码 \\
R-6 & P2 & II §6  & B-field 三次插值 / 全 div-B 闭合检查
                  & 半周代码 \\
R-7 & P2 & II §7  & cosmic-ray PDC 对位 run
                  & 1 周机时 \\
R-8 & P2 & III §5 & B-field $\pm 5\%$ 扫描 (BS-1)
                  & 80 分钟单核 \\
R-9 & P2 & III §5 & $\sigma_\mathrm{wire}$ 扫描 0.1--2.0 mm (WS-1)
                  & 3.5 小时单核 \\
R-10 & P2 & III §5 & $\alpha_\mathrm{tgt}$ 扫描 (TS-1, sanity)
                  & 2 小时单核 \\
R-11 & P2 & III §5 & ensemble 规模扫描 10/50/200/1000 (ES-1)
                  & 13 小时单核 \\
R-12 & P3 & II §3  & PDC 气体混合优化 (Ar/CF$_4$ 替代)
                  & 设计研究 \\
R-13 & P3 & II §5  & E2E 几何全材料调研
                  & 1--2 天 Geant4 配置审计 \\
R-14 & P3 & III §6 & $(-100,0)$ $\chi^2$ 地形 PNG
                  & 1 天画图 \\
\bottomrule
\end{tabular}
\end{table}

\paragraph{下一步建议。}
P0 = R-1 + R-2 同时并行；预期能把 \StatusReport\ § 已知局限 中所列两个主要
症状（$p_z$ $-10.8\,\mathrm{MeV}/c$ 系统偏差、$p_y$ Fisher 欠估 1.9$\times$）
同时关闭。完成后再做 R-3 + R-4 的 P1 二轮，并把 89k 大数据集覆盖率推到 P0
fix 后的水平做最终验证。

% ============================================================
\end{document}
